\documentclass[11pt]{article}
\usepackage[a4paper,margin=1.45cm]{geometry}
\usepackage[T1]{fontenc}
\usepackage[utf8]{inputenc}
\usepackage{lmodern}
\usepackage[spanish,es-noquoting,es-noshorthands]{babel}
\usepackage{amsmath,amssymb,mathtools,array,booktabs,pdflscape,adjustbox,diagbox}
\usepackage[protrusion=true,expansion=false]{microtype}
\setlength{\parindent}{1.25em}
\setlength{\parskip}{0pt}
\newcommand{\widebar}[1]{\overline{#1}}
\newcommand{\A}{\Delta}
\newcommand{\D}{\Delta}
\allowdisplaybreaks
\setlength{\emergencystretch}{12em}
\sloppy
\hbadness=10000
\vbadness=10000
\hfuzz=2pt
\vfuzz=10pt
\raggedbottom
\newcommand{\R}{\varrho}
\newcommand{\hata}{\widehat{a}}
\newcommand{\hatv}{\widehat{v}}
\newcommand{\modu}[1]{\;\operatorname{mod} #1}
\newcommand{\mcell}[1]{\ensuremath{\begin{gathered}#1\end{gathered}}}
\newcommand{\pair}[2]{(#1\,|\,#2)}
\newcommand{\eps}{\varepsilon}
\makeatletter
\renewcommand\footnotesize{\@setfontsize\footnotesize{7.5}{8.7}\selectfont}
\makeatother
\begin{document}

\begin{center}
{\Large\bfseries 1. Sobre la formación del sistema de formas de la forma\\ ternaria bicuadrática}\par
\vspace{1.5em}
\emph{Sitzungsberichte der Physikalisch-medizinischen Sozietät in Erlangen} 39 (1907), pp. 176--179
\end{center}

\vspace{3em}
\begin{center}
(Extracto de la disertación de la autora.)
\end{center}

Los trabajos de Gordan, Maisano y Pascal\footnote{P. Gordan, ``Über das volle Formensystem der ternären biquadratischen Form'' \(f=x_1^3x_2+x_2^3x_3+x_3^3x_1\) (Math. Annalen, vol. XVII, pp. 217--233, 1880); G. Maisano, 1. ``Sistemi completi dei primi cinque gradi della forma ternaria biquadratica e degl' invarianti, covarianti e contravarianti di sesto grado'' (Giorn. di Battaglini XIX). 2. ``Sui covarianti indipendenti di \(6^\circ\) grado nei coefficienti della forma biquadratica ternaria'' (Rend. Circ. Mat. di Palermo I, 1887); E. Pascal, ``Contributo alla teoria della forma ternaria biquadratica e delle sue varie decomposizioni in fattori'' (memoria premiada por la R. Accademia delle scienze fisiche e matematiche di Napoli, 1905).} se ocupan del sistema de formas de la forma ternaria bicuadrática. Gordan establece, para la forma automorfa especial
\[
f=x_1^3x_2+x_2^3x_3+x_3^3x_1,
\]
el sistema completo de formas, compuesto de 54 formaciones, sobre la base de principios análogos a los que había dado para los sistemas de formas en el dominio binario.

En el trabajo de Maisano, para la forma bicuadrática general, se establecen las formas hasta el quinto orden inclusive\footnote{Por ``orden'' debe entenderse la dimensión en los coeficientes, y por ``grado'' la dimensión en las variables.}, junto con algunos invariantes, covariantes y contravariantes de orden superior, según el método empleado por Gordan en el volumen I de los \emph{Mathematische Annalen} para la forma ternaria cúbica. Pascal, utilizando los resultados de Maisano, se ocupa principalmente de la cuestión de la descomposición de la forma bicuadrática en factores\footnote{En su segunda nota breve Maisano intenta demostrar una dependencia lineal entre los tres covariantes de orden 6 y grado 6. Frente a esta demostración no del todo completa, Pascal cree haber probado la independencia lineal de los tres covariantes de orden 6, así como la de los tres invariantes de orden 9. Pero que en efecto existe una relación lineal entre los tres covariantes, por una parte, y los tres invariantes, por otra, resultó de mis cálculos; en la notación de Pascal, las fórmulas explícitas son
\[
0=20\Omega_1+6\Omega_2-3\Omega_3+4fC_1-12fC_2-4A\cdot\Delta,
\]
\[
0=4A^3-15AB+30C-90D+9E.
\]}.

\textbf{El objetivo de mis investigaciones es establecer el sistema de formas de la forma ternaria bicuadrática general; por lo pronto se dan sólo los fundamentos principales y se establece un llamado ``sistema relativamente completo''\footnote{Para la terminología, véase Gordan--Kerschensteiner, \emph{Vorlesungen über Invariantentheorie}, vol. II, p. 227.}.}

La idea básica para la formación de sistemas de formas ternarias es la misma que en el dominio binario. Partiendo de un primer sistema relativamente completo, a saber, el sistema de formas tomado del caso binario, se pasa, según una ley determinada, a sistemas con módulo cada vez más alto, hasta que o bien el sistema de un módulo, tomado éste como forma fundamental, se vuelve finito y conocido, o bien un módulo puede reducirse a formas que tienen invariantes como factores. Por transvección del sistema relativamente completo con el sistema del módulo, en el primer caso se obtiene el sistema absolutamente completo; en el segundo, el sistema relativamente completo y el absolutamente completo coinciden. Por la demostración general de Hilbert de la finitud de los sistemas de formas, este procedimiento debe necesariamente llegar a término.

En nuestro caso, el módulo \((abc)\) del primer sistema relativamente completo puede reducirse a los módulos
\[
\Delta=(abc)^2a_x^2b_x^2c_x^2\quad\text{y}\quad \nu=(abu)^4;
\]
de aquí se obtiene fácilmente el sistema relativamente completo módulo \(\nu\). Como sucesión de módulos elegimos ahora las formas
\[
\nu;\qquad \nu^{(\nu)}=(\nu\nu_1x)^4=s_x^4,
\]
\[
\nu^{(s)}=(ss'u)^4\ldots,
\]
y añadimos, como módulos adicionales, dos formas cuadráticas que aparecen en la formación del sistema relativamente completo módulo \(s\), a saber, \(u_\sigma^2\) y \(t_x^2\). Entonces se puede mostrar que el módulo que sigue a \(s\), es decir \((ss'u)^4\), es reducible al módulo \((\varrho,t)\). Pero como el sistema simultáneo de dos formas cuadráticas es finito y conocido, se alcanza así el final descrito arriba. Como sistema relativamente completo módulo \((\varrho,t)\) se obtienen 331 formaciones.

Añado algunos detalles sobre el método empleado.

El proceso fundamental para producir formas, el proceso de plegamiento, puede definirse en el dominio ternario de la siguiente manera.

Si en un producto simbólico
\[
s_x^m t_x^n u_\sigma^\mu u_\tau^\nu
\]
se reemplazan los pares de factores
\[
\begin{array}{c|c|c|c|c}
 & s_xt_x & u_\sigma u_\tau & s_xu_\sigma\ \text{o}\ s_xu_\tau & t_xu_\tau\ \text{o}\ t_xu_\sigma\\
\text{respectivamente por} & (stu) & (\sigma\tau x) & s_\sigma\ \text{o}\ s_\tau & t_\tau\ \text{o}\ t_\sigma\\
\text{plegamiento} & \mathrm{I} & \mathrm{II} & \mathrm{III} & \mathrm{IV},
\end{array}
\]
entonces las formas así obtenidas han surgido de la forma original por plegamiento\footnote{Gordan, Math. Annalen, vol. XVII, p. 219.}.

Aquí todavía se puede poner siempre \(s_\sigma=0;\ t_\tau=0\). Para la relación entre los plegamientos individuales vale el siguiente teorema:

\textbf{Los plegamientos I y II son plegamientos fundamentales, a partir de los cuales, independientemente del orden de composición, pueden componerse los plegamientos III y IV. En otras palabras: para formar todas las formas que surgen de una expresión dada por plegamiento, sólo hay que aplicar los plegamientos I y II}\footnote{El teorema tiene una excepción en el caso en que \(n\) y \(\mu\), respectivamente \(m\) y \(\nu\), se anulan simultáneamente; la ``sucesión de formas'' se reduce entonces al término diagonal.}.

Por una ``sucesión de formas'' entendemos una forma inicial junto con todas las formas que surgen de ella por plegamiento en sí misma. Por el teorema, una sucesión de formas
\[
s_x^m t_x^n u_\sigma^\mu u_\tau^\nu
\]
puede disponerse en un esquema rectangular. Al avanzar en él se obtiene:
\begin{enumerate}
\item al pasar una columna hacia la derecha, todas las formas que surgen por plegamiento I de las formas adyacentes;
\item al pasar una fila hacia abajo, todas las formas que surgen por plegamiento II de las formas situadas encima, y por tanto todas las formas que surgen por plegamiento.
\end{enumerate}

Bajo estas hipótesis, los teoremas sobre reductores pueden formularse en su forma más general, con la definición
\begin{center}
``Un reductor es una sucesión de formas reducible,''
\end{center}
del siguiente modo:

\begin{center}
\textbf{Si la forma inicial de una sucesión de formas es reducible porque uno de sus términos ha surgido por plegamiento con un reductor, y si la forma final de la sucesión de formas ha surgido de ese mismo término por plegamiento, entonces toda la sucesión de formas es reducible.}
\end{center}

Como métodos adicionales de reducción hay que mencionar:
\begin{enumerate}
\item la llamada ``doble reducción'', es decir, la reducción de una forma de dos maneras para obtener una relación entre las formas superiores;
\item el ``plegamiento con formas descompuestas'', que en parte tiene el carácter de reducción por reductores y en parte el de ``doble reducción''.
\end{enumerate}

\clearpage

\begin{center}
{\Large\bfseries 2. Sobre la formación del sistema de formas de la forma\\ bicuadrática ternaria}\par
\vspace{1.5em}
Journal f. d. reine u. angew. Math. 134 (1908), pp. 23--90 y dos tablas
\end{center}

\vspace{3em}

\begin{center}
\textbf{Introducción.}
\end{center}

Del sistema de formas de la forma bicuadrática ternaria tratan trabajos de Gordan, Maisano y Pascal\footnote{Un breve extracto de la introducción y del capítulo I apareció en los \emph{Sitzungsberichte der phys.-med. Sozietät Erlangen 1907}, pp. 176--179.}\footnote{P. Gordan, sobre el sistema completo de formas de la forma bicuadrática ternaria
\[
f=x_1^3x_2+x_2^3x_3+x_3^3x_1.
\]
(\emph{Math. Annalen}, vol. XVII (1880), pp. 217--233.)
G. Maisano, 1) \emph{Sistemi completi dei primi cinque gradi della forma ternaria biquadratica e degl' invarianti, covarianti e contravarianti di sesto grado}. (Giorn. di Battaglini XIX.)
2) \emph{Sui covarianti indipendenti di $6^\circ$ grado nei coefficienti della forma biquadratica ternaria}. (Rend. Circ. Mat. di Palermo I, 1887.)
E. Pascal, \emph{Contributo alla teoria della forma ternaria biquadratica e delle sue varie decomposizioni in fattori}. (Memoria premiada por la R. Accademia delle scienze fisiche e matematiche di Napoli. 1905.)}. El señor Gordan establece el sistema completo de formas, compuesto de 54 formaciones, de la forma automorfa especial:
\[
f=x_1^3x_2+x_2^3x_3+x_3^3x_1
\]
sobre la base de principios análogos a los que había dado para los sistemas de formas en el dominio binario.

En el señor Maisano se establecen, para la forma bicuadrática general, las formas hasta el quinto orden\footnote{Por ``orden'' debe entenderse la dimensión en los coeficientes, y por ``grado'' la dimensión en las variables.} inclusive, así como algunos invariantes, covariantes y contravariantes de orden superior, según el método aplicado por el señor Gordan en el volumen I de los \emph{Math. Annalen} a la forma cúbica ternaria. El señor Pascal, usando los resultados de Maisano, se ocupa principalmente de la cuestión de la descomposición de la forma bicuadrática en factores\footnote{En la segunda nota breve, el señor Maisano intenta demostrar una dependencia lineal de los tres covariantes de sexto orden y sexto grado. Frente a esta demostración no del todo completa, el señor Pascal cree haber demostrado la independencia lineal de los tres covariantes de sexto orden, así como la de los tres invariantes de noveno orden. Pero que en efecto existe una relación lineal entre los tres covariantes, por una parte, y los tres invariantes, por otra, lo muestran las fórmulas (13), § 11, y (22), nota del § 17, del presente trabajo, donde esas relaciones se dan explícitamente. Según una comunicación del señor Pascal, la contradicción se explica por un error numérico en la expresión de su contravariante $p$ (aquí llamada $\varrho$), p. 46 de su memoria.}.

\emph{El objetivo de las investigaciones que siguen es establecer el sistema de formas de la forma bicuadrática ternaria general; más precisamente, en este trabajo sólo se dan los fundamentos principales y se establece un llamado ``sistema relativamente completo''\footnote{Para la denominación, véase § 3a.}.}

El trabajo se vincula estrechamente con el trabajo de Gordan; no obstante, los principios que allí sólo se daban en forma muy general debían primero desarrollarse en detalle. Con ayuda de un teorema sobre la conexión de los plegamientos y mediante la introducción de la ``sucesión de formas'' (§ 1), los teoremas de reducción se formulan de manera precisa y se completan (§ 3), mientras que el establecimiento recurrente de desarrollos en serie especiales (§ 2) da el medio de cálculo para realizar efectivamente las reducciones.

La idea fundamental de la formación de sistemas de formas ternarias es la misma que en el dominio binario. Partiendo de un primer sistema relativamente completo ---el sistema de las formas tomadas del caso binario--- se llega, según una ley determinada, a sistemas con módulo cada vez más alto, hasta que el sistema de un módulo ---tomado el módulo como forma fundamental--- se vuelva finito y conocido, o bien hasta que un módulo pueda reducirse a formas que tengan invariantes como factores. Por transvección del sistema relativamente completo con el sistema del módulo surge, en el primer caso, el sistema absolutamente completo, mientras que en el segundo caso el sistema relativamente completo y el sistema absolutamente completo se vuelven idénticos. A consecuencia de la finitud de los sistemas de formas, demostrada en general por el señor Hilbert, este procedimiento debe necesariamente terminar.

En nuestro caso, el módulo $(abc)$ del primer sistema relativamente completo (§ 4) se reduce a los módulos
\[
\Delta=(abc)^2a_x^2b_x^2c_x^2\quad\text{y}\quad \nu=(abu)^4
\]
(§ 5), y luego se encuentra el sistema relativamente completo módulo $\nu$ (§ 6). Estos resultados ya se dan en el trabajo de Gordan para la forma bicuadrática general; sin embargo, nuestro modo de derivación se transporta más fácilmente a formas de grado superior.

Como sucesión de módulos elegimos ahora las formas (§ 7):
\[
\nu,\qquad \nu(\nu)=(\nu\nu_1x)^4=s_x^4,\qquad \nu(s)=(ss'u)^4,\ldots .
\]

En la formación del sistema relativamente completo módulo $s$, se adjuntan como módulos dos formas cuadráticas que aparecen en el sistema, $u_\varrho^2$ y $t_x^2$ (§§ 9 y 10), y por consiguiente se forma el sistema relativamente completo módulo $(s,\varrho,t)$; después se muestra (§ 17) que el módulo $(ss'u)^4$ es reducible a los módulos $\varrho$ y $t$. De ese modo, el sistema inmediatamente superior pasa a un sistema relativamente completo módulo $(\varrho,t)$. Pero como el sistema simultáneo de dos formas cuadráticas es finito y conocido, y en el caso general consta de 20 formaciones\footnote{Cf. Clebsch--Lindemann, \emph{Vorlesungen über Geometrie}. (vol. I, p. 288 ss.) Sistema de dos formas cogredientes. Gordan, \emph{Über Büschel von Kegelschnitten}. (\emph{Math. Ann.} vol. XIX, p. 530.) Sistema de dos formas contragredientes.}, con el establecimiento del sistema relativamente completo módulo $(\varrho,t)$ se alcanza la terminación caracterizada arriba. Queda reservada la transvección de este sistema con el sistema de $(\varrho,t)$ para formar el sistema absolutamente completo.

Como sistema relativamente completo módulo $(\varrho,t)$ se encuentran 331 formaciones, ordenadas en la tabla adjunta según su grado en las variables $x$ y $u$.

Finalmente damos todavía una visión general de los símbolos introducidos:
\begin{align*}
f&=a_x^4,\\
\theta&=\theta_x^4u_\vartheta^2=(abu)^2a_x^2b_x^2,\\
K&=K_x^6u_k^3=(a\theta u)u_\vartheta^2a_x^3\theta_x^3,\\
j&=u_j^6=(a\theta u)^4u_\vartheta^2,\\
\Delta&=\Delta_x^6=a_\vartheta^2a_x^2\theta_x^4=(abc)^2a_x^2b_x^2c_x^2,\\
N&=N_x^8u_\eta=(a\Delta u)a_x^3\Delta_x^5,\\
\nu&=u_\nu^4=(abu)^4,\\
H&=H_x^2u_\eta^4=(\nu\nu_1x)^2u_\nu^2u_{\nu_1}^2,\\
L&=L_x^3u_l^6=(\nu\eta x)H_x^2u_\nu^3u_\eta^3,\\
g&=g_x^6=(\nu\eta x)^4H_x^2,\\
\sigma&=u_\sigma^6=H_\nu^2u_\nu^2u_\eta^4,\\
s&=s_x^4=(\nu\nu_1x)^4,\\
Z&=Z_x^4u_\zeta^2=(ss'u)^2s_x^2s_x'^2,\\
\varrho&=u_\varrho^2=\theta_\nu^4u_\vartheta^2,\\
t&=t_x^2=a_\eta^4H_x^2,\\
i&=a_\nu^4,\\
J&=s_\nu^4.
\end{align*}

\begin{center}
{\Large\bfseries Capítulo I. \quad Teoremas generales sobre formas ternarias.}
\end{center}

\begin{center}
\textbf{§ 1. \quad Proceso de plegamiento. \quad Sucesiones de formas.}
\end{center}

El proceso fundamental para producir formas es el proceso de plegamiento, que en el dominio ternario se define de la siguiente manera. Sea dado un producto simbólico:
\[
s_x^m t_x^n u_\sigma^\mu u_\tau^\nu .
\]
Si se reemplazan los pares de factores
\[
s_xt_x\qquad u_\sigma u_\tau\qquad s_xu_\sigma\ \text{o}\ s_xu_\tau
\qquad t_xu_\tau\ \text{o}\ t_xu_\sigma
\]
respectivamente por
\[
(stu)\qquad (\sigma\tau x)\qquad s_\sigma\ \text{o}\ s_\tau
\qquad t_\tau\ \text{o}\ t_\sigma,
\]
plegamiento
\[
\text{I}\qquad\text{II}\qquad\text{III}\qquad\text{IV},
\]
entonces las formas resultantes han surgido de la forma originaria por plegamiento\footnote{Gordan, \emph{Math. Annalen} vol. XVII, p. 219.}.

La expresión $s_x^m t_x^n u_\sigma^\mu u_\tau^\nu$ puede concebirse ya sea como el producto efectivo de dos formas $S\cdot T$, ya sea como una sola forma con varias series de símbolos:
\[
s_x^m t_x^n u_\sigma^\mu u_\tau^\nu=A_x^{m+n}u_x^{\mu+\nu}.
\]
En el primer caso hablamos del plegamiento de la forma $S$ con $T$, en el segundo del ``plegamiento de la forma $A$ en sí misma''. Para abreviar, en lo que sigue supondremos siempre (lo cual de hecho ocurre para todas las formas consideradas más adelante) que
\[
\tag{a}
s_\sigma^\lambda=0,\qquad t_\tau^\lambda=0
\]
para todo $\lambda$ no nulo y cualquier $x$, y junto con cualesquiera otros plegamientos. Esto dice que la forma $s_x^m t_y^n u_\sigma^\mu u_\tau^\nu$ es una llamada ``forma normal'', es decir, que se anula al aplicar el proceso $\Omega$, tanto con respecto a las variables $x$ y $u$ como con respecto a las variables $y$ y $v$.

Por tanto, la condición (a) no representa ninguna restricción, sino que siempre puede alcanzarse\footnote{Cf. Gordan, \emph{Math. Annalen} vol. V, p. 104.}.

Para la conexión entre los distintos plegamientos vale el siguiente teorema:

\emph{Teorema I. Los plegamientos I y II son plegamientos fundamentales a partir de los cuales, independientemente del orden de composición, pueden componerse los plegamientos III y IV. En otras palabras: para formar todas las formas que surgen por plegamiento de una expresión dada, basta aplicar los plegamientos I y II.}

Demostración: por el teorema de identidad, respectivamente por el teorema del producto para matrices, teniendo en cuenta (a), se obtiene
\[
(\widehat{st}\sigma x)=-t_\sigma,\qquad
(\widehat{\sigma\tau}su)=-s_\tau,
\]
\[
(\widehat{st}\tau x)=s_\tau,\qquad
(\widehat{\sigma\tau}tu)=t_\sigma,
\]
\[
(stu)(\sigma\tau x)=s_\tau+t_\sigma-u_x\,s_\tau t_\sigma .
\]
Así, por composición de los plegamientos I y II surgen los plegamientos III y IV, tanto cuando se aplica el plegamiento II al plegamiento I, esto es, a los factores $(stu)u_\sigma u_\tau$, como cuando se aplica el plegamiento I al plegamiento II, a los factores $(\sigma\tau x)s_xt_x$.

Llamamos \emph{sucesión de formas}\footnote{Cf. Clebsch, \emph{Über eine Fundamentalaufgabe der Invariantentheorie} (Abh. der Gött. Ges. d. Wiss. vol. XVII): § 17 y final del § 18. La ``sucesión de formas'' difiere del ``sistema equivalente propiamente reducido'' introducido allí por el principio de ordenación según formas superiores.} a una forma inicial junto con todas las formas que han surgido de ella por plegamiento en sí misma, y designamos la sucesión de formas por la forma inicial. Aquí se entenderá por ``forma superior'' toda forma con más plegamiento en sí misma; entre los plegamientos igualmente legítimos $s_\tau$ y $t_\sigma$, uno ha de normalizarse como superior. De la ley de formación de la sucesión de formas se sigue:

1) Cada forma de la sucesión de formas es lineal en los símbolos de la forma inicial.

2) La sucesión de formas está determinada unívocamente para formas iniciales que tienen dos series de símbolos contragredientes cada una; para formas iniciales con más series de símbolos debe normalizarse unívocamente distinguiendo plegamientos especiales entre aquellos que están conectados por el teorema de identidad.

Por el teorema I, una sucesión de formas $s_x^m t_x^n u_\sigma^\mu u_\tau^\nu$ ($m>n;\ \mu>\nu$) puede disponerse según el siguiente esquema rectangular:
\begingroup\small
\[
\begin{array}{c|c|c|c}
s_x^m t_x^n u_\sigma^\mu u_\tau^\nu
& (stu)
& (stu)^2
& \cdots\\[0.3em]
(\sigma\tau x)
& t_\sigma,\ s_\tau
& t_\sigma(stu),\ s_\tau(stu)
& \cdots\\[0.3em]
(\sigma\tau x)^2
& t_\sigma(\sigma\tau x),\ s_\tau(\sigma\tau x)
& t_\sigma^2,\ t_\sigma s_\tau,\ s_\tau^2
& \cdots\\[0.3em]
\vdots&\vdots&\vdots&\ddots\\[0.3em]
(\sigma\tau x)^\nu
& t_\sigma(\sigma\tau x)^{\nu-1},\ s_\tau(\sigma\tau x)^{\nu-1}
& t_\sigma^2(\sigma\tau x)^{\nu-2},\ t_\sigma s_\tau(\sigma\tau x)^{\nu-2},\ s_\tau^2(\sigma\tau x)^{\nu-2}
& \cdots\\[0.3em]
\vdots
& t_\sigma(\sigma\tau x)^\nu
& t_\sigma^2(\sigma\tau x)^{\nu-1},\ t_\sigma s_\tau(\sigma\tau x)^{\nu-1}
& \cdots
\end{array}
\]
\endgroup
y la parte inferior correspondiente, continuada como sigue\footnote{Aquí, y análogamente en lo que sigue, la parte inferior marcada con un asterisco debe colocarse en el lugar igualmente marcado, arriba a la derecha del esquema.}
\begingroup\small
\[
\begin{array}{c|c|c}
(stu)^n & \cdots & \cdots\\[0.3em]
t_\sigma(stu)^{n-1},\ s_\tau(stu)^{n-1}
& s_\tau(stu)^n & \cdots\\[0.3em]
t_\sigma^2(stu)^{n-2},\ t_\sigma s_\tau(stu)^{n-2},\ s_\tau^2(stu)^{n-2}
& t_\sigma s_\tau(stu)^{n-1},\ s_\tau^2(stu)^{n-1}
& s_\tau^2(stu)^n .
\end{array}
\]
\endgroup

Aquí se obtiene, cuando se avanza

1) una columna hacia la derecha, todas las formas que surgen por plegamiento I de las formas adyacentes,

2) una fila hacia abajo, todas las formas que surgen por plegamiento II de las formas situadas encima,

y por tanto todas las formas que surgen por plegamiento.

Una forma cualquiera de la sucesión total, $t_\sigma^\kappa s_\tau^\lambda(stu)^\mu$ o $t_\sigma^\kappa s_\tau^\lambda(\sigma\tau x)^\nu$, constituye el comienzo de una nueva sucesión de formas, compuesta por todas aquellas formas del esquema rectangular que tienen $t_\sigma^\kappa s_\tau^\lambda(stu)^\mu$ o $t_\sigma^\kappa s_\tau^\lambda(\sigma\tau x)^\nu$ como esquina superior izquierda y poseen el factor $t_\sigma^\kappa s_\tau^\lambda$.

Observación final. El teorema I tiene una excepción en el caso en que los números $n$ y $\mu$ (respectivamente $m$ y $\nu$) se vuelven nulos, de modo que los plegamientos I, II y IV ya no existen, mientras que el plegamiento III (respectivamente IV) todavía existe. En este caso designamos como sucesión de formas el miembro diagonal del esquema:
\[
\begin{array}{ccccc}
s_x^m u_\tau^\nu &&&& t_x^n u_\sigma^\mu\\
& s_\tau && t_\sigma &\\
& s_\tau^2 && t_\sigma^2 &\\
&&\ddots&&\\
&&s_\tau^\nu && t_\sigma^\mu .
\end{array}
\]
De acuerdo con la ausencia de los plegamientos I y II, el desarrollo en serie según las polares de la sucesión de formas se reduce entonces siempre al miembro inicial; sin embargo, los teoremas sobre reductores siguen siendo válidos.

\begin{center}
\textbf{§ 2. \quad Desarrollos en serie según polares de la sucesión de formas.}
\end{center}

Para formas con $n$ variables se han establecido explícitamente dos clases de desarrollos en serie\footnote{Cf. Gordan, \emph{Über Kombinanten}, §§ 2 y 5 (\emph{Math. Ann.} vol. V).}:

1) para formas con una serie cada una de variables contragredientes, $s_x^m u_\tau^\nu$, una serie que progresa según potencias de $u_x$, cuyos coeficientes son ``formas normales'';

2) para formas con dos series de variables cogredientes, $s_x^m t_x^n$, un desarrollo según las polares de los covariantes elementales (polares de la sucesión de formas $s_x^m t_x^n$), que en el caso ternario tiene la forma siguiente\footnote{Cf. Study, \emph{Methoden zur Theorie der ternären Formen} (Teubner 1889) II. § 7. A diferencia de la derivación dada allí, la nuestra proporciona el método para determinar rápidamente por cálculo los coeficientes numéricos $C_{pr\varrho}$. El desarrollo de Clebsch--Study, bajo la especialización a), a'), sería
\[
(stu)^\mu u_\tau^{\nu-\kappa}v_\tau^\kappa s_x^m t_x^{\,n-\lambda}(tuv)^\lambda
=\sum_{p,r,\varrho} C_{pr\varrho}
\bigl[s_\tau^\varrho(stu)^{\mu+r-\varrho}\bigr]_{\widehat{uv}^{\lambda-r+\varrho}v^{\nu-\varrho},\,u^p,\,(vw=\widehat{xw})},
\]
y por tanto no permite la simplificación que en nuestro desarrollo aparece ya durante el cálculo, cuando se sabe que han de omitirse todos los términos que contienen el factor $u_x$.}:
\[
\tag{I}
s_x^m t_x^{n-\lambda}t_y^\lambda
=\sum_{\varrho}
\frac{\binom{m}{\varrho}\binom{\lambda}{\varrho}}
{\binom{m+n-\varrho+1}{\varrho}}
\bigl[(st\widehat{xy})^\varrho s_x^{m-\varrho}t_x^{n-\varrho}\bigr]_{y^{\lambda-\varrho}}.
\]

Combinamos los dos desarrollos estableciendo, para formas con dos series cada una de variables contragredientes,
\[
s_x^m t_x^{n-\lambda}t_y^\lambda u_\sigma^\mu u_\tau^{\nu-\kappa}v_\tau^\kappa
\]
desarrollos según las potencias de $u_x$, cuyos coeficientes son polares compuestas de la sucesión de formas $s_x^m t_x^n u_\sigma^\mu u_\tau^\nu$, tomadas según las variables $x$ y $u$.

Basta establecer la serie para dos series contragredientes de símbolos (respectivamente de variables), pues, reuniendo por pares las series de símbolos (de variables) en una nueva serie única, las formas con más series de símbolos (de variables) se reducen a este caso.

Para que todas las formas de la sucesión de formas contengan sólo dos series de símbolos, respectivamente de variables, especializamos:
\[
\begin{array}{ll}
\text{a) } \sigma=\widehat{st}, & \text{a') } y=\widehat{uv},\\[0.3em]
\text{b) } s=\widehat{\sigma\tau}, & \text{b') } v=\widehat{xy},
\end{array}
\]
y efectuamos el desarrollo para la especialización a), a').

Por transferencia directa del desarrollo I se obtienen polares compuestas para las formas $(stu)^\mu s_x^m t_x^{n-\lambda}t_y^\lambda$, mientras que para las formas $(stu)^\mu u_\tau^{\nu-\kappa}v_\tau^\kappa s_x^m t_x^{n-\lambda}t_y^\lambda$, en lugar de polares compuestas, aparecen términos cuya evaluación hace necesaria la introducción de variables auxiliares siempre nuevas, que sólo han de identificarse al final, lo cual complica innecesariamente el cálculo.

Por eso seguimos un camino indirecto: representamos las polares compuestas de todas las formas de la sucesión de formas, es decir, las expresiones
\[
\bigl[s_\tau^\varrho(stu)^{\mu-\varrho+r}\bigr]_{y^\lambda}^{v^\kappa}
\]
por los términos iniciales de esas polares, y obtenemos después el desarrollo buscado según polares mediante inversión del sistema recurrente de ecuaciones que resulta.

Según el principio de transferencia, para la $\lambda$-ésima polar de una forma $s_x^m t_x^n:[s_x^m t_x^n]_{y^\lambda}$ ($\lambda\le n$) vale el desarrollo siguiente (el caso $\lambda>n$ se reduce al primero por el desarrollo de $[s_y^m t_y^n]_{x^{m+n-\lambda}}$)\footnote{Gordan, \emph{Die Resultante binärer Formen} (cap. I, § 5). Rend. Circ. Mat. di Palermo XXII. 1906.}:
\[
\bigl[s_x^m t_x^n\bigr]_{y^\lambda}
=\sum_r (-1)^r\,
\frac{\binom{m}{r}\binom{\lambda}{r}}{\binom{m+n}{r}}\,
(st\widehat{xy})^r s_x^{m-r}t_x^{n-\lambda}t_y^{\lambda-r},
\]
y del mismo modo para la $\kappa$-ésima polar de $u_\sigma^\mu u_\tau^\nu:[u_\sigma^\mu u_\tau^\nu]_{v^\kappa}$,
\[
\bigl[u_\sigma^\mu u_\tau^\nu\bigr]_{v^\kappa}
=\sum_\varrho (-1)^\varrho\,
\frac{\binom{\mu}{\varrho}\binom{\kappa}{\varrho}}{\binom{\mu+\nu}{\varrho}}\,
(\sigma\tau\widehat{uv})^\varrho u_\sigma^{\mu-\varrho}u_\tau^{\nu-\kappa}v_\tau^{\kappa-\varrho}.
\]

Por multiplicación se sigue de ello, introduciendo las especializaciones a) y a'),
\[
\begin{aligned}
\bigl[s_x^m t_x^n(stu)^\mu u_\tau^\nu\bigr]_{\widehat{uv}^{\lambda}}^{v^\kappa}
={}&
\sum_{r,\varrho}c_{r\varrho}\,
(st\widehat{uv})^r(st\widehat{uv})^\varrho
s_x^{m-r}t_x^{n-\lambda}(tuv)^{\lambda-r}(stu)^{\mu-\varrho}
u_\tau^{\nu-\kappa}v_\tau^{\kappa-\varrho},
\end{aligned}
\]
y de ahí, desarrollando según potencias de $u_x$, distinguiendo el primer término ($\sum'_{r,\varrho}$ significa que se omite el término $r=0$, $\varrho=0$) y teniendo en cuenta (a) p. 26,
\[
\tag{II}
\begin{aligned}
&s_x^m t_x^{n-\lambda}(tuv)^\lambda(stu)^\mu u_\tau^{\nu-\kappa}v_\tau^\kappa\\
&\quad =
\bigl[s_x^m t_x^n(stu)^\mu u_\tau^\nu\bigr]_{\widehat{uv}^{\lambda}}^{v^\kappa}\\
&\qquad
-\sum_{r,\varrho}' c_{r\varrho}\,s_\tau^\varrho(stu)^{\mu-\varrho+r}
u_\tau^{\nu-\kappa}v_\tau^{\kappa-\varrho}
s_x^{m-r}t_x^{n-\lambda}(tuv)^{\lambda-r+\varrho}v_x^r\\
&\qquad
+u_x\sum_{\varrho}\sum_{r=1}^{\lambda} c_{r\varrho}\,
s_\tau^\varrho(stu)^{\mu-\varrho+r-1}(stv)
u_\tau^{\nu-\kappa}v_\tau^{\kappa-\varrho}
s_x^{m-r}t_x^{n-\lambda}(tuv)^{\lambda-r+\varrho}v_x^{r-1}\\
&\qquad
\vdots\\
&\qquad
-(-1)^\lambda u_x^\lambda\sum_\varrho c_{\lambda\varrho}\,
s_\tau^\varrho(stu)^{\mu-\varrho}(stv)^\lambda
u_\tau^{\nu-\kappa}v_\tau^{\kappa-\varrho}
s_x^{m-\lambda}t_x^{n-\lambda}(tuv)^\varrho,
\end{aligned}
\]
donde
\[
c_{r\varrho}=(-1)^{r+\varrho}\cdot
\frac{\binom{m}{r}\binom{\lambda}{r}}{\binom{m+n}{r}}\cdot
\frac{\binom{\mu}{\varrho}\binom{\kappa}{\varrho}}{\binom{\mu+\nu}{\varrho}}
\qquad
\begin{array}{l}
\text{para }\lambda\le n,\\
\kappa\le\nu.
\end{array}
\]
y de manera análoga para las demás combinaciones de los números
\[\lambda,n;\kappa,\nu.\]
La expresión
\[s_x^m t_x^{n-\lambda}(tuv)^\lambda(stu)^\mu u_\tau^{\nu-\kappa}v_\tau^\kappa\]
queda, por tanto, reducida a una polar compuesta y, aplicando la identidad
\[
(stv)u_\tau=(stu)v_\tau-s_\tau(tuv),
\]
a una suma de expresiones formadas de manera análoga, que corresponden a formas superiores de la sucesión de formas.

Por iteración se llega al desarrollo:
\[
\tag{III}
s_x^m t_x^{n-\lambda}(tuv)^\lambda(stu)^\mu u_\tau^{\nu-\kappa}v_\tau^\kappa
=
\sum_{p,r,\varrho}u_x^p\,C_{pr\varrho}\cdot
\bigl[s_\tau^\varrho(stu)^{\mu-\varrho+r}\bigr]_{\widehat{uv}^{\lambda-r+\varrho}v^{\kappa-\varrho+p},\,v_x^{r-p}}.
\]
Los coeficientes numéricos $C_{pr\varrho}$ se calculan de manera unívoca y recursiva a partir de los coeficientes $c_{r\varrho}$, $c'_{r\varrho}$ del sistema de ecuaciones que resulta de iterar (II.) (véase el ejemplo, p. 42). Desarrollos análogos se obtienen para las demás especializaciones.

\begin{center}
\textbf{§ 3. \quad Teoremas de reducción.}
\end{center}

Los teoremas conocidos sobre la reducción de formas y sistemas de formas han de ponerse ahora en relación con los §§ 1 y 2; para ello son necesarias algunas definiciones tomadas de la teoría de las formas binarias.

a) Designamos como \emph{sistema relativamente completo}, módulo una sucesión prefijada de formas, a un sistema de formas con la propiedad de que todas las formaciones que surgen por plegamiento de un producto arbitrario de esas formas pueden expresarse como funciones racionales enteras de las formas del sistema, salvo un término aditivo formado por expresiones que han surgido por plegamiento de las formas del sistema con el sistema del módulo\footnote{Gordan--Kerschensteiner, \emph{Vorlesungen über Invariantentheorie}. vol. II, p. 227.}.

b) Llamamos \emph{reducible} a una forma cuando puede expresarse por formas que tienen invariantes como factores, o por ``formas superiores''; esto es, por formas superiores de la sucesión total de formas a la que pertenece la forma, o por formas que contienen los símbolos del módulo en orden superior.

Los teoremas sobre reductores\footnote{Gordan, \emph{Math. Annalen} vol. XVII, p. 222.} pueden entonces enunciarse en su forma más general de la siguiente manera:

Definición: \emph{Un reductor es una sucesión de formas reducible.}

\emph{Teorema II. Si la forma inicial de una sucesión de formas es reducible porque uno de sus miembros ha surgido por plegamiento con un reductor (tiene un reductor como factor), y si la forma final de la sucesión de formas ha surgido de ese mismo miembro por plegamiento, entonces toda la sucesión de formas es reducible}\footnote{La simplificación que introduce el teorema II puede verse en el siguiente ejemplo: para la forma especial $f=x_1^3x_2+x_2^3x_3+x_3^3x_1$, $a_y^2a_x^2u_y^2$ es un reductor. De ello se sigue, por el teorema II, la reducción de la sucesión de formas
\[
\theta_\nu(\vartheta\nu x)\theta_x^3u_\vartheta u_\nu^2
=
a_\nu^2(abu)-a_\nu b_\nu(aba),
\]
pues $\theta_\nu^4=a_\nu^2b_\nu^2(abu)^2$ ha surgido del miembro $a_\nu^2(abu)$ por plegamiento. En Gordan, loc. cit. p. 231, en cambio, se efectúa separadamente la reducción de las formas
\[
\theta_\nu(\vartheta\nu x),\quad \theta_\nu(\vartheta\nu x)^2,\quad
\theta_\nu^2,\quad \theta_\nu^2(\vartheta\nu x),\quad
\theta_\nu^2(\vartheta\nu x)^2
\]
.}.

Demostración: la sucesión de formas surge del miembro inicial por plegamiento en sí mismo, por tanto, por un lado, mediante plegamiento superior con el reductor y, por otro, mediante plegamiento del reductor en sí mismo. En ambos casos, según § 2, todas las formas de la sucesión de formas pueden representarse como polares (transvecciones) sobre la sucesión de formas del reductor.

La conclusión desde la forma inicial a toda la sucesión de formas ya no se da en los demás métodos de reducción. Éstos son:

1) la llamada ``reducción doble''.

Entendemos por ella la reducción, de dos maneras, de una expresión que tiene dos reductores independientes entre sí como factores, con lo cual surgen relaciones entre las formas superiores a las que se redujo.

Mediante la aplicación sistemática de la ``reducción doble'' se deben obtener, por una parte, todas las relaciones\footnote{Así se encontraron, por ejemplo, mediante ``reducción doble'': la reducción de la forma $a_\eta^4$ (§ 10), la única forma que Maisano incluyó superfluamente en el sistema; luego, las relaciones mencionadas en la nota a la introducción entre los tres covariantes y los tres invariantes.}; por otra parte, cuando la ``reducción doble'', aplicada a distintas expresiones, no da nuevas relaciones, se tiene a la vez un criterio para la independencia de las formas superiores y un control del cálculo.

2) \emph{Plegamiento con formas descomponibles.}

Por ``formas descomponibles'' entenderemos ---con una ligera generalización del concepto usual--- formas que pueden expresarse por productos de formas de grado inferior y por formas ``superiores''. Los productos de formas pasan, bajo un plegamiento simple, a productos; asimismo, los productos de formas no lineales bajo un plegamiento doble en las especializaciones a') y b') (§ 2), según el teorema del producto:
\[
a_yb_y a_xb_x=\frac12a_y^2b_x^2+\frac12b_y^2a_x^2-\frac12(ab\widehat{yx})^2.
\]
Las primeras polares (transvecciones) y ciertas segundas polares de formas descomponibles son, por tanto, de nuevo formas descomponibles. Se sigue:

Un miembro de una transvección simple (respectivamente doble) sobre una forma descomponible, surgido por plegamiento simple (respectivamente doble) con una forma descomponible, se vuelve él mismo una forma descomponible:

a) si las formas inmediatamente superiores en la sucesión de formas de la forma descomponible se descomponen o se vuelven reducibles, o también si, en un plegamiento simple, se presentan las especializaciones $y=\widehat{uv}$ o $v=\widehat{xy}$;

b) si los demás plegamientos simples conducen a formas superiores (véase § 12, B. $H_k$ y $H_k(k\eta x)$).

Un tal miembro de una transvección puede también descomponerse:

c) si los demás plegamientos simples conducen a formas inferiores que se descomponen independientemente del plegamiento con la forma descomponible, es decir, que pueden expresarse por productos y por la forma que se ha de reducir, pero verificando en cada caso por cálculo que el coeficiente numérico del miembro correspondiente no se anula idénticamente. Esto no es otra cosa que la ``reducción doble'' de la forma inferior (véase § 23, C. $H_q(\theta H u)(\theta s u)$).

Formas descomponibles surgen por todos los plegamientos simples con determinantes funcionales\footnote{cf. Gordan, loc. cit. § 3.}, y además por ciertos plegamientos superiores. Se tiene
\[
(st\widehat{yx})s_y=\frac12\,t\,s_y^2-\frac12\,s\,t_y^2+\frac12(st\widehat{yx})^2,
\]
\[
(st\widehat{yx})^2s_y^2=\frac13\,t\,s_y^3-\frac13\,s\,t_y^3+s_y(st\widehat{yx})^2-\frac13(st\widehat{yx})^3.
\]
Fórmulas análogas valen para las formas dualísticas
\[
(\sigma\tau\widehat{\nu u})v_\sigma
\quad\text{y}\quad
(\sigma\tau\widehat{\nu u})v_\sigma^2.
\]

De los teoremas de este parágrafo resulta la siguiente regla para establecer todas las formas irreducibles que surgen del plegamiento de $S$ con $T$:

Se procede, comenzando por los plegamientos más bajos, por cualquier camino fijado de antemano (por ejemplo, fila por fila o simétricamente respecto del término diagonal), en la formación de la sucesión de formas $ST$, hasta llegar a una forma que tiene un reductor como factor. La sucesión de formas definida por esta forma debe omitirse tan pronto como la forma final satisfaga la condición establecida en el teorema II. Después de excluir todas las sucesiones de formas reducibles, se han de eliminar, mediante la reducción doble, todas las demás formas reducibles, así como todas las formas descomponibles. Los lugares de la sucesión total de formas que quedan doblemente cubiertos por sucesiones de formas reducibles independientes entre sí dan lugar a la reducción de formas superiores (véase § 11, D).

\clearpage
\noindent\textbf{Teorema III.}
\emph{Las formas tomadas del dominio binario, es decir, aquellas formas que surgen de las formas del sistema de la forma binaria correspondiente, según el principio de transferencia de Clebsch, por bordado, forman un sistema relativamente completo módulo $(abc)$.}

\noindent\textbf{Demostración:}
A una relación binaria que afirma que las formas $Q'_1,\ldots,Q'_n$ forman un sistema completo de una forma fundamental binaria,
\[
Q'-F(Q')=0=\sum_\lambda\bigl\{(ab)c_x+(bc)a_x+(ca)b_x\bigr\}\varphi_\lambda^{*}
\]
corresponde por bordado la relación ternaria:
\[
Q-F(Q_i)=\sum_\lambda u_x\cdot(abc)\varphi_\lambda .
\]
Pero ésta es precisamente la ecuación definitoria de un sistema relativamente completo módulo $(abc)$. Para la forma bicuadrática, el sistema de las formas tomadas consta de las siguientes formas:
\[
\begin{gathered}
f=a_x^4,\qquad
\theta=\theta_x^4u_\vartheta^2=(abu)^2a_x^2b_x^2,\qquad
K=K_x^6u_k^3=(a\theta u)u_\vartheta^2a_x^3\theta_x^3,\\
j=u_j^6=(a\theta u)^4u_\vartheta^2,\qquad
\nu=u_\nu^4=(abu)^4,
\end{gathered}
\]
entre las cuales las cuatro primeras forman un sistema relativamente completo módulo $((abc),\nu)$.

\begin{center}
\textbf{§ 5. \quad Reducción del módulo $(abc)$ a los módulos $\A$ y $\nu$.}
\end{center}

El módulo $(abc)$, dado sólo por un factor de paréntesis, debe, de acuerdo con la definición (§ 3, a), reducirse a formas del sistema. En otras palabras: la sucesión de formas $(abc)$ ha de normalizarse unívocamente (cf. § 1).

\[
a_s^2a_xu_x=\frac13 a_s^3u_x=\frac13 S\cdot u_x. \tag{2}
\]
\[
\left.
\begin{aligned}
(a\A u)^2&=a_s-\frac{S}{6}u_x^2,\\
(a\A u)^3&=0.
\end{aligned}
\right\} \tag{3}
\]
\[
\left.
\begin{aligned}
\theta(s)=(ss,x)^2&=2a_t+\frac13 S\cdot\theta-\frac13 Tu_x^2,\\
\theta(t)=(tt,x)^2&=-\frac13 S\cdot a_t+\frac23 T\cdot a_s+\frac1{18}S^2\cdot\theta-\frac1{18}u_x^2\cdot S\cdot T.
\end{aligned}
\right\} \tag{4}
\]

Sucesión de los módulos: $(abc)$; $\A$; $s$; $t$ (el sistema del módulo $t$ consta de la única forma $t$).\footnote{Gordan--Kerschensteiner, loc. cit. p. 134.}

Mostraremos que las formas
\[
\A=(abc)^2a_x^2b_x^2c_x^2,\qquad \nu=(abu)^4
\]
bastan como módulos, es decir, que las formas
\[
\begin{gathered}
\A=(abc)^2a_x^2b_x^2c_x^2,\qquad
a_\nu=(abc)(bcu)^3a_x^3,\\
a_\nu^2=(abc)^2(bcu)^2u_x^2,\qquad
i=a_\nu^4=(abc)^4
\end{gathered}
\]
definen la sucesión de formas $(abc)$. En la sucesión de formas $a_\nu$ falta la forma $a_\nu^3$, según la identidad
\[
a_\nu^3=(abc)^3a_x(bcu)=\frac13 u_x\cdot(abc)^4=\frac13 i\cdot u_x.
\]

Para reducir un módulo a un módulo superior debemos, por definición, reducir los plegamientos más generales, o bien todos los plegamientos especiales que entren en consideración con el módulo, a plegamientos con el módulo superior.

Los plegamientos más generales surgen, en nuestro caso, de las expresiones
\[
(abc)a_x^3b_y^3c_z^3,\qquad
(abc)a_x^2b_y^2c_z^2a_zb_xc_x
\]
(tomadas de las formas cúbicas),
\[
(abc)a_xb_yc_za_x^2b_x^2c_x^2
\]
(tomada de las formas cuadráticas), y por polarización de estas expresiones.

Para calcular estas expresiones por las polares de $\A$ y de la sucesión de formas $a_\nu$, respectivamente por sus términos iniciales, seguimos, como en § 2, el camino indirecto. Desarrollamos los términos polares
\[
(abc)^2a_x^2b_y^2c_z^2,\qquad
a_\nu(\nu xy)(\nu yz)(\nu zx)a_xa_ya_z
=(abc)(bc\widehat{x}u)(bc\widehat{y}z)(bc\widehat{z}x)a_xa_ya_z\quad\text{etc.}
\]
como agregados lineales de expresiones que tienen el factor simbólico $(abc)$, y por inversión del sistema de ecuaciones obtenemos las relaciones buscadas, así como una relación entre las polares tomadas según diferentes combinaciones de las variables. Para la evaluación sirven el teorema de identidad y el teorema del producto para determinantes.

El sistema de ecuaciones para $(abc)a_x^3b_y^3c_z^3$ es:
\begingroup\small
\[
\begin{aligned}
1)\quad
\sum_3 a_\nu(\nu yz)^3a_x^3
&=6(abc)a_x^3b_y^3c_z^3
-6\sum_3(abc)a_x^3b_y^2b_zc_z^2c_y^{*},\\[0.4em]
2)\quad
3(abc)^2a_x^2b_y^2c_z^2\cdot(xyz)
-\frac12\sum_3 a_\nu^2(\nu yz)^2\nu_x^2(xyz)
&=2\sum_3(abc)a_x^3b_y^2b_zc_z^2c_y\\
&\quad-4\sum_3(abc)a_xa_ya_zb_y^2b_zc_z^2c_x,\\[0.4em]
3)\quad
a_\nu(\nu xy)(\nu yz)(\nu zx)a_xa_ya_z
&=2\sum_3(abc)a_xa_ya_zb_y^2b_zc_z^2c_x,\\[0.4em]
4)\quad
\sum_3 a_\nu(\nu yz)^3u_x^3
-\frac32\sum_3 a_\nu^2(\nu yz)^2\nu_x^2(xyz)
+\frac16 i\cdot(xyz)^3
&=6\sum_3(abc)a_xa_ya_zb_y^2b_zc_z^2c_x.
\end{aligned}
\]
\endgroup

De ello resulta para
\[
(abc)a_x^3b_y^3c_z^3,
\]
y, por un cálculo análogo, para
\[
(abc)a_x^2b_y^2c_z^2a_zb_xc_x,\qquad
(abc)a_xb_yc_za_z^2b_x^2c_x^2:
\]
\[
\begin{aligned}
(abc)a_x^3b_y^3c_z^3
&=\frac32(abc)^2a_x^2b_y^2c_z^2(xyz)
+\frac32 a_\nu(\nu xy)(\nu yz)(\nu zx)a_xa_ya_z\\
&\quad-\frac1{12}i\cdot(xyz)^3,\\[0.4em]
(\text{IV.})\quad
(abc)a_x^2b_y^2c_z^2a_zb_xc_x
&=\frac23(abc)^2a_x^2b_y^2c_z^2c_x\cdot(xyz)\\
&\quad+\frac13a_\nu(\nu xy)(\nu yz)(\nu zx)a_x^3
-\frac16 a_\nu^2(\nu xy)(\nu zx)a_x^2\cdot(xyz),\\[0.4em]
(abc)a_xb_yc_za_z^2b_x^2c_x^2
&=\frac16(abc)^2a_x^2b_z^2c_z^2\cdot(xyz).
\end{aligned}
\]

Hemos obtenido así un sistema relativamente completo módulo $(\A,\nu)$, compuesto de las cuatro formas $f,\theta,K,j$. De ello resulta que todas las transvecciones de $f$ sobre $\theta$ que no se han tomado del dominio binario, así como la transvección $(a\theta u)^2$, reducible en el binario a $(ab)^4$, pueden expresarse mediante los símbolos $\A$ y $\nu$.\footnote{$\sum_3$ se refiere a la suma, de tres términos, obtenida a partir del término inicial por permutación cíclica de las variables. Las ecuaciones valen también separadamente para cada término de la suma.}

Obtenemos las fórmulas de reducción fundamentales para lo que sigue por especialización del desarrollo IV o, más brevemente, por cálculo directo (permutación de los símbolos), para $a_\vartheta(a\theta u)^2$, $a_\vartheta^2((a\theta u)^2$, $(a\theta u)^2$ según el siguiente planteamiento:
\[
\begin{aligned}
a_\vartheta(a\theta u)^2
&=(abc)(bcu)(abu)^2-\frac13(abc)(bcu)\{(bcu)a_x-u_x(abc)\}^2,\\
a_\vartheta^2(a\theta u)^2
&=(abc)^2(abu)^2-\frac13(abc)^2\{(bcu)a_x-u_x(abc)\}^2,
\end{aligned}
\]
para $((a\theta u)^2)^{*}$:
\[
\begin{aligned}
(abu)a_y^3b_x^3&=\frac32u_\vartheta(\vartheta yx)\theta_y^2+\frac14(\nu yx)^3,\\
((abu)(acu))b_x^3&=\frac32u_\vartheta(\vartheta\widehat{a}ux)(\theta au)^2+\frac14(\nu\widehat{a}ux)^3.
\end{aligned}
\]
\[
\left.
\begin{aligned}
(a\theta u)^2&=\frac16\nu\cdot f-\frac23u_x\cdot a_\nu
+\frac23u_x^2\cdot a_\nu^2-\frac{i}{18}u_x^4,\\
a_\vartheta&=\frac13u_x\cdot\A,\\
a_\vartheta(a\theta u)&=0,\\
a_\vartheta^2&=\A,\\
(a\theta u)^3&=0,\\
a_\vartheta(a\theta u)^2&=-\frac56a_\nu+\frac76u_x\cdot a_\nu^2-\frac{i}{9}u_x^3,\\
a_\vartheta^2(a\theta u)&=0,\\
(a\theta u)^4&=j,\\
a_\vartheta(a\theta u)^3&=0,\\
a_\vartheta^2(a\theta u)^2&=\frac23a_\nu^2-\frac{i}{9}u_x^2.
\end{aligned}
\right\} \tag{1}
\]

Añadimos la fórmula, procedente igualmente de la reducción del módulo $(abc)$,
\[
a_\nu^3a_xu_\nu=\frac13 i\cdot u_x. \tag{2}
\]

\emph{De las fórmulas (1.) y (2.) y de las fórmulas construidas de manera análoga para la forma fundamental $\nu$, todas las fórmulas de reducción posteriores se deducen por polarización}, con excepción de las relaciones para las formas descomponibles. De las fórmulas (1.) vemos:

La sucesión de formas:
\[
\begin{array}{rcl}
a_\vartheta(a\theta u)^2 &\text{conduce a}& \text{símbolos }\nu\text{ solamente},\\[0.2em]
a_\vartheta &\text{conduce a}& \text{símbolos }\A\text{ y }\nu,\\[0.2em]
(a\theta u)^2 &\text{conduce a}& \text{símbolos }j,\A\text{ y }\nu,\\[0.2em]
(a\theta u)^2\text{ bajo plegamiento I} &\text{conduce a}& \text{símbolos }j\text{ y }\nu,\\[0.2em]
(a\theta u)^2\text{ bajo plegamiento II} &\text{conduce a}& \text{símbolos }\A\text{ y }\nu.
\end{array}
\]

Por tanto podemos reemplazar:
\begin{enumerate}
\item módulo $(\A,\nu)$: los plegamientos con $j$ por los plegamientos correspondientes con $(a\theta u)^2$ o $(a\theta u)^3$;
\item módulo $(\nu)$: los plegamientos con $\A$ por los plegamientos correspondientes con $a_\vartheta$ o $a_\vartheta(a\theta u)$, o también por plegamientos especiales con $(a\theta u)^2$ (que sólo conducen a un plegamiento I simple).
\end{enumerate}
En particular se tiene:
\[
\begin{aligned}
(a\theta u)^2\theta_y^2\theta_x^2&=\frac13(jyx)^2\pmod{\nu},&
(a\theta u)^2u_y\theta_y&=-\frac16(jyx)^2\pmod{\nu},\\
(a\theta u)^2\nu_y^2&=\frac13(\A u\nu)^2\pmod{\nu},&
(a\theta u)(a\theta\nu)u_\vartheta\nu_\vartheta&=-\frac16(\A u\nu)^2\pmod{\nu},\\
(a\theta u)^2u_\nu^2\vartheta_\nu^2&=\frac13(\A u\nu)^2\A_\nu\pmod{\nu}.
\end{aligned}
\]

\begin{center}
\textbf{§ 6. \quad Reducción del módulo $(\A,\nu)$ al módulo $(\nu)$.}
\end{center}

Las fórmulas de reducción del parágrafo anterior nos dan el medio para establecer directamente el sistema relativamente completo módulo $\nu$; veremos que al sistema de las formas tomadas sólo tenemos que adjuntar las formas $\A$ y $(a\A u)=N$ (de manera análoga al caso de la forma cúbica, y válida más generalmente).

Para reducir el módulo $\A$, consideramos todos los plegamientos especiales que entran en cuenta, es decir, los plegamientos del sistema relativamente completo módulo $(\A,\nu)$ con el sistema de $\A$, y partimos de las formas de menor orden en los coeficientes, los plegamientos de $f$ con $\A$.

Para reducir las formas $(a\A u)^2$, $(a\A u)^3$, $(a\A u)^4$, reemplazamos, según § 5, los símbolos $\A$ por formas inferiores de la sucesión de formas, es decir, desarrollamos las expresiones
\[
a_\vartheta b_\vartheta(b\theta u)^2,\qquad
a_\vartheta(a\theta u)b_\vartheta(b\theta u)^2,\qquad
a_\vartheta(a\theta u)b_\vartheta(b\theta u)^3
\]
1) según las polares de la sucesión de formas $a_\vartheta$ (símbolos $\A$ y $\nu$),

2) según las polares de la sucesión de formas $b_\vartheta(b\theta u)^2$ (símbolos $\nu$),

y obtenemos las fórmulas de reducción (cálculo, véase abajo):
\begin{align}
\text{a)}\quad (a\A u)^2
&=\frac12(\vartheta\nu x)^2-\frac25 f\cdot a_\nu^2
-\frac45u_x\cdot\theta_\nu(\vartheta\nu x)^2 \notag\\
&\quad+u_x^2\left\{\frac16 i\cdot f-\frac1{40}S\right\},\notag\\
\text{b)}\quad (a\A u)^3
&=-\frac3{10}\theta_\nu(\vartheta\nu x),\tag{3}\\
\text{c)}\quad (a\A u)^4
&=\frac{21}{10}\theta_\nu^2-\frac3{10}\Pi
-\frac65u_x\cdot\theta_\nu^3+\frac1{10}u_x^2\cdot\theta .\notag
\end{align}

(Los mismos resultados se obtienen también por reducción doble de las expresiones $a_\vartheta^2(b\theta u)^2$, $a_\vartheta^2(b\theta u)^3$, $a_\vartheta^2((a\theta u)^2(b\theta u)^2$ y $a_\vartheta^2((a\theta u)(b\theta u))^3$, después de eliminar $a_\vartheta^2$.)

De las fórmulas (3.) resulta que $(a\A u)^2$ es un reductor relativamente al módulo $\nu$. De ello se deduce:

1) La reducción del sistema de $\A$: la sucesión de formas $(\A\A' u)^2=a_\vartheta^2((a\A u)^2$ tiene como factor el reductor $(a\A u)^2$.

2) La reducción del sistema nacido por plegamiento con el módulo $\A$:

La sucesión de formas $(\theta\A u)^2$ tiene como factor el reductor $(a\A u)^2$.

Para las formas $(\theta\A u)$ y $\A_\vartheta$, se obtiene:
\[
(a\A u)^2(abu)=(\theta\A u)-u_x\cdot\A_\vartheta(\theta\A u),
\]
\[
(a\A u)(a\A b)=-\frac12\A_\vartheta+\frac12u_x\cdot\A_\vartheta^2.
\]

Pero del reductor $(\theta\A u)$ resulta la reducción del sistema nacido por plegamiento con el módulo $\A$.

\bigskip

El sistema relativamente completo módulo $\nu$ se compone, por tanto, de las seis formas:
\[
f,\ \theta,\ K,\ j,\ \A,\ N.
\]

Como ejemplo del desarrollo en serie del § 2, damos en detalle la deducción de la fórmula (3.)a; en los cálculos posteriores se escribirá directamente la fórmula final III. (Las polares de las formas nulas ya se han omitido durante el cálculo; la primera línea da los valores insertados según el desarrollo II, la segunda línea los valores finales encontrados recursivamente empezando por la última ecuación del sistema.) Del desarrollo II se sigue:

\begin{align*}
\text{1)}\quad&m=3,\quad n=4,\quad \lambda=2,\quad \mu=0,\quad \nu=\chi=1,\\
a_\vartheta b_\vartheta(b\theta u)^2
&=[a_\vartheta]_{b^2}b+\frac67\{a_\vartheta b_\vartheta(a\theta u)(b\theta u)-u_xa_\vartheta b_\vartheta(a\theta b)(b\theta u)\}\\
&\quad-\frac17\{a_\vartheta(a\theta u)^2b_\vartheta-2u_xa_\vartheta(a\theta u)(a\theta b)b_\vartheta+u_x^2a_\vartheta(a\theta b)^2b_\vartheta\}\\
&=\frac23(a\Delta u)^2+\frac15[a_\vartheta(a\theta u)^2]b+\frac1{30}f[a_\vartheta^2(a\theta u)^2]-\frac25u_x[a_\vartheta(a\theta u)^2]b^2\\
&\quad-\frac1{15}u_x[a_\vartheta^2(a\theta u)^2]b+\frac15u_x^2[a_\vartheta(a\theta u)^2]b^3+\frac1{30}u_x^2[a_\vartheta^2(a\theta u)^2]b^2;\\[.6em]
\text{2)}\quad&m=2,\quad n=3,\quad \lambda=1,\quad \mu=1,\quad \nu=\chi=1,\\
a_\vartheta(a\theta u)b_\vartheta(b\theta u)
&=\frac25\{a_\vartheta(a\theta u)^2b_\vartheta-u_xa_\vartheta(a\theta u)(a\theta b)b_\vartheta\}+\frac12a_\vartheta^2(b\theta u)^2\\
&\quad-\frac15\{a_\vartheta^2(b\theta u)(a\theta u)-u_xa_\vartheta^2(a\theta b)(b\theta u)\}\\
&=\frac12(a\Delta u)^2+\frac25[a_\vartheta(a\theta u)^2]b+\frac1{15}[a_\vartheta^2(a\theta u)^2]f\\
&\quad-\frac25u_x[a_\vartheta(a\theta u)^2]b^2-\frac1{10}u_x[a_\vartheta^2(a\theta u)^2]b+\frac1{30}u_x^2[a_\vartheta^2(a\theta u)^2]b^2;\\[.6em]
\text{3)}\quad&m=2,\quad n=3,\quad \lambda=1,\quad \mu=1,\quad \nu=1,\quad \chi=2,\\
a_\vartheta(a\theta b)b_\vartheta(b\theta u)
&=\frac25\{a_\vartheta(a\theta b)(a\theta u)b_\vartheta-u_xa_\vartheta(a\theta b)^2b_\vartheta\}\\
&=\frac25[a_\vartheta(a\theta u)^2]b^2+\frac1{30}[a_\vartheta^2(a\theta u)^2]b-\frac25u_x[a_\vartheta(a\theta u)^2]b^3-\frac1{30}u_x[a_\vartheta^2(a\theta u)^2]b^2;\\[.6em]
\text{4)}\quad&m=1,\quad n=2,\quad \lambda=0,\quad \mu=2,\quad \nu=\chi=1,\\
a_\vartheta(a\theta u)^2b_\vartheta
&=[a_\vartheta(a\theta u)^2]b+\frac23a_\vartheta^2(a\theta u)(b\theta u)\\
&=[a_\vartheta(a\theta u)^2]b+\frac16[a_\vartheta^2(a\theta u)^2]f-\frac16u_x[a_\vartheta^2(a\theta u)^2]b;\\[.6em]
\text{5)}\quad&m=1,\quad n=2,\quad \lambda=0,\quad \mu=2,\quad \nu=1,\quad \chi=2,\\
a_\vartheta(a\theta u)(a\theta b)b_\vartheta
&=[a_\vartheta(a\theta u)^2]b^2+\frac13a_\vartheta^2(a\theta b)(b\theta u)\\
&=[a_\vartheta(a\theta u)^2]b^2+\frac1{12}[a_\vartheta^2(a\theta u)^2]b-\frac1{12}u_x[a_\vartheta^2(a\theta u)^2]b^2;\\[.6em]
\text{6)}\quad&m=1,\quad n=2,\quad \lambda=0,\quad \mu=2,\quad \nu=1,\quad \chi=3,\\
a_\vartheta(a\theta b)^2b_\vartheta&=[a_\vartheta(a\theta u)^2]b^3;\\[.6em]
\text{7)}\quad&m=2,\quad n=4,\quad \lambda=2,\quad \mu=\nu=\chi=0,\quad (\text{según el desarrollo I}),\\
a_\vartheta^2(b\theta u)^2&=[a_\vartheta^2]_{ub^2}+\frac1{10}\{[a_\vartheta^2(a\theta u)^2]f-2u_x[a_\vartheta^2(a\theta u)^2]b+u_x^2[a_\vartheta^2(a\theta u)^2]b^2\};\\[.6em]
\text{8)}\quad&m=1,\quad n=3,\quad \lambda=1,\quad \mu=1,\quad \nu=\chi=0,\quad (\text{desarrollo I}),\\
a_\vartheta^2(a\theta u)(b\theta u)&=\frac14[a_\vartheta^2(a\theta u)^2]f-\frac14u_x[a_\vartheta^2(a\theta u)^2]b;\\[.6em]
\text{9)}\quad&m=1,\quad n=3,\quad \lambda=1,\quad \mu=\chi=1,\quad \nu=0,\quad (\text{desarrollo I}),\\
a_\vartheta^2(a\theta b)(b\theta u)&=\frac14[a_\vartheta^2(a\theta u)^2]b-\frac14u_x[a_\vartheta^2(a\theta u)^2]b^2.
\end{align*}

De 1) y 4) resulta:
\[
\begin{aligned}
(a\Delta u)^2
&=\frac65[a_\vartheta(a\theta u)^2]b+\frac15f[a_\vartheta^2(a\theta u)^2]+\frac35u_x[a_\vartheta(a\theta u)^2]b^2-\frac3{20}u_x[a_\vartheta^2(a\theta u)^2]b\\
&\quad-\frac3{10}u_x^2[a_\vartheta(a\theta u)^2]b^3-\frac1{20}u_x^2[a_\vartheta^2(a\theta u)^2]b^2,\\
(a\Delta u)^2
&=-a_\nu b_\nu+\frac35fa_\nu^2+\frac45u_xa_\nu^2b_\nu-\frac3{20}u_x^2a_\nu^2b_\nu^2-\frac1{12}u_x^2if.
\end{aligned}
\]

(Para la conversión en símbolos $\theta$, véanse § 8 y § 9.) La fórmula (3.)a podría calcularse aún más brevemente por transferencia directa del desarrollo I, introduciendo la variable auxiliar $w=x\widehat{u}b$, mientras que, ya desde las fórmulas (3.)b y (3.)c, el camino que hemos seguido se vuelve más simple; para (3.)b se sabe de antemano, por § 9, que todos los términos que contienen el factor $u_x$ deben omitirse. Para calcular (3.)b y (3.)c se obtiene:
\[
\begin{aligned}
(3.)\mathrm{b)}\quad 1)&\quad a_\vartheta(a\theta u)b_\vartheta(b\theta u)^2=\frac12(a\Delta u)^3+\frac45[a_\vartheta(a\theta u)^2]_{ub}b+\frac7{360}[a_\vartheta^2(a\theta u)^2]_{ub^2},\\
2)&\quad a_\vartheta(a\theta u)b_\vartheta(b\theta u)^2=[a_\vartheta(a\theta u)^2]_{ub}b+\frac{13}{36}[a_\vartheta^2(a\theta u)^2]_{ub^2};\\[.5em]
(3.)\mathrm{c)}\quad 1)&\quad a_\vartheta(a\theta u)b_\vartheta(b\theta u)^3=\frac12(a\Delta u)^4+\frac65[a_\vartheta(a\theta u)^2]_{ub^2}b+\frac{53}{120}[a_\vartheta^2(a\theta u)^2]_{ub^3}\\
&\qquad-\frac65u_x[a_\vartheta(a\theta u)^2]_{ub^3}b^2-\frac35u_x[a_\vartheta^2(a\theta u)^2]_{ub^2}b+\frac{19}{120}u_x^2[a_\vartheta^2(a\theta u)^2]_{ub^3}b^2,\\
2)&\quad a_\vartheta(a\theta u)b_\vartheta(b\theta u)^3=\frac34[a_\vartheta^2(a\theta u)^2]_{ub^2}.
\end{aligned}
\]

Observación final: las transvecciones de $f$ sobre $j$ que son reducibles según § 5 se calculan de manera análoga. Se tiene
\[
\tag{4}
\begin{aligned}
g_\nu&=-\theta_\nu+u_x\left(\frac92\theta_\nu^2-\frac12H\right)-3u_x^2\theta_\nu^3+\frac12u_x^3\theta,\\
g_\nu^2&=\frac{21}{10}\theta_\nu^2-\frac3{10}H-2u_x\theta_\nu^3+\frac12u_x^2\theta,\\
g_\nu^3&=-\frac7{10}\theta_\nu^3+\frac12u_x^2\theta,\qquad g_\nu^4=\frac35\theta.
\end{aligned}
\]

\[
\tag{4}
g_\nu^3=-\frac7{10}\theta_\nu^3+\frac12u_x^2\theta,
\qquad
g_\nu^4=\frac35\theta.
\]

De (3.) y (4.) resulta, por transferencia desde el dominio binario,
\[
(\theta\theta'u)^2=\frac13\,j\cdot f\pmod{\nu}.\footnotemark
\]
Las demás formas de la sucesión de formas $(\theta\theta'u)^2$ y la forma $\theta'_\nu$ son reducibles a los símbolos $\nu$. Por tanto podemos reemplazar:
\[
\bmod\ \nu:\quad \text{los plegamientos con } f\cdot j
\text{ por los plegamientos correspondientes con }(\theta\theta'u)^2.
\]
Además se tiene:
\[
(\vartheta\vartheta'x)^2=\frac43 f\cdot\A,
\qquad
\theta_{g\nu}(\vartheta\vartheta'x)=-N.
\]
\footnotetext{Gordan--Kerschensteiner, loc. cit. p. 181.}



\begin{center}
{\Large\bfseries Capítulo III.\quad El sistema relativamente completo módulo $(s,\varrho,t)$.}
\end{center}

\begin{center}
\textbf{\S\ 7.\quad Visión general de la formación del sistema.}
\end{center}

Para pasar del sistema relativamente completo módulo $\nu$ al sistema relativamente completo con el módulo inmediatamente superior, según la definición de \S\ 3 hay que formar el sistema de $\nu$ con respecto a este módulo superior y plegarlo con las formas del sistema relativamente completo módulo $\nu$. Pero $\nu=u_\nu^4$, como contravariante de cuarto grado de la forma $f$, se halla en relación dualística con ella. Si, por tanto, consideramos $\nu$ como forma fundamental, obtenemos un sistema relativamente completo de seis formas módulo $\nu(\nu)=(\nu\nu_1x)^4=s_x^4$.

Así, para formar el sistema relativamente completo módulo $s$ (sistema III), debemos transveccionar
\[
\begin{array}{ll}
\text{Sistema I:} & f,\ \theta,\ K,\ j,\ \A,\ N\quad \text{con}\vphantom{\dfrac11}\\[0.3em]
\text{Sistema II:} & \nu,\ H=(\nu\nu_1x)^2u_\nu^2u_{\nu_1}^2,\ L=(\nu\eta x)H_x^2u_\nu^3u_\eta^3,\\
& g=(\nu\eta x)^4H_x^2,\ \sigma=H_\eta^2u_\nu^2u_\eta^4,\ (\nu\sigma x).
\end{array}
\]
Se verá (fórmula (13)) que la forma $g$ es reducible, de modo que el sistema II consta sólo de cinco formas.

En el curso del cálculo se adjuntan como módulos las formas cuadráticas
\[
\varrho=u_\varrho^2=\theta_\nu^4u_\vartheta^2,
\qquad
 t=t_x^2=a_\eta^4H_x^2,
\]
de manera que se obtiene un sistema relativamente completo módulo $(s,\varrho,t)$; y el módulo $(\varrho,t)$ debe considerarse como módulo superior al módulo $s$.

La ordenación de las formas individuales debe seguir el orden en los coeficientes. Normalizamos el plegamiento
\[
\theta_\nu(K_\nu,\theta_\nu,\ldots)
\]
como superior al plegamiento
\[
H_y(H_y,L_y,\ldots).
\]
Por esto, según \S\ 1 y \S\ 3b, queda determinada de modo unívoco la sucesión de las formas individuales del mismo orden.

La formación de las formas de igual orden se llevará a cabo en forma tabular:

I. Indicación de qué formas de los sistemas I y II han de plegarse entre sí para alcanzar el orden determinado en los coeficientes.

II. Disposición de las formas irreducibles según el esquema de la sucesión de formas.

III. Reducción de las sucesiones de formas o formas reducibles o descomponibles, es decir, de los lugares en los que se interrumpe la sucesión total de formas; con ello se demuestra que todas las formaciones irreducibles quedan agotadas por las indicadas.

IV. Consecuencias: 1) indicación de los reductores recién surgidos; 2) indicación de aquellas formas que no entran, en productos con formas del mismo sistema, en plegamiento con formas del otro sistema.

Se verá que sólo aparecen:

1) plegamientos de una forma del sistema I con una forma del sistema II;

2) plegamientos de potencias de $\theta$, respectivamente de $H$, o de dos formas de uno de los sistemas con una forma del segundo sistema.

Obtendremos el siguiente esquema de plegamiento:
\[
\begin{array}{r@{\ }c@{\ }l@{\qquad}c@{\qquad}r@{\ }c@{\ }l}
\multicolumn{3}{c}{\text{Sistema I}} & \text{plegado con} & \multicolumn{3}{c}{\text{Sistema II}}\\[0.25em]
1. & \text{orden} & f & & 2. & \text{orden} & \nu\\
2. & ,, & \theta & & 4. & ,, & H\\
3. & ,, & K,j,\A & & 6. & ,, & L,\sigma\\
4. & ,, & N,\theta^2,f\cdot j & & 8. & ,, & (\nu\sigma x),H^2\\
5. & ,, & \theta\cdot K & & 10. & ,, & H\cdot L\\
6. & ,, & \theta^3,\A\cdot j & & 12. & ,, & H^3.
\end{array}
\]

Además de la ``doble reducción'', sirven como control del cálculo las dos formas especiales:

\[
\begin{array}{ll}
\text{I.}&\begin{array}{rlrlrl}
f&=\displaystyle\sum_3x_1^3x_2,& u_\nu^2&=\frac{i}{6}u_x^2,& s&=\frac{i}{3}f,\\
a_\eta^2&=-\frac19 i\theta,& H_\nu^2&=\frac{i}{6}\theta,& \sigma&=-\frac{i}{6}j,\\
g&=-\frac23i\Delta,& \varrho&=0,& t&=0,
\end{array}\\[1.2em]
\text{II.}&\begin{array}{rlrlrl}
f&=\displaystyle\sum_3x_1^4,& s&=\frac43if,& a_\eta^2&=\frac29i\theta,\\
\Delta_\nu^2&=\frac1{15}i\theta,& \sigma&=\frac43ij,& g&=\frac43i\Delta,\\
\varrho&=0,& t&=0.
\end{array}
\end{array}
\]

Observación final: a las fórmulas (1), (2), (3), (4) les corresponden, para la forma fundamental $\nu$, las siguientes:
\[
\tag{5}
\begin{array}{c|c|c}
(\nu\eta x)^2=A & H_\nu(\nu\eta x)=0 & H_\nu^2=\sigma\\[0.4em]
(\nu\eta x)^3=0 & H_\nu(\nu\eta x)^2=B & H_\nu^2(\nu\eta x)=0\\[0.4em]
(\nu\eta x)^4=g & H_\nu(\nu\eta x)^3=0 & H_\nu^2(\nu\eta x)^2=C,
\end{array}
\]
\[
A=\frac16s\nu-\frac23u_xs_\nu+\frac23u_x^2s_\nu^2-\frac{J}{18}u_x^4,
\quad B=-\frac56s_\nu+\frac76u_xs_\nu^2-\frac{J}{9}u_x^3,
\quad C=\frac23s_\nu^2-\frac{J}{9}u_x^2.
\]
\[
\tag{6}
s_\nu^3u_xs_\nu=\frac13u_xs_\nu^4=\frac13Ju_x.
\]

\[
\tag{7}
\begin{aligned}
\text{a)}\quad (\nu\sigma x)^2
&=\frac12(Hsu)^2-\frac25\nu\cdot s_\nu^2-
\frac45u_x\,s_\eta(Hsu)^2
+u_x^2\left\{\frac16J\cdot\nu-\frac1{40}(ss'u)^4\right\},\\
\text{b)}\quad (\nu\sigma x)^3&=-\frac3{10}s_\eta(Hsu),\\
\text{c)}\quad (\nu\sigma x)^4&=\frac{21}{10}s_\eta^2-\frac3{10}Z-
\frac65u_xs_\eta^3+\frac1{10}u_x^2s_\eta^4.
\end{aligned}
\]
\[
\tag{8}
\begin{aligned}
\text{a)}\quad g_\nu&=-s_\eta+u_x\left(\frac92s_\eta^2-\frac12Z\right)-3u_x^2s_\eta^3+\frac12u_x^3s_\eta^4,\\
\text{b)}\quad g_\nu^2&=\frac{21}{10}s_\eta^2-\frac3{10}Z-2u_xs_\eta^3+\frac12u_x^2s_\eta^4,\\
\text{c)}\quad g_\nu^3&=-\frac7{10}s_\eta^3+\frac12u_xs_\eta^4,\\
\text{d)}\quad g_\nu^4&=\frac35s_\eta^4.
\end{aligned}
\]

\begin{center}
\textbf{\S\ 8.\quad Formas de tercer orden (sistema III).}
\end{center}

\noindent\textbf{Plegamiento de $f$ con $\nu$.}

\noindent\emph{Formas irreducibles:}
\[
\begin{array}{cccc}
a_\nu&\cdot&\cdot&\cdot\\
\cdot&a_\nu^2&\cdot&\cdot\\
\cdot&\cdot&\cdot&\cdot\\
\cdot&\cdot&\cdot&a_\nu^4=i.
\end{array}
\]

\noindent\emph{Reducción:}
\[
a_\nu^3=\frac13 i u_x\qquad\text{(fórmula (2)).}
\]
\noindent\emph{Consecuencias:} 1) $a_\nu^3$ es reductor. 2) $f$ sólo entra en plegamiento con $\nu$ en productos con formas contragredientes $(j)$. Lo mismo vale para $\nu$ con respecto a $f$.

\begin{center}
\textbf{\S\ 9.\quad Formas de cuarto orden (sistema III).}
\end{center}

\noindent\textbf{Plegamiento de $\theta$ con $\nu$.}

\noindent\emph{Formas irreducibles:}
\[
\begin{array}{cccc}
(\theta\nu x)&\theta_\nu&\cdot&\cdot\\
(\theta\nu x)^2&\theta_\nu(\theta\nu x)&\theta_\nu^2&\cdot\\
\cdot&\theta_\nu(\theta\nu x)^2&\cdot&\theta_\nu^3.
\end{array}
\]

\noindent\emph{Reducciones:} Del reductor $a_\nu^3$, por doble reducción de las expresiones $a_\nu^3(abu)$, $a_\nu^3b_\nu$, $a_\nu^3b_\nu(abu)$ según el planteamiento:

\[
\begin{aligned}
\theta_\nu^2(\vartheta\nu x)
&=a_\nu^2(\widehat{a}b\nu x)(abu)-\frac13(\nu\nu_1x)^3=a_\nu b_\nu(\widehat{a}b\nu x)(abu)+\frac16(\nu\nu_1x)^3\\
&=a_\nu^3(abu)-a_\nu^2b_\nu(abu)=2a_\nu^2b_\nu(abu)=\frac23a_\nu^3(abu),\\[.4em]
\theta_\nu^2(\vartheta\nu x)^2
&=a_\nu^2(\widehat{a}b\nu x)^2-\frac13(\nu\nu_1x)^4=a_\nu b_\nu(\widehat{a}b\nu x)^2+\frac16(\nu\nu_1x)^4\\
&=if-2a_\nu^3b_\nu+a_\nu^2b_\nu^2-\frac13s=2a_\nu^3b_\nu-2a_\nu^2b_\nu^2+\frac16s=\frac23if-\frac23a_\nu^3b_\nu-\frac16s,\\[.4em]
\theta_\nu^3(\vartheta\nu x)&=a_\nu^2b_\nu(\widehat{a}b\nu x)(abu)=a_\nu^3b_\nu(abu).
\end{aligned}
\]

Se obtienen las fórmulas:
\[
\tag{9}
\begin{array}{c|c|c}
\theta_\nu^2(\theta\nu x)=0
& \theta_\nu^3(\theta\nu x)=0
& \theta_\nu^4u_x^2=u_\varrho^2=\varrho\quad(\text{módulo adjunto, \S\ 7}),\\[0.6em]
\theta_\nu^2(\theta\nu x)^2=\dfrac49 i f-\dfrac16s
& &
\end{array}
\]

\noindent\emph{Consecuencias:} 1) $\theta_\nu^2(\theta\nu x)^2$ es reductor. 2) $\nu$ sólo entra en plegamiento con $\theta$ en productos con formas contragredientes $(g)$.

\begin{center}
\textbf{\S\ 10.\quad Formas de quinto orden (sistema III).}
\end{center}

\[
\text{Plegamiento de }\left\{\begin{array}{l}
K,j,\A \text{ con }\nu,\\
f\text{ con }H.
\end{array}\right.
\]

\noindent\emph{Formas irreducibles:}
\[
\begin{array}{c@{\qquad}c@{\qquad}c@{\qquad}c}
\text{A)} & \begin{array}{cccc}
\cdot&K_\nu&\cdot&\cdot\\
(k\nu x)^2&\cdot&K_\nu^2&\cdot\\
(k\nu x)^3&K_\nu(k\nu x)^2&\cdot&\cdot\\
\cdot&K_\nu(k\nu x)^3&\cdot&\cdot
\end{array}
&
\text{B)}\ \begin{array}{c}(j\nu x)\\(j\nu x)^2\\(j\nu x)^3\\ \cdot\end{array}
&
\text{C)}\ \begin{array}{c}\A_\nu\\\cdot\\\cdot\\\cdot\end{array}
\end{array}
\]
\[
\text{D)}\qquad
\begin{array}{cccc}
(aHu)&(aHu)^2&\cdot&\cdot\\
a_\eta&a_\eta(aHu)&a_\eta(aHu)^2&\cdot\\
\cdot&a_\eta^2&\cdot&\cdot
\end{array}
\]

\noindent\emph{Reducciones:}

\noindent A) Sucesión de formas $K_\nu^3$: reductor $a_\nu^3$.

Forma $K_\nu^2(k\nu x)^2$: reductor $\theta_\nu^2(\vartheta\nu x)^2$.
\[
(k\nu x),\quad K_\nu(k\nu x),\quad K_\nu^2(k\nu x)
\]
son formas descomponibles según \S\ 3.

\noindent B) Forma
\[
\begin{aligned}
(j\nu x)^4&=(\hata\theta\nu x)^4
=a_\nu^4\theta+\theta_\nu^4 f-4a_\nu^3\theta_\nu
-4a_\nu\{a_\nu\theta_x-(\hata\theta\nu x)\}^3\\
&\quad+6a_\nu^2\{a_\nu\theta_x-(\hata\theta\nu x)\}^2
=\R f+\frac53 i\theta-6a_\nu^2(\hata\theta\nu x)^2
+4a_\nu(\hata\theta\nu x)^3,
\end{aligned}
\]
y por tanto
\[
(j\nu x)^4=\R f+\frac53 i\theta\modu{(\A,\nu)},\qquad (\text{cf. final de \S\ 5}).
\]

\noindent C) Forma $\A_\nu^4$: reductor $\theta_\nu^4$.

Formas $\A_\nu^2$ y $\A_\nu^3$: reductor $a_\nu^3$ o $\theta_\nu^4$ por sustitución de la forma $\A$ por formas inferiores de la sucesión de formas (final de \S\ 5), es decir, por doble reducción de las expresiones
\[
a_\vartheta a_\nu^3,
\qquad a_\vartheta a_\nu^3\theta_\nu,
\qquad a_\vartheta\theta_\nu^4.
\]
El doble cálculo de $\A_\nu^3$ da ocasión a la reducción de la forma superior $a_\eta^3$.

\noindent\emph{Fórmulas de reducción:}
\[
\tag{10}
\begin{aligned}
\text{a)}\quad \A_\nu^2
&=\frac35a_\eta^2-\frac3{10}(asu)^2+\frac13 i\theta
+\frac15u_x^2(2a_\R^2-t),\\
\text{b)}\quad \A_\nu^3
&=-\frac3{10}a_\R+\frac35u_x\left(\frac{13}{10}a_\R^2-\frac25t\right),\\
\text{c)}\quad \A_\nu^4&=a_\R^2-\frac25t.
\end{aligned}
\]

\noindent\emph{Derivación de las fórmulas:}
\[
\begin{aligned}
\text{a)}\quad
 a_\vartheta a_\nu^3
&=\frac13 i\theta
=[a_\vartheta]_{\nu^3}+\frac67[a_\vartheta^2]_{\nu^2}
+\frac65[a_\vartheta(a\theta u)^2]_{\nu^2}\,
u x^2
+\frac{11}{30}[a_\vartheta^2(a\theta u)^2]_\nu\,\nu x^2\\
&\quad-\frac67u_x[a_\vartheta^2]_{\nu^3}
-\frac16u_x[a_\vartheta^2(a\theta u)^2]_{\nu^2}\,\nu x^2,\\
\frac13 i\theta
&=\A_\nu^2-\frac23u_x\A_\nu^3-a_\nu a_{\nu_1}(\nu\nu_1x)^2
+\frac25a_\nu^2(\nu\nu_1x)^2
+\frac15u_x^2a_\nu^2a_{\nu_1}(\nu\nu_1x)^2;
\end{aligned}
\]
\[
\begin{aligned}
\text{b)}\quad
 a_\vartheta a_\nu^3\theta_\nu
&=0=[a_\vartheta]_{\nu^4}+\frac9{14}[a_\vartheta^2]_{\nu^3}
+\frac35[a_\vartheta(a\theta u)^2]_{\nu^2}\,\nu x^2
+\frac{17}{120}[a_\vartheta^2(a\theta u)^2]_{\nu}\,\nu x^2\\
&\quad-\frac9{14}u_x[a_\vartheta^2]_{\nu^4}
-\frac1{24}u_x[a_\vartheta^2(a\theta u)^2]_{\nu^2}\,\nu x^2,\\
0&=\frac56\A_\nu^3-\frac12u_x\A_\nu^4-\frac14a_\eta^3+\frac1{20}u_xa_\eta^4;
\end{aligned}
\]
\[
\begin{aligned}
\text{b')}\quad
 a_\vartheta\theta_\nu^4
&=a_\R=a_\vartheta\theta_\nu\{a_\nu-(a\theta\nu x)\}^3\\
&=-\frac32[a_\vartheta]_{\nu^3}+\frac35[a_\vartheta(a\theta u)^2]_{\nu^2}\,\nu x^2
-\frac{37}{120}[a_\vartheta^2(a\theta u)^2]_{\nu}\,\nu x^2\\
&\quad+\frac32u_x[a_\vartheta^2]_{\nu^4}
+\frac{49}{120}u_x[a_\vartheta^2(a\theta u)^2]_{\nu^2}\,\nu x^2,\\
a_\R&=-\frac32\A_\nu^3+\frac32u_x\A_\nu^4-\frac{11}{20}a_\eta^3+\frac7{20}u_xa_\eta^4;
\end{aligned}
\]
\[
\begin{aligned}
\text{c)}\quad
 a_\vartheta^2\theta_\nu^4
&=a_\R^2=[a_\vartheta^2]_{\nu^4}+\frac35[a_\vartheta^2(a\theta u)^2]_{\nu^2}\,\nu x^2,\\
a_\R^2&=\A_\nu^4+\frac25a_\eta^4.
\end{aligned}
\]

\noindent D) Reducción de $a_\eta^3$ por doble reducción de $\A_\nu^3$ (C).

Las demás reducciones se obtienen por un cálculo dualísticamente opuesto al de \S\ 9.

\noindent\emph{Fórmulas de reducción:}
\[
\tag{11}
\begin{array}{rl|rl}
a_\eta^2(aHu)&=-\dfrac16(asu)^3,&
 a_\eta^2(aHu)^2&=\dfrac49 i\nu-\dfrac16(asu)^4,\\[0.7em]
a_\eta^3&=-a_\R+\dfrac35u_xa_\R^2+\dfrac15u_xt,
& a_\eta^3(aHu)&=0,\\[0.7em]
&&a_\eta^4&=t\quad(\text{adjuntado como módulo}).
\end{array}
\]

\noindent\emph{Consecuencias:}

1) $(j\nu x)^4$, $\A_\nu^2$, $a_\eta^2(aHu)$ son reductores.

2) $\nu$ entra sólo en productos con formas contragredientes en plegamiento con $K,j,\A$; lo mismo vale para $f$ con respecto a $H$ y para $j$ con respecto a $\nu$, pues $\theta_\nu(\theta\cdot j,\nu,x^3)$ se vuelve reducible según \S\ 11 B.

\begin{center}
\textbf{\S\ 11.\quad Formas de sexto orden (sistema III).}
\end{center}
\[
\text{Plegamiento de }\left\{\begin{array}{l}
\theta^2,
 f\cdot j,
 N \text{ con }\nu,\\
\theta\text{ con }H.
\end{array}\right.
\]

\noindent\emph{Formas irreducibles:}
\[
\begin{array}{c@{\quad}c@{\quad}c}
\text{A)}\ (\vartheta^2\nu x)^3,\quad \theta_\nu(\vartheta^2\nu x)^3
&
\text{B)}\ a_\nu(j\nu x),\quad a_\nu^2(j\nu x)
&
\text{C)}\ N_\nu
\end{array}
\]
\[
\text{D)}\quad
\begin{array}{c|c|c|c}
 & (\theta Hu)&(\theta Hu)^2&\cdot\\
(\vartheta\eta x)&\theta_\eta&H_\vartheta(\theta Hu),\ \theta_\eta(\theta Hu)&\theta_\eta(\theta Hu)^2\\
(\vartheta\eta x)^2&H_\vartheta(\vartheta\eta x),\ \theta_\eta(\vartheta\eta x)&\theta_\eta^2&\theta_\eta^2(\theta Hu)\\
\cdot&\theta_\eta(\vartheta\eta x)^2&\cdot&\cdot
\end{array}
\]

\noindent\emph{Reducciones:}

\noindent A) $(\vartheta^2\nu x)^4$ se descompone según la fórmula (3)a:
\[
(a\A\vartheta x)^4=\frac12(\vartheta^2\nu x)^4-\frac35 f\cdot a_\nu^2(\nu\vartheta x)^2
=\frac13\A^2+f\A_\vartheta^2.
\]

\noindent B) $a_\nu(j\nu x)^2$ se descompone según la fórmula (9)a, sustituyendo, de acuerdo con el final de \S\ 6, $f\cdot j$ por $(\theta\theta'u)^2$ y teniendo en cuenta la reducibilidad de $\theta_\vartheta'$:
\[
\frac13a_\nu(j\nu x)^2=(\theta\theta'\nu x)^2\theta_\nu
=\theta_\nu^3\theta-\theta_\nu^2\theta_\nu'
=\theta_\nu^3\theta+\theta_\nu^2(\widehat{j}\nu\theta'u)
=\theta_\nu^3\theta-\frac23\theta_\nu^3\theta.
\]
Formas $a_\nu(j\nu x)^3$ y $a_\nu^2(j\nu x)^2$: reductores $a_j$ y $(j\nu x)^4$:
\[
a_\nu(j\nu x)^3=a_j(j\nu x)^3-(j\nu x)^3(j\nu\widehat{a}u),
\]
\[
a_\nu^2(j\nu x)^2=a_j^2(j\nu x)^2-2a_j(j\nu x)^2(j\nu\widehat{a}u)+(j\nu x)^2(j\nu\widehat{a}u)^2.
\]

\noindent C) Sucesión de formas $N_\nu^2$; reductor $\A_\nu^2$. $(n\nu x)$ y $N_\nu(n\nu x)$ se descomponen como transvecciones sobre determinantes funcionales según \S\ 3.

\noindent D) \emph{Vista general de las reducciones:}

a) Sucesión de formas $\theta_\eta^2H_\vartheta$: reductor $a_\eta^2(aHu)$. (Conduce a formas con símbolos superiores.)

b) Sucesión de formas $\theta_\eta H_\vartheta$: reductor $a_\eta^2(aHu)$ por doble reducción. (Conduce a formas con símbolos superiores y a formas superiores de la sucesión total de formas, es decir, a la sucesión de formas $\theta_\eta^2$.) En lugar de $\theta_\eta H_\vartheta$ consideramos la sucesión de formas
\[
\theta_\eta(\vartheta\eta x)(\theta Hu)=\theta_\eta H_\vartheta+\theta_\eta^2,
\]
y en ella reemplazamos los símbolos $\theta$ por la forma inferior $(abu)$, llegando así a la doble reducción de la sucesión de formas
\[
a_\eta^2(aHu)(abu)=\theta_\eta H_\vartheta+\frac32\theta_\eta^2\modu{s}
\]
hasta la forma final
\[
a_\eta^3b_\eta(bHu)(abu)=\theta_\eta^3H_\vartheta+\frac12\theta_\eta^4.
\]

c) Sucesión de formas $\theta_\eta^3$: por doble reducción de aquellas formas que pertenecen simultáneamente a las sucesiones $\theta_\eta H_\vartheta$ y $\theta_\eta^2H_\vartheta$, es decir, de las formas de la sucesión $\theta_\eta^2H_\vartheta$, hacia formas con símbolos superiores por una parte, y hacia formas con símbolos superiores y hacia la sucesión de formas superior $\theta_\eta^3$ por otra.

d) Formas $\theta_\eta^2(\vartheta\eta x)$, $\theta_\eta^2(\vartheta\eta x)^2$, $\theta_\eta^3(\vartheta\eta x)$: por doble reducción mediante el reductor $a$, de manera análoga al cálculo de \S\ 9.

e) Formas $\theta_\eta^2(\theta Hu)^2$ y $\theta_\eta^2(\vartheta\eta x)^2$: por doble reducción de las expresiones $\A_\nu^2(a\A u)^4$ y $(a\A\nu x)^4$ mediante los reductores $\A_\nu^2$ y $(a\A u)^4$.

f) Formas $H_\vartheta$ y $H_\vartheta^2$: formas descomponibles. Para $H_\vartheta^2$ esto resulta de la doble reducción de la expresión $\A_\nu^2(a\A u)^2$, o también del cálculo directo de los productos $(a_\nu^2)^2$, $a_\nu^2b_{\nu_1}$ por formas del sistema. (El cálculo análogo de $(a_\nu)^2$ da una relación entre formas descomponibles.)

g) Reducción de formas superiores porque formas de la sucesión $\theta\cdot H$ admiten dos posibilidades de reducción (pertenecen a dos sucesiones de formas reducibles), a saber:
\begin{quote}\small
$g$: por doble reducción de $\theta_\eta^2(\vartheta\eta x)^2$; admite las reducciones d) y e).\par
$s_\vartheta^2(\theta su)$: por doble reducción de $\theta_\eta^3(\vartheta\eta x)$; admite las reducciones c) y d).\par
$s_\vartheta^2(\theta su)^2$: por doble reducción de $\theta_\eta^2H_\vartheta^2$; admite la reducción a) y tiene como factor el reductor $\theta_\nu^2(\vartheta\nu x)^2$.\par
$s_\nu^2$: por doble reducción de $\theta_\eta^3H_\vartheta$; admite las reducciones a) y c).
\end{quote}
La demostración de la independencia de las formas restantes resulta por doble reducción de la sucesión de formas $\A_\nu^2(a\A u)^2=\theta_\eta^2$.

\noindent\emph{Ejecución de las reducciones:}

La existencia de las reducciones a), b), c), d) es clara a priori, pues según la ley de formación indicada las relaciones que surgen en la doble reducción son independientes entre sí. Para las reducciones e), f), g) debe verificarse por cálculo que no surjan identidades.

Damos, simétricamente respecto al término diagonal en la sucesión de formas $\theta\cdot H$ y avanzando hasta la forma $\theta_\eta$, en parte con indicación del cálculo, también las fórmulas obtenidas mediante las reducciones a), b), c), d), pues se emplean en parte en las reducciones e), f), g), y en parte más adelante.

\noindent 1) $H_\vartheta$ y $H_\vartheta^2$ según f). Con el planteamiento\footnote{El signo $\sim$ en lugar del signo de igualdad significa que se omiten productos de $u_x$ con formas superiores a las que surgen por los plegamientos III y IV.}
\[
\begin{aligned}
a_\nu b_{\nu_1}^2
&=\nu a_\nu b_\nu^2-(aHu)(bHu)b_\eta
\sim \nu a_\nu b_\nu^2-f a_\eta(aHu)^2-(abu)(aHu)b_\eta+(abu)^2b_\eta,\\
a_\nu^2b_{\nu_1}^2
&=b_\nu^2\{a_\nu u_\nu-(a\nu\nu_1u)\}^2
\sim \nu\{a_\nu^2b_\nu^2-2a_\nu b_\nu(aHu)(bHu)
-\frac16(asu)^2(bsu)^2\}\\
&\sim \nu a_\nu^2b_\nu^2-2a^2(aHu)(abu)-\frac16(asu)^2(bsu)^2+(\vartheta\eta x)^2(\theta Hu)^2
\end{aligned}
\]
surgen, al sustituir el valor de $\theta_\eta H_\vartheta$, las fórmulas:
\[
\tag{12}
\begin{aligned}
\text{a)}\quad H_\vartheta&\sim 2a_\nu a_\R^2+2\nu b_\nu(\vartheta\eta x)^2+2f a_\eta(aHu)^2-2\theta_\eta,\\
\text{b)}\quad H_\vartheta^2&\sim (a^2)^2+2\theta_\eta^2+\frac23(\theta su)^2
+\frac1{12}s\nu-\frac19if\nu-\frac16f(asu)^4.
\end{aligned}
\]

\noindent 2) $\theta_\eta H_\vartheta$ según b):
\[
\begin{aligned}
a_\eta^2(aHu)(abu)&\sim [a_\eta^2(aHu)]_{bu}+\frac12 f[a_\eta^2(aHu)^2]
=a_\eta b_\eta(aHu)(abu)\\
&\quad+a_\eta(a\widehat{b}\eta x)(abu)(aHu)
\sim \theta_\eta(\vartheta\eta x)(\theta Hu)+\frac12\theta_\eta^2+\frac14(\nu\eta x)^2.
\end{aligned}
\]
De las fórmulas (5) y (11) se sigue:
\[
\theta_\eta H_\vartheta\sim -\frac32\theta_\eta^2-\frac14(\theta su)^2-\frac1{12}s\nu+\frac29if\nu.
\]

\noindent 3) $\theta_\eta H_\vartheta(\theta Hu)$ se calcula según b) de modo análogo a $\theta_\eta H_\vartheta$:
\[
\theta_\eta H_\vartheta(\theta Hu)\sim-\theta_\eta^2(\theta Hu)-\frac16(\theta su)^3.
\]

\noindent 4) $\theta_\eta H_\vartheta(\vartheta\eta x)$ según b); $\theta_\eta^2(\vartheta\eta x)$ según d) (cf. \S\ 9):
\[
\theta_\eta H_\vartheta(\vartheta\eta x)\sim-\theta_\eta^2(\vartheta\eta x)-\frac16s_\vartheta(\theta su)
\sim-\frac13(\vartheta\R x)-\frac16s_\vartheta(\theta su),
\]
\[
\theta_\eta^2(\vartheta\eta x)\sim\frac23a_\eta^3(abu)\sim-\frac13(\vartheta\R x).
\]

\noindent 5)
\[
\begin{array}{ccc}
\theta_\eta H_\vartheta^2\text{ según b)},&
\theta_\eta^2H_\vartheta\text{ según a)},&
\theta_\eta^2H_\vartheta\text{ según b) y con ello }\theta_\eta^3\text{ según c)},\\
\text{de }a_\eta^2(aHb)(bHu);&
\text{de }a_\eta^3(Hab)(abu);&
\text{de }a_\eta^3(bHu)(abu).
\end{array}
\]
\[
\theta_\eta H_\vartheta^2\sim-\frac32\theta_\eta^2H_\vartheta-\frac14s_\vartheta(\theta su)^2+\frac16s_\nu-\frac49ia_\nu
\sim-\theta_\R-\frac16s_\nu-\frac19ia_\nu,
\]
\[
\theta_\eta^2H_\vartheta\sim\frac23\theta_\R-\frac16s_\vartheta(\theta su)^2+\frac29s_\nu-\frac29ia_\nu,
\]
\[
\theta_\eta^3H_\vartheta\sim\frac23\theta_\R-\frac23\theta_\eta^3+\frac5{36}s_\nu,
\qquad
\theta_\eta^3\sim\frac14s_\vartheta(\theta su)^2-\frac18s_\nu+\frac13ia_\nu.
\]

\noindent 6) $\theta_\eta^2(\theta Hu)^2$ según e):
\[
\A_\nu^2(a\A u)^2=[(a\A u)^4]_{\nu^2}
=-\frac{12}{5}\theta_\eta^2(\theta Hu)^2-\frac3{10}\sigma-\frac1{20}(\theta su)^4+\nu\R
\quad(\text{según la fórmula (3)c}),
\]
\[
\A_\nu^2(a\A u)^4=[\A_\nu^2]_{aa}^4=-\frac35\theta_\eta^2(\theta Hu)^2+\frac15\sigma-\frac3{10}(\theta su)^4+\frac13ij
\quad(\text{según la fórmula (10)a}),
\]
\[
\theta_\eta^2(\theta Hu)^2=-\frac16\sigma+\frac1{12}(\theta su)^4-\frac19ij+\frac13\nu\R.
\]

\noindent 7) $\theta_\eta^2(\vartheta\eta x)^2$ según d) y e). De aquí resulta la reducción de la forma superior $g$.

Por un cálculo análogo al de \S\ 9 se sigue:
\[
\begin{aligned}
\theta_\eta^2(\vartheta\eta x)^2
&=\frac12 i f-\frac12a_\eta^2b_\eta-\frac1{12}g
=\frac23tf-\frac23a_\eta^3b_\eta-\frac16g\\
&=-\frac16g-\frac13(\vartheta\R x)^2+\frac4{15}fa_\R^2+\frac8{15}ft
\quad(\text{según la fórmula (11)}).
\end{aligned}
\]
Por el cálculo correspondiente como arriba se sigue:
\[
((a\A\nu x)^4)=[(a\A u)^4]_{\nu x^4}
=\frac{21}{10}\theta_\eta^2(\vartheta\eta x)^2+\frac7{10}s_\vartheta^2-\frac3{10}g
\quad(\text{según la fórmula (3)c}),
\]
\[
\begin{aligned}
(a\A\nu x)^4&=-\frac13i\A+6[\A_\nu^2]a^2-4[\A_\nu^3]a+\A_\nu^4f\\
&= -\frac{36}{5}\theta_\eta^2(\vartheta\eta x)^2-\frac95s_\vartheta^2-\frac35g-\frac35(\vartheta\R x)^2\\
&\quad+\frac53i\A+\frac{37}{25}fa_\R^2+\frac{74}{25}ft\quad(\text{según la fórmula (10)}).
\end{aligned}
\]
Por tanto
\[
\begin{aligned}
\frac{93}{10}\theta_\eta^2(\vartheta\eta x)^2
&=-\frac3{10}g-\frac52s_\vartheta^2-\frac35(\vartheta\R x)^2\\
&\quad+\frac53i\A+\frac{37}{25}fa_\R^2+\frac{74}{25}ft,
\end{aligned}
\]
y al comparar los dos valores de $\theta_\eta^2(\vartheta\eta x)^2$ según g):
\[
\tag{13}
0=g-2s_\vartheta^2+2(\vartheta\R x)^2+\frac43i\A-\frac45fa_\R^2-\frac85ft.
\]

\noindent 8) $\theta_\eta^2H_\vartheta(\theta Hu)$ según a) y b); de aquí, $\theta_\eta^3(\theta Hu)$ según c) (las formas con el factor $u_x$ desaparecen):
\[
\theta_\eta^2H_\vartheta(\theta Hu)=-\frac16s_\vartheta(\theta su)^3-\frac13(\nu\R x),
\]
\[
\theta_\eta^2H_\vartheta(\theta Hu)=-\frac13\theta_\eta^3(\theta Hu)-\frac13(\nu\R x),
\]
\[
\theta_\eta^3(\theta Hu)=\frac12s_\vartheta(\theta su)^3.
\]

\noindent 9) $\theta_\eta^2H_\vartheta(\vartheta\eta x)$ según a) y b); de aquí, $\theta_\eta^3(\vartheta\eta x)$ según c). $\theta_\eta^3(\vartheta\eta x)$ según d); de aquí, la reducción de la forma superior $s_\vartheta^2(\theta su)$ según g) (las formas con el factor $u_x$ desaparecen):
\[
\theta_\eta^2H_\vartheta(\vartheta\eta x)=\frac13(atu),
\]
\[
\theta_\eta^2H_\vartheta(\vartheta\eta x)=\frac13(atu)+\frac13\theta_\eta^3(\vartheta\eta x)-\frac16s_\vartheta^2(\theta su),
\]
\[
\theta_\eta^3(\vartheta\eta x)=\frac12s_\vartheta^2(\theta su),
\]
\[
\theta_\eta^3(\vartheta\eta x)=\frac25\theta_\R(\vartheta\R x)+\frac15(atu),
\]
\[
\tag{14}
s_\vartheta^2(\theta su)=\frac45\theta_\R(\vartheta\R x)+\frac25(atu).
\]

\noindent 10) $\theta_\eta^2H_\vartheta^2$ según a) y mediante el reductor $\theta_\nu^2(\vartheta\nu x)^2$; $\theta_\eta^3H_\vartheta$ según a) y c); $\theta_\eta^4$ según c); de aquí, la reducción de las formas superiores $s_\vartheta^2(\theta su)^2$ y $s_\nu^2$ según g):
\[
\begin{aligned}
\text{Según a)}\quad
\theta_\eta^2H_\vartheta^2&=[a_\eta^2(aHu)^2]b^2+\frac13[a_\eta^4]_{bu^2}-\frac13H_\nu^2(\nu\eta x)^2\\
&=\frac49ia_\nu^2-\frac16s_\vartheta^2(\theta su)^2-\frac5{18}s_\nu^2+\frac13(atu)^2+\frac1{54}Ju_x^2.
\end{aligned}
\]
\[
\begin{aligned}
\text{Mediante }\theta_\nu^2(\vartheta\nu x)^2:
\quad\theta_\eta^2H_\vartheta^2
&=[\theta_\nu^2(\vartheta\nu x)^2]_{\nu_1}+\frac13[\theta_\nu^4]_{\nu_1}\hat\nu_1x^2-\frac13s_\vartheta^2(\theta su)^2\\
&=\frac49ia_\nu^2-\frac16s_\nu^2-\frac13s_\vartheta^2(\theta su)^2+\frac13(\nu\R x)^2.
\end{aligned}
\]
\[
\begin{aligned}
\text{Según a)}\quad
\theta_\eta^3H_\vartheta&=-[a_\eta^2(aHu)]_{bu}b^2-\frac12[a_\eta^2(aHu)^2]b^2-\frac13[a_\eta^3]b+\frac1{12}[a_\eta^4]_{bu^2}\\
&\quad+\frac12u_x[a_\eta^2(aHu)^2]b^3\\
&=-\frac29ia_\nu^2+\frac1{12}s_\nu^2+\frac4{15}\theta_\R^2-\frac7{90}(\nu\R x)^2+\frac1{30}(atu)^2+\frac2{27}i^2u_x^2-\frac1{36}Ju_x^2.
\end{aligned}
\]
\[
\begin{aligned}
\text{Según c)}\quad
\theta_\eta^3H_\vartheta:
 a_\eta^3(Hab)(bHu)&=\frac12(atu)^2
=\theta_\eta^2H_\vartheta(\theta Hu)(\vartheta\eta x)+\frac14H_\nu^2(\nu\eta x)^2\\
&\quad+\frac12\theta_\eta^2H_\vartheta^2
=\frac32\theta_\eta^2H_\vartheta^2+\theta_\eta^3H_\vartheta-u_x[a_\eta^2(aHu)^2]b^3\\
&\quad+\frac14H_\nu^2(\nu\eta x)^2,
\end{aligned}
\]
y por tanto
\[
\begin{aligned}
\theta_\eta^3H_\vartheta&=-\frac23ia_\nu^2+\frac5{12}s_\nu^2+\frac12(\nu\R x)^2-\frac12(atu)^2\\
&\quad+\frac4{27}i^2u_x^2-\frac1{12}Ju_x^2.
\end{aligned}
\]
Según c),
\[
\begin{aligned}
\theta_\eta^4&=\frac49ia_\nu^2-\frac16s_\nu^2+\frac{16}{15}\theta_\R^2\\
&\quad-\frac{14}{45}(\nu\R x)^2+\frac2{15}(atu)^2.
\end{aligned}
\]
Según g), a partir de los dos valores para $\theta_\eta^2H_\vartheta^2$:
\[
\tag{15}
\begin{aligned}
s_\vartheta^2(\theta su)^2
&=\frac23s_\nu^2-2(atu)^2+2(\nu\R x)^2-\frac19Ju_x^2\\
&=\frac89ia_\nu^2+\frac8{15}\theta_\R^2-\frac{14}{15}(atu)^2+\frac{38}{45}(\nu\R x)^2-\frac4{27}i^2u_x^2.
\end{aligned}
\]
Según g), a partir de los dos valores para $\theta_\eta^3H_\vartheta$:
\[
\tag{16}
s_\nu^2=\frac43ia_\nu^2+\frac45\theta_\R^2+\frac85(atu)^2-\frac{26}{15}(\nu\R x)^2+\frac16Ju_x^2-\frac29i^2u_x^2.
\]
La fórmula (16) proporciona la base para la reducción del módulo $(ss'u)^4$, y la fórmula (13) para la reducción del sistema de $s$ (\S\ 17).

\noindent\emph{Consecuencias:}

1) $a_\nu(j\nu x)^3$, $\theta_\eta H_\vartheta$, $\theta_\eta^2(\vartheta\eta x)$, $\theta_\eta^2(\theta Hu)^2$ son reductores; $a_\nu(j\nu x)^2$, $H_\vartheta$ y $H_\vartheta^2$ son formas descomponibles.

2) La forma del sistema II, $j(\nu)=g=(\nu\eta x)^4H_x^2$, es reducible.

3) Como la única forma contragrediente $g$ se vuelve reducible, $\nu$ no entra en absoluto, ni en potencias ni en productos con formas del mismo sistema, en plegamiento con el sistema I.

$\theta$ entra en plegamiento sólo considerada como contravariante en productos o potencias, y correspondientemente $H$ sólo considerada como covariante (cf. \S\ 13 B).




\begin{center}
\textbf{\S\ 12.\quad Formas de séptimo orden.}
\end{center}
\[
\text{Plegamiento de }\left\{\begin{array}{l}
\theta\cdot K\text{ con }\nu,\\
K,j,\A\text{ con }H,\\
f\text{ con }L,\sigma.
\end{array}\right.
\]

\noindent\emph{Formas irreducibles:}
\[
\text{B)}\quad
\begin{array}{c|c|c|c|c}
\cdot&\cdot&(KHu)^2&\cdot&\cdot\\
\cdot&K_\eta&\cdot&K_\eta(KHu)^2&\cdot\\
(k\eta x)^2&\cdot&K_\eta^2&\cdot&\cdot\\
(k\eta x)^3&H_k(k\eta x)^2,\ K_\eta(k\eta x)^2&\cdot&\cdot&\cdot\\
\cdot&K_\eta(k\eta x)^3&\cdot&\cdot&\cdot
\end{array}
\]
\[
\text{C)}\quad
\begin{array}{cc}
(j\eta x)&H_j\\
\cdot&H_j(j\eta x)\\
\cdot&H_j^2
\end{array}
\qquad
\text{D)}\quad
\begin{array}{c}
(\A Hu)\\\A_\eta
\end{array}
\]
\[
\text{E)}\quad
\begin{array}{cccc}
a_i&(aLu)^2&(aLu)^3&\cdot\\
\cdot&\cdot&a_i(aLu)^2&a_i(aLu)^3\\
\cdot&a_i^2&\cdot&\cdot
\end{array}
\]

\noindent\emph{Reducciones:}

\noindent A) $(\vartheta\cdot k,\nu,x)^4$ se descompone como una transvección simple sobre la forma descomponible $(\vartheta^2\nu x)^4$.

\noindent B) Sucesión de formas $H_kK_\eta$: reductor $H_\vartheta\theta_\eta$.

Sucesión de formas $K_\eta^2(KHu)$: reductor $a_\eta^2(aHu)$.

Sucesión de formas $K_\eta^2(k\eta x)$: reductor $\theta_\eta^2(\vartheta\eta x)$.

La doble reducción de la sucesión de formas $K_\eta^3$ que resulta de ahí da ocasión a la reducción de la sucesión superior $(j\eta x)^3$.

$(KHu)$, $(k\eta x)$, $K_\eta(k\eta x)$ son formas descomponibles según \S\ 3.

$H_k(KHu)$, $K_\eta(KHu)$ se descomponen como surgidas por plegamiento simple con la forma descomponible $K_\nu(k\nu x)$ y el determinante funcional
\[
K_\nu^2=\theta_\nu^2(a\theta u)\equiv a_\nu^2(a\theta u)\modu{\nu^2}.
\]
A la doble representación de $K_\nu^2$ corresponde una relación entre formas descomponibles. Se obtiene:
\[
H_k(KHu)=2K_\nu(\nu\nu_1k)=2\theta_{\nu_1}(\nu\nu_1a\theta)
=2\theta_\nu^2a_\nu-a_\nu^2\theta_\nu+f\theta_\eta(\theta Hu)^2-f\nu\theta_\vartheta^3
\modu{\nu^3},
\]
\[
\begin{aligned}
K_\eta(KHu)&=-K_\nu^2(\nu\nu_1x)
=-a_\nu^2(a\theta u)(\nu\nu_1x)\\
&=a_\eta(aHu)^2\theta+a_\nu^2\theta_\nu
=-\theta_\nu(\theta Hu)^2f-\theta_\nu^2a_\nu.
\end{aligned}
\]
Por tanto
\[
\tag{17}
0=a_\nu^2\theta_\nu+\theta_\nu^2a_{\nu_1}+f\theta_\eta(\theta Hu)^2+\theta a_\eta(aHu)^2\modu{\nu^3}.
\]
Las formas $H_k$, $H_k(k\eta x)$, $H_k^2$, $H_k^2(k\eta x)$ se descomponen como surgidas por plegamiento simple con las formas descomponibles $H_\vartheta$ y $H_\vartheta^2$ (cf. \S\ 3. 2), a) y b)).

Se tiene
\[
H_k=H_\vartheta(a\theta u)\sim [H_\vartheta]_{ua}-fH_\vartheta(\theta Hu)
\]
(especialización $y=uv$ de 2)a).
\[
H_k(k\eta x)=H_\vartheta a_\eta-H_\vartheta\theta_\eta f,
\qquad
H_\vartheta a_\eta=\frac54[H_\vartheta]a-\frac14H_\vartheta a_\vartheta
\quad(\S\ 3. 2)b),
\]
\[
H_\vartheta^2(a\theta u)=[H_\vartheta^2]_{ua},
\qquad
H_\vartheta^2a_\eta=[H_\vartheta^2]a=H_k^2(k\eta x)+H_\vartheta^2\theta_\eta f.
\]

\noindent C) Sucesión de formas $(j\eta x)^3$: reducible por doble reducción de la sucesión de formas $K_\eta^3$ según B), de acuerdo con el planteamiento:
\[
0=a_\eta^3(a\theta u)=\theta_\eta^3(a\theta u)
=\theta_\eta+(a\theta\eta x)^3(a\theta u)-\theta_\eta^3(a\theta u)
=\frac12(j\eta x)^3\modu{(\nu^3,s,\R,t,i,J)}.
\]
El planteamiento análogo vale hasta la forma final de la sucesión:
\[
0=H_\vartheta^2(a\theta\eta x)a_\eta^3
=H_\vartheta^2(a\theta\eta x)\theta_\eta^3
-\frac12H_\vartheta^2(j\eta x)^4\modu{(\nu^3,s,\R,t,i,J)}.
\]
Forma $(j\eta x)^2$: se descompone 1) por doble polarización de la relación para la forma descomponible $H_\vartheta^2$ ((12.) b);

2) por cálculo directo del producto $a_\nu\theta_\nu^3$; de ahí la descomposición de la forma superior $(\A Hu)^2$. Se obtiene:
\[
\tag{18}
\left.\begin{array}{ll}
\text{a)}&\dfrac13(\A Hu)^2+\dfrac13(j\eta x)^2=\theta_\nu^2a_{\nu_1}^2,\\[0.6em]
\text{b)}&(j\eta x)^2=\dfrac43\theta_\nu^3a_{\nu_1},
\end{array}\right\}\modu{(\nu^3,s,\R,t,i,J)}.
\]
según el planteamiento:
\[
H_\vartheta^2(a\theta u)^2-\frac13(\A Hu)^2
=2\theta_\eta^2(a\theta u)(aHu)+\theta_\nu^2a_{\nu_1}^2
=(a\theta u)(aHu)\{a_\eta-(a\theta\eta x)^2\}+\theta_\nu^2a_{\nu_1}^2,
\]
\[
\begin{aligned}
a_{\nu_1}\theta_\nu^3
&=-(a\widehat{\nu}\nu_1u)\theta_\nu^2(\theta\widehat{\nu}\nu_1u)
-\frac12(u\nu\nu_1u)\theta_\nu\theta_{\nu_1}(\vartheta\nu\nu_1u)\\
&=-\frac32\theta_\eta^2(\theta Hu)(aHu)
=-\frac32\theta_\eta^2(\theta Hu)(a\theta u)\\
&=-\frac32\{a_\eta-(a\theta\eta x)^2\}\{(aHu)-(a\theta u)\}(a\theta u).
\end{aligned}
\]

\noindent D) Sucesión de formas $\A_\eta(\A Hu)$: reductor $\A_\nu^2$.

$(\A Hu)^2$ es forma descomponible según C).

\noindent E) Sucesión de formas $a_i^2(aLu)$: reductor $a_\eta^2(aHu)$. Las formas $(aLu)$ y $a_i(aLu)$ se descomponen como transvecciones sobre determinantes funcionales según \S\ 3.

\noindent F) Sucesión de formas $a_\sigma^3$: reductor $a_\R^3$.

Formas $a_\sigma^2$ y $a_\sigma^3$: reductores $a_\nu^3$ y $a_\eta^3$ por doble reducción de las expresiones correspondientes
\[
H_\nu a_\nu^3\quad\text{y}\quad H_\nu a_\eta^3,
\qquad
H_\nu a_\nu^3a_\nu\quad\text{y}\quad H_\nu a_\eta^4
\]
(cf. \S\ 10. C), reducción de $\A_\nu^2$ y $\A_\nu^3$).

De ahí la reducción para las formas superiores $a_\nu s_\nu(asu)^2$, $a_\nu^2s_\nu(asu)^2$ y para formas en los símbolos $\R,t(a_\nu,a_\eta^2)$, teniendo en cuenta (16.):
\[
\tag{19}
\left.\begin{array}{rcl}
a_\nu s_\nu(asu)^2&=&\dfrac13 i\theta_\nu^2-\dfrac1{18}iH,\\[0.5em]
a_\nu^2s_\nu(asu)^2&=&0,
\end{array}\right\}\modu{(\R,t)}.
\]
La forma $a_\sigma$ se descompone por polarización de (12.)b o, más brevemente, por desarrollo directo del producto $a_\nu^2\theta_\nu^3$ según el planteamiento:
\[
a_\nu^2\theta_{\nu_1}^3=(a_\nu^2\theta_{\nu_1})a_{\nu_1}-(a\theta\nu_1x)^2
=-2a_\eta(aHu)(a\theta H)(a\theta\eta x)+(a\theta\nu_1x)^2a_\nu^2\theta_{\nu_1}
\]
y teniendo en cuenta la reducción de $H_j(j\eta x)^2$ (C):
\[
\tag{20}a_\sigma=2a_\nu^2\theta_{\nu_1}^3\modu{(s,\R,t,i)}.
\]

\noindent\emph{Consecuencias:}

1) $(j\eta x)^3$, $\A_\eta(\A Hu)$, $a_\sigma^2$ son reductores; $(j\eta x)^2$, $(\A Hu)^2$, $a_\sigma$ son formas descomponibles.

2) $f$ entra en plegamiento con $L$ sólo en productos con formas contragredientes; $f$, y por tanto también $K$ y $N$, no entran en absoluto en plegamiento con $\sigma$ y, por consiguiente, tampoco con $(\nu\sigma x)$. $j$ entra sólo en productos con formas contragredientes; $\A$ no entra en absoluto en productos en plegamiento con $H$ y $L$ (esto se sigue de $\A_\eta(\A Hu)$ y $(\A Hu)^2$). Del mismo modo $\sigma$, y por tanto $(\nu\sigma x)$, no entra en productos en plegamiento con ninguna forma del sistema I. En los productos del sistema II, teniendo en cuenta las consecuencias de \S\ 11, sólo quedan potencias de $H$ y productos de $H$ con $L$.

\begin{center}
\textbf{\S\ 13.\quad Formas de octavo orden (sistema III).}
\end{center}
\[
\text{Plegamiento de }\left\{\begin{array}{l}
\A\cdot j\text{ con }\nu,\\
\theta^2,\\
 f\cdot j,\\
 N\text{ con }H,\\
\theta\text{ con }L,\sigma.
\end{array}\right.
\]

\noindent\emph{Formas irreducibles:}
\[
\begin{array}{c@{\quad}l@{\qquad}c@{\quad}l}
\text{A)}&\A_\nu(j\nu x)&
\text{B)}&\begin{array}{c}(\vartheta^2\eta x)^3\\(\vartheta^2\eta x)^4\end{array}
\quad H_\vartheta(\vartheta^2\eta x)^3,
\ \theta_\eta(\vartheta^2\eta x)^3,\\[1em]
\text{C)}&(aHu)H_j,
\quad a_\eta(aHu)H_j&
\text{D)}&H_n,
\ N_\eta,
\end{array}
\]
\[
\text{E)}\quad
\begin{array}{c|c|c|c}
\cdot&(\theta Lu)^2&(\theta Lu)^3&\cdot\\
\theta_i&\cdot&L_\vartheta(\theta Lu)^2,
\ \theta_i(\theta Lu)^2&\theta_i(\theta Lu)^3\\
(\vartheta lx)^2&\theta_i^2&\cdot&\cdot\\
\cdot&\theta_i(\vartheta lx)^2&\cdot&\cdot
\end{array}
\]

\noindent\emph{Reducciones:}

\noindent A) $\A_\nu(j\nu x)^3$: reductor $a_\nu(j\nu x)^3$.

$\A_\nu(j\nu x)^2$ es reducible por la forma descomponible $a_\nu(j\nu x)^2$ según \S\ 3. 2)b, teniendo en cuenta el reductor $\theta_{\vartheta j}$. En efecto, según \S\ 11 B) se obtiene:
\[
\frac3{10}\A_\nu(j\nu x)^2+\frac35a_\nu(j\nu x)a_\vartheta(j\nu x)
+\frac1{10}a_\nu(j\nu x)^2
=\frac35\theta_\nu^3\theta_{\vartheta j}+\frac25\theta_\nu^3\theta_{\vartheta j}\theta_{\vartheta j},
\]
(reductores $a_j$ y $\theta_{\vartheta j}$).

\noindent B) Las formas $H_\vartheta(\vartheta'\eta x)^2$, $H_\vartheta^2(\vartheta'\eta x)$, $H_\vartheta^2(\vartheta'\eta x)^2$ se descomponen como surgidas por plegamiento simple sucesivo con las formas descomponibles $H_\vartheta$ y $H_\vartheta^2$, respectivamente $H_\vartheta a_\eta$ y $H_\vartheta^2a_\eta$, según \S\ 3. 2)b) y según la relación:
\[
H_\vartheta(\vartheta'\eta x)^2=2H_\vartheta a_\eta^2f-2H_\vartheta a_\eta b_\eta.
\]

\noindent C) Sucesión de formas $a_\eta(j\eta x)^2$: reductores $a_j$ y $(j\eta x)^3$.

La forma $a_\eta(j\eta x)$ se descompone: reductor $a_j$ y plegamiento simple con la forma descomponible $(j\eta x)^2$ según \S\ 3, 2)a (especialización $y=uv$).

La forma $a_\eta H_j$ se descompone: 1) por plegamiento simple con la forma descomponible $a_\nu(j\nu x)^2$;

2) por plegamiento doble especial con la forma descomponible $(j\eta x)^2$; de ahí una relación entre formas descomponibles.

A partir de $a_\nu(j\nu x)^2$:
\[
0=\theta_\nu'\theta_\nu+2a_\eta H_j\modu{(s,\R,t,i,J)}.
\]
A partir de $(j\eta x)^2$:
\[
0=\theta_\nu'\theta_\nu-\frac32a_\eta H_j\modu{(s,\R,t,i,J)}
\]
(se omiten los productos con contravariantes), según el planteamiento:
\[
0=\frac12a_\nu(abu)^2\theta_\nu^3
+\frac12u_x(ubu)\theta_{\nu_1}^3(\theta bu)
-\frac34(\widehat{j}\eta bu)^2+\frac34H_j(\widehat{j}\eta bu)
-\frac18fH_j^2
\]
y permutando los símbolos en $a_\nu(ubu)(\theta bu)\theta_{\nu_1}^3$.

\noindent D) Sucesión de formas $N_\eta(NHu)$: reductor $\A_\nu^2$.

$(NHu)$, $(n\eta x)$, $N_\eta(n\eta x)$, $H_\eta(NHu)$ son formas descomponibles (cf. \S\ 12 B); $(NHu)^2$ se descompone por plegamiento simple con la forma $(\A Hu)^2$.

\noindent E) Sucesiones de formas $\theta_iL_\vartheta$, $\theta_i^2(\theta Lu)^2$, $\theta_i^2(\vartheta lx)$:

Reductores: $\theta_\eta H_\vartheta$, $\theta_\eta^2(\theta Hu)^2$, $\theta_\eta^2(\vartheta\eta x)$.

Las formas $(\theta Lu)$, $(\vartheta lx)$, $L_\vartheta$, $L_\vartheta(\theta Lu)$, $\theta_i(\theta Lu)$, $L_\vartheta(\vartheta lx)$, $\theta_i(\vartheta lx)$, $L_\vartheta^2$, $L_\vartheta^2(\theta Lu)$, $\theta_i^2(\theta Lu)$ se descomponen. (Dualísticamente opuestas a las formas descomponibles de \S\ 12 B).)

\noindent F) Sucesión de formas $\theta_\sigma(\vartheta\sigma x)$: reductor $a_\sigma^2$.

Las formas $(\vartheta\sigma x)$ y $(\vartheta\sigma x)^2$ se descomponen, como surgidas por plegamiento simple con la forma descomponible $a_\sigma$; $\theta_\sigma$, surgida por plegamiento doble, se descompone, pues la sucesión de formas
\[
a_{\nu_2}^2(a\theta u)\theta_{\nu_1}^3\equiv H_j^2(j\eta x)\equiv0\modu{u_x(s,\R,t,i,J)}:
\]
\[
\theta_\sigma=a_\sigma(abu)^2=a_\nu^2(abu)^2\theta_\nu^3.
\]

\noindent\emph{Consecuencias:}

$\theta^3$ o $\theta^2K$ no entra en plegamiento con $H$; $\theta$, considerada sólo como contravariante, entra en potencias o productos en plegamiento con $L$, pero no entra en absoluto en plegamiento con $\sigma$ y $(\nu\sigma x)$.

$N$ no entra en productos en plegamiento con ninguna forma del sistema II, ni tampoco con potencias o productos de formas del sistema II. En los productos del sistema I, teniendo en cuenta las consecuencias de los últimos párrafos, sólo quedan los productos $f\cdot j$, $\A\cdot j$, las potencias de $\theta$ y los productos $\theta\cdot K$.

\begin{center}
\textbf{\S\ 14.\quad Formas de noveno orden (sistema III).}
\end{center}
\[
\text{Plegamiento de }\left\{\begin{array}{l}
\theta\cdot K\text{ con }H,\\
K,j,\A\text{ con }L,\\
j,\A\text{ con }\sigma,\\
f\text{ con }H^2.
\end{array}\right.
\]

\noindent\emph{Formas irreducibles:}
\[
\begin{array}{c@{\quad}l@{\qquad}c@{\quad}l}
\text{A)}&(\vartheta\cdot k,\eta,x)^4,
\ H_k(\vartheta\cdot k,\eta,x)^4&
\text{B)}&(KLu)^3,
\ K_i(KLu)^3,
\ K_i^2,
\ (klx)^3,
\ K_i(Klx)^3,\\[0.7em]
\text{C)}&L_j,
\ L_j^2&
\text{D)}&\A_i,\\[0.7em]
\text{E)}&(j\sigma x)&
\text{G)}&(aH^2u)^3,
\ (aH^2u)^4,
\ a_\eta(aH^2u)^3.
\end{array}
\]

\noindent\emph{Reducciones:}

A) Las formas $H_\vartheta(k\eta x)^3$, $H_\vartheta^2(k\eta x)^2$, $H_\vartheta^2(k\eta x)^3$ se descomponen como surgidas por plegamiento simple sucesivo con formas descomponibles.

B) Sucesiones de formas $K_iL_k$, $K_i^2(KLu)$, $K_i^2(klx)$:

Reductores: $\theta_\eta H_\vartheta$, $a_\eta^2(aHu)$, $\theta_\eta^2(\vartheta\eta x)$.

Las formas $L_i$, $K_i(KLu)$, $K_i(klx)$, $(KLu)^2$, $(klx)^2$ se descomponen como transvecciones sobre determinantes funcionales (\S\ 3).

Las formas $L_k$, $L_k(KLu)$, $L_k(KLu)^2$, $L_k(klx)$, $L_k(klx)^2$, $L_k^2$, $L_k^2(KLu)$, $L_k^2(klx)$, $L_k^3$, $K_i(KLu)^2$, $K_i(klx)^2$ se descomponen como surgidas por plegamiento simple con las formas descomponibles:
\[
L_\vartheta,
\ H_k(KHu),
\ L_\vartheta(\vartheta lx),
\ H_k^2,
\ L_\vartheta^2,
\ L_\vartheta^2(\theta Lu),
\ K_\eta(KHu),
\ \theta_i(\vartheta lx)
\]
según \S\ 3. 2), a) y b).

C) Sucesión de formas $(jlx)^3$: reductor $(j\eta x)^3$.

Las formas $(jlx)^2$ y $L_j(jlx)$ surgen por plegamiento simple con la forma descomponible $(j\eta x)^2$.

D) Sucesión de formas $\A_i(\A Lu)$: reductor $\A_\nu^2$.

Formas $(\A Lu)^2$, $(\A Lu)^3$: plegamiento simple con la forma descomponible $(\A Hu)^2$.

E) Sucesión de formas $(j\sigma x)^2$: reductor $a_\sigma^2$ por sustitución de $j$ por formas inferiores de la sucesión de formas (\S\ 5):
\[
a_\sigma^2(a\theta u)^2=\frac13(j\sigma x)^2,
\qquad
a_\sigma^2(a\theta\sigma x)^2=\frac13(j\sigma x)^4.
\]
F) Sucesión de formas $\A_\sigma^2$: reductor $a_\sigma^2$.

Forma $\A_\sigma$: reductor $a_\sigma^2$ por sustitución de $\A$ por formas inferiores de la sucesión:
\[
a_\sigma a_\R^2=\frac23\A_\sigma\modu{\nu}.
\]

\noindent\emph{Consecuencias:} Sólo la forma $j$ entra en plegamiento con $\sigma$ y $(\nu\sigma x)$.

$N$ no entra en plegamiento con $L$, pues la forma $N_i$ correspondiente a $\A_i$ se descompone.

\begin{center}
\textbf{\S\ 15.\quad Formas de décimo orden (sistema III).}
\end{center}
\[
\text{Plegamiento de }\left\{\begin{array}{l}
\theta^2,\\
\ f\cdot j\text{ con }L,\\
\theta\text{ con }H^2.
\end{array}\right.
\]

\noindent\emph{Formas irreducibles:}
\[
\begin{array}{c@{\quad}l@{\qquad}c@{\quad}l}
\text{A)}&(\vartheta^2lx)^4,
\ \theta_i(\vartheta^2lx)^4&
\text{B)}&(aLu)^2L_j,
\ a_i(aLu)^2L_j,\\[0.8em]
\text{C)}&(\theta H^2u)^3,
\ (\theta H^2u)^4,
\ H_\vartheta(\theta H^2u),
\ \theta_\eta(\theta H^2u)^3.
\end{array}
\]

\noindent\emph{Reducciones:}

A) y B): descomposición de las restantes formas por plegamiento simple con formas descomponibles.

C) Descomposición de las restantes formas según \S\ 13 B), por la consideración dualísticamente opuesta.

\noindent\emph{Consecuencias:}

$\theta^2K$ no entra en plegamiento con $L$; correspondientemente, $H^2L$ no entra en plegamiento con $K$. $H^3$ y $H^2L$ no entran en plegamiento con $\theta$. Con ello queda demostrado el esquema de plegamiento de \S\ 7.

\begin{center}
\textbf{\S\ 16.\quad Formas de orden 11, 12, 13 y 15 (sistema III).}
\end{center}

\noindent\emph{Orden 11.}
\[
\text{Plegamiento de }\left\{\begin{array}{l}
\theta\cdot K\text{ con }L,\\
K,j\text{ con }H^2,\\
j\text{ con }(\nu\sigma x),\\
f\text{ con }H\cdot L.
\end{array}\right.
\]
\noindent\emph{Formas irreducibles:}
\[
\begin{array}{c@{\quad}l@{\qquad}c@{\quad}l}
\text{A)}&(\vartheta\cdot k,l,x)^5,
\ \theta_i(\vartheta\cdot k,l,x)^5&
\text{B)}&(KH^2u)^4,
\ K_\eta(KH^2u)^4,\\[0.8em]
\text{C)}&(\nu\sigma j),&
\text{D)}&(a_\eta H\cdot L,u)^4.
\end{array}
\]

\noindent\emph{Orden 12.}
\[
\text{Plegamiento de }\left\{\begin{array}{l}
\theta^3\text{ con }L,\\
\theta\text{ con }H\cdot L.
\end{array}\right.
\]
\noindent\emph{Formas irreducibles:}
\[
\text{E)}\ (\vartheta^3lx)^6,
\qquad
\text{F)}\ (\theta,H\cdot L,u)^4,
\qquad L_\vartheta(\theta,H\cdot L,u^4).
\]

\noindent\emph{Orden 13.}
\[
\text{Plegamiento de }K\text{ con }H\cdot L.
\]
\noindent\emph{Formas irreducibles:}
\[
\text{G)}\ (K,H\cdot L,u)^5,
\qquad K_\eta(K,H\cdot L,u)^5.
\]

\noindent\emph{Orden 15.}
\[
\text{Plegamiento de }K\text{ con }H^3.
\]
\noindent\emph{Forma irreducible:}
\[
\text{H)}\ (KH^3u)^6.
\]

\noindent\emph{Reducciones:}

Las formas $H_j^2$, $H_j^2$ tienen el reductor $L_j^3$, a partir de la relación:
\[
(\nu lx)L_x^3=-\frac12H^2\modu{(\sigma,s)}.
\]
La descomposición o reducibilidad de las restantes formas sigue de las consideraciones de los párrafos anteriores.

\noindent\emph{Consecuencias:}

Según el esquema de plegamiento de \S\ 7 y las consecuencias de \S\S\ 8--15, quedan con esto agotadas las formas irreducibles del sistema III (117 formas).

\clearpage

\begin{center}
{\Large\bfseries Capítulo IV.\quad El sistema relativamente completo módulo $(\varrho,t)$.}
\end{center}

\begin{center}
\textbf{\S\ 17.\quad Reducción del módulo $(ss'u)^4$ y del sistema de $s$.}
\end{center}

Para formar el sistema relativamente completo inmediatamente superior, por analogía con \S\ 7 hay que formar el sistema de $s$ con respecto al módulo siguiente
\[
\nu(s)=(ss'u)^4
\]
y plegarlo con las formas del sistema relativamente completo módulo $s$. Se mostrará que, al introducir el módulo $(\varrho,t)$ como módulo superior,
\[
\begin{array}{ll}
\text{A)}&\text{el módulo }(ss'u)^4\text{ se vuelve reducible,}\vphantom{\dfrac11}\\
\text{B)}&\text{el sistema de }s\text{ consta de la única forma }s.
\end{array}
\]

A) La reducción del módulo $(ss'u)^4$ se obtiene inmediatamente del reductor $s_\nu^2$ encontrado mediante la fórmula (16.), \S\ 11. De
\[
 s_\nu^2=\frac43 i a_\nu^2+\frac45\theta_\varrho^2+\frac85(atu)^2-
 \frac{26}{15}(\nu\varrho x)^2+\frac16J\cdot u_x^2-
 \frac29i^2\cdot u_x^2
\]
se sigue, por polarización de $x$ hacia $\nu_1$:
\[
\tag{21}
\begin{aligned}
(ss'u)^4&=-\frac23J\cdot\nu+6s_\nu^2s_{\nu_1}^2\\
&=\frac{24}{5}\theta_\nu^2\theta_\varrho^2+
\frac{48}{5}a_\nu^2(atu)^2-
\frac{52}{5}H_\varrho^2+
\frac43i\cdot(asu)^4+\frac13J\cdot\nu-
\frac49i^2\cdot\nu,
\end{aligned}
\]
y con ello la reducción del módulo $(ss'u)^4$.\footnote{De la fórmula (21.) se sigue la relación, mencionada en la nota a la introducción, entre los tres invariantes de orden 9, aplicando la fórmula (10.)c.}

B) La reducción de
\[
Z=(ss'u)^2=\theta(s)
\]
y de ahí la reducción del sistema de $s$ se obtiene sustituyendo la forma $Z$ por la forma inferior $g_\nu^2$ según la fórmula (8.)b, \S\ 7, teniendo en cuenta la relación
\[
 s_\eta^2=s_\nu^2(\nu\nu_1x)^2-\frac13Z.
\]
La forma $g_\nu^2$ tiene, según la fórmula (13.), \S\ 11, el reductor $s_\nu^2$ como factor. Se obtiene, por las fórmulas (8.) y (16.):
\[
 g_\nu^2=-Z+\frac{14}{5}ia_\eta^2+\frac{14}{15}i(asu)^2
 \modu{(\varrho,t)},
\]
y, por las fórmulas (13.), (14.), (15.), (16.):
\[
\begin{aligned}
g_\nu^2&=2[s_\vartheta^2]_{\nu^2}-\frac43 i\A_\nu^2
      =2s_\vartheta^2s_\nu^2-\frac65[s_\vartheta^2(\theta su)^2]\widehat{\nu x}^{\,2}
        -\frac43i\A_\nu^2\\
&=\frac45i\cdot a_\eta^2-\frac25i\cdot(asu)^2+\frac13J\cdot\theta
\modu{(\varrho,t)}.
\end{aligned}
\]
De aquí se sigue
\[
\tag{23}
 Z=2i\cdot a_\eta^2+\frac43i\cdot(asu)^2-\frac13J\cdot\theta
 \modu{(\varrho,t)}.
\]

El sistema relativamente completo módulo $((ss'u)^4,\varrho,t)$ pasa por tanto a un sistema relativamente completo módulo $(\varrho,t)$; y de éste surge, por transvección sobre el sistema conocido de dos formas cuadráticas, el sistema absolutamente completo. Para formar el sistema relativamente completo módulo $(\varrho,t)$, puesto que según la fórmula (23.) el sistema de $s$ se reduce a la forma $s$, debemos transveccionar el sistema relativamente completo módulo $(s,\varrho,t)$ sobre $s$ y sobre potencias de $s$. Se verá que las potencias de $s$ sólo entran en plegamiento con el sistema I.

La relación anunciada es:
\[
\tag{22}
\begin{aligned}
(ass')^4&=\frac{24}{5}\theta_\nu^2\theta_\varrho^2a_\nu^2a_\varrho^2+
\frac{48}{5}a_\nu^2b_\nu^2(tab)^2-
\frac{52}{5}H_\varrho^2a_\nu^4+
\frac53J\cdot i-\frac49i^3,\\
(ass')^4&=\frac{24}{5}a_\varrho^2a_{\varrho'}^2-
\frac{12}{5}t_\varrho^2+\frac53J\cdot i-\frac49i^3.
\end{aligned}
\]

Como ``formas superiores'' definimos, entre las formas del mismo orden y del mismo grado, aquellas que han surgido de formas superiores del sistema relativamente completo módulo $(s,\varrho,t)$ por plegamiento con $s$, considerando el plegamiento con $s$ sólo como polarización de las formas originarias. Con ello queda determinada de modo unívoco la ordenación de las formas individuales (cf. \S\ 1, sucesiones de formas 2). Por ejemplo:
\[
(\theta_\eta(\theta Hu),s,u)^2\text{ es superior a }s_\eta((\theta Hu)^2,s,u),
\]
\[
(\theta_\eta(\theta Hu)^2,s,u)\text{ es superior a }(\theta_\eta(\theta Hu),s,u)^2.
\]

Dividimos el sistema resultante en los cuatro subsistemas:
\[
\begin{array}{ll}
\text{Sistema }s_I:&\text{surgido por plegamiento de }s\text{ y potencias de }s\text{ con el sistema I;}\vphantom{\dfrac11}\\
\text{Sistema }s_{II}:&\text{surgido por plegamiento de }s\text{ con el sistema II;}\vphantom{\dfrac11}\\
\text{Sistema }s_{III}:&\text{surgido por plegamiento de }s\text{ con las formas individuales}\\
&\text{(sin productos) del sistema III y con productos de una}\\
&\text{forma de cada uno de los sistemas I y II;}\vphantom{\dfrac11}\\
\text{Sistema }s_{IV}:&\text{surgido por plegamiento de }s\text{ con productos de una forma de cada}\\
&\text{uno de los sistemas I y III, II y III, III y III.}
\end{array}
\]
Observación: de ahora en adelante, todas las fórmulas deben entenderse módulo $(\varrho,t)$; desde la fórmula (27.) en adelante, módulo $(\varrho,t,i,J,\text{ formas superiores})$, sin que esto se indique expresamente cada vez.

Para la reducción reunimos las fórmulas obtenidas anteriormente:
\[
\tag{24}
\begin{aligned}
s_\nu^2&=\frac43ia_\nu^2+\frac16J\cdot u_x^2-\frac29i^2\cdot u_x^2,
& s_\nu^3&=\frac13J\cdot u_x,
& s_\nu^4&=J,\\[0.3em]
g&=2s_\vartheta^2-\frac43i\A,
& s_\vartheta^2(\theta su)&=0,\\[0.3em]
s_\vartheta^2(\theta su)^2&=\frac89i\cdot a_\nu^3-\frac4{27}i^2\cdot u_x^2,\\[0.3em]
Z&=2i\cdot a_\eta^2+\frac43i(asu)^2-\frac13J\cdot\theta,\\[0.3em]
g_\nu&=-s_\eta+u_x\left\{2ia_\eta^2-\frac23i(asu)^2+\frac23J\cdot\theta\right\},\\[0.3em]
g_\nu^2&=\frac45ia_\eta^2-\frac25i(asu)^2+\frac13J\cdot\theta,
&g_\nu^3&=0,
&g_\nu^4&=0.
\end{aligned}
\]

\noindent\emph{Consecuencias:}
\[
 s_\nu^3,
\quad s_\vartheta^2(\theta su),
\quad s_\vartheta^2s_\nu,
\quad s_\vartheta^2\theta_\nu
\]
son reductores.



\clearpage
\begin{center}
\textbf{\S\ 18.\quad Sistema $s_I$.}
\end{center}

\[
\text{Plegamiento de}\quad
\left\{\begin{array}{l}
s\text{ con } f;\ \theta;\ K,j,\A;\ f\cdot j,N;\ \A\cdot j,\\
s^3\text{ con }\theta,K.
\end{array}\right.
\]

\noindent\emph{Formas irreducibles.}

\noindent Orden 5:
\[
(asu),\qquad (asu)^2,\qquad (asu)^3,\qquad (asu)^4.
\]

\noindent Orden 6:
\[
\begin{array}{cccc}
(\theta su)&(\theta su)^2&(\theta su)^3&(\theta su)^4\\
s_\vartheta&s_\vartheta(\theta su)&s_\vartheta(\theta su)^2&s_\vartheta(\theta su)^3\\
\cdot&s_\vartheta^2&\cdot&\cdot
\end{array}
\]

\noindent Orden 7:
\[
\begin{array}{c@{\qquad}c@{\qquad}c@{\qquad}c}
\cdot&(Ksu)^2&(Ksu)^3&(Ksu)^4\\
\text{A)}\ s_k&s_k(Ksu)&s_k(Ksu)^2&s_k(Ksu)^3\\
\cdot&s_k^2&\cdot&\cdot
\end{array}
\qquad
\begin{array}{l}
\text{B)}\ s_j,\\[0.4em]
\text{C)}\ (\A su),\ (\A su)^2.
\end{array}
\]

\noindent Orden 8:
\[
\begin{array}{ll}
\text{A)}&s_j(asu),\ s_j(asu)^2,\ s_j(asu)^3,\\[0.4em]
\text{B)}&(Nsu)^2,\quad s_n,\quad s_n(Nsu).
\end{array}
\]

\noindent Orden 10:
\[
\text{A)}\ s_j(\A su),\qquad \text{B)}\ s_\vartheta(\theta s'u)^4.
\]

\noindent Orden 11:
\[
(Ks^2u)^6,\qquad s_k(Ks^2u)^5,
\qquad s_k(Ks^2u)^6.
\]

\noindent\emph{Reducciones:}

Las sucesiones de formas
\[
 s_\vartheta^2(\theta su),\quad s_k^2(Ksu),\quad s_j^2,\quad (\A su)^3,
 \quad (Nsu)^3
\]
tienen como factor el reductor $s_\vartheta^2(\theta su)$; $s_j^2$ y $(\A su)^3$ se obtienen retrocediendo de $j$ y de $\A$ a formas inferiores de la sucesión:
\[
\begin{aligned}
s_\vartheta^2(\theta su)(a\theta u)^3&=-\frac12s_j^2,\\
s_\vartheta^2(\theta su)(a\theta u)^2&=\frac13(\A su)^3,\\
s_\vartheta^2(\theta su)^2(a\theta u)^2&=\frac13(\A su)^4+\frac13s_j^2.
\end{aligned}
\]
Formas $s_\vartheta^2s_{\vartheta'}$ y $s_\vartheta^2s_{\vartheta'}^2$: reductores $s_\vartheta^2(\theta su)$ y $\theta_{\vartheta'}$:
\[
\begin{aligned}
s_\vartheta^2s_{\vartheta'}&=s_\vartheta^2\theta_{\vartheta'}-s_\vartheta^2(\theta s\widehat{\vartheta'}x),\\
s_\vartheta^2s_{\vartheta'}^2&=s_\vartheta^2s_{\vartheta'}\theta_{\vartheta'}-s_\vartheta^2s_{\vartheta'}(\theta s\widehat{\vartheta'}x).
\end{aligned}
\]
Sucesión de formas $s_j(\A su)^2$: reductores $\A_j$ y $(\A su)^3$.

\noindent\emph{Consecuencias:}

$s$ no entra en potencias en plegamiento con $f,j,\A,N$. Las formas $f,j,\A,N$ sólo entran en plegamiento en productos con formas contragredientes; sin embargo, en los productos con $f$ dichas formas pueden seguir siendo cuadráticas en $x$, y en los productos con $\A$ pueden seguir siendo lineales en $x$. Por tanto hay que considerar los siguientes productos de formas:
\[
 f\cdot(0,\lambda),\quad f\cdot(1,\lambda),\quad f\cdot(2,\lambda),\quad
 j\cdot(\mu,0),\quad N,\quad
 \A\cdot(0,\lambda),\quad \A\cdot(1,\lambda)\quad(\lambda>0),
\]
donde $(\mu,\lambda)$ designa una forma de grado $\mu$ en $x$ y de grado $\lambda$ en $u$.

Ni $\theta$ ni $K$ entran en potencias ni en productos con formas arbitrarias en plegamiento con $s$ o con $s^2$. En efecto, para productos con el determinante funcional $K$, $K\cdot\varphi_x$ ($\varphi\ne f$ o $\theta$), vale la relación entre formas descomponibles
\[
 K\cdot\varphi_x=(a\theta u)\varphi_x=f\cdot(\varphi\theta u)-\theta\cdot(\varphi au).
\]
Con ello queda demostrado el esquema anterior para el plegamiento del sistema $s_I$.

\begin{center}
\textbf{\S\ 19.\quad Sistema $s_{II}$.}
\end{center}

Plegamiento de $s$ con $\nu;\ H;\ L;\ H^2$.

\noindent\emph{Formas irreducibles:}
\[
\begin{array}{c@{\qquad}c@{\qquad}c@{\qquad}c}
\text{Orden 6} & \text{Orden 8} & \text{Orden 10} & \text{Orden 12}\\[0.3em]
s_\nu & (Hsu),\ (Hsu)^2 & (Lsu)^2,\ (Lsu)^3 & (H^2su)^3\\
\cdot & s_\eta & s_l & \cdot\\
\cdot & s_\nu^4=J & \cdot & \cdot
\end{array}
\]

\noindent\emph{Reducciones:}

Sucesiones de formas $s_\eta(Hsu)$ y $s_l(Lsu)$: reductor $s_\nu^2$.

Sucesión de formas $s_\sigma$: reductor $s_\nu^2$ procedente de $H$, $s_\nu^2=\dfrac23s_\sigma$.

$(H^2su)^4$ se descompone por la fórmula (7.)a (cf. \S\ 11 A). De ahí se sigue también la descomposición de $(H\cdot L,s,u)^4$.

\noindent\emph{Consecuencias:}

$s^2$ no entra en plegamiento con el sistema II. La forma $\nu$ entra sólo en productos con formas contragredientes. La forma $H$, considerada como covariante, entra en plegamiento en productos con formas $(1,\lambda)$, $(2,\lambda)$, $(3,\lambda)$, con $\lambda$ arbitrario. Las formas $\sigma$ y $(\nu\sigma x)$ no entran en absoluto, y $L$, como determinante funcional, no entra en productos en el plegamiento; las transvecciones completas sobre covariantes del sistema III se vuelven reducibles. De los productos de los sistemas I y II quedan:
\[
 f\cdot\nu,\qquad \A\cdot\nu.
\]

\begin{center}
\textbf{\S\ 20.\quad Formas de séptimo orden (sistema $s_{III}$).}
\end{center}

Plegamiento de $s$ con $f\cdot\nu$, $a_\nu$, $a_\nu^2$.

\noindent\emph{Formas irreducibles:}
\[
\begin{gathered}
s_\nu(asu),\quad a_\nu(asu),\quad s_\nu(asu)^2,\quad a_\nu(asu)^2,
\quad s_\nu(asu)^3,\quad a_\nu(asu)^3,\\
a_\nu s_\nu,
\quad a_\nu s_\nu(asu),
\quad a_\nu^2(asu),
\quad a_\nu^2(asu)^2.
\end{gathered}
\]

\noindent\emph{Reducciones:}

No mencionaremos las reducciones evidentes que surgen por plegamiento directo con los reductores $s_\nu^2$ y $s_\eta(Hsu)$, ni tampoco la descomposición de transvecciones con determinantes funcionales.

Para las formas $a_\nu s_\nu(asu)^2$ y $a_\nu^2s_\nu(asu)^2$, véase la fórmula (19.), \S\ 12; además,
\[
 a_\nu s_\nu(asu)^3=a_\nu(a\widehat{\nu}_1\nu_2u)^3(\nu_1\nu_2\nu)=0.
\]
Sucesión de formas $a_\nu^2s_\nu$: reductor $s_\vartheta^2(\theta su)$ obtenido retrocediendo de $a_\nu^2$ a formas inferiores, es decir, por doble reducción de las expresiones
\[
 s_\vartheta^2(a\theta s)(a\theta u),
 \qquad s_\vartheta^2(a\theta s)(a\theta u)^2
\]
según el planteamiento:
\[
\begin{aligned}
s_\vartheta^2(a\theta s)(a\theta u)
&=[(a\theta u)^2]s^3+\frac15[a_\vartheta(a\theta u)^2]s^2+\frac1{60}[a_\vartheta^2(a\theta u)^2]s\\
&=\frac12a_\nu^2s_\nu-\frac23a_\nu^2s_\nu^2,\\
s_\vartheta^2(a\theta s)(a\theta u)^2
&=\frac45[a_\vartheta(a\theta u)^2]_{us}s^2+\frac{23}{180}[a_\vartheta^2(a\theta u)^2]_{us}s\\
&=-\frac12a_\nu^2s_\nu(asu).
\end{aligned}
\]
Aquí vale $a_\nu^2b_\nu(abu)=0$.

Se obtienen las fórmulas:
\[
\tag{25}
\begin{array}{rcl@{\qquad}rcl}
a_\nu^2s_\nu&=&\dfrac49i\,a_\nu^2b_\nu,
& a_\nu s_\nu(asu)^2&=&\dfrac13i\theta_\nu^2-\dfrac1{18}iH,\\[0.7em]
a_\nu s_\nu(asu)^3&=&0,
& a_\nu^2s_\nu(asu)&=&0,\\[0.7em]
a_\nu^2s_\nu(asu)^2&=&0.
\end{array}
\]

\noindent\emph{Consecuencias:}

1) $a_\nu s_\nu(asu)^2$ y $a_\nu^2s_\nu$ son reductores.

2) Se obtienen los valores de las formas $(\A su)^3$, $(\A su)^4$, ya reconocidas como reducibles en \S\ 19; igualmente para la correspondiente sucesión de formas $(agu)^2$, que al plegarse con $\nu$ da lugar a una doble reducción mediante el reductor $g_\nu$:
\[
\tag{26}
\begin{aligned}
\text{a)}\quad (\A su)^3&=-\frac65a_\nu^2(asu)+3a_\nu s_\nu(asu)+2i\theta_\nu(\theta\nu x),\\
\text{b)}\quad (\A su)^4&=-\frac25a_\nu^2(asu)^2+\frac43i\theta_\nu^2+\frac19iH.
\end{aligned}
\]
De las fórmulas (24.) y (26.), módulo $(i,J)$ (cf. p. 71), se sigue
\[
\tag{27}
\begin{aligned}
\text{a)}\quad (agu)^2&=\frac23(\A su)^2-\frac43a_\nu s_\nu+\frac35s\cdot a_\nu^2,\\
\text{b)}\quad (agu)^3&=2a_\nu s_\nu(asu)-a_\nu^2(asu)^2,\\
\text{c)}\quad (agu)^4&=-a_\nu^2(asu)^2.
\end{aligned}
\]

\begin{center}
\textbf{\S\ 21.\quad Formas de octavo orden (sistema $s_{III}$).}
\end{center}

Plegamiento de $s$ con las formas de cuarto orden (sistema III) (\S\ 9).

\noindent\emph{Formas irreducibles:}
\[
\begin{array}{lll}
(\vartheta\nu x)s_\nu,& ((\vartheta\nu x),s,u)^2,\\[-0.1em]
((\vartheta\nu x)^2,s,u),& ((\vartheta\nu x)^2,s,u)^2,\ \theta s_\nu,\\[-0.1em]
(\vartheta\nu x)^2s_\nu,& ((\vartheta\nu x)^2,s,u)s_\nu,
\end{array}
\]
\[
\begin{array}{lll}
((\vartheta\nu x),s,u)^3,\ (\theta_\nu su)^2,&
((\vartheta\nu x),s,u)^4,\ (\theta_\nu su)^3,\\[-0.1em]
((\vartheta\nu x)^2,s,u)^3,\ (\theta_\nu(\vartheta\nu x),s,u)^2,\ (\theta_\nu^2su),&
((\vartheta\nu x),s,u)^3,\ (\theta_\nu^2su)^2,\ (\theta_\nu(\vartheta\nu x),s,u)^4.
\end{array}
\]

\noindent\emph{Reducciones:}

$((\vartheta\nu x),s,u)s_\nu$ se descompone como transvección sobre determinantes funcionales, por \S\ 3.

Sucesión de formas $((\vartheta\nu x),s,u)^2s_\nu$: reductores $s_\nu^2$ y $s_\vartheta^2\theta_\nu$:
\[
(\widehat\vartheta\nu su)s_\nu(\theta su)=s_\nu s_\vartheta(\theta su)=-\frac12(\widehat\vartheta\nu su)^2(\theta su),
\]
\[
\theta_\nu(\widehat\vartheta\nu su)s_\nu=\theta_\nu s_\nu s_\vartheta=-\frac12\theta_\nu(\widehat\vartheta\nu su)^2.
\]
Sucesión de formas $\theta_\nu s_\nu(\theta su)$: reductor $a_\nu s_\nu(asu)^2$ por doble reducción de las expresiones
\[
a_\nu s_\nu(asu)^2(abu)\quad\text{y}\quad a_\nu s_\nu(asu)^2(abu)(sbu);
\]
$\theta_\nu s_\nu(\theta su)^3$ por doble reducción de $(agu)^3(abu)(bgu)^3$.

Sucesión de formas $\theta_\nu(\vartheta\nu x)s_\nu$: reductor $a_\nu^2s_\nu$ a partir de $\theta_\nu(\vartheta\nu x)s_\nu=a_\nu^2(abu)s_\nu$.

Sucesión de formas $(\theta_\nu(\vartheta\nu x)^2,s,u)$: doble reducción de las formas de la sucesión $\theta_\nu(\widehat\vartheta\nu su)s_\nu$, que al mismo tiempo pertenecen a las sucesiones $((\vartheta\nu x)su)^2s_\nu$ y $\theta_\nu(\vartheta\nu x)s_\nu$.

La forma $(\theta_\nu(\vartheta\nu x),s,u)$ se descompone como transvección simple sobre el determinante funcional $\theta_\nu(\vartheta\nu x)=a_\nu^2(abu)$.

La forma $((\vartheta\nu x)^2,s,u)^4$: doble reducción de la expresión $(a\A u)^2(\A su)^4$, o también de $(\theta gu)^4$, a partir de las fórmulas (27.) y de manera análoga a la reducción de $(j\eta x)^4$ (\S\ 12 C)).

La doble reducción da las fórmulas:
\[
\tag{28}
\begin{aligned}
0&=\theta_\nu s_\nu(\theta su)-\frac16(\widehat\vartheta\nu su)^2(\theta su)-\frac1{12}(Hsu),\\
0&=\theta_\nu s(\theta su)^2-\frac14(\widehat\vartheta\nu su)^2(\theta su)^2-\frac1{24}(Hsu)^2,\\
0&=(\widehat\vartheta\nu su)^2(\theta su)^2+2\theta_\nu(\widehat\vartheta\nu su)(\theta su)^2
  +3\theta_\nu^2(\theta su)^2-\frac13H(su)^2,\\
0&=\theta_\nu s_\nu(\theta su)^3+\frac12\theta_\nu(\widehat\vartheta\nu su)(\theta su)^3,\\
0&=\theta_\nu s_\nu s_\vartheta-\frac14s_\eta,
\quad 0=\theta_\nu s_\nu s_\vartheta(\theta su),
\quad 0=\theta_\nu s_\nu s_\vartheta(\theta su)^2.
\end{aligned}
\]
El último grupo corresponde a $((\theta_\nu(\vartheta\nu x)^2,s,u)^2,\ldots)$.

\noindent\emph{Consecuencias:}

1) $\theta(\vartheta\nu x)s_\nu$, $(\widehat\vartheta\nu su)(\theta su)s_\nu$, $\theta_\nu(\widehat\vartheta\nu su)^2$, $(\widehat\vartheta\nu su)^2(\theta su)^2$ son reductores.

2) Las formas de cuarto orden $(5,\lambda)$, $(6,\lambda)$, etc., no entran en plegamiento con $s^2$.

3) Para la sucesión de formas $(\theta gu)^2$, las fórmulas (27.), teniendo en cuenta las reducciones de este parágrafo, dan las fórmulas siguientes:
\[
\tag{29}
\begin{aligned}
\text{a)}\quad (\theta gu)^2&=\frac43\theta_\nu s_\nu+\frac23s_\nu s_\vartheta-\frac13s\theta_\nu^2-\frac19sH-\frac13f\,a_\nu^2(asu)^2+\frac13a_\nu^2(asu)^2,\\
\text{b)}\quad (\theta gu)^3&=-\frac13(\widehat\vartheta\nu su)^2(\theta su)-\theta_\nu(\widehat\vartheta\nu su)(\theta su)-\theta_\nu^2(\theta su),\\
\text{c)}\quad (\theta gu)^4&=2\theta_\nu^2(\theta su)^2-\frac13(Hsu)^2,\\
g\vartheta(\theta gu)&=(\widehat\vartheta\nu su)(\vartheta\nu x)s_\nu-\theta_\nu(\vartheta\nu x)(\widehat\vartheta\nu su),\\
g\vartheta(\theta gu)^2&=\frac23s\theta_\nu^3,
\quad g\vartheta(\theta gu)^3=\frac13\theta_\nu^3(\theta su),
\quad g\vartheta(\theta gu)^4=0.
\end{aligned}
\]

\begin{center}
\textbf{\S\ 22.\quad Formas de noveno orden (sistema $s_{III}$).}
\end{center}

Plegamiento de $s$ con las formas de quinto orden (sistema III) y con el producto $\A\cdot\nu$ (\S\ 10).

\noindent\emph{Formas irreducibles:}
\[
\begin{array}{ll}
\text{A)}& (k\nu x)^2s_\nu,
((k\nu x)^3,s,u),
(k\nu x)^3s_\nu,
((k\nu x)^2,s,u)^2,
K_\nu s_\nu,\\
& ((k\nu x)^3,s,u)^2,
((k\nu x)^2,s,u)s_\nu,
(K_\nu(k\nu x)^3,s,u),\\
& (K_\nu su)^2,\ (K_\nu su)^3,\ (K_\nu su)^4,
((k\nu x)^2,s,u)^3,
(K_\nu^2su)^4,\\[0.3em]
\text{B)}& (j\nu x)s_\nu,
((j\nu x)^2,s,u),
(j\nu x)^3,s,u,
((j\nu x)^2,s,u)^2,\\[0.3em]
\text{C)}& s_\nu(\A su),\quad (\A_\nu su),\\[0.3em]
\text{D)}& (aHu)s_\eta,
(a_\eta su),
((aHu),s,u)^2,
((aHu)^2,s,u),\\
& (aHu)^2s_\eta,
(a_\eta su)^2,
(a_\eta^2su),
((aHu),s,u)^3,
((aHu)^2,s,u)^2,\\
& ((aHu),s,u)^4,
(a_\eta su)^3,
(a_\eta(aHu),s,u)^2,
(a_\eta(aHu)^2,s,u),
(a_\eta(aHu),s,u)^3.
\end{array}
\]

\noindent\emph{Reducciones:}

A) Las reducciones se siguen directamente de las de \S\ 21 por plegamiento simple. La sucesión de formas $K_\nu s_\nu(Ksu)^2$ tiene como factores los reductores $a_\nu s_\nu(asu)^2$ y $\theta_\nu s_\nu(\theta su)^2$. De ahí que las formas superiores $K_\nu^2(Ksu)^2$, $K_\nu^2(Ksu)^3$ se descompongan según el planteamiento:
\[
0=a_\nu s_\nu(asu)^2(a\theta u)=K_\nu s_\nu(Ksu)^2-\theta_\nu s_\nu(\theta su)^2(a\theta u).
\]

B) Sucesión de formas $(j\nu x)^2s_\nu$: reductor $a_\nu^2s_\nu$ a partir de
\[
(j\nu x)^2s_\nu=3a_\nu^2s_\nu(a\theta u)^2.
\]
$((j\nu x)^3,s,u)^2$: reductor $a_\nu^2s_\nu$ a partir de
\[
((j\nu x)^3,s,u)^2=-6a_\nu^2s_\nu(a\theta u)^2(\theta su).
\]

C) Formas $\A_\nu(\A su)^2$ y $s_\nu(\A su)^2$: reductor $g_\nu$ a partir de la fórmula (27.)a.

Forma $\A_\nu s_\nu$:
1) reductor $a_\nu^2s_\nu$ a partir de $a_g a_\nu^2s_\nu=\frac23\A_\nu s_\nu$;
2) se descompone por la fórmula (27.)a mediante doble reducción de la expresión $(\A s\widehat\nu x)^2$.

De ahí se descompone la forma superior $a_\eta s_\eta$.

D) Sucesión de formas $(aHu)^2s_\eta(asu)$: reductor $a_\nu s_\nu(asu)^2$ por doble reducción de la sucesión $[a_\nu s_\nu(asu)^2]_{\nu_1x^2}$; reducción de $(aHu)^2s_\eta(asu)^2$ a partir de $a_\nu s_\nu(asu)^2$ y $a_\nu s_\nu(asu)^3$. De ahí la reducción de $(a_\eta,s,u)^4$.

Sucesión de formas $a_\eta(aHu)s_\eta$: reductores $a_\nu^2s_\nu$, $a_\eta^2s_\eta$ por doble reducción de la expresión $(\A s\widehat\nu x)^3$, puesto que $a_\eta^2s_\eta$ no surge por plegamiento del término reducible $a_\nu^2s_\nu(\nu\nu_1x)$ de la forma $a_\eta(aHu)s_\eta$.

Sucesión de formas $(a_\eta^2su)^2$: reductor $a_\nu^2s_\nu$ a partir de
\[
a_\nu^2s_\nu s_{\nu_1}=-\frac12(a_\eta^2(Hsu))^2.
\]
La forma $(a_\eta(aHu),s,u)$ se descompone como transvección sobre un determinante funcional.

Las fórmulas de reducción obtenidas son:
\[
\tag{30}
\begin{aligned}
\text{a)}\quad 0&=(aHu)^2(asu)s_\eta-a_\eta(asu)(Hsu)^2-a_\eta(aHu)(asu)(Hsu),\\
\text{b)}\quad 0&=(aHu)^2(asu)^2s_\eta-a_\eta(aHu)(asu)^2(Hsu),\\
\text{c)}\quad 0&=a_\eta(asu)^2(Hsu)^2,\\
0&=a_\eta s_\eta+\frac14a_\nu^2g-\frac14s\cdot a_\eta^2,\\
0&=a_\eta(aHu)s_\eta-\frac12a_\eta^2(Hsu),
\quad 0=a_\eta^2(Hsu)^2,\ldots,
\quad 0=a_\eta^2s_\eta.
\end{aligned}
\]

\noindent\emph{Consecuencias:}

1) $(j\nu x)^2s_\nu$, $\A_\nu s_\nu$, $\A_\nu(\A su)^2$, y por consiguiente $N_\nu(Nsu)^2$, $(aHu)^2(asu)s_\eta$, $a_\eta(aHu)s_\eta$, $a_\eta^2(Hsu)^2$, son reductores; $a_\eta s_\eta$ es una forma descomponible que pasa a formas descomponibles bajo todos los plegamientos simples y dobles.

2) Las formas de quinto orden $(5,\lambda)$, $(6,\lambda)$, etc., no entran en plegamiento con $s^2$.

\begin{center}\textbf{\S\ 23.\quad Formas de décimo orden (sistema $s_{III}$).}\end{center}
Plegamiento de $s$ con las formas de sexto orden (sistema III) (\S\ 11).

\noindent\emph{Formas irreducibles:}
\[
\begin{array}{rl}
\text{A)}& (\vartheta^2\nu x)^3s_\nu,
((\vartheta^2\nu x)^3,s,u),
((\vartheta^2\nu x)^3,s,u)^2,
(\vartheta^2\nu x)^3s_\nu s_\vartheta,
\theta_\nu(\vartheta^2\nu x)^3s_\vartheta,\\[0.3em]
\text{B)}& a_\nu(j\nu x)s_\nu,
(a_\nu(j\nu x),s,u)^2,
(a_\nu(j\nu x),s,u)^3,
(a_\nu(j\nu x),s,u)^4,
(a_\nu^2(j\nu x),s,u)^2,
(a_\nu^2(j\nu x),s,u)^3,\\[0.3em]
\text{C)}& ((\vartheta\eta x)^2,s,u),
((\vartheta\eta x),s,u)^2,
(\theta Hu)s_\eta,
(\theta_\eta su),
((\theta Hu),s,u)^2,
((\theta Hu)^2,s,u),\\
& ((\theta\eta x),s,u)^3,
(\theta_\eta su)^2,
(\theta_\eta^2su),
((\theta Hu),s,u)^3,
((\theta Hu)^2,s,u)^2,
(H_\vartheta(\theta Hu),s,u)^2,\\
& (\theta_\eta(\theta Hu),s,u)^2,
(\theta_\eta(\theta Hu)^2,s,u),
((\theta Hu),s,u)^4,
(\theta_\eta su)^4,
(\theta_\eta(\theta Hu),s,u)^3,
(\theta_\eta^2(\theta Hu),s,u)^2.
\end{array}
\]
\noindent\emph{Reducciones:}
A) Las formas $((\vartheta^2\nu x)^3,s,u)^2$, $((\vartheta^2\nu x)^3,s,u)^3$ se descomponen:
\[
(\vartheta\nu x)(\widehat{\vartheta\nu su})(\widehat{\vartheta\nu su})=(\vartheta\nu x)s_\vartheta s_\nu-(\vartheta\nu x)s_\nu s_\vartheta=-2\theta(\vartheta\nu x)s_\nu s_\vartheta+(\vartheta\nu x)s_\vartheta^2.
\]
Las formas restantes tienen reductores como factores o son plegamientos simples con formas descompuestas.

B) Reductor $a_\nu^2s_\nu$ y $(\widehat{j\nu su})s_\nu$.

C) 1. La sucesión de formas $(\theta Hu)^2s_\eta$: a) reductor $(aHu)^2(asu)s_\eta$; b) doble reducción de las sucesiones $(\theta gu)^3g_\nu$ y $g_\nu(agu)^3(bgu)^3(abu)$. De ahí la reducción de las formas superiores $(\theta,s,u)^3$, $(H_\vartheta(\theta Hu),s,u)^3$ y $(a_\nu b_{\nu_1}^2,s,u)^4$ [esta última forma sustituye a $(H_\vartheta su)^4$].

2. $(\vartheta\eta x,s,u)^4$: reductor $(a_\eta su)^4$.

3. $((\vartheta\eta x)^2,s,u)^3$: reductor $s_\vartheta^2s_\nu$ a partir de $(\widehat{\vartheta\eta su})^2(Hsu)$. Las formas $(\vartheta\eta x)s_\eta$, $(\vartheta\eta x)^2s_\eta$, $((\vartheta\eta x)^2,s,u)^2$, $\theta_\eta s_\eta$ se descomponen por plegamiento con la forma descompuesta $a_\eta s_\eta$.

4. Sucesiones de formas $H_\vartheta(\theta Hu)s_\eta$, $\theta_\eta(\theta Hu)s_\eta$, $H_\vartheta(\vartheta\eta x)s_\eta$, $\theta_\eta(\vartheta\eta x)s_\eta$: reductores $a_\eta(aHu)s_\eta$ y $a_\eta^2s_\eta$.

La sucesión de formas $(\theta_\eta(\vartheta\eta x),s,u)^2$: reductor $a_\eta^2(Hsu)^2$ [excepto $(\theta_\eta^2(\theta Hu),s,u)^2$, que surge del término $a_\eta^2(\theta su)(Hsu)$ por plegamiento].

5. $(H_\vartheta(\vartheta\eta x),s,u)$, $(H_\vartheta(\vartheta\eta x),s,u)^2$: doble reducción de $a_\eta(aHu)s_\eta(abu)$ y $g_\vartheta(\theta gu)g_\nu$.

La sucesión de formas $(H_\vartheta(\vartheta\eta x),s,u)^3$: reductor $((\vartheta\nu x)^2,s,u)^2s_\nu$.

6. $(H_\vartheta(\theta Hu),s,u)$ se descompone por plegamiento simple con la forma descompuesta $(\theta_\nu(\vartheta\nu x),s,u)$ según \S\ 3. 2), c), es decir, por doble reducción de la forma descompuesta inferior $(H_\vartheta su)^2$:\footnote{El signo $\equiv$ significa que se han omitido productos de formas de grado inferior.}
\begin{align*}
0&=\theta_\nu(\vartheta\nu_1)(\theta su)+\frac23\theta_\nu^2(\vartheta\nu_1x)(\theta su)+\theta_\nu(\vartheta\nu x)(\theta su)s_{\nu_1}\\
&\equiv -\frac12H_\vartheta(\theta Hu)(\theta su)+\theta_\eta(\theta su)(Hsu)+\frac12H_\vartheta(\theta su)(Hsu)\\
&\equiv -\frac12H_\vartheta(\theta Hu)(\theta su)+\theta_\eta(\theta su)(Hsu)+\frac12[H_\vartheta]_{sx}^2-\frac12\frac35H_\vartheta(\theta Hu)(\theta su)\\
&\equiv -\frac35H_\vartheta(\theta Hu)(\theta su).
\end{align*}
Por doble reducción, según 1. y 5., se obtienen las fórmulas siguientes:
\begin{equation}
\tag{31}
\begin{aligned}
0&=\theta_\eta(\theta su)(Hsu)^2+\frac14H_\vartheta(\theta Hu)(\theta su)^2\\
&\quad-\frac34\theta_\eta(\theta Hu)(\theta su)(Hsu)-\frac54\theta_\eta(\theta Hu)^2(\theta su)^3\\
&\quad-(asu)^3b_\eta(\theta Hu)^2-a_\nu(asu)^3b_\nu^2,\\
0&=H_\vartheta(\theta Hu)(\theta su)^3-7\theta_\eta(\theta Hu)(\theta su)^2(Hsu),\\
0&=a_\nu(asu)^2b_{\nu_1}^{2}(bsu)^2-\theta_\eta(\theta su)^2(Hsu)^2\\
&\quad+6\theta_\eta(\theta Hu)(\theta su)^2(Hsu),\\
0&\equiv (H_\vartheta(\vartheta\eta x),s,u),
&0&\equiv H_\vartheta(\widehat{\vartheta\eta su})(Hsu)-4\theta_\eta^2(Hsu).
\end{aligned}
\end{equation}
\noindent\emph{Consecuencias:}
\begin{enumerate}
\item $(\theta Hu)^2s_\eta$, $H_\vartheta(\theta Hu)s_\eta$, $\theta_\eta(\theta Hu)s_\eta$, $H_\vartheta(\vartheta\eta x)s_\eta$, $\theta_\eta(\vartheta\eta x)s_\eta$, $(H_\vartheta(\vartheta\eta x),s,u)$, $(\theta_\eta(\vartheta\eta x),s,u)^2$ son reductores.
\item $s^2$ no entra en plegamiento con las formas $(5,\lambda)$, etc., de orden 6.
\end{enumerate}

\begin{center}\textbf{\S\ 24.\quad Formas de undécimo orden (sistema $s_{III}$).}\end{center}
Plegamiento de $s$ con las formas de séptimo orden (sistema III) (\S\ 12).

\noindent\emph{Formas irreducibles:}
\[
\begin{array}{rl}
\text{A)}& ((k\eta x)^3,s,u),\quad (K_\eta(k\eta x)^3,s,u),\quad (K_\eta su)^2,
\quad ((KHu)^2su)^2,\\
& ((KHu)^2,s,u)^3,
\quad ((KHu)^2,s,u)^4,
\quad (K_\eta(KHu)^2,s,u)^3,\\[0.3em]
\text{B)}& ((j\eta x),s,u)^2,
\quad (H_jsu),
\quad ((j\eta x),s,u)^3,\\[0.3em]
\text{C)}& (\D_\eta su),\\[0.3em]
\text{D)}& (aLu)^2s_\eta,
\quad (a_\eta su)^2,
\quad ((aLu)^2,s,u)^2,
\quad ((aLu)^3,s,u)^2,
\quad ((aLu)^3,s,u)^3,
\quad (a_l(aLu)^2,s,u)^2.
\end{array}
\]
\noindent\emph{Reducciones:}

A) Reducción o descomposición por plegamiento simple con las formas correspondientes en los símbolos $\theta-H$.

$(K_\eta(KHu)^2,s,u)$ y $(K_\eta(KHu)^2,s,u)^2$: descomposición según \S\ 3. 2), a), por plegamiento con $K_\nu^2(Ksu)$, $K_\nu^2(Ksu)^2$.

$(K_\eta(KHu),s,u)^4\equiv(a_\nu^2\theta_{\nu_1},s,u)^4$: descomposición según \S\ 3. 2), b), a partir de la forma descompuesta $K_\nu^2(Ksu)^3$.

B) Sucesión de formas $(j\eta x)s_\eta$: reductor $(j\nu x)^2s_\nu$.

Forma $H_j(\widehat{j\eta su})(Hsu)$: reductor $(j\nu x)(\widehat{j\nu su})^2$.

Forma $H_j(j\eta x)(Hsu)$: se descompone reduciendo la forma descompuesta $((j\eta x)^2,s,u)^2$ mediante el reductor $(j\nu x)^2s_\nu$ (fórmula (18.)a):
\[
0=-2(j\nu x)^2s_\nu s_{\nu_1}=[(j\eta x)^2]_{su^2}+H_j(j\eta x)(Hsu).
\]
C) Reductores $\D_\nu s_\nu$ y $\D_\nu(\D su)^2$.

D) Reducción y descomposición por plegamiento simple con las formas correspondientes en los símbolos $f-H$.

E) De (30.)c se obtiene la reducción de la sucesión de formas descompuesta $(a_\sigma su)^2$:
\begin{align*}
0&=a_\eta(Hsu)^2(as\widehat{\nu x})^2
=a_\eta(Hsu)^2(a\widehat{\nu\eta}u)^2+2a_\eta(a\widehat{\nu\eta}u)(\widehat{\eta\sigma u})(Hsu)^2
=\frac13a_\sigma(asu)^2,\\
0&=a_\eta a_\nu(Hsu)^2(as\widehat{\nu x})(asu)
=a_\eta(a\widehat{\nu\eta}u)^2(Hsu)^2(asu)+a_\eta(a\widehat{\nu\eta}u)(\widehat{\eta\sigma u})(Hsu)^2(asu)
=\frac13a_\sigma(asu)^3.
\end{align*}
\noindent\emph{Consecuencia:} $s^2$ no entra en plegamiento con las formas de orden 7.

\begin{center}\textbf{\S\ 25.\quad Formas de los órdenes 12, 13, 14 y 15 (sistema $s_{III}$).}\end{center}
Plegamiento de $s$ con las formas de los órdenes 8, 9, 10, 11 (sistema III) (\S\S\ 13--16).

\noindent\emph{Formas irreducibles:}

\noindent Orden 12.
\[
\begin{array}{rl}
\text{A)}& ((\vartheta^2\eta x)^4,s,u),\quad \theta_\eta(\vartheta^2\eta x)^3s_\vartheta,\\
\text{B)}& ((aHu)H_j,s,u)^2,
\quad ((aH\cdot u)H_j,s,u)^3,\\
\text{C)}& (\theta_\nu su)^2,
\quad ((\theta Lu)^2,s,u)^2,
\quad ((\theta Lu)^3,s,u),
\quad ((\theta Lu)^2,s,u)^3,
\quad (\theta_l(\theta Lu)^2,s,u)^2.
\end{array}
\]
\noindent\emph{Reducciones:} Ya no se mencionarán las reducciones que surgen por plegamiento directo con reductores o con formas descompuestas.

B) Descomposición de $((aHu)H_j,s,u)$ y $(a_\eta(aHu)H_j,s,u)$ como transvecciones sobre determinantes funcionales.

Descomposición de $(a_\eta(aHu)H_j,s,u)^2$: doble reducción de la forma descompuesta $[\theta_\nu\theta_{\nu_1}^3]_{su^3}$ (\S\ 13. C):
\begin{align*}
0&\equiv[\theta_{\nu_1}^3\theta_\nu]_{su^3}
\equiv-\frac34\theta_\nu(\theta su)^2(\theta\theta'u)\theta_{\nu_1}^3\\
&\equiv-\frac98(\theta\theta'\eta su)(\theta\theta'u)(\theta Hu)(\theta'Hu)\theta_\nu'(\theta su)
\equiv-\frac38a_\eta(aHu)H_j(asu)^2.
\end{align*}

C) De las fórmulas (31.) se sigue la reducción de las formas descompuestas
\[
 [L_\vartheta(\theta Lu)]_{su^3}\quad\hbox{y}\quad [\theta_l(\theta Lu)]_{s^2u^3},
\]
es decir, de los productos $[\theta_\nu^2H]_{su^3}$ y $[a_\nu^2(bHu)^2]_{su^3}$:
\begin{equation}
\tag{32}
\begin{aligned}
0&\equiv\theta_\nu^2(\theta su)^2(Hsu),\\
0&\equiv [L_\vartheta(\theta Lu)]_{su^4}-7[\theta_l(\theta Lu)]_{su^4},\\
0&\equiv a_\nu^2(asu)^2(bHu)^2(bsu)\\
&\quad+6\theta_l(\theta Lu)^2(\theta su)(Lsu),\\
0&\equiv\theta_\nu^3(\theta su)(Hsu)^2\\
&\quad\hbox{de }0\equiv\theta_l^2(\theta Lu)(\theta su)(Lsu)^2\\
&\quad(\hbox{reductor }a_\eta^2(Hsu)^2).
\end{aligned}
\end{equation}

\noindent Orden 13.
\[
\text{A)}\ ((KLu)^3,s,u)^2,\ ((KLu)^3,s,u)^3,\qquad
\text{B)}\ (L_jsu)^2,
\qquad
\text{C)}\ ((aH^2u)^3,s,u)^2.
\]
\noindent Orden 14.
\[
\text{A)}\ ((aLu)^2L_j,s,u)^2,
\qquad
\text{B)}\ ((\theta H^2u)^3,s,u)^2.
\]
\noindent Orden 15.
\[
 ((KH^2u)^4,s,u)^2.
\]
\noindent\emph{Reducciones:} Descomposición de $(K_l(KLu)^3,s,u)^2$, $(H_\vartheta(\theta H^2u)^3,s,u)$, $((K,H\cdot L,u)^5,s,u)$ según \S\ 3. 2), c); es decir, considerando simultáneamente las formas descompuestas $(K_l(KLu)^2,s,u)^3$, $(H_\vartheta(\theta H'u)^2,s,u)^2$, $((K,H\cdot L,u)^4,s,u)^2$.

\noindent\emph{Consecuencia:} $s^2$ no entra en plegamiento con formas individuales del sistema III.

\begin{center}\textbf{\S\ 26.\quad Sistema $s_{IV}$.}\end{center}
\noindent 1) \emph{Plegamiento de $s$ con los sistemas I--III.}

\noindent\emph{Reducciones:} Según las reducciones de los últimos parágrafos ($(a_\nu^2s_\nu(aHu)^2(asu)s_\eta$, etc.), las formas $(0,\lambda)$, $(1,\lambda)$, $(2,\lambda)$ no admiten simultáneamente plegamientos I y III; por tanto, según \S\ 18, las formas $f,\D,N$ no entran en productos en el plegamiento que forma el sistema $s_{IV}$.

La forma $j$ no entra en plegamiento, pues, por las mismas reducciones, los plegamientos contragredientes a $j$,
\[
(K_\nu(k\nu x)^3,s,u),\qquad ((\vartheta^2\eta x)^4,s,u)\quad\hbox{etc.},
\]
corresponden a los plegamientos cogredientes a $j$:
\[
K_\nu(k\nu x)^2s_k,
\qquad (\vartheta^2\eta x)^3s_\vartheta\quad\hbox{etc.}
\]

\noindent 2) \emph{Plegamiento de $s$ con los sistemas II--III}, es decir, de $s$ con $H\cdot(1,\lambda)$, $H\cdot(2,\lambda)$, $H\cdot(3,\lambda)$.

\noindent\emph{Formas irreducibles:}
\[
\begin{array}{rl}
\text{A)}&(a_\nu H,s,u)^4,
((aHu)^2H,s,u)^3,
((aHu)^2H,s,u)^4,\\
&((aLu)^2H,s,u)^4,
((aLu)^3H,s,u)^3,
((aH^2u)^3H,s,u)^4,\\[0.3em]
\text{B)}&(a_\nu^2H,s,u)^3,
(a_\eta(aHu)^2H,s,u)^3,
(a_l(aLu)^2H,s,u)^4,\\[0.3em]
\text{C)}&(\theta_\nu H,s,u)^4,
((\theta Hu)^2H,s,u)^3,
((\theta Hu)^2H,s,u)^4,\\
&((\theta Lu)^2H,s,u)^4,
((\theta Lu)^3H,s,u)^3,
((\theta H^2u)^3H,s,u)^4,\\[0.3em]
\text{D)}&(\theta_\eta(\theta Hu)^2H,s,u)^3,
(\theta_l(\theta Lu)^2H,s,u)^4,\\[0.3em]
\text{E)}&((KLu)^3H,s,u)^4,
((KH^2u)^4H,s,u)^4,\\[0.3em]
\text{F)}&((j\nu x)^2H,s,u)^3,
((j\nu x)^2H,s,u)^4,
((j\eta x)H,s,u)^4,\\
&(H_jH,s,u)^3,
(L_jH,s,u)^4.
\end{array}
\]
\noindent\emph{Reducciones:} $\nu$ no entra en plegamiento, de modo análogo a $j$.

B) y D). $(a_\nu^2H,s,u)^4$ y $(\theta_\nu^2H,s,u)^4$ se descomponen por las fórmulas (27.)c) y (29.)c), teniendo en cuenta $(gHu)^2=-\frac13s\sigma$:
\[
(a_\nu^2H,s,u)^4=\frac13(asu)^4\sigma,
\qquad
(\theta_\nu^2H,s,u)^4=-\frac16(\theta su)^4\sigma;
\]
de ahí la descomposición de $(a_\eta(aHu)H,s,u)^4$, $(\theta_\eta(\theta Hu)H,s,u)^4$.

$(\theta_\nu^2H,s,u)^3$, $(\theta_\nu^3H,s,u)^3$ y, por tanto, $(\theta_\eta^2(\theta Hu)H,s,u)^4$, se descomponen según \S\ 25, formas de orden 12, C).

\noindent 3) \emph{Plegamiento de $s$ con el sistema III--III}, es decir, de $s$ con formas del sistema III:
\[
(3,\lambda)(3,\lambda),\quad (3,\lambda)(2,\lambda),\quad (3,\lambda)(1,\lambda),\quad
(2,\lambda)(2,\lambda),\quad (2,\lambda)(1,\lambda),\quad (1,\lambda)(1,\lambda).
\]
\noindent\emph{Formas irreducibles:}
\[
\begin{array}{rl}
\text{A)}&(a_\nu^2a_\nu^2,s,u)^3,
(a_\nu^2a_\nu^2,s,u)^4,
(a^2a_\eta(aHu)^2,s,u)^3,\\[0.3em]
\text{B)}&(a_\nu H_j,s,u)^4,
((aHu)^2H_j,s,u)^3,
((aLu)^2H_j,s,u)^4,
((aLu)^3H_j,s,u)^2,
((aH^2u)^3H_j,s,u)^3.
\end{array}
\]
\noindent\emph{Reducciones:}

Los productos II--III, para el mismo orden en los símbolos $\nu$ y $f$, deben considerarse superiores a los productos III--III.

Las reducciones surgen en parte de relaciones entre productos (sizigias) y en parte de sustituir los productos por las formas descompuestas correspondientes y reducir el plegamiento de estas formas con $s$. Los productos que contienen contravariantes siempre se omiten, puesto que pasan a productos.

A) $f$ cuadrático, símbolos $\nu$ de orden arbitrario. Por \S\ 11, fórmulas (12.), y por cálculo análogo, se tiene:
\begin{align*}
1)\quad 0&=a_\nu b_{\nu_1}+\frac12f(aHu)^2-\frac14\theta H+\frac12u_xH_\vartheta-\frac14u_x^2H_\vartheta^2,\\
2)\quad 0&=a_\nu b_{\nu_1}^2+f a_\eta(aHu)^2-\theta_\eta-\frac12H_\vartheta,\\
3)\quad 0&=a_\nu^2b_{\nu_1}^2-H_\vartheta^2+2\theta_\eta^2.
\end{align*}
De ahí se sigue:
\begin{align*}
4)\quad 0&=a_\nu^2b_{\nu_1}^2(j\nu_1x),\\
5)\quad L_\vartheta(\theta Lu)&=-a_\nu b_\eta(bHu)^2+\frac34a_\nu^2(bHu)^2+\frac34\theta_\nu^2H,\\
6)\quad L_\vartheta(\theta Lu)&=-6a_\nu b_\eta(bHu)^2+2a_\nu^2(bHu)^2-\frac12\theta_\nu^2H,\\
7)\quad 0&=a_\nu(bHu)^2+f(aLu)^3+\theta_\nu H,\\
8)\quad 0&=a_\nu(bLu)^3+\frac34(\theta Hu)^2H,\\
9)\quad 0&=(aHu)^2(bHu)^2+2(\theta Hu)^2H,\\
10)\quad 0&=a_\eta(aHu)^2b_\eta(bHu)^2,\\
11)\quad 0&\equiv [a_\nu^2b_\eta(bHu)]_{\widehat{su^4}}.
\end{align*}
Las reducciones de
\[
(H_\vartheta su)^4,
\qquad (L_\vartheta(\theta Lu),s,u)^3,
\qquad (L_\vartheta(\theta Lu),s,u)^4
\]
(fórmulas (31.) y (32.)) dan, junto con las fórmulas 1) a 11), la reducción de todos los plegamientos de $s$ con productos que son cuadráticos en los símbolos $f$ y arbitrarios en $\nu$.

B) Productos de dos de los símbolos $f,\theta,j$; $\nu$ de orden arbitrario. Por las fórmulas (17.), (18.), (20.) y por cálculo directo valen las relaciones:
\begin{align*}
1)\quad 0&=a_\nu\theta_{\nu_1}+\frac14\theta(aHu)^2+\frac14f(\theta Hu)^2,\\
2)\quad 0&=\theta_\nu\theta_{\nu_1}+\frac12\theta(\theta Hu)^2.
\end{align*}
De ahí se sigue:
\begin{align*}
3)\quad 0&=3a_\nu(\theta Hu)^2+\theta(aLu)^3+2f(\theta Lu)^3,\\
4)\quad 0&=3\theta_\nu(aHu)^2+f(\theta Lu)^3+2\theta(aLu)^3,\\
5)\quad 0&=a_\nu(\theta Lu)^3,
&0&=\theta_\nu(aLu)^3,
&0&=(aHu)^2(\theta Hu)^2,\\
6)\quad 0&=a_\nu\theta_\nu^2+\theta_\nu a_\nu^2+f\theta_\eta(\theta Hu)^2+\theta a_\eta(aHu)^2,\\
7)\quad 0&=\theta_\nu\theta_\nu^2+\theta\theta_\eta(\theta Hu)^2+\frac16fH_j,\\
8)\quad 0&=a_\nu(j\nu x)^2+fH_j,\\
9)\quad 0&=(aHu)^2(j\nu x)^2-2a_\nu H_j,\\
10)\quad 0&=\theta_\nu(j\nu x)^2,
&0&=\theta H_j,\\
11)\quad 0&=a_\nu\theta_\nu^3-\frac34(j\eta x)^2,\\
12)\quad 0&=a_\nu^2\theta_\nu^2-\frac13(j\eta x)^2-\frac13(\D Hu)^2,\\
13)\quad 0&=a_\nu^2\theta_\nu^3-\frac12a_\sigma,\\
14)\quad 0&=\theta_\nu\theta_\nu^3,
&0&=a_\eta H_j.
\end{align*}
Las fórmulas se han llevado hasta el punto en que los productos se vuelven polarizables; esto es, por plegamiento sólo aparece un nuevo producto, porque a) el plegamiento del producto consigo mismo se vuelve reducible, y b) los restantes productos que surgen por plegamiento se vuelven reducibles o conducen a contravariantes (cf. \S\ 3. 2), a) y b)).

Las fórmulas 1) a 5) dan la reducción de todos los plegamientos de $s$ con productos que surgen de $a_\nu\theta_\nu$ por plegamiento. Las fórmulas 6) a 10) dan la reducción de los productos que surgen de $a_\nu\theta_\nu^2$ por plegamiento. Según \S\ 24 A), teniendo en cuenta las fórmulas 6) a 10), se descomponen los plegamientos de $s$ con productos que surgen de $a_\nu^2\theta_\nu$.

Se tiene:
\begin{align*}
0&\equiv a_\nu^2(asu)^2\theta_{\nu_1}(\theta su)^2,
&0&\equiv a_\nu^2(asu)\theta_{\nu_1}(\theta su)^3-K_\eta(KHu)^2(Ksu)^3,\\
0&\equiv a_\nu^2(asu)^2\theta_{\nu_1}(\theta su),
&0&\equiv a_\nu^2(asu)\theta_{\nu_1}(\theta su)^2,\\
0&\equiv a_\nu^2(asu)\theta_{\nu_1}(\theta su).
\end{align*}
Los reductores
\[
 ((j\eta x)^2,s,u)^3\ [\hbox{de }(\widehat{\vartheta\eta su})^2(Hsu),\ \S\ 23.\ C3],
\]
\[
(H_j(j\eta x),s,u)^2\ [\S\ 24,\ B],
\quad ((\D Hu)^2,s,u)^2\ [\S\ 24,\ C],
\quad (a_\sigma su)^2\ [\S\ 24,\ E]
\]
dan, junto con las fórmulas 11) a 14), la reducción de todos los productos que surgen de $a_\nu\theta_{\nu_1}^3$, $a_\nu^2\theta_{\nu_1}^2$, $a_\nu^2\theta_{\nu_1}^3$.

Nuestros cálculos anteriores se resumen en la siguiente \emph{conclusión}:

El sistema relativamente completo módulo $(\R,t)$ consta de las siguientes 331 formas y subsistemas, que reproducimos también a continuación ordenados en tablas.
\clearpage
\begin{landscape}
\thispagestyle{plain}
\begin{center}\bfseries Tabla I (véase p. 90)\end{center}
\vspace{-0.6em}
\begingroup
\setlength{\tabcolsep}{2pt}
\renewcommand{\arraystretch}{1.35}
\tiny
\begin{center}
\begin{adjustbox}{max width=0.98\linewidth,max totalheight=0.78\textheight,center}
\begin{tabular}{|c|c|c|c|c|c|c|c|c|c|c|c|c|c|c|c|c|c|c|c|c|c|c|}
\hline
\diagbox{$u$}{$x$} & 0 & 1 & 2 & 3 & 4 & 5 & 6 & 7 & 8 & 9 & 10 & 11 & 12 & 13 & 14 & 15 & 16 & 17 & 18 & 19 & 20 & 21 \\ \hline
0 & \mcell{i} &  &  &  & \mcell{f} &  & \mcell{\A} &  & \mcell{K_{\nu}(k\nu x)^3} &  & \mcell{K_{\eta}(k\eta x)^3} &  &  &  & \mcell{(\vartheta^3\eta x)^4} & \mcell{H_k(k\vartheta,\eta x)^4} &  & \mcell{\theta_l(\vartheta k,lx)^5} &  &  &  & \mcell{(\vartheta^2lx)^6} \\ \hline
1 &  & \mcell{u_x} &  &  &  & \mcell{\theta_\nu(\vartheta\nu x)^2} &  & \mcell{\theta_\eta(\vartheta\eta x)^2} & \mcell{N} & \mcell{(k\nu x)^3} & \mcell{\theta_\nu(\vartheta^2\nu x)^3} & \mcell{(k\eta x)^3} & \mcell{H_{\vartheta}(\vartheta^2\eta x)^3\\ \theta_\eta(\vartheta^2\eta x)^3} &  & \mcell{\theta_l(\vartheta^2lx)^4} &  & \mcell{(\vartheta k,\eta,x)^4} &  & \mcell{(\vartheta k,l,x)^5} &  &  &  \\ \hline
2 &  &  & \mcell{a_\nu^2} &  & \mcell{\theta\\ a_\eta^2} &  & \mcell{(\vartheta\nu x)^2} & \mcell{K_\nu(k\nu x)^2} & \mcell{(\vartheta\eta x)^2} & \mcell{H_k(k\eta x)^2\\K_\eta(k\eta x)^2} &  & \mcell{(\vartheta^2\nu x)^3\\ K_l(klx)^3} &  & \mcell{(\vartheta^2\eta x)^3} &  & \mcell{(\vartheta^2lx)^4} &  &  &  &  &  &  \\ \hline
3 &  & \mcell{\theta_\nu^3} &  & \mcell{a_\nu} & \mcell{\theta_\nu(\vartheta\nu x)} & \mcell{\A_\nu\\a_\eta} & \mcell{K\\H_{\vartheta}(\vartheta\eta x)\\\theta_\eta(\vartheta\eta x)} & \mcell{\A_\eta} & \mcell{(k\nu x)^2\\\theta_l(\vartheta lx)^2} &  & \mcell{(k\eta x)^2} &  & \mcell{(klx)^3} &  &  &  &  &  &  &  &  &  \\ \hline
4 & \mcell{\nu} &  & \mcell{H\\\theta_\nu^2} & \mcell{(j\nu x)^3\\a_\eta(aHu)} & \mcell{\theta_\eta^2} & \mcell{(\vartheta\nu x)\\a_l^2} &  & \mcell{N_\nu\\(\vartheta\eta x)} &  & \mcell{H_\eta\\N_\eta\\(\vartheta lx)^2} &  &  &  &  &  &  &  &  &  &  &  &  \\ \hline
5 &  & \mcell{a_\eta(aHu)^2} & \mcell{\theta_\eta^2(\theta Hu)} & \mcell{\theta_\nu} & \mcell{\frac{K_\nu^2}{(aHu)}} & \mcell{\theta_\eta} & \mcell{K_\eta^2\\(\A Hu)\\a_l} &  & \mcell{\A_l} &  &  &  &  &  &  &  &  &  &  &  &  &  \\ \hline
6 & \mcell{j\\\sigma} &  & \mcell{(j\nu x)^2\\(aHu)^2} & \mcell{L\\a_\nu^2(j\nu x)\\H_{\vartheta}(\theta Hu)\\\theta_\eta(\theta Hu)} & \mcell{\frac{K_\nu}{\theta_l^2}} &  &  & \mcell{K_\eta} &  &  &  &  &  &  &  &  &  &  &  &  &  &  \\ \hline
7 &  & \mcell{\theta_\eta(\theta Hu)^2} & \mcell{H_j(j\eta x)\\a_l(aLu)^2} &  & \mcell{a_\nu(j\nu x)\\(\theta Hu)} &  & \mcell{\A_\nu(j\nu x)\\\theta_l} & \mcell{K_l^2} &  &  &  &  &  &  &  &  &  &  &  &  &  &  \\ \hline
8 & \mcell{H_j^2\\a_l(aLu)^3} & \mcell{(\nu\sigma x)\\(j\nu x)} & \mcell{(\theta Hu)^2} & \mcell{K_\eta(KHu)^2\\(j\eta x)\\(aLu)^2} &  &  &  &  &  &  &  &  &  &  &  &  &  &  &  &  &  &  \\ \hline
9 &  & \mcell{H_j\\(aLu)^3} & \mcell{a_\eta(aHu)H_j\\L_{\vartheta}(\theta Lu)^2\\\theta_l(\theta Lu)^2} &  & \mcell{(KHu)} &  &  &  &  &  &  &  &  &  &  &  &  &  &  &  &  &  \\ \hline
10 & \mcell{\theta_l(\theta Lu)^3} & \mcell{L_j^2\\(j\sigma x)\\a_\eta(aH^2u)^3} &  & \mcell{H_j(aHu)\\(\theta Lu)^2} &  &  &  &  &  &  &  &  &  &  &  &  &  &  &  &  &  &  \\ \hline
11 &  & \mcell{(\theta Lu)^3} & \mcell{K_l(KLu)^3\\L_j\\(aH^3u)^3} &  &  &  &  &  &  &  &  &  &  &  &  &  &  &  &  &  &  &  \\ \hline
12 & \mcell{(aH^2u)^4} & \mcell{a_\eta(aLu)^2L_j\\H_{\vartheta}(\theta H^2u)^3\\\theta_\eta(\theta H^2u)^3} &  & \mcell{(KLu)^3} &  &  &  &  &  &  &  &  &  &  &  &  &  &  &  &  &  &  \\ \hline
13 & \mcell{(\nu\sigma j)} &  & \mcell{(aLu)^2L_j\\(\theta H^3u)^3} &  &  &  &  &  &  &  &  &  &  &  &  &  &  &  &  &  &  &  \\ \hline
14 & \mcell{(\theta H^2u)^4} & \mcell{K_\eta(KH^2u)^4\\(a_\eta H\!\cdot L,u)^4} &  &  &  &  &  &  &  &  &  &  &  &  &  &  &  &  &  &  &  &  \\ \hline
15 & \mcell{L_9(\theta,H\!\cdot L,u)^4} &  & \mcell{(KH^2u)^4} &  &  &  &  &  &  &  &  &  &  &  &  &  &  &  &  &  &  &  \\ \hline
16 & \mcell{(\theta,H\!\cdot L,u)^5} &  &  &  &  &  &  &  &  &  &  &  &  &  &  &  &  &  &  &  &  &  \\ \hline
17 & \mcell{K_\eta(K,H\!\cdot L,u)^5} &  &  &  &  &  &  &  &  &  &  &  &  &  &  &  &  &  &  &  &  &  \\ \hline
18 &  & \mcell{(K,H\!\cdot L,u)^5} &  &  &  &  &  &  &  &  &  &  &  &  &  &  &  &  &  &  &  &  \\ \hline
19 &  &  &  &  &  &  &  &  &  &  &  &  &  &  &  &  &  &  &  &  &  &  \\ \hline
20 &  &  &  &  &  &  &  &  &  &  &  &  &  &  &  &  &  &  &  &  &  &  \\ \hline
21 & \mcell{(KH^3u)^6} &  &  &  &  &  &  &  &  &  &  &  &  &  &  &  &  &  &  &  &  &  \\ \hline
\end{tabular}
\end{adjustbox}
\end{center}
\vspace{-0.2em}
\begin{center}\footnotesize (Los números de la fila horizontal superior indican el grado en las $x$; los de la primera fila vertical, el grado en las $u$.)\end{center}
\endgroup
\end{landscape}
\clearpage
\begin{landscape}
\thispagestyle{plain}
\begin{center}\bfseries Tabla II (véase p. 90)\end{center}
\vspace{-0.6em}
\begingroup
\setlength{\tabcolsep}{2pt}
\renewcommand{\arraystretch}{1.35}
\tiny
\begin{center}
\begin{adjustbox}{max width=0.98\linewidth,max totalheight=0.78\textheight,center}
\begin{tabular}{|c|c|c|c|c|c|c|c|c|c|c|c|c|c|c|c|c|c|}
\hline
\diagbox{$u$}{$x$} & 0 & 1 & 2 & 3 & 4 & 5 & 6 & 7 & 8 & 9 & 10 & 11 & 12 & 13 & 14 & 15 & 16 \\ \hline
0 & \mcell{J} &  &  &  & \mcell{s} &  & \mcell{s_j^3} &  &  &  &  & \mcell{s_\nu} & \mcell{(k\nu x)^3s_\nu} & \mcell{(\vartheta^3\nu x)^2s_{\vartheta}\\(\vartheta^2\nu x)^2s_{\vartheta}} &  & \mcell{\theta_\eta(\vartheta^2\eta x)^2s_\eta} &  \\ \hline
1 &  &  &  &  &  &  & \mcell{(asu)} & \mcell{s_\eta} & \mcell{s_l^2\\(\A su)} & \mcell{s_\eta(Nsu)\\(\vartheta\nu x)^2s,u} & \mcell{((k\nu x)^3,s,u)_s\\(K_\nu(k\nu x)^3,s,u)} &  & \mcell{K_\eta(k\eta x)^3,s,u} &  & \mcell{(\vartheta^2\nu x)^2s_\nu} &  & \mcell{(\vartheta^2\eta x)^4,s,u} \\ \hline
2 &  &  &  &  & \mcell{(asu)^2} & \mcell{s_\theta(\theta su)} & \mcell{(\A su)^2/a_\nu s_\nu} & \mcell{(\vartheta\nu x)^2,s,u\\(\theta_l(\vartheta\nu x)^2,s,u)} &  & \mcell{\theta_\eta(\vartheta^2\eta x)^2,s,u} &  & \mcell{(k\eta x)^2s_\nu\\((k\nu x)^2,s,u)} &  & \mcell{(k\eta x)^3,s,u} &  &  &  \\ \hline
3 &  &  & \mcell{(asu)^3} & \mcell{s_\vartheta(\theta su)^2\\s_\nu} & \mcell{a_\nu s_\nu(asu)\\a_l^2(asu)} & \mcell{s_\eta} & \mcell{(\theta su)\\a_l^2su} &  & \mcell{(Nsu)^2\\(\vartheta\eta x)s_\nu\\((\vartheta\nu x)^2,s,u)} & \mcell{((k\nu x)^2,s,u)^2} & \mcell{(\vartheta^2\eta x)^2,s,u} &  &  &  &  &  &  \\ \hline
4 & \mcell{(asu)^4} & \mcell{s_\vartheta(\theta su)^3} & \mcell{a_\nu^2(asu)^2} & \mcell{(\theta_\nu^2su)} & \mcell{(\theta su)^2} & \mcell{s_k(Ksu)^3\\a_\nu(asu)\\s_\nu(asu)} & \mcell{((\vartheta\nu x)^2,s,u)^2} & \mcell{s_\nu(\A su)\\(\A su)\\(aHu)s_\eta\\(a_\nu su)} &  & \mcell{(\A_\eta su)} &  &  &  &  &  &  &  \\ \hline
5 &  &  & \mcell{(\theta su)^3} & \mcell{s_k(Ksu)^3,s_j\\s_\vartheta(\theta s^\prime u)^4\\s_\nu(asu)^2\\a_\nu(asu)^2} & \mcell{(Hsu)\\((\vartheta\nu x)^2,s,u)\\(\theta_\nu(\vartheta\nu x)^2,s,u)\\(\theta_\nu^2su)} & \mcell{((j\nu x)^2,s,u)\\(aHu)^2s_\nu\\(a_\eta su)^2} & \mcell{(Ksu)^2\\s_\eta\\(\theta_\eta^2su)} &  & \mcell{((k\nu x)^2,s,u)^2\\K_\nu s_\nu} &  &  &  &  &  &  &  &  \\ \hline
6 & \mcell{(\theta su)^4} & \mcell{s_\nu(asu)^3\\a_\nu(asu)^3} & \mcell{(Hsu)^2\\(\theta_\nu(\vartheta\nu x),s,u)^3\\(\theta_\nu^2su)^2} & \mcell{a_\eta s_\nu^2\\(a_\eta(aHu),s,u)^2\\(a_\eta(aH^2u),s,u)} & \mcell{(Ksu)^3} & \mcell{s_\eta(asu)\\((\vartheta\nu x),s,u)^3\\(\theta_\eta su)} & \mcell{((k\nu x)^2,s,u)^3} & \mcell{s_\nu(\A su)\\a_\nu(j\nu x)s_\nu\\(\theta H u)s_\eta\\(\theta_\eta su)} &  &  &  &  &  &  &  &  &  \\ \hline
7 & \mcell{\theta_\nu(\vartheta\nu x),s,u)^4} & \mcell{(a_\eta(aHu),s,u)^3} & \mcell{(Ksu)^4\\(\theta_\eta^2(\theta Hu),s,u)^2\\(a_\eta^2-a_l^2,s,u)^3} & \mcell{s_j(asu)^3\\s_k(K^2u)^3\\((j\nu x),s,u)^2\\(\theta_\eta su)^2} & \mcell{((j\nu x),s,u)\\((j\eta x),s,u)\\(aHu),s,u)\\(aH^2u),s,u)} & \mcell{((\vartheta\eta x),s,u)^3} & \mcell{(aLu)^3s_l/(asu)^3} &  &  &  &  &  &  &  &  &  &  \\ \hline
8 & \mcell{(a^2-a_l^2,s,u)^4} & \mcell{s_j(asu)^3\\s_k(K^2u)^3\\((j\nu x),s,u)^4\\(\theta_\eta su)^3} & \mcell{((j\nu x)^2,s,u)^2\\(aHu),s,u)^3\\((aHu)^2,s,u)^3} & \mcell{(Lsu)^2\\(a_l^2(j\nu x),s,u)^3\\H_\vartheta(\theta Hu),s,u)^3\\\theta_l(\theta Hu),s,u)^3} & \mcell{(asu)^3} & \mcell{(K,su)^3} &  & \mcell{(K_\vartheta su)^2} &  &  &  &  &  &  &  &  &  \\ \hline
9 & \mcell{K_j^2su^4\\((aHu),s,u)^4} & \mcell{(Lsu)^4\\(a_l^2(j\nu x),s,u)^4\\(\theta_\nu(\theta Hu),s,u)^2} & \mcell{(K_\nu^3u)^6\\(a_\eta(aLu)^2,s,u)^3\\(a_l^2H,s,u)^3} & \mcell{(K,su)^3} & \mcell{(a_\nu(j\nu x),s,u)^3\\(\theta Hu),s,u\\(\theta H u)^2,s,u} &  & \mcell{(\theta_lsu)^3} &  &  &  &  &  &  &  &  &  &  \\ \hline
10 &  &  & \mcell{(K,su)^4\\(a_l^2a_\eta(aHu)^2,s,u)^3} & \mcell{((j\eta x),s,u)^3\\(H_j,su)\\((aLu)^2,s,u)^2\\((aLu)^3,s,u)} &  &  &  &  &  &  &  &  &  &  &  &  &  \\ \hline
11 & \mcell{(a_\nu(j\nu x),s,u)^4\\((\theta Hu),s,u)^4} & \mcell{(K_l(KH u)^2,s,u)^3\\((j\eta x),s,u)^3\\((aLu)^2,s,u)^4\\(a_\eta H,s,u)^3} & \mcell{(H^2su)^3\\\theta_l(\theta Lu)^2,s,u)^3} &  & \mcell{((KHu)^2,s,u)^3} &  &  &  &  &  &  &  &  &  &  &  &  \\ \hline
12 &  & \mcell{(a_\eta(aHu)^2H,s,u)^3} & \mcell{((KH u)^3,s,u)^3} & \mcell{((aHu)H_j,s,u)^2\\((\theta Lu)^2,s,u)^3} &  &  &  &  &  &  &  &  &  &  &  &  &  \\ \hline
13 & \mcell{((KH u)^3,s,u)^4} & \mcell{((aHu)H_j,s,u)^3\\((\theta Lu)^2,s,u)^4\\(\nu\cdot H,s,u)^3} & \mcell{(L_jsu)^2\\((aH^3u)^3,s,u)^3\\((j\eta x)H,s,u)^3\\((aHu)^2H,s,u)^3} &  &  &  &  &  &  &  &  &  &  &  &  &  &  \\ \hline
14 & \mcell{((j\nu x)^2H,s,u)^4\\((aHu)^2H,s,u)^4} & \mcell{(\theta_\eta(\theta Hu)^2H,s,u)^3} &  & \mcell{((KLu)^3,s,u)^2} &  &  &  &  &  &  &  &  &  &  &  &  &  \\ \hline
15 & \mcell{(a_l(aLu)^3H,s,u)^4} & \mcell{((KLu)^3,s,u)^3} & \mcell{((aLu)^2L_j,s,u)^2\\((\theta Hu)^3H,s,u)^3\\((\theta Hu)^2H,s,u)^2} &  &  &  &  &  &  &  &  &  &  &  &  &  &  \\ \hline
16 & \mcell{((\theta Hu)^2H,s,u)^4\\(\sigma\cdot H_j,s,u)^4} & \mcell{((j\eta x)H,s,u)^4\\(H_jH,s,u)^3\\((aLu)^2H,s,u)^4\\((aLu)^3H,s,u)^3} &  &  &  &  &  &  &  &  &  &  &  &  &  &  &  \\ \hline
17 & \mcell{(\theta(\theta Lu)^2H,s,u)^4} &  & \mcell{((KH^3u)^3,s,u)^3} &  &  &  &  &  &  &  &  &  &  &  &  &  &  \\ \hline
18 &  & \mcell{((\theta Lu)^2H,s,u)^4\\((\theta Lu)^3H,s,u)^3\\(aHu)^2H,s,u)^4} &  &  &  &  &  &  &  &  &  &  &  &  &  &  &  \\ \hline
19 & \mcell{((aH^3u)^3H,s,u)^4\\(L_jH,s,u)^5} &  &  &  &  &  &  &  &  &  &  &  &  &  &  &  &  \\ \hline
20 &  & \mcell{((KLu)^3H,s,u)^5} & \mcell{((aLu)^3H,s,u)^5} &  &  &  &  &  &  &  &  &  &  &  &  &  &  \\ \hline
21 & \mcell{((\theta H^3u)^3H,s,u)^4\\((aLu)^2H_j,s,u)^4} &  &  &  &  &  &  &  &  &  &  &  &  &  &  &  &  \\ \hline
22 &  &  &  &  &  &  &  &  &  &  &  &  &  &  &  &  &  \\ \hline
23 & \mcell{((KH^3u)^4H,s,u)^4} & \mcell{((aH^3u)^4H,s,u)^3} &  &  &  &  &  &  &  &  &  &  &  &  &  &  &  \\ \hline
\end{tabular}
\end{adjustbox}
\end{center}
\vspace{-0.2em}
\begin{center}\footnotesize (Los números de la fila horizontal superior indican el grado en las $x$; los de la primera fila vertical, el grado en las $u$.)\end{center}
\endgroup
\end{landscape}




\clearpage

\providecommand{\pair}[2]{(#1\,|\,#2)}
\providecommand{\eps}{\varepsilon}
\providecommand{\rhoidx}{\varrho}

\begin{center}
{\Large\bfseries 3. Sobre la teoría de invariantes de las formas en $n$ variables\footnote{Conferencia pronunciada en la reunión de naturalistas de Salzburgo, 1909.}}\par
\vspace{1.2em}
Jahresbericht der DMV 19 (1910), pp. 101--104
\end{center}

\vspace{1em}
En la teoría proyectiva de invariantes de las formas en $n$ variables está resuelto el problema principal: se ha demostrado la finitud del sistema de formas. Sin embargo, no se ha tratado una serie de cuestiones ulteriores, a saber, las que se refieren a métodos para investigar la conexión entre las formas. Quisiera comunicar brevemente los resultados que he obtenido respecto de tales cuestiones; pero, por claridad, esbozaré primero los planteamientos en el dominio ternario, donde son conocidos.

Pienso ante todo en los dos teoremas fundamentales del método simbólico: el teorema de que todas las formaciones invariantes pueden representarse simbólicamente, y el segundo, de que todas las relaciones existentes entre invariantes se obtienen por aplicación sucesiva de un pequeño número de identidades simbólicas; es decir, que el método simbólico da todas las relaciones sin salir del dominio de los invariantes.\footnote{Study: \emph{Methoden zur Theorie der ternären Formen}. II. \S\ 6. (Leipzig, Teubner, 1889.)} Sobre esta base vienen luego los procesos de generación de las formas, el proceso simbólico de plegamiento y los procesos de diferenciación no simbólicos, la teoría de los desarrollos en serie y análogos. En este contexto quisiera mencionar además un teorema demostrado en mi disertación,\footnote{\emph{Über die Bildung des Formensystems der ternären biquadratischen Form}. \S\ 1. (Crelle's Journal Bd. 134. 1908.)} a saber, que todos los plegamientos pueden generarse por aplicación sucesiva de los dos plegamientos ``cogredientes'' posibles; por ello entiendo, para un producto simbólico $s_x t_x u_\sigma u_\tau$, los dos plegamientos $(stu)$ y $(\sigma\tau x)$. Sobre este teorema puede construirse la teoría de los reductores y de la reducción de los sistemas de formas, análogamente al dominio binario, como he mostrado en mi disertación.

Si pasamos ahora al dominio $n$-ario, ya el primer teorema del método simbólico sólo vale en medida condicionada. Como es sabido, ya en las formas cuaternarias aparecen, además de las coordenadas de punto y de plano $x$ y $u$, coordenadas de recta $p_{ik}$, los subdeterminantes de dos filas de coordenadas de punto. Análogamente, en el dominio $n$-ario tenemos filas de variables $p_\rho$, cuyos elementos individuales son los subdeterminantes de $\rho$ filas $x$, y filas simbólicas $S^\rho$, cuyos elementos se comportan como subdeterminantes de $\rho$ filas $s$ cogredientes con las $u$. La representación simbólica por determinantes, como es sabido, sólo vale cuando las filas simbólicas $S^\rho$ se resuelven en las filas $s$ y éstas se distribuyen entre distintos determinantes; entonces sólo tienen significado real aquellos agregados de determinantes que dependen de las $S^\rho$. Pero no se puede reconocer de conjunto qué agregados de determinantes logran esto; y con ello se pierden todos los demás teoremas mencionados en el dominio ternario.

Quisiera ahora indicar un principio simple para obtener una representación simbólica explícita en las filas $S^\rho$, y con ello recuperar también los restantes teoremas. Se trata de una aplicación del conocido teorema sobre las matrices correspondientes: ``A cada fila de variables $p_\rho$ se le puede asociar biunívocamente otra fila $q_{n-\rho}$, cuyos elementos individuales son subdeterminantes de $n-\rho$ filas $u$ unidas a las $\rho$ filas $x$ por ecuaciones $(u\mid x)=0$.'' Correspondientemente, a cada fila simbólica $S^\rho$ se le asocia otra $S_{n-\rho}$, que se comporta como si estuviera compuesta de $n-\rho$ filas cogredientes con las $x$.

El teorema admite todavía una segunda formulación, apropiada para las aplicaciones: ``A cada producto matricial $(v^{(1)}\cdots v^{(\rho)}\mid y^{(1)}\cdots y^{(\rho)})$,\footnote{Es decir, la suma de los productos de determinantes correspondientes, formados a partir de la matriz de las $\rho$ filas $v$ y de la de las $\rho$ filas $y$.} donde las filas de la segunda matriz son contragredientes con las de la primera, se le asocia un determinante; y, recíprocamente, todo determinante puede transformarse en un producto matricial con un número prefijado $\rho$ de filas.'' Así, determinante y producto matricial aparecen como completamente equivalentes, y el primer teorema del método simbólico toma la forma: ``Todas las formaciones invariantes pueden representarse simbólicamente por productos matriciales.'' A partir de dos filas simbólicas $S^\sigma$, $T^\tau$ se puede ahora, con ayuda de filas de variables, formar muy fácilmente un producto matricial que contenga explícitamente esas filas simbólicas y que, por tanto, represente una formación invariante de la forma $(S^\sigma\mid p_\sigma)(T^\tau\mid p_\tau)$, a saber:
\begin{equation}
 \pair{q_{n-\sigma-\lambda}S^\sigma}{T_{n-\tau}p_{\tau-\lambda}},
 \qquad (\lambda=1,2,\ldots,\tau\ \text{o}\ n-\sigma).
\tag{1}
\end{equation}

Más importante es la recíproca: que todas las formaciones invariantes puedan representarse explícitamente por productos matriciales, de modo que el proceso anterior significa la ampliación del proceso de plegamiento binario y ternario. Para demostrar esta recíproca es necesario trasladar el segundo teorema del método simbólico al dominio $n$-ario, es decir, establecer la totalidad de las identidades independientes e irreducibles en las que todas las filas $S^\rho$, $p_\rho$ se consideran irreducibles. Estas identidades se obtienen de las conocidas, que sólo contienen filas $x$ y $u$,\footnote{E. Pascal en \emph{Memorie d. R. Acc. d. Lincei}. (4) 5 (1888), p. 375.} sustituyendo el teorema del producto para determinantes por un sistema de teoremas del producto para matrices, y contrayendo distintos productos matriciales en uno solo precisamente mediante estos teoremas del producto. Se llega así a dos fórmulas duales, de las cuales sólo se cita una:
\begin{equation}
\begin{aligned}
 \pair{S_1^{\rho_1}S_2^{\rho_2}\cdots S_k^{\rho_k}}{x p_{\rho-1}}
 &=\sum_{i=1}^{k}\eps\,\pair{S_1^{\rho_1}\cdots S_{i-1}^{\rho_{i-1}}S_{i+1}^{\rho_{i+1}}\cdots S_k^{\rho_k}q_{n-\rho_i+1}}{S_{n-\rho_i}x},\\
 \rho&=\sum_{i=1}^{k}\rho_i<n,\qquad \eps=\pm1.
\end{aligned}
\tag{2}
\end{equation}

La única especialización contenida en estas fórmulas, que entre una fila $x$, no puede evitarse; para filas simbólicas y de variables completamente generales se llega a una ``identidad de descomposición'', una identidad en la que una fila $p_\rho$ se descompone en las filas individuales $x$. Clebsch ya utilizó el caso especial más simple de la identidad de descomposición en su ``problema fundamental de la teoría de invariantes'' para derivar los covariantes elementales en los desarrollos en serie. Con ayuda de estas identidades, que median el paso de las expresiones determinantes no explícitas al producto matricial, se puede probar efectivamente que con la representación matricial se obtiene la representación simbólica explícita deseada.

Para el proceso de plegamiento vale, sin embargo, un teorema por completo análogo al del dominio ternario, sobre el cual también pueden construirse análogamente los teoremas de reducción. Si en (1) se designa el número $\lambda$ como defecto del plegamiento, y se llaman cogredientes los plegamientos para los que $\sigma=\tau$, entonces todos los plegamientos pueden generarse por aplicación sucesiva de los plegamientos cogredientes de defecto 1. La prueba de este teorema se basa esencialmente en que las formas más generales pueden sustituirse por ``formas normales'', es decir, por formas que se anulan idénticamente bajo la aplicación de todas las operaciones diferenciales invariantes. En el dominio cuaternario esta última circunstancia fue demostrada por Mertens;\footnote{Über invariante Gebilde quaternärer Formen. Wiener Berichte. Bd. 98. 1889.} nuestro conocimiento de las operaciones diferenciales más generales --que, como es sabido, proceden de los procesos de plegamiento al sustituir las filas simbólicas por símbolos diferenciales-- nos permite, usando la identidad de descomposición, transferir casi literalmente la demostración de Mertens al dominio $n$-ario.

\clearpage
\begin{center}
{\Large\bfseries 4. Sobre la teoría de invariantes de las formas en $n$ variables}\par
\vspace{1.2em}
Journal für die reine und angewandte Mathematik 139 (1911), pp. 118--154
\end{center}

\section*{Introducción.}
En la teoría proyectiva de invariantes de las formas en $n$ variables, las cuestiones que se adjuntan al problema principal --la demostración de la finitud del sistema de formas-- apenas han sido tratadas cuando se sale de los dominios binario y ternario. Se trata de las cuestiones relativas a la dependencia mutua de las formas y a los métodos para dominar esta conexión entre formas. Estos métodos descansan esencialmente --como se ha demostrado por completo en los dominios binario y ternario-- en dos teoremas designados por Study como los ``teoremas fundamentales del método simbólico'', a saber, el teorema sobre la representabilidad simbólica de las formas y el teorema sobre las identidades simbólicas.\footnote{E. Study, \emph{Methoden zur Theorie der ternären Formen} (Leipzig, Teubner, 1889). II. \S\ 6. Sobre la aparición de estos teoremas fundamentales también en la teoría de invariantes de grupos especiales de transformaciones, cf. Study, \emph{Das Apollonische Problem} (Math. Ann. Bd. 49 [1897]) y \emph{Zur Differentialgeometrie der analytischen Kurven} (Transactions of the American Math. Society, Bd. 1 [1909]).} El último teorema afirma que todas las relaciones existentes entre formaciones invariantes enteras se obtienen por aplicación sucesiva de un pequeño número de identidades simbólicas; así, el método simbólico, y sólo él, da conocimiento de la conexión entre formas sin salir del dominio de los invariantes.

Gordan, en su programa de Erlangen,\footnote{P. Gordan, \emph{Über das Formensystem binärer Formen}. S. 46. (Leipzig, Teubner, 1875).} ya señaló la razón de la dificultad del paso a $n$ variables: reside en que el primer teorema del método simbólico ya no es válido en su forma general. En el dominio $n$-ario aparecen, como es sabido, filas de variables de nivel superior $p_\rho$\footnote{Para la notación véase \S\ 1.} y, análogamente, filas simbólicas de nivel superior $S^\rho$; se llega a estas filas tanto cuando, como originalmente en Clebsch,\footnote{A. Clebsch, über eine Fundamentalaufgabe der Invariantentheorie. (Abh. der Ges. d. Wiss. zu Göttingen, Bd. XVII, 1872).} se parte de formas con un número arbitrario de filas $p_\rho$, como cuando se parte, siguiendo a F. Deruyts y A. Capelli,\footnote{Cf. A. Capelli, \emph{Lezioni sulla teoria delle forme algebriche} (Napoli 1902), \S\ 22--24, y la literatura allí citada.} de formas que contienen sólo un número arbitrario de filas cogredientes $x$. No existe una representación simbólica que contenga explícitamente las filas $p_\rho$, $S^\rho$;\footnote{Definición más precisa de ``representación explícita'' en \S\ 2.} es necesario resolver $p_\rho$, $S^\rho$ en las filas individuales $x$ y en las $s$ contragredientes con ellas, y distribuir éstas entre distintos determinantes. Es posible, ciertamente, verificar mediante condiciones diferenciales (cf. fórmula (19.)) que un agregado dado de determinantes tiene la propiedad de depender sólo de las filas $p_\rho$, $S^\rho$; pero de este modo no se obtiene una visión de conjunto de la totalidad de las formaciones invariantes ni de los procesos para generarlas.\footnote{Una ampliación del simbolismo debida a Study y muy valiosa para el cálculo (comunicada por F. Engel, Ber. d. K. Sächs. Ges. d. Wiss., math.-phys. Kl., 1900, p. 126, y para el dominio cuaternario por E. Study, \emph{Geometrie der Dynamen}, p. 135) tampoco es una representación explícita en nuestro sentido. Por ejemplo, difícilmente conduciría a los procesos de diferenciación no simbólicos para generar las formas.}

Mostraremos ahora en \S\S\ 2 y 6 cómo, con ayuda del ``teorema sobre las matrices correspondientes'',\footnote{Aparte de la \emph{Ausdehnungslehre} de Grassmann de 1862, n.º 112, primero en A. Brill, über zwei Berührungsprobleme, \S\ 2 (Math. Ann. Bd. 4. 1871).} se recupera en su generalidad el primer teorema fundamental sustituyendo la representación por determinantes por la representación mediante ``productos matriciales''. En \S\ 2 se muestra que todo producto matricial que contiene explícitamente las filas $S^\rho,p_\rho$ es una formación invariante de la forma fundamental; la recíproca dada en \S\ 6, que toda formación invariante puede representarse explícitamente por productos matriciales, sólo se logra mediante la transferencia del segundo teorema fundamental (\S\S\ 3--5).

Este teorema es conocido para formas que contienen sólo filas de variables $x$.\footnote{E. Pascal, Giorn. di mat. 26 (1888), pp. 33, 102; Rendic. d. Accad. d. Linc. (4) 4 (1888), p. 119; ib. Mem. (4) 5 (1888), p. 375.} El problema general, establecer la totalidad de las identidades irreducibles en las que todas las filas $S^\rho,p_\rho$ se consideran irreducibles, fue formulado por Study;\footnote{``Methoden \dots'', p. 204, nota 17).} al mismo tiempo ve en el número de identidades que aparecen una medida de la dificultad de los métodos en el dominio $n$-ario. Puesto que ahora es posible reunir la totalidad de las identidades en una sola fórmula (36.), esta dificultad queda reducida al menos al mínimo. En las aplicaciones, sin embargo, se presentarán a menudo ciertos casos especiales distinguidos de (36.), como (20.), (21.), caracterizados por admitir representación explícita sin sustitución de nuevas filas; o la fórmula (31.), designada como ``identidad de descomposición'', porque una fila $p_\rho$ parece descomponerse en las filas $y$.

Sobre los dos teoremas fundamentales puede ahora construirse fácilmente la teoría ulterior. El primer teorema proporciona directamente los procesos para generar las formas, el proceso simbólico de plegamiento y los procesos de diferenciación no simbólicos,\footnote{En el trabajo ``Invariante Gebilde von Nullsystemen'' (Sitzungsber. d. Ak. d. Wiss. Wien, Bd. 97, Abt. IIa, 1888), Mertens llegó por otro camino, no directamente generalizable, al proceso de plegamiento cuaternario. El invariante alternante de 6 filas $S^2$ que allí aparece difiere del que surge aquí, por una sola aplicación de la sustitución (11.), en un agregado de invariantes pares. Para un caso especial, el proceso de plegamiento cuaternario también se encuentra en P. Gordan, \emph{Das simultane System von zwei quadratischen quaternären Formen} (Math. Ann. Bd. 56. 1903).} mientras que la identidad de descomposición elimina entre esos procesos todos los equivalentes (\S\ 6). El conocimiento de los procesos de diferenciación proporciona además el traslado del concepto de ``forma normal'' al dominio $n$-ario y, por medio de un procedimiento de demostración dado por Mertens\footnote{F. Mertens, Über invariante Gebilde quaternärer Formen. (Sitzber. d. Ak. d. Wiss. Wien. Bd. 98. Abt. IIa. 1889.)} en el dominio cuaternario, el teorema de que un sistema de formaciones invariantes puede sustituirse por un sistema equivalente de formas normales (\S\ 7). Bajo esta última hipótesis mostré en mi disertación\footnote{Über die Bildung des Formensystems der ternären biquadratischen Form. (Dieses Journal, Bd. 134. 1908.)} que el proceso de plegamiento ternario puede generarse por aplicación sucesiva de dos plegamientos fundamentales. Sobre la base de este teorema y de la teoría de los desarrollos en serie desarrollé luego, introduciendo el concepto de ``fila de formas'', los principales teoremas de reducción para sistemas de formas. De modo completamente análogo, en el dominio $n$-ario el proceso de plegamiento puede generarse a partir de $n-1$ plegamientos fundamentales (\S\ 8), y pueden establecerse los teoremas de reducción (\S\ 9).

\paragraph{Nota final.} Con ocasión de la comunicación de Salzburgo, el Prof. E. Müller, de Viena, llamó mi atención sobre el hecho de que el cálculo con productos matriciales que subyace a este trabajo coincide, en principio, con el cálculo de la teoría de la extensión de Grassmann.\footnote{Hermann Graßmanns gesammelte mathematische und physikalische Werke. Ersten Bandes zweiter Teil: Die Ausdehnungslehre von 1862. In Gemeinschaft mit H. Graßmann d. Jüngeren herausgegeben v. Fr. Engel (Leipzig, Teubner 1896).} En efecto, el ``producto combinatorio'' de Grassmann $[S^{\rho_1}S^{\rho_2}\cdots S^{\rho_i}]$ puede considerarse como una matriz dada por esas filas (cf. \S\ 2): concretamente, el ``producto progresivo'' como la matriz de las propias filas, el ``regresivo'' como la matriz de las filas asociadas $S_{n-\rho_1},S_{n-\rho_2},\ldots,S_{n-\rho_i}$, mientras que la matriz correspondiente al ``producto mixto'' surge cuando varias filas $S$ se reúnen en nuevas filas $P,Q$ mediante la sustitución (9.). Nuestra ``fila asociada'' corresponde al ``complemento'' de Grassmann; nuestro ``producto matricial'', a un ``producto mixto de número de nivel 0''.

Bajo esta correspondencia, el desarrollo dado en la \emph{Ausdehnungslehre} de 1862, n.º 113, se vuelve idéntico a una fórmula dual de nuestra fórmula (32.). En el caso especial $\rho=1$, (32.) pasa a (17.), y la fórmula dual a (16.). A partir de estos casos especiales del desarrollo de Grassmann, Müller\footnote{E. Müller, Beiträge zur Graßmannschen Ausdehnungslehre. 1. Mitteilung: Einige allgemeine Sätze. (Sitzungsber. d. Ak. d. Wiss. Wien; Bd. 118. Abt. IIa. Juli 1909), Satz 16 und 18.} ha deducido recientemente las identidades (20.), (21.) mencionadas arriba.

En contraste con la derivación de Grassmann--Müller, nuestra derivación da al mismo tiempo la prueba de completitud de las identidades, lo cual parece ser el punto esencial para las cuestiones aquí consideradas; y, partiendo de (20.), (21.), avanza además hasta las identidades irreducibles más generales (36.), que no parecen haber sido advertidas hasta ahora.

\section*{\S\ 1. Notaciones y definiciones. Resumen de resultados conocidos.}
Denotamos, como de costumbre, por la fila de variables $x$ la fila de elementos $x_1,x_2,\ldots,x_n$, y análogamente por la fila de variables $p_\rho$ ($\rho=1,2,\ldots,n-1$) la fila de los $\binom n\rho$ elementos $p_{i_1i_2\ldots i_\rho}$, donde $p_{i_1i_2\ldots i_\rho}$ designa los $\binom n\rho$ determinantes formados a partir de la matriz de las $\rho$ filas $x^{(1)},x^{(2)},\ldots,x^{(\rho)}$,
\[
\left|\begin{array}{cccc}
 x_{i_1}^{(1)}&x_{i_2}^{(1)}&\cdots&x_{i_\rho}^{(1)}\\
 x_{i_1}^{(2)}&x_{i_2}^{(2)}&\cdots&x_{i_\rho}^{(2)}\\
 \vdots&\vdots&&\vdots\\
 x_{i_1}^{(\rho)}&x_{i_2}^{(\rho)}&\cdots&x_{i_\rho}^{(\rho)}
\end{array}\right|,
\qquad (i_1<i_2<\cdots<i_\rho=1,2,\ldots,n).
\]
Según el teorema sobre las matrices correspondientes,\footnote{Cf. por ejemplo Gordan--Kerschensteiner, \emph{Vorlesungen über Invariantentheorie}, Bd. I, p. 99.} a cada fila $p_\rho$ se le asocia biunívocamente una segunda fila $q_{n-\rho}$ por la relación, válida para cada elemento de la fila,
\begin{equation}
 p_{i_1i_2\ldots i_\rho}=(-1)^{\frac{\rho(\rho+1)}{2}+i_1+i_2+\cdots+i_\rho}q_{j_1j_2\ldots j_{n-\rho}},
\tag{1}
\end{equation}
donde los $i$ y los $j$ están cada uno en buen orden y juntos forman una permutación completa de los números 1 a $n$. La fila $q_{n-\rho}$ consta de los $\binom n\rho$ determinantes de grado $n-\rho$, $q_{j_1j_2\ldots j_{n-\rho}}$, formados a partir de la matriz de $n-\rho$ filas $u^{(1)},u^{(2)},\ldots,u^{(n-\rho)}$ contragredientes con las $x$. Estas filas satisfacen las $\rho(n-\rho)$ ecuaciones de condición
\begin{equation}
 (u^{(k)}\mid x^{(l)})=\sum_{\alpha=1}^{n}u_\alpha^{(k)}x_\alpha^{(l)}=0;
 \qquad (k=1,2,\ldots,n-\rho;\ l=1,2,\ldots,\rho).
\tag{2}
\end{equation}
Las supondremos normalizadas (con Clebsch\footnote{a. a. O. \S\ 5.}) de modo que el factor de proporcionalidad que en otro caso entraría en (1.) se ponga igual a uno.

Según Clebsch\footnote{a. a. O. \S\ 12.} (y asimismo según las investigaciones posteriores de F. Deruyts y A. Capelli, cf. la Introducción), las formas más generales pueden reducirse a aquellas que contienen a lo sumo una de cada una de las $n-1$ filas $p_\rho$, y por tanto pueden representarse simbólicamente como sigue:
\begin{equation}
 (S^1\mid p_1)^{m_1}(S^2\mid p_2)^{m_2}\cdots(S^{n-1}\mid p_{n-1})^{m_{n-1}},
\tag{3}
\end{equation}
donde, en analogía con (2.), hemos puesto
\[
 (S^\rho\mid p_\rho)=\sum S^{i_1i_2\ldots i_\rho}p_{i_1i_2\ldots i_\rho}.
\]
Como consecuencia de las relaciones idénticas no lineales existentes entre los elementos de una fila $p_\rho$, las filas simbólicas $S^\rho$ pueden normalizarse de modo que satisfagan exactamente las mismas identidades; es decir, de modo que se comporten como los determinantes de una matriz de $\rho$ filas cuyas filas $s^{(1)},s^{(2)},\ldots,s^{(\rho)}$ son contragredientes con las $x$.\footnote{Clebsch, a. a. O. \S\ 4.} Así, a cada fila $S^\rho$ puede asociársele biunívocamente otra $S_{n-\rho}$ mediante la relación
\begin{equation}
 S_{i_1i_2\ldots i_{n-\rho}}=(-1)^{\frac{(n-\rho)(n-\rho+1)}2+i_1+i_2+\cdots+i_{n-\rho}}S^{j_1j_2\ldots j_\rho},
\tag{4}
\end{equation}
donde de nuevo los $i$ y los $j$ están cada uno en buen orden y juntos agotan exactamente los números 1 a $n$. Los elementos $S_{i_1i_2\ldots i_{n-\rho}}$ de la fila $S_{n-\rho}$ se comportan como los determinantes formados de $(n-\rho)$ filas $\sigma^{(1)},\sigma^{(2)},\ldots,\sigma^{(n-\rho)}$ cogredientes con las $x$, y las filas $s$ y $\sigma$ están enlazadas por las $\rho(n-\rho)$ condiciones
\begin{equation}
 (s^{(k)}\mid \sigma^{(l)})=0,\qquad(k=1,2,\ldots,\rho;\ l=1,2,\ldots,n-\rho).
\tag{5}
\end{equation}
En consecuencia de las relaciones (1.) y (4.), cada factor de (3.) admite la representación simbólica cuádruple equivalente
\begin{equation}
 (S^\rho\mid p_\rho)=(S^\rho q_{n-\rho})=(S_{n-\rho}p_\rho)=(-1)^{\rho(n-\rho)}(q_{n-\rho}\mid S_{n-\rho}),
\tag{6}
\end{equation}
donde, como siempre en lo sucesivo, $(S^\rho q_{n-\rho})$ y $(S_{n-\rho}p_\rho)$ son abreviaturas de los determinantes $(s^{(1)}\cdots s^{(\rho)}u^{(1)}\cdots u^{(n-\rho)})$ y $(\sigma^{(1)}\cdots\sigma^{(n-\rho)}x^{(1)}\cdots x^{(\rho)})$; por desarrollo de Laplace éstos resultan depender sólo de las filas $S^\rho,q_{n-\rho}$ y, respectivamente, $S_{n-\rho},p_\rho$, como pretende indicar la notación anterior. Las expresiones $(S^\rho\mid p_\rho)$ y $(q_{n-\rho}\mid S_{n-\rho})$ significan, según la discusión precedente, los productos de las matrices de las filas $s$ y $x$, respectivamente de las filas $u$ y $\sigma$; por producto matricial se entiende la suma de los productos de determinantes correspondientes, es decir, el valor del determinante de la matriz producto.

El número característico $\rho(n-\rho)$ perteneciente a una fila simbólica (respectivamente, a una fila de variables) se llamará, siguiendo a Clebsch, el peso de la fila;\footnote{a. a. O. \S\ 11.} el número $\rho$, que da el número de filas $s$, el índice de peso superior; el número $n-\rho$, que da el número de filas $\sigma$, el índice de peso inferior. De las relaciones (4.) se sigue que una fila simbólica puede aparecer en una misma relación una vez con índice de peso superior y otra con índice de peso inferior; lo mismo vale para la aparición de las filas correspondientes $p_\rho$ y $q_{n-\rho}$. Queda por señalar que, en general, denotamos por $p$ las filas de variables de nivel superior cuya matriz asociada consta de filas $x$, y por $q$ aquellas cuya matriz consta de filas $u$; por letras griegas minúsculas $\sigma,\tau,\alpha,\beta$, las filas simbólicas cogredientes con las $x$; por letras latinas minúsculas $s,t,a,b$, las filas simbólicas cogredientes con las $u$; y todas las demás filas simbólicas por letras latinas mayúsculas con índice de peso: $S^\rho$, $T^\tau$, $A^\rho$ o $S_{n-\sigma}$, $T_{n-\tau}$, $A_{n-\rho}$. De acuerdo con un índice de peso superior o inferior, escribimos también los elementos individuales de una fila con índices superiores o inferiores, como por ejemplo en (4.).
\clearpage
\section*{\S\ 2. Representación mediante productos matriciales.}
En lo que sigue entenderemos siempre por ``producto matricial'' la suma de los productos de determinantes correspondientes de dos matrices con el mismo número de filas, donde las filas de la segunda matriz son contragredientes con las de la primera. Las filas de cada matriz pueden ser: 1. dadas individualmente; 2. reunidas, como en (6.), en una fila $S^\rho$; 3. reunidas en distintas filas $S^\rho,S^\sigma\ldots$. El producto matricial se denotará con el símbolo $(\ \mid\ )$ usado en (6.). La representación (6.) es una aplicación especial de una segunda formulación del teorema sobre las matrices correspondientes: ``A cada producto matricial de $\rho$ filas $(v^{(1)}v^{(2)}\cdots v^{(\rho)}\mid y^{(1)}\cdots y^{(\rho)})$ se le asigna un determinante; y, recíprocamente, a cada determinante se le puede asignar un producto matricial con número prescrito $\rho$ de filas ($\rho=1,2,\ldots,n-1$).''\footnote{Para $\rho>n$ el producto matricial, como se sabe, se anula; para $\rho=n$ se convierte en el producto de dos determinantes.} En efecto, basta reunir las filas $y$ en una fila $p_\rho$ (respectivamente, las filas $v$ en una fila $q_\rho$) y efectuar las sustituciones dadas en (1.) para pasar de un producto matricial dado al determinante, y recíprocamente. Sin embargo, el valor del determinante no está dado por sus elementos individuales, sino que resulta sólo del desarrollo de Laplace. Puesto que determinante y producto matricial se muestran así equivalentes, el teorema sobre la representabilidad simbólica de las formas (el primer teorema fundamental) se expresa también de la manera siguiente:

\emph{Teorema I.} ``Todas las formaciones invariantes pueden representarse simbólicamente por productos matriciales, y recíprocamente todos los productos matriciales son formaciones invariantes.''\footnote{Las formaciones invariantes de formas fundamentales prescritas son, naturalmente, sólo aquellos agregados de productos matriciales que dependen de las filas simbólicas $S^\rho\ldots$.}

La representación mediante productos matriciales permite superar la dificultad esencial de la representación no explícita causada por la representación determinantal.\footnote{Cf. la Introducción.} Por ``representación explícita'' entendemos una representación en la que sólo entran las propias filas $S^\rho,T^\tau,\ldots$, o filas $R^\rho$ derivables de manera única de éstas, pero en la que ya no aparecen las filas componentes individuales $s,t,\ldots$. En \S\ 6 obtendremos, para todas las formaciones invariantes que sólo dependen de las filas simbólicas $S^\rho,T^\tau,\ldots$, una representación explícita en dichas filas simbólicas, y con ello el traslado al dominio $n$-ario de los procesos de generación, del proceso de plegamiento y de los procesos de diferenciación. La posibilidad de la representación explícita es patente a partir del hecho de que sólo a los productos matriciales más simples pueden asignarse precisamente los determinantes que aparecen en la representación determinantal, de modo que todos los demás productos matriciales son combinaciones de distintos determinantes en una sola expresión. Pues si se resuelven las filas simbólicas y de variables en las filas individuales $s,u,x$, entonces, según el primer teorema fundamental en su forma usual, debe surgir necesariamente un agregado de determinantes. La demostración de que, recíprocamente, todos los agregados de determinantes que representan formaciones invariantes de formas fundamentales prescritas pueden reunirse en productos matriciales explícitos se obtiene mediante el segundo teorema fundamental (\S\ 6).

Para obtener, a partir de dos filas simbólicas $S^\sigma,T^\tau$, con ayuda de filas de variables, una formación invariante de la forma fundamental $(S^\sigma\mid p_\sigma)(T^\tau\mid p_\tau)$ que contenga explícitamente todas las filas (el proceso de plegamiento), sólo necesitamos formar un producto matricial de la tercera clase según la definición dada al principio del parágrafo:
\begin{equation}
 \pair{q_{n-\sigma-\lambda}S^\sigma}{T_{n-\tau}p_{\tau-\lambda}},
 \qquad (\lambda=1,2,\ldots,\tau,n-\sigma).
\tag{7}
\end{equation}
Desarrollado, (7.) da:
\begin{equation}
\begin{aligned}
&\pair{v^{(1)}\cdots v^{(n-\sigma-\lambda)}s^{(1)}\cdots s^{(\sigma)}}
       {\alpha^{(1)}\cdots\alpha^{(n-\tau)}y^{(1)}\cdots y^{(\tau-\lambda)}} \\
&\qquad =\sum_{i,k,l}\eps\,q_{k_1\ldots k_{n-\sigma-\lambda}}
 S^{k_{n-\sigma-\lambda+1}\ldots k_{n-\lambda}}
 T_{l_1\ldots l_{n-\tau}}
 p^{l_{n-\tau+1}\ldots l_{n-\lambda}},
\end{aligned}
\tag{8}
\end{equation}
es decir, una expresión que sólo depende de las filas $S,T,q,p$. Aquí $\eps=\pm1$\footnote{Esta notación se conservará en todo lo que sigue.} y la suma debe entenderse de tal modo que los índices de cada fila $S\ldots$ estén en buen orden, y que, para cada combinación $i_1,i_2,\ldots,i_{n-\lambda}$, tanto los $k$ como los $l$ recorran individualmente todas las permutaciones posibles de los números $i$. El número $\lambda$ que entra en (7.) se llamará el defecto del plegamiento en el caso $\sigma>\tau$; la razón de esta última restricción aparecerá en \S\ 6.\footnote{El número $\lambda$ llega siempre hasta el menor de los dos últimos números indicados.}

Los determinantes de las dos matrices que entran en (7.) pueden considerarse como elementos de nuevas filas $Q^{n-\lambda},P_{n-\lambda}$, cuyas filas asociadas son $Q_\lambda,P^\lambda$. Según (8.), estas filas $Q,P$ pueden expresarse por las originales $S$ y $q$, respectivamente $T$ y $p$. Indicamos esto mediante las ecuaciones
\begin{equation}
\begin{aligned}
 q_{n-\sigma-\lambda}S^\sigma&=Q^{n-\lambda}\sim Q_\lambda, &\quad\text{o también}\quad Q_\lambda&=[q_{n-\sigma-\lambda}S^\sigma]_\lambda,\\
 T_{n-\tau}p_{\tau-\lambda}&=P_{n-\lambda}\sim P^\lambda, &\quad\text{o también}\quad P^\lambda&=[T_{n-\tau}p_{\tau-\lambda}]^\lambda.
\end{aligned}
\tag{9}
\end{equation}
Teniendo en cuenta (4.), se tiene entonces, análogamente a (6.), la representación simbólica cuádruple equivalente para el producto matricial (7.):
\begin{equation}
\pair{q_{n-\sigma-\lambda}S^\sigma}{T_{n-\tau}p_{\tau-\lambda}}
=(q_{n-\sigma-\lambda}S^\sigma P^\lambda)=(Q_\lambda T_{n-\tau}p_{\tau-\lambda})=(-1)^{\lambda(n-\lambda)}\pair{P^\lambda}{Q_\lambda}.
\tag{10}
\end{equation}
Sobre (7.) puede efectuarse todavía otra sustitución de nuevas filas poniendo
\begin{equation}
 \pair{q_{n-\sigma-\lambda}S^\sigma}{T_{n-\tau}p_{\tau-\lambda}}
 =\pair{R^{\sigma+\lambda}}{p_{\sigma+\lambda}}(R^{\tau-\lambda}p_{\tau-\lambda}).
\tag{11}
\end{equation}
También aquí (8.) muestra que los productos de los elementos de las filas $R$ pueden expresarse de manera única por las filas $S$ y $T$. Las sustituciones (9.) y (11.) se usarán especialmente cuando con (7.) se efectúe otro plegamiento con otras filas.

La representación matricial nos permite formular invariantemente las relaciones cuadráticas entre los elementos de una fila $p_\rho$ (de las cuales, como se sabe, se siguen todas las demás), y por tanto, según \S\ 1, las relaciones lineales en los coeficientes de la forma fundamental que satisfacen las filas $S^\rho$. Éstas son las identidades
\begin{equation}
 \pair{q_{n-\rho-\lambda}S^\rho}{S_{n-\rho}p_{\rho-\lambda}}=0,
 \qquad (\lambda=1,2,\ldots,\rho,n-\rho),
\tag{12}
\end{equation}
donde las $p$ y las $q$ han de considerarse como cantidades arbitrarias. Veremos en \S\ 5 que las identidades con número de defecto impar $\lambda$ son consecuencias de las de defecto superior. Las relaciones (12.) son consecuencia directa de las relaciones (5.); pues se tiene
\[
 \pair{q_{n-\rho-\lambda}S^\rho}{S_{n-\rho}p_{\rho-\lambda}}
 =(v^{(1)}\cdots v^{(n-\rho-\lambda)}s^{(1)}\cdots s^{(\rho)}\mid \sigma^{(1)}\cdots\sigma^{(n-\rho)}y^{(1)}\cdots y^{(\rho-\lambda)}),
\]
y la aplicación del teorema del producto da, como consecuencia de (5.), un determinante de $(n-\lambda)$ filas que se anula. Las identidades (12.) dicen que, para hallar todos los tipos de plegamientos, basta plegar el producto de formas lineales en cada fila de variables; o también, por \S\ 6, que la forma normalizada $(S^\rho\mid p_\rho)^m$ no posee ninguna formación invariante lineal en los coeficientes.\footnote{Cf. un añadido de Study en las obras de Grassmann, primer tomo, segunda parte, p. 510 (condición de que una cantidad $S^\rho$ sea simple).}

Las condiciones (12.) son también suficientes para caracterizar las filas $p_\rho$. Pues la propiedad de que los elementos individuales de $p_\rho$ sean determinantes de una matriz es una propiedad invariante, y por tanto debe poder expresarse invariantemente; pero, por el Teorema VI, las condiciones (12.) dan la mayor restricción invariante posible para las $p_\rho$, a saber, que la forma lineal formada con la fila $p_\rho$ como coeficientes no posea ninguna formación invariante.

\section*{\S\ 3. Las identidades simbólicas.}
Debemos ahora trasladar el segundo teorema fundamental del método simbólico, es decir, establecer la totalidad de las identidades irreducibles e indecomponibles en las que todas las filas $S^\rho,p_\rho$ se consideran indecomponibles. Estipulamos además lo siguiente. De acuerdo con el significado de las filas simbólicas y de variables, deberán considerarse compuestas aquellas identidades que surgen al especializar la fila de variables $p_\rho$ en $p_{\rho-1}y$ y aplicar una de las identidades del sistema a las expresiones resultantes; por otra parte, las identidades que contienen $k$ filas simbólicas deben contarse como indecomponibles aun cuando surjan de identidades con menos filas simbólicas mediante el proceso anterior. Las filas $S^\rho,p_\rho$ no tienen por qué entrar necesariamente de manera explícita en las identidades; para interpretarlas como cantidades indecomponibles basta mostrar que los productos matriciales que aparecen dependen sólo de estas filas. Sin embargo, mostraremos que en este caso, en parte aplicando la sustitución (11.), siempre puede obtenerse una representación explícita. Las identidades generales se obtienen, mediante el principio de la representación matricial, a partir de las identidades especiales dadas por Pascal, que sólo contienen filas $x$ y $u$ y constituyen una transferencia directa de las dadas por Study para el dominio ternario:\footnote{Para la formulación y la bibliografía del problema, cf. la Introducción.}
\begin{align}
\sum_{i=1}^n(-1)^{i+1}\pair{a^{(i)}}{x}(a^{(1)}\cdots a^{(i-1)}a^{(i+1)}\cdots a^{(n)}u)&=(-1)^{n+1}\pair{u}{x}(a^{(1)}\cdots a^{(n)}),\tag{13}\\
\sum\pm\pair{a^{(1)}}{x^{(1)}}\pair{a^{(2)}}{x^{(2)}}\cdots\pair{a^{(n)}}{x^{(n)}}&=(a^{(1)}\cdots a^{(n)})(x^{(1)}\cdots x^{(n)}),\tag{14}\\
\sum_{i=1}^n(-1)^{i+1}\pair{u}{a^{(i)}}(a^{(1)}\cdots a^{(i-1)}a^{(i+1)}\cdots a^{(n)}x)&=(-1)^{n+1}\pair{u}{x}(a^{(1)}\cdots a^{(n)}).\tag{15}
\end{align}
De (13.) y (15.) surgen otras dos identidades por las sustituciones $(a^{(i)}\mid x)=(a^{(i)}q_{n-1})$ y, respectivamente, $(u\mid a^{(i)})=(p_{n-1}a^{(i)})$; por tanto, según nuestra interpretación, no deben considerarse esencialmente distintas. (14.) representa el teorema del producto para determinantes; mostraremos cómo (13.) y (14.) (y, dualísticamente, (15.) y (14.)) pueden sustituirse por el sistema de teoremas del producto convenientemente formados para matrices. El teorema del producto, formado una vez para matrices de $\rho$ filas y otra para matrices de $(\rho-1)$ filas, dice
\[
\pair{a^{(1)}a^{(2)}\cdots a^{(\rho)}}{xp_{\rho-1}}
=(a^{(1)}a^{(2)}\cdots a^{(\rho)}\mid xy^{(2)}\cdots y^{(\rho)})
=\sum\pm\pair{a^{(1)}}{x}\pair{a^{(2)}}{y^{(2)}}\cdots\pair{a^{(\rho)}}{y^{(\rho)}},
\]
\[
\sum\pm\pair{a^{(2)}}{y^{(2)}}\cdots\pair{a^{(\rho)}}{y^{(\rho)}}
=(a^{(2)}\cdots a^{(\rho)}\mid y^{(2)}\cdots y^{(\rho)})
=(a^{(2)}\cdots a^{(\rho)}p_{\rho-1})=(a^{(2)}\cdots a^{(\rho)}q_{n-\rho+1});
\]
y de aquí se sigue el sistema de teoremas del producto, donde $\rho=2,3,\ldots,n$:
\begin{equation}
\begin{aligned}
\pair{a^{(1)}a^{(2)}\cdots a^{(\rho)}}{xp_{\rho-1}}
 &=\sum_{i=1}^{\rho}(-1)^{i+1}\pair{a^{(i)}}{x}(a^{(1)}\cdots a^{(i-1)}a^{(i+1)}\cdots a^{(\rho)}\mid p_{\rho-1})\\
 &=\sum_{i=1}^{\rho}(-1)^{i+1}\pair{a^{(i)}}{x}(a^{(1)}\cdots a^{(i-1)}a^{(i+1)}\cdots a^{(\rho)}q_{n-\rho+1}).
\end{aligned}
\tag{16}
\end{equation}
Para $\rho=n$, (16.) pasa a (13.); si además especializamos $p_{n-1}=y^{(2)}p_{n-2}$, $p_{n-2}=y^{(3)}p_{n-3}$, etc., y aplicamos las identidades (16.) a las expresiones $(a^{(1)}\cdots a^{(i-1)}a^{(i+1)}\cdots a^{(n)}\mid y^{(2)}p_{n-2})$, $(a^{(1)}\cdots a^{(i-1)}a^{(i+1)}\cdots a^{(h-1)}a^{(h+1)}\cdots a^{(n)}\mid y^{(3)}p_{n-3})$, etc., pasamos de (13.) a (14.). Por tanto, la identidad (14.) es, según nuestra definición, compuesta a partir de las identidades (16.); equivalentemente, el sistema (16.) es equivalente a las identidades (13.) y (14.). El sistema (16.) satisface una de las condiciones requeridas para las identidades generales: contiene la fila $p_{\rho-1}$ como cantidad indecomponible y, al mismo tiempo, en forma explícita. Es el único sistema que satisface esta condición, y sólo ésta. Pues no podemos formar otros determinantes o productos matriciales a partir de las filas $a,x,p$; y entre éstos no pueden existir relaciones independientes de (16.), ya que de lo contrario, al eliminar $(a^{(1)}\cdots a^{(\rho)}\mid xp_{\rho-1})$, surgiría una relación entre $\rho$ ($\rho<n$) expresiones $(a^{(i)}\mid x)(a^{(1)}\cdots a^{(i-1)}a^{(i+1)}a^{(\rho)}q_{n-\rho+1})$, mientras que por (13.) al menos $n+1$ de tales expresiones deben entrar en una relación. De consideraciones análogas se obtiene el sistema equivalente a (15.) y (14.):
\begin{equation}
\begin{aligned}
\pair{u q_{\rho-1}}{a^{(1)}\cdots a^{(\rho)}}
 &=\sum_{i=1}^{\rho}(-1)^{i+1}\pair{u}{a^{(i)}}(q_{\rho-1}\mid a^{(1)}\cdots a^{(i-1)}a^{(i+1)}\cdots a^{(\rho)})\\
 &=\sum_{i=1}^{\rho}(-1)^{i+1}\pair{u}{a^{(i)}}(p_{n-\rho+1}a^{(1)}\cdots a^{(i-1)}a^{(i+1)}\cdots a^{(\rho)}).
\end{aligned}
\tag{17}
\end{equation}
Para pasar de (16.) a identidades entre filas arbitrarias $A^{\rho_1},B^{\rho_2},\ldots,x,p$, observamos que estas identidades deben convertirse en (16.) tan pronto como resolvamos
$A^{\rho_1}=a^{(1)}a^{(2)}\cdots a^{(\rho_1)}$, $B^{\rho_2}=b^{(1)}b^{(2)}\cdots b^{(\rho_2)}$, etc., y que también sólo pueden combinarse en expresiones finales mediante las operaciones determinadas por (16.) tan pronto como las expresiones individuales sean del tipo que entra en (16.). Así obtenemos, si además ponemos
\begin{equation}
 \rho=\rho_1+\rho_2+\cdots+\rho_k\le n,
\tag{18}
\end{equation}
\begin{align*}
&\pair{a^{(1)}\cdots a^{(\rho_1)}b^{(1)}\cdots b^{(\rho_2)}\cdots k^{(1)}\cdots k^{(\rho_k)}}{xp_{\rho-1}}
  =\pair{A^{\rho_1}B^{\rho_2}\cdots K^{\rho_k}}{xp_{\rho-1}} \\
&\quad=\sum_{i=1}^{\rho_1}(-1)^{i+1}\pair{a^{(i)}}{x}(a^{(1)}\cdots a^{(i-1)}a^{(i+1)}\cdots a^{(\rho_1)}B^{\rho_2}\cdots K^{\rho_k}q_{n-\rho+1})\\
&\qquad+(-1)^{\rho_1\rho_2}\sum_{i=1}^{\rho_2}(-1)^{i+1}\pair{b^{(i)}}{x}(b^{(1)}\cdots b^{(i-1)}b^{(i+1)}\cdots b^{(\rho_2)}A^{\rho_1}\cdots K^{\rho_k}q_{n-\rho+1})\\
&\qquad+\cdots\\
&\qquad+(-1)^{(\rho_1+\cdots+\rho_{k-1})\rho_k}\sum_{i=1}^{\rho_k}(-1)^{i+1}\pair{k^{(i)}}{x}(k^{(1)}\cdots k^{(i-1)}k^{(i+1)}\cdots k^{(\rho_k)}A^{\rho_1}\cdots K^{\rho_{k-1}}q_{n-\rho+1}).
\end{align*}
Cada suma individual (pero todavía no una parte de la suma) contiene realmente las filas $A^{\rho_1},B^{\rho_2},\ldots$ como cantidades indecomponibles; pues satisface las condiciones necesarias y suficientes\footnote{Capelli, a. a. O. \S\ 23, da condiciones suficientes: $\bigl(a^{(k)}\mid \frac{\partial}{\partial a^{(i)}}\bigr)=0$, $(k>i,\ i=1,2,\ldots,\rho_1-1)$. Éstas implican del modo más simple las condiciones necesarias y suficientes por aplicación del proceso de corchetes de Poisson según Wellstein, Von den Differentialgleichungen der projektiven Invarianten, Math. Ann. 67 (1909) \S\ 3.}
\begin{equation}
 \left(a^{(i+1)}\middle|\frac{\partial}{\partial a^{(i)}}\right)=0;\ \ldots\ ;
 \left(k^{(j+1)}\middle|\frac{\partial}{\partial k^{(j)}}\right)=0.
 \quad (i=1,2,\ldots,\rho_1-1;\ j=1,2,\ldots,\rho_k-1)
\tag{19}
\end{equation}
Mostramos que también contiene explícitamente todas las filas. En efecto, si ponemos $B^{\rho_2}\cdots K^{\rho_k}q_{n-\rho+1}=q_{n-\rho_1+1}$, entonces por (16.) y (10.) obtenemos
\begin{align*}
\sum(-1)^{i+1}\pair{a^{(i)}}{x}(a^{(1)}\cdots a^{(i-1)}a^{(i+1)}\cdots a^{(\rho_1)}q_{n-\rho+1})
&=\pair{A^{\rho_1}}{xp_{\rho_1-1}}\\
&=(A_{n-\rho_1}xp_{\rho_1-1})=(-1)^{(\rho_1+1)(n+1)}\pair{q_{n-\rho_1+1}}{A_{n-\rho_1}x}\\
&=(-1)^{(\rho_1+1)(n+1)}\pair{B^{\rho_2}\cdots K^{\rho_k}q_{n-\rho+1}}{A_{n-\rho_1}x}.
\end{align*}
Si calculamos las restantes sumas de modo análogo y luego ponemos, puesto que las filas ya están caracterizadas como distintas por los índices de peso, $B^{\rho_2}=A^{\rho_2}$, $C^{\rho_3}=A^{\rho_3},\ldots$, obtenemos las identidades deseadas:
\begin{equation}
\begin{aligned}
\pair{A^{\rho_1}A^{\rho_2}\cdots A^{\rho_k}}{xp_{\rho-1}}
&=\sum_{i=1}^{k}(-1)^{\delta_i}\pair{A^{\rho_1}\cdots A^{\rho_{i-1}}A^{\rho_{i+1}}\cdots A^{\rho_k}q_{n-\rho+1}}{A_{n-\rho_i}x},\\
\delta_i&=(\rho_1+\rho_2+\cdots+\rho_{i-1})\rho_i+(\rho_i+1)(n+1),
\end{aligned}
\tag{20}
\end{equation}
y de (17.) las identidades duales:
\begin{equation}
\begin{aligned}
\pair{u q_{\sigma-1}}{A_{\sigma_1}A_{\sigma_2}\cdots A_{\sigma_k}}
&=\sum_{i=1}^{k}(-1)^{\eps_i}\pair{u A^{n-\sigma_i}}{p_{n-\sigma+1}A_{\sigma_1}\cdots A_{\sigma_{i-1}}A_{\sigma_{i+1}}\cdots A_{\sigma_k}},\\
\eps_i&=(\sigma_1+\sigma_2+\cdots+\sigma_{i-1})\sigma_i+(\sigma_i+1)(n+\sigma).
\end{aligned}
\tag{21}
\end{equation}
\emph{Teorema II.} Con (20.) y (21.) queda agotada la totalidad de las identidades indecomponibles entre las filas $A^{\rho_i},x,p$, respectivamente $A_{\sigma_i},u,q$, y al mismo tiempo la totalidad de aquellas identidades indecomponibles que admiten representación explícita sin efectuar la sustitución (11.).

La primera parte del teorema se sigue de las consideraciones que conducen a (20.) y (21.); la segunda, de la imposibilidad de una identidad entre dos filas simbólicas:
\[
\pair{q_\rho A^\sigma}{B_{\rho+\sigma-\tau}p_\tau}
=\eps_1\pair{q_\rho B^{n-\rho-\sigma+\tau}}{A_{n-\sigma}p_\tau}
+\eps_2\pair{q_\rho q_{n-\tau}}{B_{\rho+\sigma-\tau}A_{n-\sigma}},
\]
donde $\rho\le\tau\le\sigma$ y ninguna de las filas tiene peso $1\cdot(n-1)$. (Otras hipótesis sobre $\rho,\tau,\sigma$ pueden reducirse a ésta permutando las designaciones de las filas.) Esta imposibilidad se sigue, por ejemplo, desarrollando (8.) bajo la especialización $q_\rho=lM^{\rho-1}$, $q_{n-\tau}=lN^{n-\tau-1}$, si se pone $l_1=1$, $M^{2,3,\ldots,\rho}=1$, $A^{\rho+1,\ldots,\rho+\sigma}=1$, todos los demás elementos de las filas $l,M,A$ iguales a cero, mientras que $N$ y $B$ son arbitrarios. Puesto que las identidades entre filas arbitrarias, al resolver una fila $p_\tau$ en $y^{(1)}y^{(2)}\cdots y^{(\tau)}$, pasan a identidades compuestas a partir de (20.), la segunda parte de la afirmación queda probada tan pronto como (20.) pueda generarse a partir de identidades con sólo dos filas simbólicas. De hecho, a partir de
\begin{equation}
\pair{A^{\rho_1}A^{\rho_2}}{xp_{\rho-1}}
=\eps\pair{A^{\rho_1}q_{n-\rho+1}}{A_{n-\rho_2}x}
+\pair{A^{\rho_1}}{xp_{\rho_1-1}},
\quad \text{donde }p_{\rho_1-1}\sim A^{\rho_2}q_{n-\rho+1},
\tag{20a}
\end{equation}
por especialización $A^{\rho_1}=B^{\sigma_1}B^{\sigma_2}$, es decir, por
\[
\pair{B^{\sigma_1}B^{\sigma_2}}{xp_{\rho_1-1}}
=\eps\pair{B^{\sigma_1}q_{n-\rho_1+1}}{B_{n-\sigma_2}x}
+\pair{B^{\sigma_1}}{xp_{\sigma_1-1}},
\quad \text{donde } p_{\sigma_1-1}\sim B^{\sigma_2}q_{n-\rho_1+1},
\]
se obtiene una identidad (20.) con tres filas simbólicas; y por nuevas especializaciones de $B^{\sigma_1}=C^{\tau_1}C^{\tau_2}$, etc., las identidades generales (20.). La inexistencia de una representación explícita de una identidad entre dos filas simbólicas y dos filas de variables arbitrarias implica así la inexistencia en el caso de un número arbitrario de filas simbólicas, siempre que entre las filas de variables no esté contenida ninguna fila $x$ o $u$.

\section*{\S\ 4. Las identidades de descomposición.}
Consideramos ahora el caso más general de identidades entre filas simbólicas y de variables arbitrarias, en cuyas expresiones finales, sin efectuar la sustitución (11.), aparece una fila $p_\tau$ resuelta en sus filas individuales $y^{(1)}\cdots y^{(\tau)}$; llamamos a estas identidades ``identidades de descomposición''. Siguiendo la observación del final del parágrafo precedente, primero las determinamos para el caso de sólo dos filas simbólicas, es decir, desarrollamos la expresión $\pair{q_\rho A^\sigma}{B_{\rho+\sigma-\tau}p_\tau}$. Ponemos:
\begin{equation}
\begin{aligned}
q_{\rho-\alpha}^{(i_1\ldots i_\alpha)}&\sim p_{n-\rho}\,y^{(i_1)}\cdots y^{(i_\alpha)},\qquad
p_{\tau-\alpha}^{(i_1\ldots i_\alpha)}=y^{(i_{\alpha+1})}\cdots y^{(i_\tau)},\qquad
p_\tau=y^{(1)}\cdots y^{(\tau)}.
\end{aligned}
\tag{22}
\end{equation}
Respecto de los números $\rho,\sigma,\tau$ deben distinguirse cuatro casos:
\[
\begin{array}{llll}
1)&\rho+\sigma\le n,& \tau\le\rho,& \tau\le\sigma;\\[2pt]
2)&\rho+\sigma\le n,& \tau\ge\rho,& \tau\le\sigma;\\[2pt]
3)&\rho+\sigma\le n,& \tau\le\rho,& \tau>\sigma;\\[2pt]
4)&\rho+\sigma\le n,& \tau\ge\rho,& \tau\ge\sigma.
\end{array}
\]
Destacaremos el primer caso y luego caracterizaremos brevemente los otros a partir de él. Por (20.), (10.) y (9.) obtenemos, cuando descomponemos la fila $y^{(1)}y^{(2)}\cdots y^{(\tau)}B_{\rho+\sigma-\tau}$ en $y^{(1)}$ y $y^{(2)}\cdots y^{(\tau)}B_{\rho+\sigma-\tau}$:
\begin{equation}
\begin{aligned}
\pair{q_\rho A^\sigma}{y^{(1)}y^{(2)}\cdots y^{(\tau)}B_{\rho+\sigma-\tau}}
&=\pair{q_\rho^{(1)}A^\sigma}{y^{(2)}\cdots y^{(\tau)}B_{\rho+\sigma-\tau}}\\
&\quad+(-1)^\rho\pair{q_\rho[A_{n-\sigma}y^{(1)}]^{\sigma-1}}{y^{(2)}\cdots y^{(\tau)}B_{\rho+\sigma-\tau}}.
\end{aligned}
\tag{23}
\end{equation}
El último término de (23.) tiene ahora exactamente la forma de la expresión original --sólo que $\sigma$ se sustituye por $\sigma-1$, y $\tau$ por $\tau-1$-- y por tanto admite de nuevo la ecuación (23.). Así, puesto que $\tau\le\sigma$, podemos aplicar $\tau$ veces la operación (23.) y, teniendo en cuenta (10.), obtenemos la relación:
\begin{equation}
\begin{aligned}
\pair{q_\rho A^\sigma}{p_\tau B_{\rho+\sigma-\tau}}
&=(-1)^{\rho\sigma+(\sigma+\tau)(n+1)}
  \pair{q_\rho B^{n-\rho-\sigma+\tau}}{A_{n-\sigma}p_\tau}\\
&\quad+\sum_{i_1=1}^{\tau}(-1)^{\rho(i_1-1)}
 \pair{q_{\rho-1}^{(i_1)}[A_{n-\sigma}y^{(1)}\cdots y^{(i_1-1)}]^{\sigma-i_1+1}}
 {y^{(i_1+1)}\cdots y^{(\tau)}B_{\rho+\sigma-\tau}}.
\end{aligned}
\tag{24}
\end{equation}
Cada término bajo el signo de suma vuelve a ser del tipo de la expresión original; $\rho$ se sustituye por $\rho-1$, $\sigma$ por $\sigma-i_1+1$ y $\tau$ por $\tau-i_1$, de modo que la condición 1) queda satisfecha aún más fuertemente. Como $\tau\le\rho$, podemos por tanto aplicar $\tau$ veces la operación (24.) y obtener la identidad de descomposición buscada:
\begin{equation}
\begin{aligned}
\pair{q_\rho A^\sigma}{p_\tau B_{\rho+\sigma-\tau}}
&=\eps_{i_0}\pair{q_\rho B^{n-\rho-\sigma+\tau}}{A_{n-\sigma}p_\tau}\\
&\quad+\sum_{i_1=1}^{\tau}\eps_{i_1}
 \pair{q_{\rho-1}^{(i_1)}B^{n-\rho-\sigma+\tau}}{A_{n-\sigma}p_{\tau-1}^{(i_1)}}\\
&\quad+\sum_{i_2=i_1+1}^{\tau}\eps_{i_2}
 \pair{q_{\rho-2}^{(i_1i_2)}B^{n-\rho-\sigma+\tau}}{A_{n-\sigma}p_{\tau-2}^{(i_1i_2)}}+\cdots\\
&\quad+\eps_{i_\tau}\pair{q_{\rho-\tau}^{(i_1\ldots i_\tau)}B^{n-\rho-\sigma+\tau}}{A_{n-\sigma}}\\
&=\sum_{\alpha=0}^{\tau}\sum_{i_\alpha=i_{\alpha-1}+1}^{\tau}\eps_{i_\alpha}
 \pair{q_{\rho-\alpha}^{(i_1\ldots i_\alpha)}B^{n-\rho-\sigma+\tau}}{A_{n-\sigma}p_{\tau-\alpha}^{(i_1\ldots i_\alpha)}}.
\end{aligned}
\tag{25}
\end{equation}
Para el factor de signo $\eps$ se obtiene de (24.):
\begin{equation}
\begin{aligned}
\eps_{i_\alpha}&=(-1)^{\delta_{i_\alpha}},\qquad
\delta_{i_\alpha}=\sum_{\beta=1}^{\alpha}(\rho-\beta+1)(i_{\beta-1}+i_\beta+1)\\
&\quad +(\rho-\alpha)(\sigma+i_\alpha+\alpha)+(\sigma+\tau+\alpha)(n+1),
\end{aligned}
\tag{26}
\end{equation}
donde debe ponerse $i_0=0$. El signo de suma en (25.) debe entenderse así: a cada combinación $i_1i_2\ldots i_{\alpha-1}$, donde $i_1<i_2<\cdots<i_{\alpha-1}$, se adjuntan todos los valores $i_\alpha$ para los que $i_{\alpha-1}<i_\alpha\le\tau$. Teniendo en cuenta (26.), se verifica fácilmente que las sumas individuales en (25.) satisfacen las ecuaciones diferenciales análogas a (19.):
\begin{equation}
 \left(\frac{\partial}{\partial y^{(i)}}\middle|y^{(i+1)}\right)=0,
 \qquad (i=1,2,\ldots,\tau-1),
\tag{27}
\end{equation}
y por tanto dependen sólo de la fila $p_\tau$. Así, la identidad de descomposición es una identidad indecomponible entre las filas $A,B,q,p$; volveremos a su representación explícita en el parágrafo siguiente.

En el caso 2), la operación (23.) también puede aplicarse $\tau$ veces, pero la operación (24.) se detiene después de la $\rho$-ésima aplicación. Así, en la expresión final de (25.) el primer signo de suma debe reemplazarse por $\sum_{\alpha=0}^{\rho}$. En los casos 3) y 4), la operación (23.) sólo puede realizarse $\sigma$ veces; en lugar de (24.) se tiene entonces la relación:
\begin{equation}
\begin{aligned}
\pair{q_\rho A^\sigma}{p_\tau B_{\rho+\sigma-\tau}}
&=(-1)^{\rho\sigma}
 \pair{q_\rho^{(\sigma+1\ldots\tau)}B^{n-\rho+\tau-\sigma}}{A_{n-\sigma}p_{\tau-\sigma}^{(\sigma+1\ldots\tau)}}\\
&\quad+\sum_{i_1=1}^{\sigma}(-1)^{\rho(i_1-1)}
 \pair{q_{\rho-1}^{(i_1)}[A_{n-\sigma}y^{(1)}\cdots y^{(i_1-1)}]^{\sigma-i_1+1}}
 {y^{(i_1+1)}\cdots y^{(\tau)}B_{\rho+\sigma-\tau}};
\end{aligned}
\tag{28}
\end{equation}
al repetir (28.) $\tau-\sigma$ veces, donde el signo de suma se convierte en $\sum_{i_\beta=i_{\beta-1}+1}^{\sigma+\beta-1}$ --como muestra la aplicación $(\sigma-i_1+1)$ veces de (23.) a los términos bajo el signo de suma en (28.)-- se obtienen todos los términos de la suma individual de (25.) correspondientes al índice $\alpha=\tau-\sigma$, con el signo correspondiente. Entonces los casos 3) y 4) pasan a los casos 1) y 2), pues los números correspondientes a $\sigma,\tau$ se convierten en $\sigma'=\sigma-i_{\tau-\sigma}+\tau-\sigma$, $\tau'=\tau-i_{\tau-\sigma}=\sigma'$; y a medida que continúa el procedimiento, $\tau'<\sigma'$. Así, en la expresión final de (25.) el primer signo de suma debe reemplazarse por $\sum_{\alpha=\tau-\sigma}^{\tau,\rho}$.

Análogamente, de (21.) se obtiene la identidad de descomposición dual, en la que una fila $q_\rho$ se resuelve en las filas individuales $v^{(1)},\ldots,v^{(\rho)}$. Si ponemos
\begin{equation}
 \dot p_{\tau-\beta}^{(i_1\ldots i_\beta)}\sim v^{(i_1)}\cdots v^{(i_\beta)}q_{n-\tau};\qquad
 \dot q_{\rho-\beta}^{(i_1\ldots i_\beta)}=v^{(i_{\beta+1})}\cdots v^{(i_\rho)};\qquad
 q_\rho=v^{(1)}\cdots v^{(\rho)},
\tag{29}
\end{equation}
entonces
\begin{equation}
 \pair{q_\rho A^\sigma}{p_\tau B_{\rho+\sigma-\tau}}
 =\sum_{\beta=\max(0,\tau-\sigma)}^{\min(\tau,\rho)}
   \sum_{i_\beta=i_{\beta-1}+1}^{\rho}\eps
\pair{\dot q_{\rho-\beta}^{(i_1\ldots i_\beta)}B^{n-\rho-\sigma+\tau}}{A_{n-\sigma}\dot p_{\tau-\beta}^{(i_1\ldots i_\beta)}} ,
\tag{30}
\end{equation}
donde el índice de suma $\beta$ recorre desde el mayor de los valores inferiores hasta el menor de los valores superiores. Todos los términos de una suma individual en (25.) y (30.) son polares de una sola forma, producidos por los procesos polares
\[
\left(\frac{\partial}{\partial p_{\tau-\alpha}}\middle|p_{\tau-\alpha}^{(i_1\ldots i_\alpha)}\right),\qquad
\left(q_{\rho-\alpha}^{(i_1\ldots i_\alpha)}\middle|\frac{\partial}{\partial q_{\rho-\alpha}}\right).
\]
Para expresar esto, denotamos por $\Pi$ un agregado de tales operaciones polares, multiplicado por constantes, y obtenemos, en lugar de (25.) y (30.), la relación
\begin{equation}
 \pair{q_\rho A^\sigma}{p_\tau B_{\rho+\sigma-\tau}}
 =\sum_{\alpha=\max(0,\tau-\sigma)}^{\min(\tau,\rho)}
 \Pi\pair{q_{\rho-\alpha}B^{n-\rho-\sigma+\tau}}{A_{n-\sigma}p_{\tau-\alpha}}.
\tag{31}
\end{equation}

\section*{\S\ 5. Las identidades de descomposición en forma explícita.}
Si en (25), en el caso 4), especializamos $\tau$ a $\sigma+\rho$, obtenemos
\begin{equation}
 \pair{q_\rho A^\sigma}{y^{(1)}\cdots y^{(\rho+\sigma)}}
 =\sum_{1\le i_1<\cdots<i_\rho\le \rho+\sigma}
 \eps_{i_\rho}\pair{q_\rho}{y^{(i_1)}\cdots y^{(i_\rho)}}
 \pair{A^\sigma}{y^{(i_{\rho+1})}\cdots y^{(i_{\rho+\sigma})}},
\tag{32}
\end{equation}
donde $\eps_{i_\rho}=(-1)^{\delta_{i_\rho}}$ y $\delta_{i_\rho}=\rho(\rho+1)/2+i_1+\cdots+i_\rho$. Ésta es una forma más general del teorema del producto para matrices que (16.) y (17.), y se sigue por desarrollo de Laplace del determinante de la matriz producto
\[
 \sum\pm(v^{(1)}y^{(1)})\cdots(v^{(\rho)}y^{(\rho)})(a^{(1)}y^{(\rho+1)})\cdots(a^{(\sigma)}y^{(\rho+\sigma)}).
\]
Así, el teorema del producto más general se compone de las identidades (20.) y (21.), respectivamente, y con ello se explica la elección de la forma especial (16.), (17.).

La relación (32.) proporciona ahora el medio para representar explícitamente las identidades de descomposición. Para ello consideramos las formas que aparecen en (31.) como formas fundamentales; es decir, efectuamos la sustitución (11.), que no cambia nada en la representación explícita en las filas $A,B$, puesto que por (8.) las nuevas filas $R$ introducidas pueden expresarse mediante las propias filas $A,B$ y no mediante sus filas componentes $a^{(1)}a^{(2)}\ldots$, $b^{(1)}b^{(2)}\ldots$. Así, sea
\begin{equation}
 \pair{q_{\rho-\alpha}B^{n-\rho-\sigma+\tau}}{A_{n-\sigma}p_{\tau-\alpha}}
 =\pair{R_{\rho-\alpha}p_{n-\rho+\alpha}}{R^{\tau-\alpha}p_{\tau-\alpha}}.
\tag{33}
\end{equation}
Entonces, para la suma individual de (25.) correspondiente al índice $\alpha$, obtenemos
\begin{equation}
\begin{aligned}
&\sum \eps_{i_\alpha}\pair{R_{\rho-\alpha}p_{n-\rho}y^{(i_1)}\cdots y^{(i_\alpha)}}{R^{\tau-\alpha}y^{(i_{\alpha+1})}\cdots y^{(i_\tau)}}\\
&\quad=\sum \eps_{i_\alpha}\pair{[R_{\rho-\alpha}p_{n-\rho}]^\alpha y^{(i_1)}\cdots y^{(i_\alpha)}}{R^{\tau-\alpha}y^{(i_{\alpha+1})}\cdots y^{(i_\tau)}}\\
&\quad=\eps\pair{[R_{\rho-\alpha}p_{n-\rho}]^\alpha R^{\tau-\alpha}}{p_\tau}
 =\eps\pair{R^{\tau-\alpha}q_{n-\tau}}{R_{\rho-\alpha}p_{n-\rho}}.
\end{aligned}
\tag{34}
\end{equation}
Así se obtiene efectivamente una expresión final explícita, y la forma simétrica en que aparecen $q_\rho$ y $p_\tau$ muestra que (30.) conduce a la misma expresión final, que por tanto es independiente de la descomposición inicial. Entre (25.) y (30.) sólo hay diferencia mientras las filas $y$ y $v$ conservan significado independiente, de modo que la identidad de descomposición pasa, según nuestra definición, a ser una identidad descomponible. Para llegar a identidades con más filas simbólicas, basta, según \S\ 3 (final), resolver una fila $q_\rho$ en $q_{\rho_1}C^{\sigma_1}$, o $p_\tau$ en $p_{\tau_1}C_{\tau-\tau_1}$. Las expresiones resultantes
\begin{equation}
 \pair{q_{\rho_1}C^{\sigma_1}}{[R^{\tau-\alpha}q_{n-\tau}]_\lambda R_{\rho+\sigma_1-\alpha-\lambda}},
 \qquad
 \pair{[R_{\rho-\alpha}p_{n-\rho}]^\lambda R^{\tau-\alpha}}{p_{\tau_1}C_{\tau-\tau_1}}
\tag{35}
\end{equation}
son exactamente del tipo que aparece para dos filas simbólicas y, por tanto, admiten nuevamente representación explícita después de una sustitución correspondiente a (33.). La repetición proporciona así representación explícita para un número arbitrario de filas. También puede notarse que la primera y la última suma en (25.) y (30.), respectivamente, dan la representación explícita sin aplicar (33), sólo mediante la fórmula (32), o bien se reducen por completo a un solo término que contiene todas las filas explícitamente. La relación obtenida de (32.) al resolver $q_\rho$ en $q_{\rho_1}C^{\sigma_1}$ y así sucesivamente puede considerarse como el teorema del producto en su forma más general.

En resumen mostraremos:

\emph{Teorema III:} Con las identidades
\begin{equation}
 \pair{q_\rho A^\sigma}{p_\tau B_{\rho+\sigma-\tau}}
 =\sum_{\alpha=\max(0,\tau-\sigma)}^{\min(\tau,\rho)}
 \eps\pair{R^{\tau-\alpha}q_{n-\tau}}{R_{\rho-\alpha}p_{n-\rho}},
\tag{36}
\end{equation}
donde las filas $R$ están determinadas por (33), junto con las que se obtienen de ellas al resolver $q_\rho,p_\tau$ en $q_{\rho_1}C^{\sigma_1}$, etc., queda agotada la totalidad de las identidades indecomponibles.

En efecto, las identidades así obtenidas son las más generales de su especie, puesto que las filas individuales ya no están sometidas a ninguna restricción de peso o número, como aún ocurría en (20.) y (21.). Tampoco existen otras identidades entre estas filas; pues cuando una fila $p_\tau$ se resuelve en $y^{(1)}\cdots y^{(\tau)}$, las identidades más generales se componen de (20.), y por tanto, por \S\ 3 (final), de (20a.); y la composición discutida después de (16.) es precisamente la efectuada en \S\ 4. El paso a un número arbitrario de filas lo proporciona, como muestra la discusión de (20a.), aplicar siempre la misma operación a las expresiones que surgen de resolver $q_\rho$ en $q_{\rho_1}C^{\sigma_1}$, y así sucesivamente. También se sigue que el paso directo desde (20.) en lugar de desde (20a.) no aporta nada nuevo. Esto se verifica fácilmente por cálculo, especializando una vez una fila $q_\rho$ en (25.), y otra vez efectuando el paso desde (20.) por el método de \S\ 4. En ambos casos surgen las mismas sumas, que por aplicación del teorema general del producto (véase arriba) pasan a las identidades obtenidas de (36.). Las identidades (20.), (21.) corresponden a los casos especiales $\tau=1$, $\rho=1$; entonces sólo aparecen dos sumas individuales y por tanto, de acuerdo con una observación anterior, una representación explícita sin la sustitución (33.), coincidente con lo encontrado antes.

Finalmente mencionemos brevemente dos casos especiales de la identidad de descomposición. Póngase en (31): $\tau=\sigma$, $\rho=n-\sigma$, y denótese la fila $p_\tau$ por $p'_\sigma\sim q'_{n-\sigma}$. Entonces
\begin{equation}
\begin{aligned}
\pair{A^\sigma}{p_\sigma}\pair{B^\sigma}{p'_\sigma}
-\pair{A^\sigma}{p'_\sigma}\pair{B^\sigma}{p_\sigma}
=\sum_{\alpha=1}^{\min(\sigma,n-\sigma)}
\Pi\pair{q_{n-\sigma-\alpha}B^\sigma}{A_{n-\sigma}p_{\sigma-\alpha}}.
\end{aligned}
\tag{37}
\end{equation}
Ésta es la fórmula de Clebsch\footnote{Clebsch, loc. cit., pp. 25--32. Clebsch demuestra que los términos individuales de los lados derechos dependen sólo de las filas $A,B$; la representación explícita se obtendría aplicando (32.) a sus expresiones $\Phi_h$.} para desarrollar una forma con dos filas cogredientes $p_\sigma,p'_\sigma$ en covariantes elementales. Los covariantes elementales pertenecientes a una forma $\pair{A^\sigma}{p_\sigma}^m\pair{B^\sigma}{p'_\sigma}^n$ son entonces los siguientes:
\begin{equation}
 \left\{\sum_{\alpha=1}^{\min(\sigma,n-\sigma)}
 \eps\pair{q_{n-\sigma-\alpha}B^\sigma}{A_{n-\sigma}p_{\sigma-\alpha}}
 \right\}^{\rho}
 \pair{A^\sigma}{p_\sigma}^{m-\rho}\pair{B^\sigma}{p'_\sigma}^{n-\rho},
 \qquad (\rho=0,1,\ldots,m,n).
\tag{38}
\end{equation}
Proceden, como diremos brevemente, de la forma fundamental por ``plegamiento cogrediente de defecto $\rho$'' (cf. \S\ 2).

En segundo lugar, póngase $\tau=\sigma-\lambda$, $\rho=n-\sigma-\lambda$, $B^\sigma=A^\sigma$. Entonces
\begin{equation}
\begin{aligned}
\pair{q_{n-\sigma-\lambda}A^\sigma}{A_{n-\sigma}p_{\sigma-\lambda}}
&=(-1)^\lambda\pair{q_{n-\sigma-\lambda}A^\sigma}{A_{n-\sigma}p_{\sigma-\lambda}}\\
&\quad+(-1)^{(n-\sigma)(\sigma-\lambda)}
 \sum_{\alpha=1}^{\min(\sigma-\lambda,n-\sigma-\lambda)}
 \Pi\pair{q_{n-\sigma-\lambda-\alpha}A^\sigma}{A_{n-\sigma}p_{\sigma-\lambda}}.
\end{aligned}
\tag{39}
\end{equation}
Así, como se observó en \S\ 2, las condiciones (12.) que caracterizan las filas simbólicas para defecto impar $\lambda$ son consecuencias de las condiciones de defecto superior. Para una forma lineal $\pair{A^\sigma}{p_\sigma}=\pair{B^\sigma}{p_\sigma}$, (39.) implica que sus formaciones invariantes irreducibles cuadráticas en los coeficientes se reducen a las de defecto par; esto concuerda con el hecho conocido de que una forma lineal en los dominios binario y ternario no tiene formaciones invariantes.

Debe observarse que las identidades generales se derivaron originariamente bajo la hipótesis de que las filas individuales satisfacen las condiciones (12.). Sin embargo, puesto que estas filas individuales entran en las identidades sólo linealmente, éstas valen también para las filas más generales que no satisfacen las condiciones (12.).

\section*{\S\ 6. Representación explícita de las formaciones invariantes.\\Procesos de plegamiento y de diferenciación.}
Los resultados de los parágrafos precedentes nos permiten completar el teorema enunciado en \S\ 2 sobre la representabilidad simbólica (el primer teorema fundamental) añadiendo la representabilidad explícita; a saber, mostramos:

\emph{Teorema IV:} ``Todas las formaciones invariantes de formas fundamentales arbitrarias pueden representarse explícitamente por productos matriciales; es decir, aparte de las filas de variables $p_\rho$, sólo entran en las expresiones finales las filas $S^\sigma,\ldots,T^\tau$ de las formas originales, o las filas $P,Q,R$ derivadas de ellas por sustitución (9.) o (11.).''

Para la prueba, supongamos que una formación invariante depende, además de filas de variables, de las filas $S^\sigma,\ldots,T^\tau$; y supongamos que esas filas se han resuelto suficientemente en filas individuales $S^\sigma=S^{\sigma_1}\cdots S^{\sigma_i},\ldots,T^\tau=T^{\tau_1}\cdots T^{\tau_k}$ para hacer posible la representación por determinantes. Ordenemos las filas en una sucesión $S^\sigma,\ldots,T^\tau$ arbitraria en sí misma pero fijada de una vez por todas. Extraemos un agregado de determinantes que depende de la primera fila total $S^\sigma=S^{\sigma_1}\cdots S^{\sigma_i}$, no sólo de filas componentes individuales. Esta dependencia, sin embargo, es por \S\ 3--\S\ 5 una consecuencia de las identidades generales. Si resolvemos ahora una fila de variables $p_\rho$ en las filas $y,v$, o una de las filas posteriores $T$ en las filas individuales $t$, respectivamente $\tau$, entonces las identidades más generales se componen de (20.) y (21.). Pero estas últimas identidades conducen siempre a una expresión que contiene explícitamente la fila $S^\sigma=S^{\sigma_1}\cdots S^{\sigma_i}$, a saber, el lado izquierdo de (20.) y (21.); pues por (19.) se verifica que sólo la suma completa de la derecha --no una parte de ella correspondiente a un término de (20a.)-- posee la propiedad de depender sólo de $S^\sigma$. A medida que continúa el procedimiento, todas las filas que han entrado una vez explícitamente permanecen explícitas; la aplicación de las identidades puede requerir a lo sumo la reunión de varias filas en una sola, es decir, la sustitución (9.). Si finalmente sólo queda por llevar a forma explícita una única fila $T^\tau$, hay dos casos posibles. Puede aparecer todavía una fila de variables libre, esto es, no conectada con otras filas mediante la sustitución (9.); resolverla en filas componentes $y$ o $v$ completa el procedimiento y da, como en (31.), polares de una o más formas que contienen explícitamente todas las filas. Si no queda ninguna fila de variables libre, entonces por el modo en que surgieron reconocemos, en la totalidad de las expresiones que contienen las filas individuales $t$ o $\tau$ pero dependen de la fila $T^\tau$, los términos de una identidad de descomposición; por \S\ 5, sin embargo, éstas siempre admiten representación explícita en todas las filas después de aplicar la sustitución (11.). Esto prueba el teorema en toda generalidad. También debe señalarse que cuando sólo aparecen dos filas simbólicas, la sustitución (11.) nunca aparece. Pues, dado que $\sigma,n-\sigma,\tau,n-\tau<n$, sólo pueden faltar filas de variables cuando las dos filas simbólicas se unen en un único determinante y por tanto aparecen explícitamente; pero tan pronto como entran filas de variables, (34.) y (33.) muestran que la sustitución (11.) se vuelve innecesaria.

De lo recién dicho se sigue que, para generar todas las formaciones invariantes, basta unir las filas simbólicas, tras adjuntar filas de variables, en todos los productos matriciales posibles. El producto matricial más general tiene ahora la forma
\begin{equation}
 \pair{P^{\mu_1}\cdots P^{\mu_i}}{Q_{\nu_1}\cdots Q_{\nu_k}},\qquad \sum\mu=\sum\nu,
\tag{40}
\end{equation}
donde $P$ y $Q$ pueden ser filas de las formas originales, o pueden haber surgido de las filas originales por una sucesión de sustituciones (9.) o (11.), o, finalmente, pueden ser filas de variables. Pero las expresiones (40.) pueden generarse uniendo sucesivamente dos filas simbólicas cada vez en un producto matricial, con ayuda de filas de variables; pues de
\[
 \pair{q_{n-\mu-\lambda}P^\mu}{Q_{\nu_1}p_{n-\mu-\nu_1-\lambda}}=\pair{R_\lambda Q_{\nu_1}}{p_{n-\mu-\nu_1-\lambda}}
\]
se obtiene, al unir con la fila $Q_{\nu_2}$, una expresión
\[
 \pair{[R_\lambda Q_{\nu_1}]^{n-\lambda-\nu_1}}{Q_{\nu_2}p_{n-\mu-\nu_1-\nu_2-\lambda}}
 =\pair{q_{n-\mu-\lambda}P^\mu}{Q_{\nu_1}Q_{\nu_2}p_{n-\mu-\nu_1-\nu_2-\lambda}},
\]
y por consiguiente, continuando el procedimiento, una expresión (40.). Si $P$ y $Q$ no son filas de las formas originales, entonces (9.) y (11.), junto con lo que acaba de mostrarse, implican que también ellas surgieron por una sucesión de uniones de dos filas simbólicas cada vez. Así:

\emph{Teorema V:} El proceso para generar todas las formaciones invariantes, el proceso de plegamiento, consiste en unir sucesivamente dos filas simbólicas cada vez en un producto matricial, con ayuda de filas de variables.

A partir de dos filas simbólicas $S^\sigma,T^\tau$, donde $\sigma\ge\tau$, pueden formarse los siguientes productos matriciales:
\begin{align}
&\pair{q_{n-\sigma-\lambda}S^\sigma}{T_{n-\tau}p_{\tau-\lambda}}, &&(\lambda=1,2,\ldots,\min(\tau,n-\sigma)),\tag{41}\\
&\pair{q_{n-\tau-\mu}T^\tau}{S_{n-\sigma}p_{\sigma-\mu}}, &&(\mu=\sigma-\tau+\lambda),\tag{42}\\
&\pair{q_{n-\tau-\nu}T^\tau}{S_{n-\sigma}p_{\sigma-\nu}}, &&(\nu=1,2,\ldots,\sigma-\tau).
\tag{43}
\end{align}
De (31.), sin embargo, se obtienen los siguientes valores para (42.) y (43.):
\begin{align}
&\pair{q_{n-\sigma-\lambda}S^\sigma}{T_{n-\tau}p_{\tau-\lambda}}
 +\sum_{\alpha}\Pi\pair{q_{n-\sigma-\lambda-\alpha}S^\sigma}{T_{n-\tau}p_{\tau-\lambda-\alpha}},\tag{42a}\\
&\Pi(q_{n-\sigma}S^\sigma)(T_{n-\tau}p_\tau)
 +\sum_{\beta}\Pi\pair{q_{n-\sigma-\beta}S^\sigma}{T_{n-\tau}p_{\tau-\beta}}.
\tag{43a}
\end{align}
Así, las formas (42.) pueden expresarse linealmente por (41.) y polares de (41.), y las formas (43.) por polares de la forma fundamental y de las formas (41.). Asimismo (31.) da la dependencia lineal de las formas (41.) respecto de (42.) y polares de (42.). Por tanto, según la definición de Clebsch,\footnote{Clebsch, loc. cit., \S\ 7.} las formas (42.) son equivalentes a las formas (41.), mientras que (43.) conduce de vuelta a la forma fundamental. El proceso generador queda entonces dado con igual derecho por (41.) o por (42.); definiremos el proceso (41.), más simple en lo que se refiere al número $\lambda$, como el proceso de plegamiento. El número $\lambda$, el defecto del plegamiento (cf. \S\ 2), da el número de filas que faltan en el producto determinantal, a saber, en aquel producto matricial que, para el plegamiento dado, contiene el número máximo de filas. Como $\lambda$ sólo llega hasta el menor de los dos valores $\tau,n-\sigma$, donde $\sigma\ge\tau$, tenemos
\begin{equation}
 \lambda\le\frac n2,\ n\text{ par};\qquad
 \lambda\le\frac{n-1}{2},\ n\text{ impar}.
\tag{44}
\end{equation}
La reducción de (42.) y (43.) a (41.) sigue siendo válida aun cuando, por plegamientos posteriores, las filas $p,q$ hayan pasado a ser filas simbólicas; pues por (36.) se tiene
\[
 \pair{A^{n-\tau-\varkappa}T^\tau}{S_{n-\sigma}B_{\sigma-\varkappa}}
 =\sum_{\alpha=\max(0,\varkappa-\sigma+\tau)}^{\min(\tau,n-\sigma)}
 \eps\pair{R^{\tau-\alpha}B^{n-\sigma+\varkappa}}{A_{\tau+\varkappa}R_{n-\sigma-\alpha}},
\]
donde las filas $R$ han surgido de $S$ y $T$ según (33.). Por tanto, estas formas que sólo contienen filas simbólicas se obtienen plegando $\pair{q_\tau A^{n-\tau-\varkappa}}{B_{\sigma-\varkappa}p_{n-\sigma}}$ o $\pair{q_{\tau-\alpha}B^{n-\sigma+\varkappa}}{A_{\tau+\varkappa}p_{n-\sigma-\alpha}}$ --según que $\sigma-\tau\gtreqless 2\varkappa$-- con la forma fundamental y las formas (41.).

Del proceso simbólico de plegamiento se obtienen, de la manera usual, los procesos invariantes de diferenciación si simplemente se reemplazan las filas simbólicas $S^\sigma,T_{n-\tau}$ por las filas de símbolos de diferenciación $\frac{\partial}{\partial p_\sigma}$, $\frac{\partial}{\partial q_{n-\tau}}$; como es sabido, el producto $\frac{\partial}{\partial p_{i_1\ldots i_\sigma}}\frac{\partial}{\partial q_{j_1\ldots j_{n-\tau}}}$ tiene entonces el significado real $\frac{\partial^2}{\partial p_{i_1\ldots i_\sigma}\partial q_{j_1\ldots j_{n-\tau}}}$. Como las filas $\frac{\partial}{\partial p_\sigma}$ y $\frac{\partial}{\partial q_{n-\tau}}$ son cogredientes con las filas $S^\sigma,T_{n-\tau}$, satisfacen las mismas identidades que éstas, en particular también las que existen entre (42.) y (42a.), y entre (43.) y (43a.). En resumen obtenemos:

\emph{Teorema VI:} Todas las formaciones invariantes de formas fundamentales arbitrarias se obtienen de éstas por aplicación repetida del proceso simbólico de plegamiento (41.), o también de los procesos invariantes de diferenciación
\begin{equation}
 \left(q_{n-\sigma-\lambda}\frac{\partial}{\partial p_\sigma}\middle|\frac{\partial}{\partial q_{n-\tau}}p_{\tau-\lambda}\right),
 \qquad(\sigma\ge\tau,\ \lambda=1,2,\ldots,\min(\tau,n-\sigma)).
\tag{45}
\end{equation}
Queda por notar que en cada plegamiento (41.) el peso total de la forma, es decir, la suma de los pesos de las filas de variables individuales (cf. \S\ 1), disminuye en el número positivo $2(\lambda^2+\lambda(\sigma-\tau))$. Esto implica, independientemente de la demostración general de finitud, la finitud de la fila de formas, es decir, de la totalidad de las formaciones invariantes de una forma fundamental dada que son lineales en los coeficientes.\footnote{Cf. Clebsch, loc. cit., \S\ 11 y 12: finitud del sistema de covariantes elementales.}

Además debe observarse que el Teorema VI sigue siendo válido para formas que no satisfacen las condiciones (12.). Pues cada factor en (3), $\pair{S^\rho}{p_\rho}^{m_\rho}$, puede reemplazarse por un producto de $m_\rho$ factores lineales $\pair{A^\rho}{p_\rho}\pair{B^\rho}{p_\rho}\cdots\pair{M^\rho}{p_\rho}$; las formaciones invariantes pasan así a las de un sistema de formas lineales, que sólo después deben ponerse iguales entre sí. Para cada forma lineal individual, sin embargo, las condiciones (12.) --que, al introducir símbolos diferenciales, contienen sólo cocientes diferenciales de segundo orden-- se satisfacen siempre.

\section*{\S\ 7. Desarrollo de las formas generales según formas normales.}
En lo que sigue (\S\ 8) derivaremos una propiedad principal del proceso de plegamiento (41.), a saber, su composibilidad a partir de plegamientos fundamentales. Para ello necesitamos trasladar al dominio $n$-ario un importante teorema dado por Mertens\footnote{F. Mertens, \emph{"Uber invariante Gebilde quatern"arer Formen} (véase la Introducción), citado como ``Mertens''.} en el dominio cuaternario. Este teorema dice que todo sistema finito de formas puede reemplazarse por un sistema equivalente de formas normales. Por ``forma normal'' entendemos --generalizando la definición binaria y ternaria\footnote{Study, loc. cit., p. 54.}-- una forma cuyas formaciones invariantes lineales en los coeficientes se anulan idénticamente, con excepción de la propia forma fundamental. Así, por el Teorema VI (\S\ 6), una forma normal $\varphi$ satisface el sistema de ecuaciones diferenciales
\begin{equation}
 \left(q_{n-\sigma-\lambda}\frac{\partial}{\partial p_\sigma}\middle|\frac{\partial}{\partial q_{n-\tau}}p_{\tau-\lambda}\right)\varphi=0,
 \qquad(\sigma\ge\tau,\ \lambda=1,2,\ldots,\min(\tau,n-\sigma)),
\tag{46}
\end{equation}
donde las filas $q_{n-\sigma-\lambda},p_{\tau-\lambda}$ han de considerarse como cantidades arbitrarias. En el dominio cuaternario --teniendo en cuenta (39.) para $\sigma=\tau=2$-- estas ecuaciones diferenciales coinciden por completo con las diez ecuaciones diferenciales dadas por Mertens sin introducir las filas arbitrarias $p,q$ (loc. cit., fórmula 28). Su conocimiento, junto con el Teorema VI, nos permite trasladar su demostración casi palabra por palabra a $n$ variables. Lo único que debe añadirse es la verificación, obtenida mediante la identidad de descomposición (30.), de que las formas construidas son formas normales; en el dominio cuaternario esta verificación se seguía todavía sin transformación ulterior. Por ello basta dar una breve exposición de la prueba, y por simplicidad en el caso relevante más abajo, a saber, el de una sola forma fundamental $f$ y su fila de formas (cf. \S\ 6, final); para formaciones invariantes de grado arbitrario en los coeficientes de un número arbitrario de formas fundamentales, la demostración sigue siendo la misma.

La idea básica de la demostración es reemplazar, mediante la introducción de los coeficientes transformados, una forma con $n-1$ filas $p_\rho$ por formas con $n$ filas $\xi$ cogredientes con las $x$, a saber, las filas de las cantidades de transformación. De estas formas se obtienen directamente las formas normales buscadas mediante procesos invariantes de diferenciación. El desarrollo de las formas generales según formas normales es mediado por un desarrollo en serie para formas con $\rho$ ($\rho<n$) filas cogredientes $\xi$; los procesos invariantes de diferenciación conducen de vuelta a las formas en las filas originales $p_\rho$.

Sean $\xi^{(1)},\ldots,\xi^{(n)}$ las filas de cantidades de transformación, de modo que
\begin{equation}
\begin{aligned}
x_i&=\xi_i^{(1)}x'_1+\xi_i^{(2)}x'_2+\cdots+\xi_i^{(n)}x'_n,\\
p_{i_1\ldots i_\rho}&=\sum_j(\xi_{i_1}^{(j_1)}\cdots \xi_{i_\rho}^{(j_\rho)})p'_{j_1\ldots j_\rho}.
\end{aligned}
\tag{47}
\end{equation}
Los coeficientes transformados de las formas $f,\varphi,\ldots$ se denotarán por $(f),(\varphi),\ldots$, y sea
\[
 f=\pair{S^1}{p_1}^{m_1}\pair{S^2}{p_2}^{m_2}\cdots\pair{S^{n-1}}{p_{n-1}}^{m_{n-1}}.
\]
Entonces
\begin{equation}
\begin{aligned}
(f)&=\pair{S^1}{\xi^{(1)}}^{m_{11}}\cdots\pair{S^1}{\xi^{(n)}}^{m_{1n}}
\pair{S^2}{\xi^{(1)}\xi^{(2)}}^{m_{21}}\cdots
\pair{S^{n-1}}{\xi^{(2)}\cdots\xi^{(n)}}^{m_{n-1,n}}\\
&\quad\cdot\pair{s^{(1)}}{\xi^{(1)}}^{k_1}\pair{s^{(2)}}{\xi^{(2)}}^{k_2}\cdots\pair{s^{(n)}}{\xi^{(n)}}^{k_n},
\end{aligned}
\tag{48}
\end{equation}
donde $m_{11}+\cdots+m_{1n}=m_1$, y así sucesivamente. De cada coeficiente $(f)$ para el cual
\begin{equation}
 k_1\ge k_2\ge k_3\ge\cdots\ge k_n,
\tag{49}
\end{equation}
el proceso de diferenciación que agota exactamente las filas $\xi$,
\begin{equation}
\left(\frac{\partial}{\partial\xi^{(1)}}\middle|p_1\right)^{k_1-k_2}
\left(\frac{\partial}{\partial\xi^{(1)}}\frac{\partial}{\partial\xi^{(2)}}\middle|p_2\right)^{k_2-k_3}
\cdots
\left(\frac{\partial}{\partial\xi^{(1)}}\cdots\frac{\partial}{\partial\xi^{(n-1)}}\middle|p_{n-1}\right)^{k_{n-1}-k_n}
\left(\frac{\partial}{\partial\xi^{(1)}}\cdots\frac{\partial}{\partial\xi^{(n)}}\right)^{k_n}
\tag{50}
\end{equation}
da la forma
\begin{equation}
 \varphi=\bigl(s^{(1)}\mid p_1\bigr)^{k_1-k_2}
 \bigl(s^{(1)}s^{(2)}\mid p_2\bigr)^{k_2-k_3}\cdots
 \bigl(s^{(1)}\cdots s^{(n-1)}\mid p_{n-1}\bigr)^{k_{n-1}-k_n}
 \bigl(s^{(1)}s^{(2)}\cdots s^{(n)}\bigr)^{k_n}.
\tag{51}
\end{equation}
Las formas $\varphi$ son las formas normales buscadas. En efecto, poseen las siguientes propiedades que definen las formas normales.

1. Son formaciones invariantes de la forma fundamental $f$ lineales en los coeficientes. Esto se sigue del hecho de que surgen de $f$ por los procesos invariantes de diferenciación $\left(\frac{\partial}{\partial p_\rho}\middle|\xi^{(i_1)}\xi^{(i_2)}\cdots\xi^{(i_\rho)}\right)$ y (50.). Así, (por \S\ 6, final) pueden expresarse linealmente por polares de la fila finita de formas de $f$. Estos polares se obtienen considerando las variables $p_\rho$ que entran en $\varphi$ como coeficientes de una forma descompuesta en formas lineales con las filas de variables $p_\rho$:
\begin{equation}
 H=(q_{n-1}\mid p_{n-1})^{k_1-k_2}(q_{n-2}\mid p_{n-2})^{k_2-k_3}\cdots(q_1\mid p_1)^{k_{n-1}-k_n}.
\tag{52}
\end{equation}
Los invariantes simultáneos de las filas de formas $f$ y $H$ dan todos los polares pertinentes, puesto que éstos deben tener la misma dimensión en las filas individuales $p_\rho$ que la propia $\varphi$.

2. Satisfacen las ecuaciones diferenciales (46.). Si se aplica a $\varphi$ el proceso diferencial (45.), entonces
\[
 \varphi'=\bigl(q_{n-\sigma-\lambda}s^{(1)}\cdots s^{(\sigma)}\mid [s^{(1)}\cdots s^{(\tau)}]_{n-\tau}p_{\tau-\lambda}\bigr),\qquad (\sigma>\tau).
\]
De la identidad de descomposición (30.) se sigue, si las filas $q_\sigma=s^{(1)}\cdots s^{(\sigma)}$, $p_{n-\tau}=[s^{(1)}\cdots s^{(\tau)}]_{n-\tau}$ formadas a partir de las filas $s$ se identifican con las $q_\rho,p_\tau$ que allí aparecen, que
\[
 \varphi'=\sum_{\beta=\sigma-\tau+\lambda}^{n-\tau,\sigma}\sum_{i_\beta}\eps\,
 \pair{q_{\sigma-\beta}^{(i_1\ldots i_\beta)}q_{n-\tau+\lambda}}{p_{\sigma+\lambda}p_{n-\tau-\beta}^{(i_1\ldots i_\beta)}}.
\]
Aquí
\[
 p_{n-\tau-\beta}^{(i_1\ldots i_\beta)}\sim s^{(1)}\cdots s^{(\tau)}s^{(i_1)}\cdots s^{(i_\beta)},
\]
donde $i_1,\ldots,i_\beta$ designan cualesquiera $\beta$ números distintos elegidos entre los números 1 a $\sigma$. Como $\beta>\sigma-\tau$, se sigue que entre estos números $i_1,\ldots,i_\beta$ debe haber alguno entre 1 y $\tau$; por tanto $p_{n-\tau-\beta}^{(i_1\ldots i_\beta)}$ se anula idénticamente, y con ello se anula cada producto matricial individual y, con él, $\varphi'$.

Para la equivalencia del sistema de formas normales con la fila de formas $f$, sólo queda mostrar que todas las formas de la fila de formas pueden desarrollarse en polares de las formas normales en el sentido fijado bajo 1. Sea $\Pi$ un agregado de operaciones polares
\begin{equation}
 \left(\frac{\partial}{\partial\xi^{(i)}}\middle|\xi^{(k)}\right),
 \qquad (i\ne k;\ i=1,2,\ldots,\rho;\ k=1,2,\ldots,\rho),
\tag{53}
\end{equation}
y sea $\Pi^\sigma$ un agregado de las operaciones especiales (53.) para las que $i=\sigma$ y $k<\sigma$. Para una forma $F$ de las $\rho$ filas $\xi^{(1)},\ldots,\xi^{(\rho)}$ que tiene grado $k_\rho$ en la fila $\xi^{(\rho)}$, vale el desarrollo\footnote{F. Mertens, \emph{"Uber eine Formel der Determinantentheorie}. Sitzungsberichte der Akademie der Wissenschaften, Viena, Math.-Naturw. Kl., vol. 91, 2.ª sección (1885), fórmula 20).}
\begin{equation}
 F=\sum_{\alpha=0}^{k_\rho}\Pi F_\alpha,
 \qquad(\rho\le n),
\tag{54}
\end{equation}
donde $F_\alpha$ tiene grado $\alpha$ en la última fila $\xi^{(\rho)}$ y surge de $F$ por las operaciones diferenciales invariantes
\begin{equation}
 \left(\frac{\partial}{\partial\xi^{(1)}}\cdots\frac{\partial}{\partial\xi^{(\rho)}}\middle|\xi^{(1)}\cdots\xi^{(\rho)}\right)^\alpha
\tag{55}
\end{equation}
y $\Pi^\rho$. Si al construir $F_\alpha$ se reemplaza el proceso (55.) por
\begin{equation}
 \left(\frac{\partial}{\partial\xi^{(1)}}\cdots\frac{\partial}{\partial\xi^{(\rho)}}\middle|p_\rho\right)^\alpha,
\tag{56}
\end{equation}
entonces se obtiene una forma $F_\alpha(p_\rho)$ con sólo $\rho-1$ filas $\xi^{(1)},\ldots,\xi^{(\rho-1)}$, a la que puede aplicarse de nuevo (54.). La continuación del procedimiento conduce a formas $G(p_\rho,p_{\rho-1},\ldots,p_1)$. Para obtener de éstas el desarrollo de $F$ según (54.), se sustituyen las filas individuales $p_\rho$ por $\xi^{(1)},\ldots,\xi^{(\rho)}$ y se aplica a las formas resultantes el proceso polar $\Pi$ (53.); esto da (cf. (48.)) una suma de coeficientes transformados $(G)$. Para $\rho=n$, el proceso (55.) permanece en el primer paso. Si $c$ denota constantes numéricas y ponemos, para abreviar,
\begin{equation}
 (\xi^{(1)}\xi^{(2)}\cdots\xi^{(n)})\sim\Delta;
 \qquad
 \left(\frac{\partial}{\partial\xi^{(1)}}\frac{\partial}{\partial\xi^{(2)}}\cdots\frac{\partial}{\partial\xi^{(n)}}\right)=\nabla,
\tag{57}
\end{equation}
entonces, en el caso $\rho=n$,
\begin{equation}
 F=\sum c\Delta^\alpha(G),
\tag{58}
\end{equation}
mientras que para $\rho<n$ sólo aparece en el lado derecho $\alpha=0$, y por tanto $\sum c(G)$.

Aplicando (58.) a un coeficiente transformado $(f)$ (48.), y usando la conmutatividad de los procesos polares $\Pi^\sigma$ con todas las operaciones (56.) para las que $\rho\ge\sigma$, así como el hecho de que los polares de coeficientes transformados vuelven a dar coeficientes transformados, se obtiene
\[
 G=\left(\frac{\partial}{\partial\xi^{(1)}}\middle|p_1\right)^{\beta_1}
 \left(\frac{\partial}{\partial\xi^{(1)}}\frac{\partial}{\partial\xi^{(2)}}\middle|p_2\right)^{\beta_2}
 \cdots
 \left(\frac{\partial}{\partial\xi^{(1)}}\cdots\frac{\partial}{\partial\xi^{(n-1)}}\middle|p_{n-1}\right)^{\beta_{n-1}}
 \nabla^{\beta_n}\Pi(f)=\sum c\varphi,
\]
y de (58.) se sigue que
\begin{equation}
 (f)=\sum c\Delta^\alpha(\varphi)\qquad\text{(Mertens, fórmula 23)).}
\tag{59}
\end{equation}
Para pasar de (59.) al desarrollo de la propia $f$, consideramos de nuevo las variables $p_\rho$ como coeficientes de una forma $H$ (52.) con los grados $m_1,\ldots,m_{n-1}$ correspondientes a $f$, y ponemos $m=m_1+\cdots+m_{n-1}$. Entonces, por la definición de los invariantes,
\[
 f\cdot\Delta^m=\sum(f)(H)=\sum c\Delta^\alpha(\varphi)(H),
\]
y la aplicación del proceso $\nabla^m$, que agota exactamente las filas $\xi$, da
\begin{equation}
 f=\sum c\Phi,
\tag{60}
\end{equation}
donde las formas $\Phi$ se derivan de $\varphi$ y $H$ por procesos invariantes de diferenciación respecto de las variables, y por tanto son invariantes simultáneos de las filas de formas $\varphi$ y $H$.\footnote{En lugar de la conclusión directa, Mertens, \S\ 7, introduce una serie de procesos diferenciales ulteriores que sirven esencialmente para formar la fila de formas $H$; esto lo hace nuestro Teorema VI.} Pero como $\varphi$ es una forma normal, la fila de formas $\varphi$ se reduce a la propia forma $\varphi$; la fila de formas $H$ contiene todas las combinaciones invariantes posibles de variables. Así, los invariantes simultáneos de $\varphi$ y $H$ son los polares de $\varphi$ en el sentido indicado bajo 1.; en otras palabras, (60.) da el desarrollo de $f$ en polares de las formas normales. Queda verificar el mismo desarrollo para una forma arbitraria $\theta$ de la fila de formas $f$. Sean $\Gamma$ las formas normales derivadas de $(\theta)$ por (50.), de modo que
\begin{equation}
 (\theta)=\sum c\Delta^\beta(\Gamma).
\tag{61}
\end{equation}
Las formas $\theta$ se caracterizan por la relación, válida para todo coeficiente transformado,
\[
 (\theta)\Delta^\sigma=\sum c(f)\qquad\text{(Mertens, fórmula 9)).}
\]
Pero toda forma $F$ de las $n$ filas $\xi$ que contiene el factor $\Delta^\sigma$ contiene también la segunda $\nabla^\sigma\Pi F$ (Mertens, fórmula 20)); por tanto
\[
 (\theta)=\sum c\nabla^\sigma(f),\qquad\text{por consiguiente}\qquad \Gamma=\sum c\varphi,
\]
y de (61.) obtenemos la relación correspondiente a (59.),
\[
 (\theta)=\sum c\Delta^\beta(\varphi),
\]
que implica para $\theta$ el desarrollo (60.) en polares de las formas normales $\varphi$. Así:

\emph{Teorema VII.} Toda fila de formas $f$ puede sustituirse por un sistema equivalente de formas normales.

\section*{\S\ 8. Reducción del proceso de plegamiento a plegamientos fundamentales.}
Después del Teorema VII, llevamos ahora a cabo la reducción del proceso de plegamiento a plegamientos fundamentales mencionada al comienzo de \S\ 7. Llamamos (cf. \S\ 5) a todo plegamiento de dos filas cogredientes $S^\rho,T^\rho$ un ``plegamiento cogrediente'', y en particular al plegamiento cogrediente de defecto 1 de dos filas $S^\rho,T^\rho$ un ``plegamiento fundamental de tipo $\rho$''. Demostramos entonces, completando el Teorema VI:

\emph{Teorema VIII:} El proceso general de plegamiento (41.) puede componerse a partir de los $(n-1)$ plegamientos fundamentales de tipo $\rho$, donde $\rho=1,2,\ldots,n-1$. En otras palabras, para generar todas las formaciones invariantes basta aplicar sucesivamente los $n-1$ plegamientos fundamentales.

En el dominio binario el proceso de plegamiento (41.) coincide directamente con el único plegamiento fundamental posible; en el dominio ternario probé el Teorema VIII en \S\ 1 de mi disertación (cf. la Introducción). Allí, a causa de (44.), los plegamientos cogredientes generales son todavía idénticos a los plegamientos fundamentales, y es notable que el Teorema VIII para $n$ variables dé la reducción a plegamientos fundamentales, no sólo la reducción mucho más débil a plegamientos cogredientes. La prueba se modela en principio sobre la prueba ternaria: los plegamientos más generales se construyen a partir de los plegamientos fundamentales, usando las identidades simbólicas. Generaremos:
\begin{enumerate}
\item plegamientos arbitrarios de defecto 1 mediante plegamientos cogredientes de defecto 1, es decir, mediante plegamientos fundamentales (el análogo del caso ternario),
\item plegamientos arbitrarios de defecto $\lambda$ mediante plegamientos cogredientes de defecto $\lambda$ y plegamientos arbitrarios de defecto 1,
\item plegamientos cogredientes de defecto $\lambda$ mediante plegamientos cogredientes de defecto $\lambda-1$ y plegamientos arbitrarios de defecto 1.
\end{enumerate}
Como es esencial, suponemos por el Teorema VII que las dos formas $S$ y $T$ que han de plegarse son formas normales. Así, por (46.), valen las relaciones:
\begin{equation}
\begin{aligned}
\pair{q_{n-\sigma_1-\lambda}S^{\sigma_1}}{S_{n-\sigma_2}p_{\sigma_2-\lambda}}&=0,
&& (\sigma_1\ge\sigma_2,\ \lambda=1,2,\ldots,\sigma_2,n-\sigma_1),\\
\pair{q_{n-\tau_1-\lambda}T^{\tau_1}}{T_{n-\tau_2}p_{\tau_2-\lambda}}&=0,
&& (\tau_1\ge\tau_2,\ \lambda=1,2,\ldots,\tau_2,n-\tau_1).
\end{aligned}
\tag{62}
\end{equation}
Además, por \S\ 2 (final), basta suponer que las formas $S$ y $T$ son lineales en todas las filas de variables; es decir, las formas construidas pertenecen a la fila de formas
\[
 f=\pair{S^1}{p_1}\pair{S^2}{p_2}\cdots\pair{S^{n-1}}{p_{n-1}}
   \pair{T^1}{p_1}\pair{T^2}{p_2}\cdots\pair{T^{n-1}}{p_{n-1}}.
\]
1) Generación de $\pair{q_{n-\sigma-1}S^\sigma}{T_{n-\tau}p_{\tau-1}}$, donde $\sigma>\tau$, a partir de plegamientos fundamentales de tipo $\tau+\alpha$, $\alpha=0,1,\ldots,\sigma-\tau$. Del plegamiento fundamental de tipo $\tau$
\[
 \pair{q_{n-\tau-1}S^\tau}{T_{n-\tau}p_{\tau-1}}
\]
obtenemos, mediante la sustitución $w\sim T_{n-\tau}p_{\tau-1}$, una forma de la fila de formas $f$ con tres filas de índice de peso superior (cf. \S\ 1) $\tau+1$:
\begin{equation}
 \pair{wS^\tau}{p_{\tau+1}}\pair{S^{\tau+1}}{p_{\tau+1}}\pair{T^{\tau+1}}{p_{\tau+1}}.
\tag{63}
\end{equation}
Un plegamiento fundamental de tipo $\tau+1$ aplicado a (63.) da
\[
 \pair{q_{n-\tau-2}wS^\tau}{S_{n-\tau-1}p_\tau};
\]
por (31.), con la sustitución $q_{n-\tau-2}w=q'_{n-\tau-1}$, éste tiene el valor
\[
 \eps\pair{q_{n-\tau-2}wS^{\tau+1}}{S^\tau p_\tau}
 +\sum_{\alpha=1}^{\tau,n-\tau-1}\Pi\pair{q'_{n-\tau-1-\alpha}S^{\tau+1}}{S_{n-\tau}p_{\tau-\alpha}}.
\]
La expresión bajo el signo de suma se anula idénticamente por (62.); así hemos alcanzado una forma
\[
 \pair{wS^{\tau+1}}{p_{\tau+2}}\pair{S^{\tau+2}}{p_{\tau+2}}\pair{T^{\tau+2}}{p_{\tau+2}},
\]
que se obtiene de (63.) cambiando $\tau$ por $\tau+1$. Repitiendo el proceso $\sigma-\tau$ veces se llega por tanto al plegamiento deseado:
\[
 \pair{wS^\sigma}{p_{\sigma+1}}=\eps\pair{q_{n-\sigma-1}S^\sigma}{T_{n-\tau}p_{\tau-1}}.
\]
Indicamos el proceso que acabamos de efectuar mediante la fórmula
\begin{equation}
 T^\tau S^\tau,
\quad S^{\tau+1},\quad S^{\tau+2},\ldots,S^\sigma;
\tag{64}
\end{equation}
y vemos que sólo se usa la propiedad de forma normal de $S$, no la de $T$. Puede efectuarse un proceso dual a (64.), usando sólo la propiedad de forma normal de $T$, en el orden inverso, a saber:
\begin{equation}
 S^\sigma T^\sigma,
\quad T^{\sigma-1},\quad T^{\sigma-2},\ldots,T^\tau,
\tag{65}
\end{equation}
según las relaciones:
\begin{align*}
\pair{q_{n-\sigma-1}S^\sigma}{T_{1-\sigma}p_{\sigma-1}}&=\eps\pair{T_{n-\sigma}y}{p_{\sigma-1}},\quad y\sim q_{n-\sigma-1}S^\sigma,\\
\pair{q_{n-\sigma}T^{\sigma-1}}{T_{n-\sigma}yp_{\sigma-2}}&=\eps\pair{T_{n-\sigma+1}y}{p_{\sigma-2}},\\
&\vdots\\
\pair{q_{n-\tau-1}T^\tau}{T_{n-\tau-1}yp_{\tau-1}}&=\eps\pair{q_{n-\sigma-1}S^\sigma}{T_{n-\tau}p_{\tau-1}}.
\end{align*}
También debemos señalar la conexión con la disminución del peso total dada en \S\ 6 (final): para plegamientos de defecto 1 esta disminución tiene el valor $2(1+(\sigma-\tau))$, y para plegamientos fundamentales el valor $2\cdot1$. Por tanto, si la composición a partir de plegamientos fundamentales es posible en absoluto, se requieren necesariamente exactamente $\sigma-\tau+1$ tales plegamientos, como se llevó a cabo arriba.

2) Generación de plegamientos generales a partir de plegamientos cogredientes y fundamentales. La disminución de peso para un plegamiento general (41) es $2(\lambda^2+\lambda(\sigma-\tau))=2[\lambda^2+(\sigma-\tau)(1+(\lambda-1))]$. Esto da la posibilidad de descomponerlo en un plegamiento cogrediente de defecto $\lambda$ (disminución de peso $2\lambda^2$) y $\sigma-\tau$ plegamientos de defecto 1 con diferencia de índices de peso $\lambda-1$ (disminución de peso $2\{1+(\lambda-1)\}$). En el caso $\lambda=1$ esta descomposición coincide con la dada bajo 1); mostramos que efectivamente puede realizarse.

Del plegamiento cogrediente de defecto $\lambda$
\[
 \pair{q_{n-\tau-\lambda}S^\tau}{T_{n-\tau}p_{\tau-\lambda}}
\]
obtenemos, mediante la sustitución $R^\lambda\sim T_{n-\tau}p_{\tau-\lambda}$, una forma con dos filas de índices de peso superior $\tau+\lambda$ y $\tau+1$:
\begin{equation}
 \pair{R^\lambda S^\tau}{p_{\tau+\lambda}}\pair{S^{\tau+1}}{p_{\tau+1}}.
\tag{66}
\end{equation}
La fila con el índice de peso inferior $\tau+1$ pertenece a una forma normal; por tanto, (66.) admite un plegamiento de defecto 1 compuesto de $\lambda$ plegamientos fundamentales, en el orden (65):
\[
 (R^\lambda S^\tau)S^{\tau+\lambda},\quad S^{\tau+\lambda-1},\quad S^{\tau+\lambda-2},\ldots,S^{\tau+1}.
\]
Teniendo en cuenta (31.) y (62.), esto da
\[
 \pair{q_{n-\tau-\lambda-1}R^\lambda S^\tau}{S_{n-\tau-1}p_\tau}=\eps\pair{R^\lambda S^{\tau+1}}{p_{\tau+\lambda+1}},
\]
es decir, una forma obtenida de (66.) sustituyendo $\tau$ por $\tau+1$, exactamente como en 1) para (63.). Repitiendo el proceso $\sigma-\tau$ veces se llega al plegamiento deseado:
\[
 \pair{R^\lambda S^\sigma}{p_{\sigma+\lambda}}=\eps\pair{q_{n-\sigma-\lambda}S^\sigma}{T_{n-\tau}p_{\tau-\lambda}}.
\]
En todo el proceso sólo se usó la propiedad de forma normal de $S$. Considerando $T$ como forma normal se obtiene el proceso dual, con inversión completa del orden, según las fórmulas
\begin{align*}
\pair{q_{n-\sigma-\lambda}S^\sigma}{T_{n-\sigma}p_{\sigma-\lambda}}&=\eps\pair{T_{n-\sigma}R_\lambda}{p_{\sigma-\lambda}},\quad R_\lambda\sim q_{n-\sigma-\lambda}S^\sigma,\\
\pair{q_{n-\sigma}T^{\sigma-1}}{T_{n-\sigma}R_\lambda p_{\sigma-\lambda-1}}&=\eps\pair{T_{n-\sigma+1}R_\lambda}{p_{\sigma-\lambda-1}},\\
&\vdots\\
\pair{q_{n-\tau-1}T^\tau}{T_{n-\tau-1}R_\lambda p_{\tau-\lambda}}&=\eps\pair{q_{n-\sigma-\lambda}S^\sigma}{T_{n-\tau}p_{\tau-\lambda}}.
\end{align*}
3) Generación de plegamientos cogredientes de defecto $\lambda$ a partir de plegamientos cogredientes de defecto $\lambda-1$ y plegamientos fundamentales. La disminución de peso para los plegamientos cogredientes de defecto $\lambda$ es $2\lambda^2=2((\lambda-1)^2+2\lambda-1)$, lo que da la posibilidad de una descomposición en un plegamiento cogrediente de defecto $\lambda-1$ y $2\lambda-1$ plegamientos fundamentales. Esto se ve del modo más sencillo poniendo primero $n=2\lambda$. El único plegamiento cogrediente posible de defecto $\lambda$ es $\pair{S^\lambda}{T_\lambda}=\pair{S^\lambda}{T_{n/2}}$; contiene una fila $u$ y una fila $x$ menos que el plegamiento de defecto $\lambda-1$ de las mismas filas, $\pair{uS^\lambda}{T_\lambda x}$. Recíprocamente, componemos el plegamiento de defecto $\lambda$ a partir de plegamientos fundamentales que efectúan la desaparición de una fila $u$ y una fila $x$ (disminución de peso $2(n-1)=2(2\lambda-1)$) y al mismo tiempo generan una nueva fila simbólica $R^\lambda$, junto con un plegamiento cogrediente de defecto $\lambda-1$. El plegamiento más simple que hace esto es el siguiente de defecto 1:
\[
 \pair{sT^\lambda}{\sigma p_\lambda}=\eps\pair{S^{2\lambda-1}}{R_{\lambda-1}p_\lambda}=\eps\pair{R^\lambda}{p_\lambda},
\]
donde
\[
 R^{\lambda+1}=sT^\lambda,
 \quad s=S^1,
 \quad \sigma=S_1,
 \quad R^\lambda\sim\sigma R_{\lambda-1}.
\]
Por (65.) y (64.) se compone de los $2\lambda-1$ plegamientos fundamentales
\begin{equation}
 T^\lambda S^\lambda,S^{\lambda-1},\ldots,S^1;
 \qquad R^{\lambda+1}S^{\lambda+1},S^{\lambda+2},\ldots,S^{2\lambda-1}.
\tag{67}
\end{equation}
Un plegamiento de defecto $\lambda-1$, teniendo en cuenta (21.) y (62.), da
\[
 \pair{uS^\lambda}{\sigma R_{\lambda-1}x}=\eps\pair{R^{\lambda+1}}{S_\lambda x}
 =\eps\pair{sT^\lambda}{S^\lambda x}=\eps\pair{S^\lambda}{T_\lambda}.
\]
El caso general, para dos filas $S^\rho,T^\rho$ y $n$ arbitrario, se obtiene directamente del caso especial. En lugar de (67.) se tienen los siguientes $2\lambda-1$ plegamientos fundamentales:
\begin{equation}
 T^\rho S^\rho,S^{\rho-1},\ldots,S^{\rho-\lambda+1};
 \qquad R^{\rho+1}S^{\rho+1},S^{\rho+2},\ldots,S^{\rho+\lambda-1}.
\tag{68}
\end{equation}
La primera mitad de los plegamientos da la forma
\[
 \pair{q_{n-\rho-1}T^\rho}{S_{n-\rho+\lambda-1}p_{\rho-\lambda}},
\]
de la cual, mediante las sustituciones
\[
 S_{n-\rho+\lambda-1}p_{\rho-\lambda}\sim w,
 \qquad T^\rho w=R^{\rho+1},
\]
y la aplicación de la segunda mitad de los plegamientos, se obtiene
\[
 \pair{q_{n-\rho-\lambda}S^{\rho+\lambda-1}}{R_{n-\rho-1}p_\rho}.
\]
Sustituyendo
\[
 q_{n-\rho-\lambda}S^{\rho+\lambda-1}\sim y,
\]
se obtiene, mediante un plegamiento cogrediente de defecto $\lambda-1$,
\begin{equation}
 \pair{q_{n-\rho-\lambda+1}S^\rho}{yR_{n-\rho-1}p_{\rho-\lambda+1}}.
\tag{69}
\end{equation}
Mostramos que éste es el plegamiento deseado de defecto $\lambda$. Sustitúyase
\[
 R_{n-\rho-1}p_{\rho-\lambda+1}=R_{n-\lambda}
\]
y obsérvese que resolver $y$ da
\[
 \pair{q_{n-\rho-\lambda+1}R^\lambda}{S_{n-\rho}y}
 =\eps\pair{q_{n-\rho-\lambda}S^{\rho+\lambda-1}}{S_{n-\rho}p'_{\rho-1}}=0.
\]
Así, por (20.), el valor de (69.) es
\[
 \eps\pair{S^\rho R^\lambda}{p_{\rho+\lambda-1}y}
 =\eps\pair{q_{n-\rho-\lambda}S^{\rho+\lambda-1}}{[S^\rho R^\lambda]_{n-\rho-\lambda}p_{\rho+\lambda-1}}.
\]
Este plegamiento es de tipo (43.), y por tanto equivalente a
\[
 \pair{q_{n-\rho-\lambda}S^\rho}{R_{n-\rho-1}p_{\rho-\lambda+1}}
 =\eps\pair{wT^\rho}{R_\lambda p_{\rho-\lambda+1}},
 \qquad R_\lambda\sim q_{n-\rho-\lambda}S^\rho.
\]
Aquí $R_\lambda$ ya está en la forma final deseada; la expresión corresponde a la forma $\pair{sT^\lambda}{S_\lambda x}$ del caso especial. Ahora, observando que la resolución de $w$ y $R_\lambda$ da
\[
 \pair{wR^{n-\lambda}}{T_{n-\rho}p_{\rho-\lambda+1}}
 =\eps\pair{q'_{n-1}R^{n-\lambda}}{S_{n-\rho+\lambda-1}p_{\rho-\lambda}}
 =\eps\pair{q''_{n-\sigma-1}S^\rho}{S_{n-\rho+\lambda-1}p_{\rho-\lambda}}=0,
\]
obtenemos por (21.) el valor
\[
 \pair{wq_{n-\rho+\lambda-1}}{T_{n-\rho}R_\lambda}
 =\eps\pair{q_{n-\rho+\lambda-1}[T_{n-\rho}R_\lambda]^{n-\lambda}}{S_{n-\rho+\lambda-1}p_{\rho-\lambda}}
 \quad\text{equivalente a}\quad
 \pair{q_{n-\rho-\lambda}S^\rho}{T_{n-\rho}p_{\rho-\lambda}}.
\]
En todo el proceso se usó sólo la propiedad de forma normal de $S$, no la de $T$. Por tanto, el plegamiento cogrediente de defecto $\lambda-1$ usado para generar (69.) puede sustituirse por $2\lambda-3$ plegamientos fundamentales y un plegamiento cogrediente de defecto $\lambda-2$, que a su vez admite una descomposición en plegamientos fundamentales, y así sucesivamente. Análogamente, la generación se obtiene usando la propiedad de forma normal de $T$ y aplicando el proceso dualístico, esto es, los plegamientos fundamentales
\[
 S^\rho T^\rho,\quad T^{\rho-1},\ldots,T^{\rho-\lambda+1};
 \qquad R^{\rho-1}T^{\rho-1},\quad T^{\rho-2},\ldots,T^{\rho-\lambda+1},
\]
y un plegamiento dual a (69.). Entonces se obtiene el plegamiento
\[
\pair{q_{n-\rho-\lambda}T^\rho}{S_{n-\rho}p_{\rho-\lambda}},
\]
equivalente al anterior.

Añadiendo la descomposición obtenida bajo 2), hemos obtenido en efecto una generación de los plegamientos más generales a partir de plegamientos fundamentales. Resta mostrar que esta generación es única, salvo permutación del orden, es decir, que a cada plegamiento le corresponde uno y sólo un número $\alpha_\rho$ de plegamientos fundamentales de tipo $\rho$, donde $\rho=1,2,\ldots,n-1$.

Sean $m_\rho$ los grados en las filas de variables $p_\rho$ de la forma inicial $f$, y $l_\rho$ los grados de la forma producida por plegamiento. Cada plegamiento fundamental de tipo $\rho$ transforma los grados $m_{\rho-1},m_\rho,m_{\rho+1}$ respectivamente en $m_{\rho-1}+1,m_\rho-2,m_{\rho+1}+1$. Así, los $\alpha$ satisfacen el sistema de ecuaciones
\begin{equation}
 m_\rho-l_\rho=-\alpha_{\rho-1}+2\alpha_\rho-\alpha_{\rho+1},
 \qquad (\rho=1,2,\ldots,n-1),
\tag{70}
\end{equation}
donde debe ponerse $\alpha_0=\alpha_n=0$, y cuyo determinante tiene valor $n$.\footnote{Si $\Delta_n$ denota el determinante de $n$ filas, entonces vale la fórmula de recurrencia: $\Delta_n=2\Delta_{n-1}-\Delta_{n-2}$; es decir, $\Delta_n-\Delta_{n-1}=\Delta_{n-1}-\Delta_{n-2}=\cdots=\Delta_2-\Delta_1=1$, puesto que $\Delta_2=3$, $\Delta_1=2$.} Los $\alpha$ quedan así determinados de manera única, y en efecto, por el Teorema VIII, como enteros positivos; esto da condiciones para las diferencias $m_\rho-l_\rho$.

Los resultados de este parágrafo valen en virtud de las relaciones (62.); sin embargo, si las relaciones (62.) no se cumplen, las expresiones finales no son en modo alguno equivalentes a las expresiones iniciales. La definición de equivalencia dada en \S\ 6 admite también la formulación:\footnote{Cf. la consideración al final del parágrafo siguiente.} ``Dos formas son equivalentes si y sólo si difieren por formas de plegamiento superior, es decir, por formas que exhiben una mayor disminución de peso.'' Esta mayor disminución de peso ocurre siempre cuando las dos filas $q_{n-\sigma-i},p_{\tau-\lambda}$ que entran en el plegamiento son realmente filas de variables, como sucede en (41.), (42.), en los plegamientos que aparecen en este parágrafo como equivalentes, y también en las formas equivalentes del Teorema VII. Pero no tiene por qué ocurrir cuando esas filas se derivan de filas simbólicas por sustitución, como fue el caso de las formas de este parágrafo que pudieron expresarse unas en términos de otras en virtud de (62.). Se ve fácilmente que las formas que se anulan tienen incluso, en parte, un peso total mayor.

\section*{\S\ 9. Filas de formas. Teoremas de reducción.}
Definimos ahora con algo más de precisión la fila de formas introducida en \S\ 6 (final) como ``la totalidad de las formaciones invariantes de una forma inicial dada que son lineales en los coeficientes, es decir, la totalidad de las formas que surgen de la forma inicial por plegamiento en sí misma, ordenadas según formas superiores''.

Para explicar las formas superiores, supongamos por el Teorema VII que la forma inicial está dada como producto de dos formas normales $S\cdot T$. Por forma superior entendemos formas con plegamiento superior en sí mismas, y entre formas con el mismo plegamiento en sí mismas, unas especialmente normalizadas. Más precisamente: supongamos que la forma $A$ ha surgido mediante $\alpha_\rho$ plegamientos fundamentales de tipo $\rho$, y la forma $B$ mediante $\beta_\rho$ tales plegamientos. Entonces $B$ es superior a $A$:
\begin{enumerate}
\item si, en el sistema de desigualdades $\beta_\rho\ge\alpha_\rho$, $\rho=1,2,\ldots,n-1$, el signo de desigualdad es estricto para al menos un valor de $\rho$;\footnote{Si en el sistema de desigualdades el signo $>$ aparece en parte y el signo $<$ en parte, las formas no son comparables, pues una forma que surge de otra por plegamiento tiene siempre un sistema de valores positivo de plegamientos fundamentales, y por tanto en nuestro caso valores positivos o nulos $\beta_\rho-\alpha_\rho$.}
\item si el signo de igualdad vale para todos los valores $\rho$, pero se han combinado menos filas simbólicas en un único producto matricial. Esta normalización resulta necesaria por las reducciones de \S\ 8. Así, por ejemplo, $\pair{S^{n-1}}{T_{n-1}}$ es superior a $\pair{S^{n-1}}{[S^1T^\rho]_{n-\rho-1}p_\rho}$;
\item si el signo de igualdad vale para todos los valores $\rho$ y el número de filas simbólicas combinadas en un producto matricial es el mismo, entonces $B$ se hace superior a $A$ por una normalización arbitraria en sí misma pero fijada de una vez por todas. Esta normalización es siempre posible por el número finito de formas pertenecientes a un sistema de valores $\alpha$; puede realizarse, por ejemplo, privilegiando la única forma $S$ (mayor número de filas simbólicas $S$, mayor suma de los índices de peso superiores de las filas $S$, etc.).
\end{enumerate}
Conviene imaginar la fila de formas como dispuesta en una variedad de puntos reticulares de dimensión $(n-1)$, donde las $n-1$ direcciones corresponden a los $n-1$ plegamientos fundamentales. A cada forma le corresponde entonces uno y sólo un sistema de valores $\alpha$, es decir, un punto reticular entero positivo, mientras que las distintas formas pertenecientes a un sistema de valores $\alpha$ quedan determinadas de manera única en su orden por la normalización anterior. Una forma arbitraria de la fila de formas da la forma inicial de una nueva fila de formas, contenida como parte de la fila total de formas; en efecto, está contenida como la parte que consta de las formas obtenidas a partir de la nueva forma inicial al avanzar en las $n-1$ direcciones positivas.

En virtud del Teorema VIII y de la teoría general de los desarrollos en serie, valen para las filas de formas exactamente los teoremas de reducción de los dominios binario y ternario (cf. disertación, \S\ 3). Derivaremos brevemente el teorema principal. Primero enunciamos la definición:

``En un sistema dado de formas, supóngase que cierto número finito de formas se ha reunido en un módulo. Llamamos reducible a una forma si puede expresarse mediante formas que tengan invariantes como factores, o mediante `formas superiores'; es decir, mediante formas superiores de la fila total de formas a la que pertenece la forma, o mediante formas que contengan los símbolos del módulo en orden superior. Una fila de formas reducible se llama reductor.'' Entonces:

\emph{Teorema IX:} Si la forma inicial de una fila de formas es reducible por haber surgido mediante plegamiento con un reductor, entonces toda la fila de formas que surge de esta forma inicial es reducible.

Para la prueba, obsérvese que una forma $h$ surgida por plegamiento de una forma $f$ con una segunda forma $g$ puede, a causa de (45.), considerarse como una forma $f$ que contiene varias filas de variables cogredientes $p_\rho,p'_\rho$. Tales formas pueden representarse, por la teoría general de los desarrollos en serie, como polares de los covariantes elementales de $f$, aplicando sucesivamente el desarrollo en serie a cada par de filas de variables cogredientes $p_\rho,p'_\rho$. Por (38.), los covariantes elementales que aparecen son las formas generadas por plegamiento cogrediente. Sin embargo, como el orden en que se aplica el desarrollo en serie sigue siendo arbitrario, podemos escoger en particular un orden correspondiente a la generación de los plegamientos generales a partir de plegamientos fundamentales. El sistema resultante de covariantes elementales es entonces idéntico a la fila de formas $f$; las formas que surgen por plegamiento cogrediente de defecto superior se ordenan entre las que surgen por plegamientos fundamentales. Pero también se obtienen polares de la fila de formas $f$ para cualquier orden del desarrollo en serie, al que puede corresponder un segundo sistema de covariantes elementales. Como la fila de formas agota la totalidad de las formaciones invariantes, toda forma del segundo sistema puede expresarse linealmente por polares de la fila de formas, y en efecto por polares de aquellas formas que muestran la misma o una mayor disminución de peso. Esto último se sigue porque (cf. \S\ 7) los polares pueden considerarse como invariantes simultáneos de una fila de formas $H$ (52.), con los números de grado correspondientes a la forma $h$, y de la fila de formas $f$. Cada plegamiento de $H$ en sí misma corresponde a una disminución de peso que, después de efectuar la sustitución (11.), se transfiere a las filas simbólicas y por tanto a las formas de $f$.\footnote{Esta consideración subyace también a la segunda formulación del concepto de equivalencia (\S\ 8, final). Otras explicaciones de los diferentes sistemas de covariantes elementales en el dominio ternario se encuentran en Study, loc. cit., II, \S\ 7, Teoremas 7) y 8). Por nuestro Teorema VII los mismos resultados valen también para $n$ variables.}

Una forma arbitraria de la fila de formas $h$ ha surgido por plegamiento de $h$ en sí misma; es decir, por una parte mediante un plegamiento superior de $g$ con $f$, y por otra mediante plegamiento de $f$ o $g$ en sí mismas. En ambos casos el desarrollo en serie conduce a polares de la fila de formas $f$, tal como se mostró antes para la forma inicial $h$. Si en particular $f$ es un reductor, resulta un desarrollo en polares del reductor; por tanto la fila de formas $h$ se vuelve reducible.

Las aplicaciones del teorema a la reducción de sistemas completos de formas se siguen análogamente a los dominios binario y ternario.

\clearpage
% ---- Continuation appended 2026-06-02: Paper 05 complete and Paper 06 through Section 6, Spanish ----
\providecommand{\Kfield}{\mathfrak{K}}
\providecommand{\Ssys}{\mathfrak{S}}
\providecommand{\Mfield}{\mathfrak{M}}
\providecommand{\tq}{\tau}
\providecommand{\Lsys}{\mathfrak{L}}
\providecommand{\Jdom}{\mathfrak{J}}
\providecommand{\ds}{\displaystyle}

\begin{center}
{\Large\bfseries 5. Cuerpos de funciones racionales}\par
\vspace{1.2em}
\emph{Jahresbericht der DMV} 22 (1913), pp. 316--319
\end{center}

\vspace{1.5em}

Las cuestiones siguientes se remontan originalmente a conversaciones con E. Fischer. Algunas de ellas, por lo demás, ya habían sido planteadas y resueltas por E. Steinitz para el caso especial de una indeterminada, e incluso bajo hipótesis más generales sobre el dominio de coeficientes.\footnote{E. Steinitz: \emph{Algebraische Theorie der Körper}, especialmente \S\ 24. Crelle's Journal 137, 1910.}

Por ``cuerpo de funciones racionales'' entiendo un cuerpo cuyos elementos son funciones racionales de $n$ indeterminadas, con coeficientes en un cuerpo numérico dado, que en particular puede comprender también todos los números complejos. Ejemplos de tales cuerpos son el cuerpo de las funciones simétricas de $n$ cantidades, o, más generalmente, la totalidad de las funciones racionales de $n$ indeterminadas que admiten las permutaciones de cierto grupo (los \emph{Gattungsbereiche} de Lagrange); además, el cuerpo de invariantes.

Entre los dos primeros cuerpos mencionados y el último existe una diferencia característica: los primeros contienen $n$ funciones algebraicamente independientes, las funciones simétricas elementales; mientras que el número de invariantes algebraicamente independientes es siempre menor que el número de indeterminadas. Por eso es importante que, en general, se puedan asociar de manera biunívoca tales cuerpos de la segunda especie con cuerpos de la primera especie, sustituyendo algunas de las indeterminadas por números. Por consiguiente, en lo que sigue me limitaré a cuerpos de la primera especie, es decir, a cuerpos con $n$ funciones algebraicamente independientes; en virtud de esa asociación, los resultados valen entonces en general.

Las cuestiones se agrupan alrededor de tres conceptos de base: base racional, base mínima y base de integridad.

1. Por una ``base racional'' entiendo un número finito de funciones del cuerpo tal que toda función del cuerpo puede representarse como combinación racional de ese número finito, con coeficientes en el dominio de coeficientes prescrito. Mediante consideraciones sencillas, que en el caso de una indeterminada se encuentran también en E. Steinitz, loc. cit., se demuestra la existencia de una base racional para todo cuerpo de funciones racionales. Por ejemplo, según el teorema de Lagrange, las funciones simétricas elementales y una función perteneciente al grupo forman una base racional para los \emph{Gattungsbereiche} de Lagrange antes mencionados.

2. Toda base racional debe contener, para cuerpos de la primera especie, al menos $n$ funciones; si contiene exactamente $n$ funciones, que entonces han de ser algebraicamente independientes, la llamaré ``base mínima''. La cuestión de la existencia de una base mínima puede responderse parcialmente mediante teoremas de Lüroth, Castelnuovo y Enriques.\footnote{J. Lüroth: Beweis eines Satzes über rationale Kurven. \emph{Math. Ann.} 9, 1875. --- G. Castelnuovo: Sulla razionalità delle involuzioni piane. \emph{Math. Ann.} 44, 1893. --- F. Enriques: Sopra una involuzione non razionale dello spazio. \emph{Rend. Acc. Linc.} vol. 21, 21 de enero de 1912.} Según ellos, para cuerpos de una indeterminada existe siempre una base mínima: la función de Lüroth del cuerpo, determinada de manera única salvo transformación lineal fraccionaria (cf. Steinitz \S\ 24). Para cuerpos de dos indeterminadas existe siempre una base mínima, y en general sólo entonces, si se permite una extensión algebraica del cuerpo de coeficientes; mientras que para tres o más indeterminadas, en general, no existe en absoluto una base mínima. Todavía no se ha intentado caracterizar los cuerpos especiales que tienen base mínima.

He seguido estudiando la cuestión de la base mínima para los \emph{Gattungsbereiche} de Lagrange. Aquí adquiere el significado siguiente: ``si el \emph{Gattungsbereich} de Lagrange perteneciente al grupo $G$ posee una base mínima, entonces se puede construir racionalmente, por representación paramétrica, la totalidad de las ecuaciones con grupo prescrito $G$; y esto para todo cuerpo numérico que contenga el dominio de coeficientes de la base mínima, por tanto en particular para cualquier cuerpo numérico arbitrario si dicho dominio de coeficientes es el cuerpo de los números racionales.''

El ejemplo más sencillo lo proporciona el grupo simétrico; aquí las funciones simétricas elementales forman una base mínima con coeficientes numéricos racionales. De ello se obtiene como representación paramétrica simplemente la ecuación con coeficientes indeterminados. El teorema de irreducibilidad de Hilbert muestra cómo se pasa de ésta a arbitrariamente muchas ecuaciones no afectas sobre todo cuerpo numérico.

Ahora bien, de los teoremas de Lüroth y Castelnuovo se deduce directamente la existencia de la base mínima para todos los grupos que aparecen en las ecuaciones de grados 3 y 4; al mismo tiempo se ve además que esta base mínima tiene coeficientes numéricos racionales. Por tanto, sobre cualquier cuerpo numérico arbitrario, se puede construir racionalmente la totalidad de las ecuaciones de grados 3 y 4 con grupo prescrito. Para los grupos cíclicos y diedrales existe una base mínima con el cuerpo de las raíces $n$-ésimas de la unidad como dominio de coeficientes; lo mismo vale para algunos otros grupos metacíclicos. La cuestión de si la base mínima es siquiera característica de ciertos grupos debe quedar indecisa.

3. Para llegar al último concepto de base, considero la totalidad de los polinomios contenidos en el cuerpo. Por una ``base de integridad'' entiendo un número finito de esos polinomios tal que todo polinomio del cuerpo pueda representarse como combinación racional entera de ese número finito, con coeficientes en el dominio de coeficientes prescrito. La cuestión de la existencia de una base de integridad, que por ejemplo contiene como caso especial la finitud del sistema de invariantes, fue planteada por Hilbert en sus ``Problemas matemáticos''\footnote{D. Hilbert: Mathematische Probleme. Conferencia de París, 1900. Problema 14. \emph{Göttinger Nachr.} 1900.} en una forma algo distinta, como problema de las funciones relativamente enteras.

Por ahora sólo puedo indicar una clase de cuerpos para los cuales existe una base de integridad. En efecto, si el cuerpo contiene un sistema de $n$ polinomios tal que la resultante homogénea de los términos de dimensión más alta de esos polinomios es distinta de cero, entonces existe una base de integridad; está dada por los coeficientes de todos los $u$ en la ecuación irreducible, dentro del cuerpo, de la forma lineal:
\[
  u_0=u_1x_1+u_2x_2+\cdots+u_nx_n.\footnotemark
\]
\footnotetext{Esta ecuación irreducible aparece, en el caso de una indeterminada, en la demostración de Steinitz del teorema de Lüroth.}
Si, en particular, el valor de la resultante es igual a una unidad del dominio de coeficientes, queda asegurada también la integralidad de la representación. Un ejemplo de esta clase de cuerpos lo dan los \emph{Gattungsbereiche} de Lagrange, pues la resultante de las funciones simétricas elementales tiene el valor $\pm 1$. Otro ejemplo (sin integralidad) lo dan todos los cuerpos que contienen sólo un polinomio algebraicamente independiente; aquí la base de integridad es idéntica a la función de Lüroth del subcuerpo más pequeño que contiene todos los polinomios. Así, como es sabido, todos los invariantes proyectivos racionales enteros de una forma cuadrática en $n$ variables son funciones enteras de la discriminante.

La condición indicada es sólo suficiente, de ningún modo necesaria, para la existencia de la base de integridad. Se pueden construir fácilmente cuerpos en los que la resultante de cualesquiera $n$ polinomios se anula y que, no obstante, poseen una base de integridad dada por los coeficientes de aquella ecuación irreducible. Surge la pregunta de si quizá también en el caso más general los coeficientes de aquella ecuación suministran la base de integridad.

\clearpage

\begin{center}
{\Large\bfseries 6. Cuerpos y sistemas de funciones racionales}\par
\vspace{1.2em}
\emph{Mathematische Annalen} 76 (1915), pp. 161--196
\end{center}

\vspace{1.5em}

El presente trabajo trata cuestiones de base para sistemas arbitrarios de funciones racionales y racionales enteras. Los métodos de la teoría de cuerpos aquí empleados hacen que la resolución de estas cuestiones para cuerpos de funciones racionales --- cuerpos de funciones racionales\footnote{Cf. una comunicación preliminar con ocasión de la reunión de naturalistas de Viena: \emph{Rationale Funktionenkörper}, \emph{Jahresber. d. D. Math.-Ver.} 22 (1913), donde se da una visión de conjunto de las cuestiones y de los resultados referentes a los cuerpos de funciones.} --- aparezca como lo esencial, mientras que la generalización de los resultados a sistemas arbitrarios se obtiene como consecuencia.

Entre las cuestiones de base para sistemas generales, hasta ahora sólo se conocía la existencia de una base de módulo, garantizada por el teorema de Hilbert (\emph{Math. Ann.} 36). En lo que sigue, la cuestión de la representabilidad racional se responde completamente mediante la existencia de una base racional para todo sistema arbitrario (\S\ 7); la base racional de los cuerpos se obtiene ya en \S\ 4. Esta existencia de la base racional permite partir siempre del cuerpo o sistema definido abstractamente, y evitar así dificultades que son causadas sólo por la elección especial de la base racional, y no por el sistema mismo, tales como la aparición de denominadores especiales o de puntos fundamentales de las funciones de la base.

Para cuerpos se trata además la cuestión de la base mínima, es decir, de una base racional compuesta de funciones algebraicamente independientes (\S\ 6). Los métodos de la teoría de cuerpos conducen además a una base racional distinguida, la base de involución\footnote{El nombre se eligió porque, en su interpretación geométrica, el cuerpo de funciones racionales representa la involución más general en un espacio lineal.} (\S\ 5), que resulta especialmente importante para dominios de integridad compuestos de polinomios (\S\ 8). En ese contexto da una representación con denominador fijo (una potencia de una función determinada por el dominio de integridad), tal como se conoce, por ejemplo, en el caso especial de la representación típica de invariantes.

De la representabilidad racional extraemos ahora, en la medida de lo posible, consecuencias para la representabilidad racional entera, es decir, para la cuestión de finitud en sentido estricto.

Los teoremas de finitud conocidos hasta ahora descansan todos en el hecho de que el teorema de Hilbert sobre bases de módulos garantiza la finitud para todos los sistemas en los que una representación
\[
  F=A_1f_1+A_2f_2+\cdots+A_kf_k
\]
entraña una segunda representación
\[
  F=B_1f_1+B_2f_2+\cdots+B_kf_k,
\]
donde los $B_i$ pertenecen al sistema y son de grado menor que $F$.

En cambio, el teorema sobre la base racional y los métodos de la teoría de cuerpos permiten inferir la finitud bajo hipótesis de una especie completamente distinta.

Así, como respuesta parcial a un problema de Hilbert (Problemas matemáticos 14), se obtiene una clase de funciones relativamente enteras (dominios relativamente enteros de primera especie) que son dominios de integridad finitos y cuya base de integridad está dada por la base de involución antes mencionada (\S\ 10). En analogía con la demostración de Hilbert del teorema de Kronecker sobre el sistema fundamental de enteros algebraicos (\emph{Complete invariant systems}, \S\ 2; \emph{Math. Ann.} 42), se puede mostrar además que todos los sistemas regulares de polinomios, es decir, todos los sistemas que pueden aplicarse en otros sin puntos fundamentales, poseen una base de integridad (\S\ 12), mientras que se pueden especificar fácilmente sistemas no regulares sin base de integridad. Como ejemplo sencillo de este hecho puede mencionarse el teorema: ``en todo sistema arbitrario de infinitos polinomios en una indeterminada existe un número finito de esos polinomios tal que todo polinomio del sistema puede representarse como combinación racional entera de ese número finito.'' Otros ejemplos se encuentran en \S\ 13.

Finalmente, los últimos parágrafos (\S\S\ 14 y 15) dan un refinamiento de los teoremas de base con respecto a la integralidad.

Un instrumento esencial de la investigación es el principio de transferencia de \S\ 3, según el cual basta considerar todas las cuestiones de base aquí planteadas para aquellos sistemas en los que el número de indeterminadas coincide con el número de funciones algebraicamente independientes del sistema.

El impulso para el presente trabajo vino de conversaciones con el señor E. Fischer, en particular de una pregunta suya sobre la base mínima de los \emph{Gattungsbereiche} de Lagrange. Las investigaciones relativas a ello, que condujeron a la construcción de ecuaciones con grupo prescrito (cf. la comunicación citada sobre ``Cuerpos de funciones racionales'', parte 2), quedan reservadas para una publicación ulterior.

La ampliación esencial del presente trabajo respecto de la comunicación original consiste en que los teoremas de base allí enunciados sólo para cuerpos o dominios de integridad especiales se trasladan aquí a sistemas arbitrarios. Esta transferencia se logra de manera sencilla por medio del concepto del menor cuerpo que contiene un sistema arbitrario, como generalización del cuerpo de fracciones de un dominio de integridad (\S\ 7). Me condujo a formar este concepto el hecho de que el señor K. Hentzelt, después de conocer la base racional de los cuerpos, pudo demostrar la existencia de la base racional para sistemas arbitrarios dando un rodeo por la base de módulo de Hilbert. También me llevó al teorema sobre la base de integridad de los sistemas regulares la observación de K. Hentzelt de que los razonamientos que yo había aplicado a dominios de integridad eran válidos para sistemas arbitrarios sujetos a las mismas hipótesis; asimismo debo al señor Hentzelt algunas observaciones aisladas de menor tamaño.

Finalmente, se menciona que en el caso de una indeterminada la base racional y la base mínima de los cuerpos se encuentran en la \emph{Algebraische Theorie der Körper} de E. Steinitz, \S\ 24 (Crelle's Journal 137); allí el dominio de coeficientes se toma como cuerpo en el sentido más general. La cuestión de hasta qué punto los teoremas dados aquí siguen siendo válidos bajo esas hipótesis más generales no se trata en lo que sigue; en cualquier caso, las demostraciones requieren modificación, ya que, por ejemplo, los medios de la teoría de los determinantes funcionales fallan para cuerpos de característica $p$.

\section*{\S\ 1. Cuerpos $\Kfield_{n\rho}$ y sistemas $\Ssys_{n\rho}$.}

En lo que sigue se trata de investigar sistemas de funciones racionales de $n$ indeterminadas, en particular cuerpos de funciones racionales.

Por $f(x_1,\ldots,x_n)$, $g(x_1,\ldots,x_n)$, \ldots, o abreviadamente $f(x)$, $g(x)$, \ldots, se entenderán siempre funciones racionales de $x_1,\ldots,x_n$; se suponen en representación reducida, es decir, con numerador y denominador primos entre sí. Las funciones racionales enteras (polinomios) se designarán con letras mayúsculas, $F(x)$, $G(x)$, \ldots. El dominio de coeficientes $\Omega$ se supone algún cuerpo numérico; en particular puede contener también todos los números complejos. El cuerpo $\Omega$ puede contener también un número finito de parámetros.

Las cantidades de $\Omega$, por tanto todas las funciones de grado cero, se cuentan siempre como pertenecientes al sistema.

Con estas estipulaciones, el cuerpo formado por funciones racionales --- el cuerpo de funciones racionales --- puede definirse abstractamente.

\emph{Definición I: Un sistema de funciones racionales se llama cuerpo si satisface las condiciones:}
\begin{enumerate}
\item \emph{Junto con $f(x)$, y para cada cantidad $c$ de $\Omega$, contiene también $c\cdot f(x)$.}
\item \emph{Junto con $f(x)$ y $g(x)$, contiene siempre también $f(x)+g(x)$, $f(x)\cdot g(x)$ y, para $g(x)\ne0$, el cociente $f(x):g(x)$.}\footnote{Aquí $f(x)$ y $g(x)$ pueden ser también de grado cero, es decir, cantidades de $\Omega$.}
\end{enumerate}
Los cuerpos especiales son los que surgen por adjunción de un número finito de funciones racionales
\[
  f_1(x),\ f_2(x),\ldots, f_k(x)
\]
a $\Omega$; se denotarán por
\[
  \Omega\bigl(f_1(x),\ldots,f_k(x)\bigr)
  \quad\text{o}\quad
  \Omega(f_1,\ldots,f_k).
\]
\footnote{Que estos ``cuerpos especiales'' agotan ya la totalidad de todos los cuerpos se deduce en \S\ 4.}
Entre estos cuerpos especiales mencionamos también el cuerpo que surge por adjunción de $x_1,\ldots,x_n$ a $\Omega$,
\[
  \Omega(x_1,\ldots,x_n),
\]
que es idéntico a la totalidad de las funciones racionales de $x_1,\ldots,x_n$ con coeficientes en $\Omega$.

Introducimos además, siguiendo a E. Steinitz, el concepto de cuerpo intermedio:
\begin{quote}
``un cuerpo $\Mfield$ está entre $\Omega_1$ y $\Omega_2$ si contiene a $\Omega_1$ y está contenido en $\Omega_2$;''
\end{quote}
entonces la Definición I admite también la siguiente formulación:

El cuerpo de funciones racionales de $n$ indeterminadas es un cuerpo intermedio entre $\Omega$ y $\Omega(x_1,\ldots,x_n)$ que contiene efectivamente las $n$ indeterminadas en al menos una función.

Junto al número de indeterminadas aparece para sistemas de funciones racionales otro número característico, el número de funciones algebraicamente independientes, o rango algebraico, mediante la siguiente definición.

\emph{Definición II: Un sistema de funciones racionales tiene rango algebraico $\rho$ si se pueden indicar $\rho$ funciones del sistema que sean algebraicamente independientes, pero cualesquiera $(\rho+1)$ funciones del sistema son algebraicamente dependientes.}\footnote{Se llaman algebraicamente independientes $\rho$ funciones $f_1(x),\ldots,f_\rho(x)$ si de toda relación
\[
F\bigl(f_1(x),\ldots,f_\rho(x)\bigr)=0\quad\text{idénticamente en }x_1,\ldots,x_n
\]
se sigue que
\[
F(\lambda_1,\ldots,\lambda_\rho)=0\quad\text{idénticamente en }\lambda_1,\ldots,\lambda_\rho.
\]
En el caso contrario se llaman algebraicamente dependientes.}

Los sistemas de (finitas o infinitas) funciones racionales en $n$ indeterminadas y de rango algebraico $\rho$ se denotarán mediante el doble índice
\[
  \Ssys_{n\rho}.
\]
En particular, por
\[
  \Kfield_{n\rho}
\]
se entenderá un cuerpo de $n$ indeterminadas y de rango algebraico $\rho$. Se tiene la desigualdad
\[
  1\leq \rho\leq n.
\]

Damos aún algunos ejemplos de cuerpos $\Kfield_{nn}$ y $\Kfield_{n\rho}$:

1. Cuerpos $\Kfield_{nn}$ son los \emph{Gattungsbereiche} de Lagrange, que pueden definirse como
\begin{quote}
``la totalidad de las funciones racionales (y racionales enteras) de $x_1,\ldots,x_n$ que admiten las permutaciones de un grupo de permutaciones $G$ de $x_1,\ldots,x_n$.''
\end{quote}
Siempre contienen el cuerpo de las funciones simétricas, y por tanto también las $n$ funciones simétricas elementales algebraicamente independientes.

2. Cuerpos $\Kfield_{n\rho}$ son los cuerpos de invariantes proyectivos, formados por la totalidad de los invariantes racionales (y racionales enteros) de una o varias formas fundamentales.\footnote{Conviene considerar sólo transformaciones de determinante $1$; entonces toda combinación racional de arbitrariamente muchos invariantes admite de nuevo la transformación, y efectivamente se obtiene un cuerpo. Los invariantes homogéneos e isobáricos tomados por sí solos no forman un cuerpo, pero sí un sistema $\Ssys_{n\rho}$.} Aquí $\rho<n$, pues el número de invariantes algebraicamente independientes es siempre menor que el número de coeficientes.

Lo correspondiente vale para los cuerpos de invariantes respecto de subgrupos del grupo proyectivo.

\section*{\S\ 2. Lema auxiliar sobre funciones algebraicamente dependientes.}

El lema auxiliar de este parágrafo mostrará en \S\ 3 que el rango algebraico de un sistema es su número propiamente característico. En efecto, el lema auxiliar muestra que mediante una especialización numérica sucesiva de algunas indeterminadas, de tal modo que $\rho$ funciones algebraicamente independientes sigan siendo algebraicamente independientes, una función arbitraria algebraicamente dependiente de esas $\rho$ funciones no puede ni anularse idénticamente ni volverse indeterminada.

Sea, pues, el sistema
\begin{equation}
  g_1(x),\ g_2(x),\ldots,g_\rho(x)
\tag{1}
\end{equation}
de rango algebraico $\rho$, donde $\rho<n$, y sea $h(x)$ una función racional algebraicamente dependiente de las funciones (1). Por teoremas conocidos, el rango de la matriz funcional de (1) es igual a $\rho$. Por tanto las indeterminadas pueden llamarse, por ejemplo, de manera que
\begin{equation}
 \frac{\partial(g_1,g_2,\ldots,g_\rho)}{\partial(x_1,x_2,\ldots,x_\rho)}
 =\frac{\Gamma(x)}{G(x)^{\rho+1}}\ne0,
\tag{2}
\end{equation}
donde $G(x)$, como es usual, denota el denominador común de las funciones (1).

\emph{El número $a$ se elige ahora de modo que el producto $G(x)\cdot\Gamma(x)$ no sea divisible por $(x_i-a)$}, entendiendo por $i$ uno de los números $\rho+1,\rho+2,\ldots,n$. Así, para $x_i=a$, el denominador común de las funciones (1) no puede anularse, y las $\rho$ funciones
\begin{equation}
  [g_1(x)]_{x_i=a},\ldots,[g_\rho(x)]_{x_i=a}
\tag{3}
\end{equation}
son igualmente algebraicamente independientes. Como diremos, el rango algebraico del sistema (1) se conserva bajo la sustitución $x_i=a$.

La ecuación irreducible que satisface $h(x)$ respecto de las funciones (1) sea
\begin{equation}
\begin{aligned}
&A_0(g_1,\ldots,g_\rho)\,h(x)^\tau
 +A_1(g_1,\ldots,g_\rho)\,h(x)^{\tau-1}
 +\cdots \\
&\hspace{6em}+A_\tau(g_1,\ldots,g_\rho)=0,
\end{aligned}
\tag{4}
\end{equation}
donde $A_0(g_1,\ldots,g_\rho)$ y $A_\tau(g_1,\ldots,g_\rho)$ no se anulan idénticamente. Para pasar de aquí a una identidad entre polinomios, ponemos
\begin{equation}
  g_i(x)=G_i(x):G(x);\qquad h(x)=H_1(x):H_2(x),
\tag{5}
\end{equation}
y observamos que, por la representación reducida supuesta de $h(x)$, $H_1(x)$ y $H_2(x)$ son primos entre sí. En virtud de (5), (4) se transforma en
\begin{equation}
\begin{aligned}
&B_0(G,G_1,\ldots,G_\rho)\,G(x)^{\sigma_0}H_1(x)^\tau \\
&\quad+B_1(G,G_1,\ldots,G_\rho)\,G(x)^{\sigma_1}H_1(x)^{\tau-1}H_2(x)+\cdots \\
&\quad+B_\tau(G,G_1,\ldots,G_\rho)\,G(x)^{\sigma_\tau}H_2(x)^\tau=0,
\end{aligned}
\tag{6}
\end{equation}
donde los $B_k$ son formas homogéneas en $G,G_i$, y al menos uno de los exponentes no negativos $\sigma_k$ debe ser cero.

Supongamos ahora que $H_1(x)$ fuese divisible por $(x_i-a)$. Entonces, por (6),
\[
  B_\tau(G,G_1,\ldots,G_\rho)\,G(x)^{\sigma_\tau}H_2(x)^\tau
\]
sería también divisible por $(x_i-a)$. Esto está excluido por hipótesis para $G(x)$; por tanto, de la divisibilidad de $B_\tau$ se seguiría también la divisibilidad de
\[
  A_\tau=B_\tau:G^\lambda,
\]
y entre las funciones (3) existiría la relación $A_\tau=0$, contra la independencia algebraica supuesta. Finalmente $H_2(x)$, por ser primo con $H_1(x)$, tampoco puede ser divisible por $(x_i-a)$;\footnote{Esta última conclusión, basada en la representación reducida, ya no sería posible para una sustitución $x_i=a$, $x_k=b$ en la que se conservaran las restantes hipótesis. Entonces $H_1(x)$ y $H_2(x)$ podrían anularse simultáneamente, y el cociente se volvería indeterminado bajo la sustitución, igual a $\frac{0}{0}$; por ejemplo
\[
 h(x)=\frac{x_3(\alpha x_1+\beta x_2)+x_4(\gamma x_1+\delta x_2)}
 {x_3(\alpha' x_1+\beta' x_2)+x_4(\gamma' x_1+\delta' x_2)}
\]
bajo la sustitución $x_3=0$, $x_4=0$.} así, la hipótesis de que $H_1(x)$ sea divisible por $(x_i-a)$ conduce a una contradicción. Del mismo modo se muestra que $H_2(x)$ tampoco puede ser divisible por $(x_i-a)$. Por consiguiente, la función $[h(x)]_{x_i=a}=h'(x)$ no puede ni anularse idénticamente ni volverse indeterminada.

La conclusión precedente puede aplicarse ahora de nuevo al sistema (3) y a la función $h'(x)$, puesta otra vez en representación reducida; continuando el procedimiento se llega finalmente al siguiente lema auxiliar.

\emph{Lema auxiliar: Los dos sistemas}
\[
  g_1(x),\ g_2(x),\ldots,g_\rho(x),
\]
\[
  g_1(x),\ g_2(x),\ldots,g_\rho(x),\ h(x)
\]
\emph{tienen cada uno rango algebraico $\rho$, donde $\rho<n$, y supóngase, por ejemplo, que}
\[
 \frac{\partial(g_1,\ldots,g_\rho)}{\partial(x_1,\ldots,x_\rho)}
 =\frac{\Gamma(x)}{G(x)^{\rho+1}}\ne0.
\]
\emph{Los números}
\[
  a_n,\ a_{n-1},\ldots,a_{\rho+1}
\]
\emph{deben elegirse de tal manera que, bajo la sustitución}
\[
  x_n=a_n,\quad x_{n-1}=a_{n-1},\ldots,\quad x_{\rho+1}=a_{\rho+1},
\]
\emph{el producto $G(x)\cdot\Gamma(x)$ no se anule, es decir, que el rango algebraico del sistema (1) se conserve. En las sustituciones sucesivas}
\[
\begin{aligned}
 [h(x)]_{x_n=a_n}&=h^{(n-1)}(x),\qquad
 [h^{(n-1)}(x)]_{x_{n-1}=a_{n-1}}=h^{(n-2)}(x),\ldots,\\
 [h^{(\rho+1)}(x)]_{x_{\rho+1}=a_{\rho+1}}&=h^*(x_1,x_2,\ldots,x_\rho),
\end{aligned}
\]
\emph{pónganse las funciones $h(x)$, $h^{(n-1)}(x)$, \ldots, $h^{(\rho+1)}(x)$ cada vez en representación reducida. Bajo estas hipótesis, en la función $h^*(x_1,\ldots,x_\rho)$ ni el numerador ni el denominador pueden anularse idénticamente, y la función tampoco puede convertirse en $\frac{0}{0}$.}\footnote{Por ejemplo, en la función $h(x)$ de la nota precedente, la sustitución sucesiva da
\[
 [h(x)]_{x_4=0}=h^{(3)}(x)=\frac{\alpha x_1+\beta x_2}{\alpha' x_1+\beta' x_2},\qquad
 h^*(x_1,x_2)=\frac{\alpha x_1+\beta x_2}{\alpha' x_1+\beta' x_2}.
\]}

\section*{\S\ 3. Aplicación biunívoca de sistemas $\Ssys_{n\rho}$ sobre sistemas $\Ssys^*_{\rho\rho}$.}

Del lema auxiliar obtendremos de inmediato una aplicación de los sistemas $\Ssys_{n\rho}$ sobre sistemas $\Ssys^*_{\rho\rho}$, sustituyendo algunas indeterminadas por números; con ello el rango algebraico $\rho$ se presenta como el número propiamente característico.

Puesto que $\Ssys_{n\rho}$ tiene rango algebraico $\rho$, existen $\rho$ funciones de $\Ssys_{n\rho}$,
\begin{equation}
  g_1(x),\ldots,g_\rho(x),
\tag{1}
\end{equation}
que son algebraicamente independientes. Además, si $f(x)$ denota una función arbitraria de $\Ssys_{n\rho}$, entonces el sistema
\begin{equation}
  g_1(x),\ldots,g_\rho(x),\ f(x)
\tag{2}
\end{equation}
tiene también rango algebraico $\rho$, y los sistemas (1) y (2) satisfacen las condiciones del lema auxiliar.\footnote{Si es necesario, deben renumerarse las indeterminadas.}

Si, por tanto, se determinan los números
\[
  a_n,\ a_{n-1},\ldots,a_{\rho+1}
\]
según el lema auxiliar, de manera que el rango algebraico $\rho$ del sistema (1) se conserve por la sustitución, lo cual es posible de arbitrariamente muchas maneras, entonces en la función definida por
\[
\begin{aligned}
 [f(x)]_{x_n=a_n}&=f^{(n-1)}(x),\qquad
 [f^{(n-1)}(x)]_{x_{n-1}=a_{n-1}}=f^{(n-2)}(x);\ldots,\\
 [f^{(\rho+1)}(x)]_{x_{\rho+1}=a_{\rho+1}}&=f^*(x_1,\ldots,x_\rho),
\end{aligned}
\]
a saber,
\begin{equation}
  f^*(x_1,\ldots,x_\rho),
\tag{3}
\end{equation}
ni el numerador ni el denominador pueden anularse idénticamente.

Hagamos ahora que $f(x)$ recorra todas las funciones, finitas o infinitas, de $\Ssys_{n\rho}$. Consideramos el sistema que consta de la totalidad de todas las funciones (3); tiene las siguientes propiedades:

1. Tiene rango algebraico $\rho$; pues contiene las funciones provenientes de (1),
\[
  g_1^*(x_1,\ldots,x_\rho),\ldots,g_\rho^*(x_1,\ldots,x_\rho),
\]
que, por la elección de $a_n,a_{n-1},\ldots,a_{\rho+1}$, son algebraicamente independientes. Por tanto el sistema debe denotarse por
\[
  \Ssys^*_{\rho\rho}.
\]
\footnote{Correspondientemente, por $f^*(x)$, $g^*(x)$, \ldots, se entenderán en todo momento funciones de aplicación, es decir, funciones de $x_1,\ldots,x_\rho$.}

2. La correspondencia de las funciones $f(x)$ y $f^*(x)$ es sin excepción biunívoca.

En efecto, si $f_1(x)$ y $f_2(x)$ pertenecen a $\Ssys_{n\rho}$, entonces su diferencia
\[
  d(x)=f_1(x)-f_2(x)
\]
es también algebraicamente dependiente del sistema (1); además,
\[
  d^*(x)=f_1^*(x)-f_2^*(x).
\]
Como $d(x)$, junto con el sistema (1), satisface las condiciones del lema auxiliar, de
\[
  d(x)\ne0
\]
se sigue necesariamente
\[
  d^*(x)\ne0;
\]
en otras palabras, a funciones distintas de $\Ssys_{n\rho}$ corresponden funciones distintas de $\Ssys^*_{\rho\rho}$.

3. De toda relación entre funciones de $\Ssys^*_{\rho\rho}$,
\[
  F\bigl(f_1^*(x),f_2^*(x),\ldots,f_k^*(x)\bigr)=0
  \quad\text{idénticamente en }x_1,\ldots,x_\rho,
\]
se sigue, para las funciones unívocamente correspondientes de $\Ssys_{n\rho}$, que
\[
  F\bigl(f_1(x),f_2(x),\ldots,f_k(x)\bigr)=0
  \quad\text{idénticamente en }x_1,\ldots,x_n.
\]
Pues con $f_1(x),f_2(x),\ldots,f_k(x)$ también toda combinación racional de estas funciones es algebraicamente dependiente del sistema (1).

Además, esto se sigue de
\begin{equation}
\begin{aligned}
 h(x)&=F\bigl(f_1(x),\ldots,f_k(x)\bigr):\\
 h^*(x)&=F\bigl(f_1^*(x),\ldots,f_k^*(x)\bigr).
\end{aligned}
\tag{4}
\end{equation}
Por el lema auxiliar, por tanto,
\[
  h(x)\ne0\quad\text{implica}\quad h^*(x)\ne0;
\]
q.e.d.

En resumen obtenemos:

\emph{Teorema I: A todo sistema $\Ssys_{n\rho}$ de funciones racionales (finitas o infinitas) se le puede asignar, de arbitrariamente muchas maneras, un sistema biunívoco $\Ssys^*_{\rho\rho}$, de tal modo que todas las relaciones existentes entre funciones de $\Ssys^*_{\rho\rho}$ se conservan también en $\Ssys_{n\rho}$ y recíprocamente. En particular, por (4), a un cuerpo $\Kfield_{n\rho}$ corresponde de nuevo un cuerpo $\Kfield^*_{\rho\rho}$.}

Del Teorema I extraemos la consecuencia que simplifica todas las demostraciones posteriores: \emph{las propiedades de los sistemas $\Ssys_{n\rho}$ que se caracterizan por finitas o infinitas relaciones entre funciones del sistema valen para los sistemas $\Ssys_{n\rho}$ tan pronto como valen para un sistema de aplicación $\Ssys^*_{\rho\rho}$.}

Llamamos brevemente a esta consecuencia el ``principio de transferencia''.

\providecommand{\Lsys}{\mathfrak{L}}
\providecommand{\Jdom}{\mathfrak{J}}

\clearpage
El ejemplo más sencillo de esta aplicación lo proporciona la teoría de los invariantes algebraicos. En efecto, mediante una transformación lineal se pueden fijar numéricamente algunos coeficientes de la forma fundamental de tal modo que todas las relaciones válidas para la forma especial sigan siendo válidas en general.

\section*{\S\ 4. Base racional de los cuerpos $\Kfield_{n\rho}$.}

En las investigaciones siguientes, agrupadas alrededor de conceptos de base, utilizaremos de manera constante el ``principio de transferencia'' de \S\ 3. Primero demostramos la existencia de la base para cuerpos, pero enunciamos las definiciones en general:

\emph{Definición III: Un número finito de funciones de $\Ssys_{n\rho}$ se llama base racional si toda función de $\Ssys_{n\rho}$ puede representarse como combinación racional de ese número finito, con coeficientes de $\Omega$.}

La base racional debe contener $\rho$ funciones algebraicamente independientes, puesto que el rango algebraico de un sistema derivado de una base racional es igual al rango algebraico de la base racional.

Mostramos primero que la demostración de existencia de la base racional admite el principio de transferencia. En efecto, la Definición III puede formularse como sigue:

Las funciones de $\Ssys_{n\rho}$
\[
  g_1(x),\ g_2(x),\ldots, g_k(x)
\]
forman una base racional si y sólo si se satisfacen las infinitas relaciones
\begin{equation}
  f(x)=\gamma\bigl(g_1(x),\ldots,g_k(x)\bigr),
\tag{1}
\end{equation}
donde $f(x)$ recorre todas las funciones de $\Ssys_{n\rho}$ y $\gamma$ denota una función racional de $k$ argumentos, variable con $f$.

Así, si en el sistema de aplicación $\Ssys^*_{\rho\rho}$ se da una base racional por
\[
  g_1^*(x),\ g_2^*(x),\ldots,g_k^*(x),
\]
lo cual entraña las relaciones
\[
  f^*(x)=\gamma\bigl(g_1^*(x),\ldots,g_k^*(x)\bigr),
\]
entonces, por el Teorema I, las relaciones (1) se cumplen para $\Ssys_{n\rho}$. Así, el \emph{principio de transferencia para la base racional} dice:

\emph{De la existencia de una base racional para un sistema de aplicación $\Ssys^*_{\rho\rho}$ se sigue la base racional para el sistema original $\Ssys_{n\rho}$.}\footnote{El principio de transferencia se formula correspondientemente para toda cuestión de base tratada aquí.}

Para demostrar la existencia de la base racional de todos los cuerpos $\Kfield_{n\rho}$ podemos, por tanto, restringirnos a cuerpos $\Kfield^*_{\rho\rho}$; aquí una generalización directa de un argumento empleado por E. Steinitz conduce al objetivo.\footnote{E. Steinitz, \emph{Algebraische Theorie der Körper}, \S\ 24, Crelle's Journal 137 (1909).}

Sea
\begin{equation}
  g_1^*(x_1,\ldots,x_\rho),\ldots,
  g_\rho^*(x_1,\ldots,x_\rho)
\tag{2}
\end{equation}
un sistema de $\rho$ funciones algebraicamente independientes de $\Kfield^*_{\rho\rho}$. Eliminando $x_1,\ldots,x_\rho$ se obtienen $\rho$ relaciones:
\begin{equation}
\begin{aligned}
  F_1\bigl(g_1^*(x),\ldots,g_\rho^*(x),x_1\bigr)&=0,\ldots,\\
  F_\rho\bigl(g_1^*(x),\ldots,g_\rho^*(x),x_\rho\bigr)&=0,
\end{aligned}
\tag{3}
\end{equation}
donde en cada relación $F_i=0$ aparece efectivamente la indeterminada $x_i$; de lo contrario, contra la hipótesis, habría una dependencia algebraica entre las funciones (2). Las relaciones (3) dicen que
\[
  x_1,x_2,\ldots,x_\rho
\]
son elementos algebraicos respecto de
\[
  \Omega\bigl(g_1^*(x),\ldots,g_\rho^*(x)\bigr).
\]
Por consiguiente, el cuerpo que surge por adjunción de $x_1,\ldots,x_\rho$,
\[
  \Omega\bigl(g_1^*,\ldots,g_\rho^*,x_1,
  \ldots,x_\rho\bigr)=\Omega(x_1,
  \ldots,x_\rho),
\]
es una extensión algebraica finita de $\Omega(g_1^*,\ldots,g_\rho^*)$. Por un teorema conocido,\footnote{Compárese, por ejemplo, Weber, \emph{Lehrbuch d. Algebra} 1, \S\ 144 (1.ª ed.) o \S\ 55 (edición pequeña).} $\Kfield^*_{\rho\rho}$, como cuerpo intermedio entre $\Omega(g_1^*,\ldots,g_\rho^*)$ y $\Omega(x_1,\ldots,x_\rho)$, es él mismo una extensión algebraica de $\Omega(g_1^*,\ldots,g_\rho^*)$, mientras que $\Omega(x_1,\ldots,x_\rho)$ es una extensión algebraica de $\Kfield^*_{\rho\rho}$. Así, $\Kfield^*_{\rho\rho}$ se obtiene de $\Omega(g_1^*,\ldots,g_\rho^*)$ por adjunción de un número finito de elementos de $\Kfield^*_{\rho\rho}$, o también por adjunción de una función convenientemente elegida; es decir, existe una base racional. Por el principio de transferencia se obtiene entonces

\emph{Teorema II: Todo cuerpo $\Kfield_{n\rho}$ posee una base racional.}

\section*{\S\ 5. Forma de involución y base de involución de los cuerpos $\Kfield_{n\rho}$.}

Del argumento precedente obtenemos además una base racional determinada por el cuerpo mismo. La indicamos primero para los cuerpos $\Kfield^*_{\rho\rho}$.

Sea $\Kfield^*_{\rho\rho}$ un cuerpo entre $\Omega$ y $\Omega(x_1,\ldots,x_\rho)$. Como en \S\ 4, se eligen $\rho$ funciones algebraicamente independientes de $\Kfield^*_{\rho\rho}$,
\[
  g_1^*(x),\ldots,g_\rho^*(x),
\]
y se forma la extensión algebraica finita $\Omega(x_1,\ldots,x_\rho)$ de $\Omega(g_1^*,\ldots,g_\rho^*)$. Sea $i$ el grado de $\Omega(x_1,\ldots,x_\rho)$ respecto de $\Kfield^*_{\rho\rho}$.

Introducimos una forma lineal
\[
  u(x)=u_1x_1+\cdots+u_\rho x_\rho
\]
con coeficientes indeterminados $u_1,\ldots,u_\rho$, y sea
\begin{equation}
 \Phi^*(z,u)=z^i+\sum_{\alpha+\alpha_1+\cdots+\alpha_\rho=i}
 g^*_{\alpha\alpha_1\ldots\alpha_\rho}(x)
 z^\alpha u_1^{\alpha_1}\cdots u_\rho^{\alpha_\rho}
\tag{1}
\end{equation}
la ecuación irreducible de $u(x)$ respecto de $\Kfield^*_{\rho\rho}$.\footnote{Noether escribe la suma sobre todos los sistemas de índices con suma total $i$; el término principal $z^i$ se ha separado.} Los coeficientes
\[
  g^*_{\alpha\alpha_1\ldots\alpha_\rho}(x)
\]
pertenecen a $\Kfield^*_{\rho\rho}$ y están determinados por el cuerpo. Llamaremos a $\Phi^*(z,u)$ la forma de involución del cuerpo.

Afirmamos que estos coeficientes forman una base racional. En efecto, la ecuación (1), vista como ecuación de $z$, es irreducible respecto de $\Kfield^*_{\rho\rho}$ y por tanto tanto más irreducible respecto del subcuerpo
\[
  \Omega\bigl(g^*_{\alpha\alpha_1\ldots\alpha_\rho}(x)\bigr).
\]
Por consiguiente $\Phi^*(z,u)$ da la ecuación irreducible de $u(x)$ respecto de ese subcuerpo. De ello se sigue que $\Omega(x_1,\ldots,x_\rho)$ es una extensión algebraica --- de grado $i$ --- de $\Omega(g^*_{\alpha\alpha_1\ldots\alpha_\rho}(x))$. Como $\Kfield^*_{\rho\rho}$ está entre este subcuerpo y $\Omega(x_1,\ldots,x_\rho)$, la igualdad de grados implica, como es sabido, la igualdad de los cuerpos. Por el principio de transferencia, los coeficientes de $\Phi(z,u)$ suministran correspondientemente una base racional de $\Kfield_{n\rho}$. Así tenemos:

\emph{Teorema III: Los coeficientes de la forma de involución constituyen una base racional determinada por el cuerpo, la base de involución; para cuerpos $\Kfield^*_{\rho\rho}$ esta base está determinada de manera única.}\footnote{Para cuerpos $\Kfield_{n\rho}$ la forma de involución puede depender de la aplicación; la unicidad puede obtenerse mediante una aplicación con coeficientes indeterminados.}

Añadimos una interpretación irracional de este hecho, que establece la conexión con la teoría geométrica de las involuciones.

La factorización de $\Phi^*(z,u)$ en factores lineales viene dada en un cuerpo de extensión por
\begin{equation}
\begin{aligned}
 \Phi^*(z,u)&=(z-u_1x_1-\cdots-u_\rho x_\rho)
 (z-u_1x_1^{(1)}-\cdots-u_\rho x_\rho^{(1)})\cdots\\
 &\qquad\cdot
 (z-u_1x_1^{(i-1)}-\cdots-u_\rho x_\rho^{(i-1)}).
\end{aligned}
\tag{4}
\end{equation}
La comparación con (1) muestra que las funciones
\[
  g^*_{\alpha\alpha_1\ldots\alpha_\rho}(x_1,
  \ldots,x_\rho)\quad(\alpha+\alpha_1+\cdots+\alpha_\rho=i)
\]
son idénticas a las funciones simétricas elementales de las filas de cantidades
\begin{equation}
\begin{array}{cccc}
  x_1 & x_2 & \cdots & x_\rho\\
  x_1^{(1)} & x_2^{(1)} & \cdots & x_\rho^{(1)}\\
  \vdots & \vdots & & \vdots\\
  x_1^{(i-1)} & x_2^{(i-1)} & \cdots & x_\rho^{(i-1)}.
\end{array}
\tag{5}
\end{equation}
Así, toda función de $\Kfield^*_{\rho\rho}$ es una combinación racional de estas funciones elementales, simétricas en las filas (5), y recíprocamente toda función simétrica de las filas (5) pertenece a $\Kfield^*_{\rho\rho}$.

Así, cuando
\[
  f^*(x)=\frac{F^*(x)}{H^*(x)}
\]
recorre todas las funciones de $\Kfield^*_{\rho\rho}$, existe un sistema de infinitas relaciones:
\begin{equation}
  \frac{F^*(x)}{H^*(x)}=
  \frac{F^*(x^{(1)})}{H^*(x^{(1)})}=\cdots=
  \frac{F^*(x^{(i-1)})}{H^*(x^{(i-1)})},
\tag{6}
\end{equation}
o, en forma resumida,
\[
  F^*(x^{(k)})H^*(x^{(l)})-
  F^*(x^{(l)})H^*(x^{(k)})=0
  \quad(k,l=0,1,\ldots,i-1),
\]
donde ni $F^*(x^{(k)})$ ni $H^*(x^{(k)})$ se anula idénticamente, pues son conjugadas de $F^*(x)$ y $H^*(x)$ y formalmente simétricas.

Recíprocamente, para toda fila de cantidades $\xi$ que satisfaga las infinitas relaciones
\begin{equation}
  F^*(x)H^*(\xi)-H^*(x)F^*(\xi)=0,
  \qquad H^*(\xi)\ne0,
\tag{7}
\end{equation}
se obtiene la pertenencia a las filas de cantidades (5).

En efecto, si $G^*(x)$ denota el denominador común de los coeficientes de $\Phi^*(z,u)$, entonces por la hipótesis (7) la función
\[
  G^*(\xi)\cdot\Phi^*(z,u;\xi),
\]
que también es entera en $\xi$, no se anula idénticamente, es decir, los coeficientes de $\Phi^*(z,u;\xi)$ no se vuelven indeterminados, y tiene el cero
\[
  z=u_1\xi_1+u_2\xi_2+\cdots+u_\rho\xi_\rho.
\]
Por (7), $\Phi^*(z,u;x)$ tiene por tanto también el cero $z=u(\xi)$; es decir, $\xi$ pertenece a las filas de cantidades (5). La afirmación correspondiente vale, por (6), si $\xi$ satisface un sistema de ecuaciones respecto de una fila conjugada de cantidades:
\[
  F^*(x^{(k)})H^*(\xi)-H^*(x^{(k)})F^*(\xi)=0;
  \qquad H^*(\xi)\ne0.
\]
Así, las filas de cantidades (5) se definen como las soluciones $\xi$ del sistema de ecuaciones
\[
  F^*(x)H^*(\xi)-H^*(x)F^*(\xi)=0;
  \qquad H^*(\xi)\ne0,
\]
donde $F^*(x):H^*(x)$ recorre todas las funciones de $\Kfield^*_{\rho\rho}$; aquí la fila $x$ puede reemplazarse también por cualquier fila conjugada.

Geométricamente esto dice lo siguiente, interpretando cada fila $x$ como un punto: las soluciones de (7) representan, en un espacio lineal de dimensión $\rho$, $\infty^\rho$ grupos de $i$ puntos cada uno, de tal manera que cualquier punto del grupo determina unívocamente el grupo individual. Tal sistema de grupos de puntos se llama una ``\emph{involución en un espacio lineal}''.\footnote{Las soluciones excluidas de (7), para las cuales $F^*(\xi)=0$, $H^*(\xi)=0$, y asimismo las soluciones excluidas del sistema de ecuaciones correspondiente a la base de involución, se llaman puntos fundamentales de la involución; esto incluye también los que surgen por homogenización, los puntos ``situados en el infinito''.}

Correspondientemente, un cuerpo $\Kfield_{n\rho}$ caracteriza una involución en un espacio lineal compuesto de curvas, superficies, etcétera.

Por lo precedente, el grado $i$ de $\Omega(x_1,\ldots,x_\rho)$ respecto de $\Kfield^*_{\rho\rho}$ es igual al número de sistemas de soluciones de (7) que satisfacen la condición lateral $H^*(\xi)\ne0$, a los que podemos llamar los sistemas de soluciones de $\Kfield^*_{\rho\rho}$. Por tanto, el argumento empleado para demostrar la existencia de la base de involución permite también la siguiente formulación irracional:

``Si $\mathfrak{L}^*$ es un divisor de $\Kfield^*_{\rho\rho}$ y todo sistema de soluciones del cuerpo $\mathfrak{L}^*$ es también un sistema de soluciones de $\Kfield^*_{\rho\rho}$, entonces $\mathfrak{L}^*$ y $\Kfield^*_{\rho\rho}$ son idénticos.''

El enunciado correspondiente vale, por el principio de transferencia, para divisores $\mathfrak{L}$ de $\Kfield_{n\rho}$.

\section*{\S\ 6. Base mínima de los cuerpos $\Kfield_{n\rho}$.}

Este parágrafo da una visión de conjunto de teoremas conocidos, pero en una formulación ampliada por el principio de transferencia y por la existencia de la base racional.

\emph{Definición IV: Una base racional de $\Kfield_{n\rho}$ se llama base mínima si contiene sólo el número mínimo $\rho$ de funciones, que entonces deben ser algebraicamente independientes.}

De los teoremas de Lüroth, Castelnuovo y Enriques\footnote{J. Lüroth: Beweis eines Satzes über rationale Kurven, \emph{Math. Ann.} 9 (1875). --- G. Castelnuovo: Sulla razionalità delle involuzioni piane, \emph{Math. Ann.} 44 (1893). --- F. Enriques: Sopra una involuzione non razionale dello spazio, \emph{Rend. Acc. Linc.}, vol. 21, 21 de enero de 1912.} se obtienen los hechos:

I. Todo cuerpo $\Kfield_{n1}$ posee una base mínima, la función de Lüroth del cuerpo, determinada de modo único salvo transformación lineal fraccionaria.

II. Para cuerpos $\Kfield_{n2}$ existe siempre una base mínima, y en general sólo entonces, si se permite una extensión algebraica del cuerpo de coeficientes $\Omega$.

III. Para cuerpos $\Kfield_{n\rho}$ ($\rho\ge3$), en general no existe base mínima. Todavía no se ha intentado una caracterización de los cuerpos especiales $\Kfield_{n\rho}$ que tienen base mínima.

Indicamos brevemente la demostración del teorema de Lüroth para cuerpos $\Kfield^*_{11}$ dada por E. Steinitz:\footnote{Loc. cit., \S\ 24.}

Sea $\Phi^*(z,u)$ la forma de involución de $\Kfield^*_{11}$, y sea $\psi^*(x)$ cualquier coeficiente de $\Phi^*(z,u)$ que contenga efectivamente la indeterminada $x$. Resulta que el grado de $\Omega(x)$ respecto de $\Omega(\psi^*(x))$ es igual al grado de $\Omega(x)$ respecto de $\Kfield^*_{11}$, de donde se sigue la identidad de los cuerpos. Además, si $\Omega(\chi^*(x))$ es igual a $\Omega(\psi^*(x))$, entonces la dependencia racional mutua de $\chi^*(x)$ y $\psi^*(x)$ implica la dependencia lineal fraccionaria de las dos funciones.

Por el principio de transferencia, el teorema de Lüroth admite por tanto también la formulación:

\emph{La base mínima de los cuerpos $\Kfield_{n1}$ consta de cualquier función de la base de involución; y cualesquiera dos funciones de la base de involución de $\Kfield_{n1}$ dependen una de otra por una transformación lineal fraccionaria.}

G. Castelnuovo demuestra el Hecho II, por medio de la teoría geométrica de las involuciones, para cuerpos $\Kfield^*_{22}$ derivados de una base racional de tres funciones. Puesto que el principio de transferencia sigue siendo válido también bajo una extensión algebraica finita del cuerpo de coeficientes $\Omega$, la existencia de la base racional completa la demostración para todos los cuerpos $\Kfield_{n2}$.

Respecto de III, debe observarse que Enriques construye una involución no racional en el espacio tridimensional; es decir, un cuerpo $\Kfield_{33}$ que no tiene base mínima.


\providecommand{\Gdom}{\mathfrak{G}}
\providecommand{\Omegaint}{[\Omega]}
\providecommand{\eps}{\varepsilon}

\clearpage
\section*{\S\ 7. Sistema arbitrario $\Ssys_{n\rho}$, familia lineal $\Lsys_{n\rho}$ y dominio de integridad $\Jdom_{n\rho}$ de funciones racionales. Existencia de la base racional.}

A partir de la existencia de la base racional de los cuerpos $\Kfield_{n\rho}$ obtendremos ahora la base racional para sistemas arbitrarios de funciones racionales y racionales enteras. Ordenamos estos sistemas según el modo en que surgen sucesivamente unos de otros mediante las operaciones elementales de adición, multiplicación y división.

I. Como siempre, $\Ssys_{n\rho}$ designa un sistema arbitrario de funciones racionales (o racionales enteras) de $n$ indeterminadas y de rango algebraico $\rho$. Las cantidades de $\Omega$ se cuentan como pertenecientes al sistema.

II. El sistema $\Lsys_{n\rho}$ se llama una \emph{familia lineal} si satisface las condiciones siguientes:

1. Junto con $f(x)$, $\Lsys_{n\rho}$ contiene siempre también $c\cdot f(x)$, donde $c$ es una cantidad arbitraria de $\Omega$;

2. junto con $f(x)$ y $g(x)$, $\Lsys_{n\rho}$ contiene siempre también $f(x)+g(x)$.

III. La familia lineal $\Jdom_{n\rho}$ se llama un \emph{dominio de integridad} bajo la condición ulterior:

3. junto con $f(x)$ y $g(x)$, $\Jdom_{n\rho}$ contiene siempre también $f(x)\cdot g(x)$.

IV. El dominio de integridad $\Kfield_{n\rho}$ se llama un \emph{cuerpo} bajo la condición ulterior:

4. junto con $f(x)$ y $g(x)$ ---donde $g(x)\ne0$---, $\Kfield_{n\rho}$ contiene siempre también el cociente $f(x):g(x)$.

Las condiciones 1 a 4 son las indicadas en \S\ 1 para el cuerpo.\footnote{$f(x)$ y $g(x)$ pueden ser también de grado cero, es decir, cantidades de $\Omega$.}

Estos dominios deben derivarse ahora sucesivamente unos de otros.

a) Un sistema arbitrario $\Ssys_{n\rho}$ se amplía a una familia lineal $\Lsys_{n\rho}$ mediante la exigencia siguiente:

$\Lsys_{n\rho}$ contiene, además de $\Ssys_{n\rho}$, todas las funciones $h(x)$, y sólo éstas, que pueden representarse linealmente con coeficientes de $\Omega$ por medio de un número finito de funciones de $\Ssys_{n\rho}$:
\[
  h(x)=c_1f_1(x)+c_2f_2(x)+\cdots+c_kf_k(x),
\]
donde $f_1(x),\ldots,f_k(x)$ recorren todas las funciones de $\Ssys_{n\rho}$, y $c_1,
\ldots,c_k$ todas las cantidades de $\Omega$.

En efecto, al añadir combinaciones racionales de funciones del sistema se conserva el rango algebraico de un sistema. $\Lsys_{n\rho}$ satisface las condiciones 1 y 2, y por tanto es una familia lineal. Además, toda familia lineal que contiene $\Ssys_{n\rho}$ contiene también $\Lsys_{n\rho}$ por 1 y 2; por consiguiente $\Lsys_{n\rho}$ es la menor familia lineal que contiene $\Ssys_{n\rho}$.

b) Correspondientemente, de una familia lineal $\Lsys_{n\rho}$ surge el menor dominio de integridad que la contiene, $\Jdom_{n\rho}$, mediante la exigencia siguiente:

$\Jdom_{n\rho}$ contiene, además de $\Lsys_{n\rho}$, todas las funciones $h(x)$, y sólo éstas, que pueden representarse como una \emph{combinación racional entera} ---con coeficientes de $\Omega$--- de un número finito de funciones de $\Lsys_{n\rho}$, $f_1(x),\ldots,f_k(x)$, donde $f_1(x),\ldots,f_k(x)$ recorren todas las funciones de $\Lsys_{n\rho}$.

c) Finalmente, de un dominio de integridad $\Jdom_{n\rho}$ surge el menor cuerpo que lo contiene, $\Kfield_{n\rho}$, mediante la exigencia siguiente: $\Kfield_{n\rho}$ contiene, además de $\Jdom_{n\rho}$, todas las funciones $h(x)$, y sólo éstas, que pueden representarse como cocientes de dos funciones de $\Jdom_{n\rho}$.

Si se reúnen las tres ampliaciones en un solo paso, resulta:

\emph{Todo sistema $\Ssys_{n\rho}$ puede ampliarse al menor cuerpo que lo contiene, $\Kfield_{n\rho}$, mediante la exigencia siguiente:}

\emph{$\Kfield_{n\rho}$ contiene, además de $\Ssys_{n\rho}$, todas las funciones $h(x)$, y sólo éstas, que pueden representarse como combinaciones racionales ---con coeficientes de $\Omega$--- de un número finito de funciones de $\Ssys_{n\rho}$, $f_1(x),\ldots,f_k(x)$, donde $f_1(x),\ldots,f_k(x)$ recorren todas las funciones de $\Ssys_{n\rho}$.}

De este último hecho se sigue inmediatamente la existencia de la base racional para dominios arbitrarios:

Sea $\Kfield_{n\rho}$ el menor cuerpo que contiene $\Ssys_{n\rho}$, y sea una base racional de $\Kfield_{n\rho}$ dada por
\begin{equation}
  h_1(x),\ h_2(x),\ldots,h_\sigma(x).
\tag{1}
\end{equation}
Por la definición de $\Kfield_{n\rho}$ existen relaciones
\begin{equation}
\begin{aligned}
 h_1(x)&=r_1\bigl(g_1(x),\ldots,g_\lambda(x)\bigr);\ldots;\\
 h_\sigma(x)&=r_\sigma\bigl(g_\mu(x),\ldots,g_\nu(x)\bigr),
\end{aligned}
\tag{2}
\end{equation}
donde $r_1,
\ldots,r_\sigma$ son funciones racionales de sus argumentos, y
\begin{equation}
  g_1(x),\ldots,g_\lambda(x),\ldots,
  g_\mu(x),\ldots,g_\nu(x)
\tag{3}
\end{equation}
pertenecen todas a $\Ssys_{n\rho}$. A causa de las relaciones (2), (3) representa, además de (1), también una base racional de $\Kfield_{n\rho}$; y, puesto que las funciones (3) pertenecen a $\Ssys_{n\rho}$, es al mismo tiempo una base racional de $\Ssys_{n\rho}$. Tenemos por tanto:

\emph{Teorema IV: Todo sistema arbitrario $\Ssys_{n\rho}$ de funciones racionales (o racionales enteras) posee una base racional, compuesta por un número finito de funciones del sistema, de rango algebraico $\rho$.}

Sea ahora, en particular, $\Lsys_{n\rho}$ una familia lineal, y sea
\begin{equation}
  g_1(x),\ldots,g_\rho(x),\ g_{\rho+1}(x),\ldots,g_\nu(x)
\tag{4}
\end{equation}
una base racional de $\Lsys_{n\rho}$. Supongamos, por ejemplo, que
\[
  g_1(x),\ldots,g_\rho(x)
\]
son algebraicamente independientes. Si se pone entonces
\[
  g(x)=c_1g_{\rho+1}(x)+c_2g_{\rho+2}(x)+\cdots+c_{\nu-\rho}g_\nu(x),
\]
los $c_i$ pueden elegirse como números de $\Omega$ de tal modo que
\begin{equation}
  g_1(x),\ldots,g_\rho(x),\ g(x)
\tag{5}
\end{equation}
sea una base racional del menor cuerpo que contiene la familia, $\Kfield_{n\rho}$. Pero como $\Lsys_{n\rho}$ es una familia lineal, $g(x)$ pertenece a $\Lsys_{n\rho}$, y (5) representa una base racional de $\Lsys_{n\rho}$; es decir:

\emph{Teorema V: La base racional de toda familia lineal $\Lsys_{n\rho}$ de funciones racionales (o racionales enteras) ---por tanto, en particular, la base racional de todo dominio de integridad $\Jdom_{n\rho}$ y de todo cuerpo $\Kfield_{n\rho}$--- puede elegirse de modo especial para que contenga a lo sumo $\rho+1$ funciones (y al menos $\rho$).}

La cota para el número de funciones en la base racional encontrada en \S\ 4 para los cuerpos descansa, por tanto, en la propiedad del cuerpo de ser una familia lineal, mientras que la restricción a $\rho$ funciones en caso de existir una base mínima utiliza todas las hipótesis del cuerpo.

Queda por señalar que, por \S\ 3 (4), todas las propiedades características de los sistemas y de los sistemas mínimos que los contienen se conservan bajo la aplicación.

\section*{\S\ 8. Base de involución de los dominios de integridad $\Jdom_{n\rho}$ formados por polinomios.}

En los parágrafos siguientes tratamos sistemas $\Ssys_{n\rho}$ cuyos elementos son funciones racionales enteras (polinomios); primero, dominios de integridad $\Jdom_{n\rho}$ formados por polinomios. Para estos dominios $\Jdom_{n\rho}$ deben fijarse ante todo los conceptos de divisibilidad.

Sean $F(x),G(x)$ dos polinomios de $\Jdom_{n\rho}$, y sea $L(x)$ un divisor común de estos polinomios, de modo que se tenga
\[
  F(x)=L(x)\cdot F_1(x),\qquad G(x)=L(x)\cdot G_1(x).
\]
$L(x)$ se llama un \emph{divisor común en $\Jdom_{n\rho}$} si y sólo si los tres polinomios $L(x),F_1(x),G_1(x)$ pertenecen a $\Jdom_{n\rho}$.

Dos polinomios de $\Jdom_{n\rho}$ se llaman \emph{primos entre sí en $\Jdom_{n\rho}$} si no poseen ningún divisor común en $\Jdom_{n\rho}$.

Sea $\Kfield_{n\rho}$ el menor cuerpo que contiene $\Jdom_{n\rho}$. Entonces, por la definición, toda función de $\Kfield_{n\rho}$ posee una representación
\[
  f(x)=F_1(x):F_2(x),
\]
donde $F_1(x),F_2(x)$ pertenecen a $\Jdom_{n\rho}$ y son primos entre sí en $\Jdom_{n\rho}$. Correspondientemente, varias funciones de $\Kfield_{n\rho}$ pueden ponerse sobre un denominador común en $\Jdom_{n\rho}$:
\[
  f_1(x)=\frac{F_1(x)}{F(x)},\ldots,
  f_k(x)=\frac{F_k(x)}{F(x)}.
\]
$F(x)$ se llama \emph{un mínimo denominador común en $\Jdom_{n\rho}$} si y sólo si los polinomios de $\Jdom_{n\rho}$
\[
  F(x),F_1(x),\ldots,F_k(x),
\]
son primos entre sí en $\Jdom_{n\rho}$.

Tal representación puede ser posible de maneras distintas, pues en $\Jdom_{n\rho}$ no valen en general las leyes de la descomposición única en funciones irreducibles.

Investigamos ahora el significado de la base de involución de $\Kfield_{n\rho}$ para $\Jdom_{n\rho}$, transfiriendo los conceptos de forma de involución y base de involución a los dominios de integridad. Sea
\[
  \Phi(z,u)=z^i+\sum{}' g_{\alpha\alpha_1\ldots\alpha_\rho}(x)
  z^\alpha u_1^{\alpha_1}\cdots u_\rho^{\alpha_\rho},
  \qquad
  (\alpha+\alpha_1+\cdots+\alpha_\rho=i)
\]
una forma de involución de $\Kfield_{n\rho}$, y sea $G(x)$ un mínimo denominador común en $\Jdom_{n\rho}$ de los coeficientes de $\Phi(z,u)$. Por multiplicación por $G(x)$, $\Phi(z,u)$ pasa a una forma $\Psi(z,u)$ cuyos coeficientes son polinomios de $\Jdom_{n\rho}$ y son primos entre sí en $\Jdom_{n\rho}$:
\[
  G(x)\cdot\Phi(z,u)=\Psi(z,u)
  =G(x)z^i+
  \sum{}' G_{\alpha\alpha_1\ldots\alpha_\rho}(x)
  z^\alpha u_1^{\alpha_1}\cdots u_\rho^{\alpha_\rho},
  \qquad
  (\alpha+\alpha_1+\cdots+\alpha_\rho=i).
\]

Toda forma $\Psi(z,u)$ que surge de una forma de involución $\Phi(z,u)$ de $\Kfield_{n\rho}$ por multiplicación por un polinomio de $\Jdom_{n\rho}$, de tal manera que los coeficientes de $\Psi(z,u)$ sean polinomios de $\Jdom_{n\rho}$ y primos entre sí en $\Jdom_{n\rho}$, se llamará una \emph{forma de involución} de $\Jdom_{n\rho}$; y los coeficientes de $\Psi(z,u)$ se llamarán una \emph{base de involución} correspondiente.

Ante todo es claro, por el significado de la base de involución de $\Kfield_{n\rho}$, que toda base de involución de $\Jdom_{n\rho}$ es al mismo tiempo una base racional. Para la investigación ulterior consideramos un dominio de aplicación $\Jdom^*_{\rho\rho}$ para cuyo menor cuerpo contenedor $\Kfield^*_{\rho\rho}$ la ecuación irreducible con cero
\[
  z=u_1x_1+\cdots+u_\rho x_\rho
\]
está dada por $\Phi^*(z,u)=0$.

Según \S\ 5, los coeficientes de
\[
  \Phi^*(z,u)=z^i+\sum{}' g^*_{\alpha\alpha_1\ldots\alpha_\rho}(x)
  z^\alpha u_1^{\alpha_1}\cdots u_\rho^{\alpha_\rho},
  \qquad
  (\alpha+\alpha_1+\cdots+\alpha_\rho=i)
\]
representan las funciones simétricas elementales de las filas de cantidades conjugadas con respecto a $\Kfield^*_{\rho\rho}$,
\[
  x,\quad x^{(1)},\ldots,x^{(i-1)}.
\]
Sea $F^*(x)$ un polinomio cualquiera de $\Kfield^*_{\rho\rho}$. Puesto que
\[
  F^*(x)=\frac1i\Bigl(F^*(x)+F^*(x^{(1)})+
  \cdots+F^*(x^{(i-1)})\Bigr),
\]
se sigue que $F^*(x)$ es una función simétrica entera de estas filas de cantidades; por consiguiente puede expresarse como combinación racional entera de las funciones simétricas elementales, es decir, de las funciones
\[
  g^*_{\alpha\alpha_1\ldots\alpha_\rho}(x).
\]
Si ahora
\[
  G^*(x)\cdot\Phi^*(z,u)=\Psi^*(z,u)
  =G^*(x)z^i+
  \sum{}'G^*_{\alpha\alpha_1\ldots\alpha_\rho}(x)
  z^\alpha u_1^{\alpha_1}\cdots u_\rho^{\alpha_\rho},
  \qquad
  (\alpha+\alpha_1+\cdots+\alpha_\rho=i)
\]
es una forma de involución de $\Jdom^*_{\rho\rho}$, entonces todo polinomio de $\Kfield^*_{\rho\rho}$, y en particular todo polinomio de $\Jdom^*_{\rho\rho}$, es una combinación racional entera de los cocientes
\[
  G^*_{\alpha\alpha_1\ldots\alpha_\rho}(x):G^*(x).
\]

Teniendo en cuenta el principio de transferencia, esto da:

\emph{Teorema VI: Para todo dominio de integridad de polinomios $\Jdom_{n\rho}$ existe al menos una base racional distinguida, la base de involución, tal que en la representación racional de todos los polinomios de $\Jdom_{n\rho}$ por medio de las funciones de la base de involución sólo aparece en el denominador una potencia de una función fija (el coeficiente principal de la forma de involución correspondiente).}

\section*{\S\ 9. Dominios relativamente enteros $\Gdom_{n\rho}$ formados por polinomios.}

En lo que sigue consideramos dominios de integridad especiales formados por polinomios, que constituyen una etapa intermedia entre los dominios de integridad y los cuerpos.

Un dominio de integridad de polinomios $\Gdom_{n\rho}$ se llama ``dominio relativamente entero'' bajo la condición:

\emph{4a) $\Gdom_{n\rho}$ contiene, junto con $F(x)$ y $G(x)$ ---donde $G(x)\ne0$---, el cociente $F(x):G(x)$ siempre y sólo cuando este cociente es entero en las indeterminadas, es decir, es un polinomio.}

Primero debe demostrarse la identidad de los dominios relativamente enteros con la totalidad de los polinomios de un cuerpo $\Kfield_{n\sigma}$ ($\sigma\ge\rho$) de funciones racionales.

Sea $\Kfield_{n\rho}$ el menor cuerpo que contiene $\Gdom_{n\rho}$. Para toda función $f(x)$ de $\Kfield_{n\rho}$ hay una representación como cociente de dos polinomios de $\Gdom_{n\rho}$:
\[
  f(x)=F_1(x):F_2(x),
\]
donde se admiten también polinomios de grado cero. Tan pronto como $f(x)$ sea un polinomio, pertenece a $\Gdom_{n\rho}$ por la condición 4a); y, puesto que $\Gdom_{n\rho}$ no puede contener otros polinomios, \emph{todo dominio relativamente entero $\Gdom_{n\rho}$ es idéntico con la totalidad de los polinomios del menor cuerpo que lo contiene.}

Recíprocamente, sea $\Kfield_{n\sigma}$ un cuerpo cualquiera de funciones racionales (por tanto, no necesariamente el menor cuerpo que contiene un sistema dado), y sea $\Jdom$ la totalidad de los polinomios contenidos en $\Kfield_{n\sigma}$. $\Jdom$ satisface las condiciones 1, 2, 3 de un dominio de integridad y la condición 4a), y por tanto es un dominio relativamente entero; es decir, \emph{la totalidad de los polinomios contenidos en un cuerpo $\Kfield_{n\sigma}$ forma un dominio relativamente entero $\Gdom_{n\rho}$ ($\rho\le\sigma$).}\footnote{También puede ocurrir que $\rho=0$; es decir, que la totalidad de los polinomios de $\Kfield_{n\sigma}$ consista en elementos del dominio de coeficientes $\Omega$.}

Mostramos además la identidad de los dominios relativamente enteros con los ``dominios de integridad de funciones relativamente enteras'' de Hilbert.\footnote{D. Hilbert: Mathematische Probleme. Conferencia de París, 1900. Problema 14 (\emph{Göttinger Nachrichten} 1900).}

Sea
\begin{equation}
  G_1(x),\ldots,G_k(x)
\tag{1}
\end{equation}
una base racional de $\Gdom_{n\rho}$. Por las condiciones 1, 2, 3, 4a, toda combinación racional de los polinomios (1) que sea entera en las indeterminadas, es decir, que sea un polinomio, pertenece a $\Gdom_{n\rho}$. Recíprocamente, todo polinomio de $\Gdom_{n\rho}$ puede expresarse racionalmente, por el concepto de base racional, mediante los polinomios (1); de ahí que $\Gdom_{n\rho}$ sea \emph{idéntico con la totalidad de aquellas combinaciones racionales de los polinomios (1) que son enteras en las indeterminadas, es decir, que son polinomios.} Así tenemos la definición de Hilbert a partir de la base racional especial, mientras que nuestra definición usa sólo propiedades abstractas del dominio.

Para dominios relativamente enteros, la definición de divisor común dada en \S\ 8 para dominios de integridad generales se simplifica:

Sean $F(x),G(x)$ dos polinomios de $\Gdom_{n\rho}$, y sea $L(x)$ un divisor común de estos polinomios. Entonces $L(x)$ se llama un \emph{divisor común en $\Gdom_{n\rho}$} si y sólo si $L(x)$ pertenece a $\Gdom_{n\rho}$.

En efecto, la pertenencia de los cocientes $F(x):L(x)$ y $G(x):L(x)$ se sigue entonces de la condición 4a).

Otra propiedad esencial de los dominios relativamente enteros es que el concepto de ser ``algebraicamente entero con respecto a $\Gdom_{n\rho}$'' puede definirse por medio de cualquier ecuación exactamente como se conoce para los cuerpos de números algebraicos o para el cuerpo $\Omega(x_1,\ldots,x_n)$.

\emph{Definición V: Una cantidad $\xi$ se llama algebraicamente entera con respecto a $\Gdom_{n\rho}$ ---o también relativamente algebraicamente entera--- si satisface alguna ecuación cuyo coeficiente principal es la unidad y cuyos restantes polinomios pertenecen a $\Gdom_{n\rho}$.}

Supóngase, a saber, que la cantidad relativamente algebraicamente entera $\xi$ satisface, por ejemplo, la ecuación
\[
  L(z)=z^\lambda+G_1(x)z^{\lambda-1}+\cdots+G_\lambda(x)=0.
\]
$L(z)$ es divisible por la función entera irreducible con respecto al menor cuerpo contenedor $\Kfield_{n\rho}$,
\[
  M(z)=z^\mu+H_1(x)z^{\mu-1}+\cdots+H_\mu(x).
\]
Pero de esta divisibilidad se sigue, por la extensión conocida del teorema de Gauss,\footnote{Compárese, por ejemplo, J. König: \emph{Einleitung in die allgemeine Theorie der algebraischen Größen} (Leipzig, B. G. Teubner, 1903), cap. IX, \S\ 2, o Weber (edición menor), \S\ 20.} que las funciones racionales $H_1(x),\ldots,H_\mu(x)$ son polinomios de $\Kfield_{n\rho}$, y por tanto pertenecen a $\Gdom_{n\rho}$. Así, la definición de la cantidad relativamente algebraicamente entera $\xi$ se obtiene de la ecuación irreducible. En particular, todo polinomio $F(x)$ de $\Gdom_{n\rho}$ es también algebraicamente entero con respecto a $\Gdom_{n\rho}$, pues satisface la ecuación $z-F(x)=0$.

Debe mencionarse además que, para un dominio de integridad $\Jdom_{n\rho}$ de polinomios, análogamente al menor cuerpo contenedor, también se define el menor dominio relativamente entero que lo contiene, $\Gdom_{n\rho}$, mediante la exigencia siguiente:

$\Gdom_{n\rho}$ contiene, junto con $\Jdom_{n\rho}$, todos los polinomios $H(x)$, y sólo éstos, que pueden representarse como cocientes de dos polinomios de $\Jdom_{n\rho}$.

De esta definición se sigue además: de toda forma de involución y base de involución de $\Jdom_{n\rho}$ surge una correspondiente de $\Gdom_{n\rho}$ escindiendo, a lo sumo, un divisor común en $\Gdom_{n\rho}$ de todos los polinomios de la base.

En efecto, $\Jdom_{n\rho}$ y $\Gdom_{n\rho}$ poseen el mismo menor cuerpo contenedor $\Kfield_{n\rho}$; y toda forma de involución de $\Jdom_{n\rho}$ surge de una forma de involución de $\Kfield_{n\rho}$ por multiplicación por un polinomio de $\Jdom_{n\rho}$, y tiene por tanto coeficientes en $\Gdom_{n\rho}$, los cuales sin embargo pueden tener un divisor común en $\Gdom_{n\rho}$.\footnote{Por ejemplo, para el dominio de integridad $\Jdom_{22}$ formado por todas las combinaciones racionales enteras de $x_1^2$ y $x_1x_2$, se obtiene como forma de involución:
\[
\Psi(z,u)=x_1^2z^2-x_1^4u_1^2-2x_1^3x_2u_1u_2-x_1^2x_2^2u_2^2.
\]
Puesto que $x_1^2$ pertenece a $\Jdom_{22}$, y por tanto con mayor razón al menor dominio relativamente entero que lo contiene, $\Gdom_{22}$, y es en consecuencia un divisor común en $\Gdom_{22}$, resulta como forma de involución de $\Gdom_{22}$:
\[
\Psi(z,u)=z^2-x_1^2u_1^2-2x_1x_2u_1u_2-x_2^2u_2^2.
\]}

Finalmente debe señalarse que el \emph{dominio de aplicación} de un dominio relativamente entero $\Gdom_{n\rho}$, que por \S\ 7 vuelve a ser un dominio de integridad $\Jdom^*_{\rho\rho}$ de polinomios, no necesita ser él mismo un dominio relativamente entero. Pues si $F(x)$ y $G(x)$ son dos polinomios de $\Gdom_{n\rho}$ cuyo cociente no es un polinomio y por tanto no pertenece a $\Gdom_{n\rho}$, entonces tampoco pertenece a $\Jdom^*_{\rho\rho}$ el cociente $F^*(x):G^*(x)$. Pero este cociente, obtenido por especialización de algunas indeterminadas, puede ser perfectamente un polinomio, y entonces para $\Jdom^*_{\rho\rho}$ no se cumple la condición 4a).\footnote{Sea $\Gdom_{22}$, por ejemplo, la totalidad de los polinomios que son combinaciones racionales de $F=x_1^2x_2$ y $G=x_1^2(1+x_3)$. La especialización admisible $x_3=0$ da $F^*:G^*=x_2$, mientras que $F:G$ no es un polinomio. Por tanto, $\Jdom^*_{22}$ no es un dominio relativamente entero.}

\section*{\S\ 10. Dominios relativamente enteros de primera especie y su base de integridad.}

El concepto introducido en \S\ 9 (Definición V), ``relativamente algebraicamente entero'', permite dar una clase de dominios relativamente enteros que son dominios de integridad finitos, y responder así positivamente, al menos para esta clase, una cuestión planteada por D. Hilbert (\emph{Problemas matemáticos}, problema 14). Damos la siguiente

\emph{Definición VI: Un número finito de polinomios de $\Gdom_{n\rho}$ se llama una base de integridad si todo polinomio de $\Gdom_{n\rho}$ puede representarse como combinación racional entera de ese número finito, con coeficientes de $\Omega$. Un dominio de integridad de polinomios que posee una base de integridad se llama dominio de integridad finito.}

La clase de dominios relativamente enteros que debe considerarse, para la cual la base de involución resultará ser una base de integridad, la designamos como dominios relativamente enteros de primera especie según la siguiente

\emph{Definición VII: Un dominio relativamente entero $\Gdom_{n\rho}$ se llama de primera especie si, como dominio de aplicación, posee un dominio relativamente entero $\Gdom^*_{\rho\rho}$ tal que todo polinomio arbitrario en $x_1,\ldots,x_\rho$ (con coeficientes de $\Omega$) es algebraicamente entero sobre $\Gdom^*_{\rho\rho}$.}

Según esta definición, las indeterminadas $x_1,\ldots,x_\rho$, y por consiguiente también la forma lineal $u(x)=u_1x_1+\cdots+u_\rho x_\rho$, son algebraicamente enteras sobre $\Gdom^*_{\rho\rho}$. La función entera irreducible con cero $z=u(x)$ ---la forma de involución del menor cuerpo contenedor $\Kfield^*_{\rho\rho}$--- tiene por tanto la unidad como coeficiente principal y, como restantes coeficientes, polinomios de $\Gdom^*_{\rho\rho}$; por consiguiente es al mismo tiempo una forma de involución de $\Gdom^*_{\rho\rho}$, y en $\Gdom^*_{\rho\rho}$ no puede existir ninguna otra forma de involución. Por el principio de transferencia, también $\Gdom_{n\rho}$ posee entonces una forma de involución cuyo coeficiente principal es la unidad; y, puesto que por el Teorema VI en la representación racional de todos los polinomios de $\Gdom_{n\rho}$ por la base de involución correspondiente sólo entra en el denominador ese coeficiente principal, la base de involución se convierte en una base de integridad. Así:

\emph{Teorema VII. Todo dominio relativamente entero de primera especie es un dominio de integridad finito; la base de integridad está dada por la base de involución obtenida del dominio de aplicación que lo caracteriza. En particular, la base de integridad de los dominios relativamente enteros de primera especie $\Gdom_{\rho\rho}$ está dada por la base de involución unívocamente determinada de $\Gdom_{\rho\rho}$.}

Añadimos un criterio para dominios relativamente enteros de primera especie. Sea $\Gdom_{n\rho}$ de primera especie, y sea una base de integridad de $\Gdom_{n\rho}$ dada por
\begin{equation}
  F_1(x),F_2(x),\ldots,F_k(x).
\tag{1}
\end{equation}
Entonces, por un lema auxiliar de Hilbert,\footnote{D. Hilbert: Über die vollen Invariantensysteme, \emph{Math. Ann.} 42 (1893), \S\ 1. El lema auxiliar allí enunciado sólo para polinomios homogéneos vale igualmente para los no homogéneos.} existen $\rho$ polinomios algebraicamente independientes de $\Gdom_{n\rho}$,
\begin{equation}
  G_1(x),G_2(x),\ldots,G_\rho(x),
\tag{2}
\end{equation}
tales que los polinomios (1), y por tanto todo polinomio de $\Gdom_{n\rho}$, son algebraicamente enteros sobre los polinomios (2). Lo mismo vale para todo polinomio del dominio de aplicación $\Gdom^*_{\rho\rho}$ con respecto a los polinomios correspondientes
\begin{equation}
  G_1^*(x),G_2^*(x),\ldots,G_\rho^*(x).
\tag{3}
\end{equation}

Por la Definición VII, todo polinomio arbitrario en $x_1,\ldots,x_\rho$ es entonces también algebraicamente entero sobre los polinomios (3).

Recíprocamente, supóngase que un dominio de aplicación $\Jdom^*_{\rho\rho}$ de un dominio relativamente entero arbitrario $\Gdom_{n\rho}$ contiene $\rho$ polinomios (3) sobre los cuales todo polinomio en $x_1,\ldots,x_\rho$ es algebraicamente entero. Mostramos entonces que $\Jdom^*_{\rho\rho}$ se convierte por ello en un dominio relativamente entero $\Gdom^*_{\rho\rho}$, y además que $\Gdom^*_{\rho\rho}$, y por consiguiente $\Gdom_{n\rho}$, son de primera especie.

En efecto, sea $F^*(x)$ un polinomio del menor cuerpo $\Kfield^*_{\rho\rho}$ que contiene $\Jdom^*_{\rho\rho}$. Por hipótesis, es algebraicamente entero sobre los polinomios (3). Como la función correspondiente de $\Kfield_{n\rho}$ es entonces algebraicamente entera sobre $\rho$ polinomios de $\Gdom_{n\rho}$, es un polinomio $F(x)$ y por tanto pertenece a $\Gdom_{n\rho}$; por consiguiente $F^*(x)$ pertenece a $\Jdom^*_{\rho\rho}$. El dominio de aplicación $\Jdom^*_{\rho\rho}$ contiene por tanto todos los polinomios de $\Kfield^*_{\rho\rho}$ y es un dominio relativamente entero $\Gdom^*_{\rho\rho}$. Para $\Gdom^*_{\rho\rho}$, y por tanto también para $\Gdom_{n\rho}$, se sigue inmediatamente de las Definiciones V y VII y de la hipótesis que son de primera especie; es decir:

\emph{Los dominios relativamente enteros de primera especie $\Gdom_{n\rho}$ se caracterizan por el hecho de que en un dominio de aplicación $\Jdom^*_{\rho\rho}$ existen $\rho$ polinomios sobre los cuales todo polinomio arbitrario en $x_1,\ldots,x_\rho$ es algebraicamente entero. De esta hipótesis resulta que el propio dominio de aplicación es un dominio relativamente entero.}\footnote{Por ejemplo, en el dominio $\Gdom_{22}$ mencionado en la observación final de \S\ 9, la especialización admisible $x_1=1$ da los dos polinomios $F^*=x_2$, $G^*=1+x_3$, sobre los cuales todo polinomio arbitrario en $x_2,x_3$ es algebraicamente entero. Por tanto, $\Gdom_{22}$ es de primera especie, de hecho con $F(x)$ y $G(x)$ como base de integridad.}

\section*{\S\ 11. Lema auxiliar sobre dependencia algebraicamente entera.}

Damos ahora un criterio para que todo polinomio arbitrario en $x_1,\ldots,x_\rho$ sea algebraicamente entero sobre $\rho$ polinomios en $x_1,\ldots,x_\rho$,
\begin{equation}
  G_1^*(x),\ldots,G_\rho^*(x).
\tag{1}
\end{equation}

Si los polinomios (1) son, en particular, formas homogéneas, entonces, para todo sistema especial de valores que haga anularse (1), se anulan simultáneamente todas las formas algebraicamente enteramente dependientes; por tanto, en particular, también $x_1,\ldots,x_\rho$.

Las formas (1), por consiguiente, no pueden anularse para ningún sistema de valores distinto de
\[
  x_1=0,
  \ldots,
  x_\rho=0;
\]
su resultante debe ser distinta de cero.

Recíprocamente, esta condición es también suficiente y admite una extensión a polinomios no homogéneos según el siguiente

\emph{Lema auxiliar: Una condición suficiente para que todo polinomio arbitrario en $x_1,\ldots,x_\rho$ sea algebraicamente entero sobre $\rho$ polinomios en $x_1,\ldots,x_\rho$}
\[
  G_1^*(x),\ldots,G_\rho^*(x)
\]
\emph{es que no se anule la resultante, tomada con respecto a $x_1,\ldots,x_\rho$, de los términos de dimensión más alta}
\[
  H_1^*(x),\ldots,H_\rho^*(x)
\]
\emph{de los polinomios (1):}
\begin{equation}
  R_0=\left[
  \begin{array}{ccc}
    H_1^*&\cdots&H_\rho^*\\
    x_1&\cdots&x_\rho
  \end{array}
  \right]\ne0.
\tag{2}
\end{equation}
\emph{Para formas homogéneas (1), la condición (2) es necesaria.}\footnote{Para la teoría de resultantes, compárese F. Mertens, Zur Theorie der Elimination (I. y II. parte), \emph{Sitzungsber. d. k. Ak. d. Wiss. Wien}, Math.-naturw. Kl. Abt. IIa, Bd. 108 (1899). Reproducido parcialmente en J. König, \emph{Theorie d. algebr. Größen}, cap. VI.}

Para la demostración formamos los polinomios
\begin{equation}
\begin{gathered}
  \bar G_1^*(x)=G_1^*(x)-\lambda_1,
  \ldots,
  \bar G_\rho^*(x)=G_\rho^*(x)-\lambda_\rho,\\
  \bar\Phi^*(x)=\Phi^*(x)-\mu,
\end{gathered}
\tag{3}
\end{equation}
donde $\lambda_1,\ldots,\lambda_\rho,\mu$ designan indeterminadas y $\Phi^*(x)$ un polinomio arbitrario en $x_1,\ldots,x_\rho$. Sea la resultante de los polinomios (3) con respecto a $x_1,\ldots,x_\rho,1$\footnote{Es decir, la resultante homogénea con respecto a $x_1,\ldots,x_\rho,x_0$, si los polinomios (3) se homogeneizan mediante una indeterminada adicional $x_0$ de modo que se conserve el grado en las $x$.} dada por
\begin{equation}
  R(\lambda,\mu)=\left[
  \begin{array}{cccc}
   \bar G_1^*&\cdots&\bar G_\rho^*&\bar\Phi^*\\
   x_1&\cdots&x_\rho&1
  \end{array}
  \right];
\tag{4}
\end{equation}
es una función racional entera de $\lambda_1,\ldots,\lambda_\rho,\mu$. Según F. Mertens, en el desarrollo de (4) en potencias de $\mu$ puede indicarse explícitamente el coeficiente principal;\footnote{A. a. O. \S\ 3 (demostración del Teorema I), y explícitamente \S\ 6 (determinación de la totalidad de todos los términos de la resultante $R$ que contienen como factor un producto de potencias dado $a_{\lambda\lambda}^{p_\lambda}a_{\mu\mu}^{p_\mu}\cdots a_{\sigma\sigma}^{p_\sigma}$). En nuestro caso se toma $p_\lambda=0$, $p_\mu=0,\ldots,p_\rho=\tau$, $a_{\sigma\sigma}=(-\mu)$. Reproducido en König, cap. VI, \S\ 9.} se obtiene
\begin{equation}
  (-1)^\tau R(\lambda,\mu)
  =R_0^\sigma\mu^\tau+A_1(\lambda)\mu^{\tau-1}+\cdots+A_\tau(\lambda),
\tag{5}
\end{equation}
donde $A_1(\lambda),\ldots,A_\tau(\lambda)$ designan funciones enteras con valores enteros de $\lambda_1,\ldots,\lambda_\rho$ y de los coeficientes de los polinomios (3). Este desarrollo muestra primero que bajo la hipótesis (2), $R_0\ne0$, tampoco $R(\lambda,\mu)$ puede anularse idénticamente. La identidad en $x_1,\ldots,x_\rho,\lambda_1,\ldots,\lambda_\rho,\mu$,
\[
  R(\lambda,
\mu)\equiv 0\pmod{\bar G_1^*(x),\ldots,
\bar G_\rho^*(x),\bar\Phi^*(x)},
\]
pasa, por la sustitución
\[
  \lambda_i=G_i^*(x),\qquad \mu=\Phi^*(x),
\]
a
\begin{equation}
\begin{aligned}
  R(G_i^*(x),\Phi^*(x))
  &=R_0^\sigma\Phi^{*\tau}(x)
    +A_1(G_i^*(x))\Phi^{*\tau-1}(x)+\cdots \\
  &\qquad +A_\tau(G_i(x))=0,
\end{aligned}
\tag{6}
\end{equation}
y esto demuestra el lema auxiliar, pues $R_0$ es, por hipótesis, un número de $\Omega$.

\section*{\S\ 12. La base de integridad de sistemas regulares $\Ssys_{n\rho}$ formados por polinomios.}

El lema auxiliar de \S\ 11, junto con la conclusión de \S\ 10, da inmediatamente el hecho siguiente:

\emph{Si en un dominio de aplicación $\Jdom^*_{\rho\rho}$ de $\Gdom_{n\rho}$ existen $\rho$ polinomios tales que la resultante homogénea de sus términos de dimensión más alta no se anula idénticamente ---$R_0\ne0$---, entonces $\Gdom_{n\rho}$ es de primera especie y, por consiguiente, es un dominio de integridad finito, con la base de involución como base de integridad.}

La condición $R_0\ne0$ permite ahora, mediante el lema auxiliar, una segunda demostración de existencia de la base de integridad, debida esencialmente a D. Hilbert;\footnote{Über die vollen Invariantensysteme, \S\ 2. En el caso de Hilbert, bajo la hipótesis de la existencia obtenida de otro modo de la base de integridad del cuerpo de invariantes, se trata sólo de una interpretación de esta base mediante el sistema fundamental de Kronecker de las cantidades algebraicamente enteras de un cuerpo.} ciertamente esta demostración no hace reconocible como tal la base de involución, pero vale uniformemente para todos los sistemas $\Ssys_{n\rho}$ con la propiedad indicada.

En efecto, sean dados los $\rho$ polinomios de $\Jdom^*_{\rho\rho}$ para los cuales $R_0\ne0$, y que en consecuencia son algebraicamente independientes,
\begin{equation}
  G_1^*(x),\ldots,G_\rho^*(x).
\tag{1}
\end{equation}
Por el lema auxiliar, todo polinomio arbitrario en $x_1,\ldots,x_\rho$, y por tanto en particular todo polinomio del menor cuerpo $\Kfield^*_{\rho\rho}$ que contiene $\Jdom^*_{\rho\rho}$, es algebraicamente entero sobre los polinomios (1). Si los polinomios correspondientes en $\Ssys_{n\rho}$ son por tanto
\begin{equation}
  G_1(x),\ldots,G_\rho(x),
\tag{2}
\end{equation}
entonces, por el principio de transferencia, todo polinomio del menor cuerpo $\Kfield_{n\rho}$ que contiene $\Ssys_{n\rho}$, y por tanto en particular todo polinomio de $\Ssys_{n\rho}$, es algebraicamente entero sobre los polinomios (2). Recíprocamente, toda función de $\Kfield_{n\rho}$ que sea algebraicamente entera sobre los polinomios (2) es también un polinomio. Así, si se considera $\Kfield_{n\rho}$ (según \S\ 4) como un cuerpo algebraico sobre $\Omega(G_1,\ldots,G_\rho)$, la totalidad de los polinomios de $\Kfield_{n\rho}$ ---el dominio relativamente entero $\Gdom_{n\rho}$--- es idéntica con las cantidades algebraicamente enteras de este cuerpo. Sea $G(x)$ un polinomio de $\Kfield_{n\rho}$ que completa (2) a una base racional, y sea $D(x)$ el discriminante de la ecuación irreducible de grado $k$ que $G(x)$ satisface con respecto a $\Omega(G_1,\ldots,G_\rho)$. Entonces todo polinomio de $\Kfield_{n\rho}$, y en particular todo polinomio $F(x)$ de $\Ssys_{n\rho}$, como cantidad algebraicamente entera, admite la representación
\begin{equation}
  F(x)=\frac{\Gamma_1(G_i)G(x)^{k-1}+\Gamma_2(G_i)G(x)^{k-2}+\cdots+\Gamma_k(G_i)}{D(x)}.
\tag{3}
\end{equation}

Hagamos recorrer a $F(x)$ todos los polinomios de $\Ssys_{n\rho}$. Entonces la aplicación del Teorema I de Hilbert sobre la base de módulo (\emph{Math. Ann.} 36) al sistema formado por las funciones correspondientes
\[
  v_1\Gamma_1(G_i)+v_2\Gamma_2(G_i)+\cdots+v_k\Gamma_k(G_i)
\]
muestra la existencia de un número finito de polinomios de $\Ssys_{n\rho}$,
\begin{equation}
  F_1(x),F_2(x),\ldots,F_\lambda(x),
\tag{4}
\end{equation}
tal que todo polinomio de $\Ssys_{n\rho}$ admite una representación
\begin{equation}
  F(x)=A_1(G_i)F_1(x)+A_2(G_i)F_2(x)+\cdots+A_\lambda(G_i)F_\lambda(x),
\tag{5}
\end{equation}
donde $A_1,\ldots,A_\lambda$ designan polinomios en los $G_i(x)$. Los polinomios (2) y (4) forman por tanto una base de integridad de $\Ssys_{n\rho}$; es decir:

\emph{Teorema VIII: Si en un sistema de aplicación $\Jdom^*_{\rho\rho}$ del sistema de polinomios $\Ssys_{n\rho}$ existen $\rho$ polinomios tales que la resultante de los términos de dimensión más alta no se anula idénticamente ---$R_0\ne0$---, entonces $\Ssys_{n\rho}$ posee una base de integridad.}

Este teorema admite una ligera generalización. Basta que los $\rho$ polinomios (1) para los cuales $R_0\ne0$ pertenezcan al dominio de integridad $\Jdom^*_{\rho\rho}$ que es el menor dominio contenedor de $\Ssys^*_{\rho\rho}$. En efecto, los argumentos que conducen al Teorema VIII muestran que la representación (5) sigue siendo válida. Pero, por la definición, todo polinomio $L(x)$ del menor dominio de integridad contenedor $\Jdom_{n\rho}$ puede representarse de manera racional entera mediante un número finito ---dependiente de $L(x)$--- de polinomios de $\Ssys_{n\rho}$; por consiguiente existen $\sigma$ polinomios de $\Ssys_{n\rho}$,
\begin{equation}
  K_1(x),K_2(x),\ldots,K_\sigma(x),
\tag{6}
\end{equation}
tales que los polinomios $G_i(x)$ se convierten en combinaciones racionales enteras de (6). Así, (6) y (4) dan una base de integridad. Introduciendo la terminología siguiente:

\emph{Definición VIII: Un sistema $\Ssys_{n\rho}$ de polinomios se llama regular si en un dominio de aplicación $\Jdom^*_{\rho\rho}$ del menor dominio de integridad que lo contiene existen $\rho$ polinomios tales que la resultante de sus términos de dimensión más alta no se anula idénticamente ---$R_0\ne0$---}

tenemos por tanto:

\emph{Teorema IX: Todo sistema regular $\Ssys_{n\rho}$ de polinomios posee una base de integridad.}\footnote{Que hay sistemas no regulares sin base de integridad fue mostrado por Hilbert con un ejemplo (\emph{Gött. Nachr.} 1891, p. 233). El siguiente es aún más sencillo:
\[
  \Ssys_{22}=x_1,
  x_1x_2,
  x_1x_2^2,
  \ldots,
  x_1x_2^\nu,
  \ldots\quad\text{in inf.}
\]
La existencia de la base de integridad aquí equivaldría a la existencia de un entero positivo $\nu$ tal que se verifiquen las identidades
\[
  x_1x_2^{\nu+\sigma}=\sum C\cdot x_1^{i_0}(x_1x_2)^{i_1}(x_1x_2^2)^{i_2}\cdots(x_1x_2^\nu)^{i_\nu}
  \quad(\sigma=1,2,\ldots).
\]
La comparación de las dimensiones en $x_1,x_2$, ya para $\sigma=1$, conduce a las dos ecuaciones condicionales para los exponentes
\[
  1=i_0+i_1+\cdots+i_\nu,
  \qquad
  \nu+1=i_1+2i_2+\cdots+\nu i_\nu,
\]
que claramente no tienen solución en enteros no negativos. Por tanto $\Ssys_{22}$ no posee base de integridad.}

Damos además, análogamente a \S\ 5 para la involución, una interpretación geométrica de los sistemas regulares. Para ello consideramos, como allí en (7), el sistema de ecuaciones
\begin{equation}
  L^*(x)-L^*(\xi)=0,
\tag{7}
\end{equation}
donde $L^*(x)$ recorre todos los polinomios de $\Jdom^*_{\rho\rho}$. Al homogenizar mediante $x_0,\xi_0$, (7) pasa al sistema de ecuaciones entre formas homogéneas
\begin{equation}
  \xi_0^{m}M^*(x)-x_0^{m}M^*(\xi)=0,
\tag{8}
\end{equation}
donde, si $H^*(x)$ designa los términos de dimensión más alta de $L^*(x)$, vale la relación
\begin{equation}
  H^*(x)=M^*(x)\pmod{x_0}.
\tag{9}
\end{equation}

Como \emph{puntos fundamentales} de $\Jdom^*_{\rho\rho}$ (compárese la nota a \S\ 5), designamos las soluciones de (8) que son independientes de $x$, es decir, distintas de las filas conjugadas a $x$. Los puntos fundamentales son, por tanto, los sistemas de valores distintos de $\xi_1=0,\ldots,\xi_\rho=0$ para los cuales se cumple simultáneamente
\[
  \xi_0^m=0,
  \qquad M^*(\xi)=0,
\]
o, por (9),
\[
  \xi_0=0,
  \qquad H^*(\xi)=0.
\]

Si existen ahora $\rho$ formas $H^*(\xi)$ para las cuales $R_0\ne0$, y por tanto sin sistema común de soluciones, entonces $\Jdom^*_{\rho\rho}$ no puede poseer ningún punto fundamental. Si, en particular, $\Ssys^*_{\rho\rho}$ consiste en formas homogéneas, entonces para $R_0\ne0$ tampoco $\Ssys^*_{\rho\rho}$ puede poseer ningún punto fundamental, puesto que tal punto sería al mismo tiempo un punto fundamental de $\Jdom^*_{\rho\rho}$. Pero también vale el recíproco. Para ello consideramos el sistema $\Jdom(H^*)$ formado por la totalidad de las formas $H^*(x)$; éste vuelve a ser un dominio de integridad.

Del Teorema I de Hilbert sobre la base de módulo y del lema auxiliar de Hilbert (compárese la cita en \S\ 10) se sigue la existencia de $\rho$ formas de $\Jdom(H^*)$ cuya anulación tiene como consecuencia la anulación de todas las formas de $\Jdom(H^*)$. Si ahora $\Jdom^*_{\rho\rho}$ no posee ningún punto fundamental, estas $\rho$ formas no pueden anularse simultáneamente; su resultante $R_0\ne0$.

Así:

\emph{Los sistemas regulares $\Ssys_{n\rho}$ de polinomios se caracterizan por el hecho de que un dominio de aplicación $\Jdom^*_{\rho\rho}$ del menor dominio de integridad que los contiene no posee ningún punto fundamental. Para sistemas regulares $\Ssys^*_{\rho\rho}$ de formas homogéneas, el propio sistema no posee ya ningún punto fundamental.}\footnote{El sistema no regular $\Ssys_{22}$ dado en la nota precedente posee el punto fundamental
\[
  \xi_1:\xi_2:\xi_0=0:1:0.
\]}

\section*{\S\ 13. Ejemplos de dominios relativamente enteros de primera especie y de sistemas regulares.}

1. Como primer ejemplo de un dominio relativamente entero de primera especie $\Gdom_{nn}$, consideramos la \emph{totalidad de los polinomios en $x_1,\ldots,x_n$ que admiten un grupo de permutaciones dado $G$}; es decir, la totalidad de los polinomios de un dominio de género de Lagrange. Puesto que $x_1,\ldots,x_n$, en virtud de la identidad
\begin{equation}
  x_k^n-\sigma_1(x)x_k^{n-1}+\sigma_2(x)x_k^{n-2}+\cdots\pm\sigma_n(x)=0
  \qquad(k=1,2,
\ldots,n),
\tag{1}
\end{equation}
son algebraicamente enteros sobre las funciones simétricas elementales pertenecientes a $\Gdom_{nn}$, y por tanto, por la Definición V, algebraicamente enteros sobre $\Gdom_{nn}$, $\Gdom_{nn}$ es de primera especie. Además $\Gdom_{nn}$, como dominio de primera especie, puesto que sus elementos son formas homogéneas y $n=\rho$, forma un sistema regular (por la conclusión de \S\ 10). En efecto, por (1), la resultante $R_0$ de las funciones simétricas elementales no puede anularse idénticamente; por las reglas de cálculo de la teoría de resultantes, se calcula $R_0=\pm1$. La \emph{forma de involución} de $\Gdom_{nn}$ está determinada unívocamente ---pues $\Gdom_{nn}$ es de primera especie y $n=\rho$--- como la forma de involución del correspondiente dominio de género de Lagrange. Como las cantidades conjugadas a $x_1,\ldots,x_n$ con respecto a este dominio surgen por las permutaciones del grupo $G$, esta forma de involución está dada por
\begin{equation}
\begin{aligned}
  \Psi(z,u)&=\prod_i
  (z-u_1x_{i_1}-u_2x_{i_2}-\cdots-u_nx_{i_n})\\
  &=z^i+\sum{}'G_{\alpha\alpha_1\ldots\alpha_n}(x)
  z^\alpha u_1^{\alpha_1}\cdots u_n^{\alpha_n},
\end{aligned}
\tag{2}
\end{equation}
donde el producto debe extenderse sobre todas las permutaciones de $G$. La fórmula (2) representa la ``forma de Galois del grupo $G$'', y por tanto vale el hecho siguiente:

\emph{Todos los polinomios en $x_1,\ldots,x_n$ que admiten un grupo de permutaciones $G$ son combinaciones racionales enteras de los coeficientes $G_{\alpha\alpha_1\ldots\alpha_n}(x)$ de la forma de Galois del grupo.}

2. Como ejemplo de \emph{sistemas regulares}, considérense los sistemas $\Ssys_{n1}$ de polinomios, es decir, sistemas que contienen sólo un polinomio algebraicamente independiente.

En efecto, sea $\Ssys^*_{11}$ un sistema de aplicación de $\Ssys_{n1}$ y sea
\[
  F^*(x)=a_0x^k+a_1x^{k-1}+\cdots+a_k\qquad(a_0\ne0)
\]
un polinomio de $\Ssys^*_{11}$. Entonces
\begin{equation}
  R_0=a_0\ne0.
\tag{3}
\end{equation}
Así $\Ssys_{n1}$ es regular por la Definición VIII, y en consecuencia:

\emph{Todo sistema $\Ssys_{n1}$ de polinomios, en particular todo sistema de infinitos polinomios en una indeterminada, posee una base de integridad.}\footnote{Los sistemas no regulares más sencillos posibles sin base de integridad son por tanto sistemas $\Ssys_{22}$. Que tales sistemas existen efectivamente lo muestran el ejemplo de Hilbert y el ejemplo dado en la nota a \S\ 12.}

3. Especializamos ahora los sistemas $\Ssys_{n1}$ a dominios relativamente enteros $\Gdom_{n1}$, que, como sistemas regulares por \S\ 12, son dominios relativamente enteros de primera especie. Su base de integridad está por tanto dada por una base de involución, que al mismo tiempo es la base de involución del menor cuerpo contenedor $\Kfield_{n1}$. Pero por \S\ 6, cualquier función de la base de involución de $\Kfield_{n1}$ es al mismo tiempo una base mínima ---la llamamos función de Lüroth del cuerpo--- y toda función de la base de involución depende de esta función de Lüroth mediante una transformación lineal fraccionaria. En el caso presente la función de Lüroth, puesto que pertenece a $\Gdom_{n1}$, es un polinomio; todo otro polinomio de la base de involución depende por tanto lineal-fraccionariamente, y por (3) algebraicamente enteramente, luego entera y linealmente, de esta función de Lüroth $L(x)$, que así se convierte en la base de integridad. La forma de involución de $\Gdom^*_{11}$, poniendo también $(u_1=1)$, es
\[
  \Psi^*(z)=L^*(z)-L^*(x),
\]
de donde $L(x)$ aparece de nuevo como base de integridad. Así:

\emph{Los dominios relativamente enteros $\Gdom_{n1}$ son de primera especie; su base de integridad consta de un polinomio, la función de Lüroth del menor cuerpo contenedor.}

4. Como ejemplo de un dominio relativamente entero $\Gdom_{n1}$, mencionemos especialmente los invariantes proyectivos racionales enteros de una forma cuadrática en $s$ variables. De acuerdo con lo precedente, es sabido que pueden representarse como combinaciones racionales enteras del determinante $|a_{ik}|$ de la forma; y este determinante es al mismo tiempo la función de Lüroth del cuerpo de invariantes correspondiente.

\section*{\S\ 14. Sistemas formados por polinomios enteros.}

Los teoremas de base obtenidos deben ahora afinarse con respecto a la integralidad sobre los enteros.

En lo que sigue, el dominio de coeficientes $\Omega$ se supondrá por tanto un cuerpo de números algebraicos (finito), y $[\Omega]$ designará la totalidad de los enteros algebraicos de $\Omega$, que llamaremos brevemente enteros. Por polinomio entero se entiende un polinomio con coeficientes en $[\Omega]$; de ahora en adelante, $F(x),G(x),\ldots$ designarán sólo polinomios enteros.

En analogía con \S\ 7, consideramos los distintos sistemas compuestos de polinomios enteros:

Un sistema arbitrario de polinomios enteros se designará como un sistema entero $\Ssys_{n\rho}$.

La familia lineal entera $\Lsys_{n\rho}$ es un sistema entero que, junto con $F_1(x)$ y $F_2(x)$, contiene siempre $c_1F_1(x)+c_2F_2(x)$, donde $c_1,c_2$ son enteros arbitrarios de $[\Omega]$.

El dominio de integridad entero $\Jdom_{n\rho}$ es una familia lineal entera que, junto con $F(x)$ y $G(x)$, contiene siempre $F(x)\cdot G(x)$.

El dominio relativamente entero entero $\Gdom_{n\rho}$ es un dominio de integridad que, junto con $F(x)$ y $G(x)$ ---donde $G(x)\ne0$---, contiene siempre y sólo contiene el cociente $F(x):G(x)$ cuando este cociente es un polinomio entero.

Correspondientemente a \S\ 7, el menor dominio de integridad entero $\Jdom_{n\rho}$ que contiene un sistema entero $\Ssys_{n\rho}$ consta de todos los polinomios $H(x)$ que son combinaciones racionales enteras, en el sentido entero, de un número finito de polinomios de $\Ssys_{n\rho}$;

el menor dominio relativamente entero entero $\Gdom_{n\rho}$ que lo contiene consta de todos los polinomios enteros $H(x)$ que son combinaciones racionales de un número finito de polinomios de $\Ssys_{n\rho}$.

Pasamos ahora a la transferencia de los teoremas de base. Con respecto a la base racional, debe observarse que los conceptos de cuerpo, menor cuerpo contenedor y base racional, que descansan sólo en la ``racionalidad'', no sufren cambio. Además, el sistema de aplicación $\Ssys^*_{\rho\rho}$ de un sistema entero $\Ssys_{n\rho}$ puede elegirse de nuevo, de maneras arbitrariamente numerosas, como sistema entero. Puesto que también para la familia lineal entera siguen siendo válidos los razonamientos que llevan a la cota del número de generadores, se obtiene:

\emph{Teorema V$'$: Todo sistema entero $\Ssys_{n\rho}$ posee una base racional; la base racional de una familia lineal entera $\Lsys_{n\rho}$ puede elegirse en particular de modo que contenga a lo sumo $\rho+1$ polinomios.}

Investigamos ahora el análogo de la base de involución de los dominios de integridad de polinomios (\S\ 8). Para ello, el teorema sobre las \emph{funciones simétricas de filas de cantidades} debe extenderse con respecto a la \emph{integralidad sobre los enteros} por medio del siguiente

\emph{Lema auxiliar. Las funciones simétricas enteras, en el sentido de los enteros, de $n$ filas de cantidades $(yz\cdots t)$:}
\begin{equation}
\begin{array}{cccc}
  y_1&z_1&\cdots&t_1\\
  y_2&z_2&\cdots&t_2\\
  \vdots&\vdots&&\vdots\\
  y_n&z_n&\cdots&t_n
\end{array}
\tag{1}
\end{equation}
\emph{son funciones enteras, en el sentido de los enteros, de un número finito de entre ellas, a saber, de las funciones simétricas elementales}
\begin{equation}
\begin{gathered}
  \sigma_1(y),\sigma_2(y),\ldots,\sigma_n(y);\quad
  \sigma_1(z),\sigma_2(z),\ldots,\sigma_n(z);\quad\ldots;\\
  \sigma_1(t),\ldots,\sigma_n(t);
\end{gathered}
\tag{2}
\end{equation}
\emph{y de las funciones}
\begin{equation}
  \chi_1(y,z,\ldots,t),\ldots,\chi_k(y,z,\ldots,t),
\tag{3}
\end{equation}
\emph{que son algebraicamente enteras y enteras sobre las funciones (2).}

En efecto, las $n$ cantidades $y_1,\ldots,y_n$ dependen algebraicamente entera e integralmente sobre los enteros de $\sigma_1(y),\ldots,\sigma_n(y)$; la afirmación correspondiente vale para $z_1,\ldots,z_n$ respecto de $\sigma_1(z),\ldots,\sigma_n(z)$, y así sucesivamente. El cuerpo de funciones simétricas de las filas de cantidades (1) forma por tanto un cuerpo algebraico sobre el cuerpo determinado por las funciones (2), de tal manera que las cantidades algebraicamente enteras e integrales sobre los enteros de este cuerpo son idénticas con las funciones simétricas enteras, en el sentido de los enteros, de las filas de cantidades (1). Si se toma entonces (3) como sistema fundamental de Kronecker, el lema auxiliar queda demostrado.

Sea ahora $\Jdom_{n\rho}$ un dominio de integridad entero, $\Jdom^*_{\rho\rho}$ un dominio de aplicación entero, y
\[
  \Psi^*(z,u)=G^*(x)z^i+
  \sum{}'G^*_{\alpha\alpha_1\ldots\alpha_\rho}(x)
  z^\alpha u_1^{\alpha_1}\cdots u_\rho^{\alpha_\rho}
\]
una forma de involución de $\Jdom^*_{\rho\rho}$; aquí $G^*(x)$ puede también tener dimensión cero, es decir, ser un número de $[\Omega]$. Las funciones simétricas elementales correspondientes a las funciones (2) están entonces dadas por ciertas funciones entre los cocientes
\begin{equation}
  G^*_{\alpha\alpha_1\ldots\alpha_\rho}(x):G^*(x).
\tag{4}
\end{equation}
Como por el lema auxiliar las funciones (3) dependen algebraicamente entera e integralmente sobre los enteros de (2), las funciones de $\Kfield^*_{\rho\rho}$ correspondientes a (3) están dadas ---por la extensión del teorema de Gauss a polinomios con coeficientes algebraicamente enteros--- por
\begin{equation}
  K_\nu^*(x):(G^*(x))^\sigma,
\tag{5}
\end{equation}
donde $K_\nu^*(x)$ son polinomios enteros de $\Jdom^*_{\rho\rho}$. Así, $K_\nu^*(x)$ pertenece al dominio si se trata de dominios relativamente enteros; en general son cocientes de polinomios de $\Jdom^*_{\rho\rho}$.

Si ahora $F^*(x)$ es un polinomio entero arbitrario de $\Jdom^*_{\rho\rho}$, entonces se tiene
\[
  F^*(x)=\frac1i\Bigl(F^*(x)+F^*(x^{(1)})+\cdots+F^*(x^{(i-1)})\Bigr),
\]
y por tanto $i\cdot F^*(x)$ se convierte en una función simétrica entera, en el sentido de los enteros, de las filas de cantidades $x,x^{(1)},\ldots,x^{(i-1)}$.

Por el lema auxiliar y el principio de transferencia, por consiguiente:

\emph{Teorema VI$'$: Para dominios de integridad enteros $\Jdom_{n\rho}$, cada forma de involución determina una base racional distinguida de tal manera que todo polinomio $F(x)$ de $\Jdom_{n\rho}$ admite una representación}
\[
  F(x)=\frac{\Gamma(H(x),H_1(x),\ldots,H_s(x))}{i\cdot H(x)^i},
\]
\emph{donde $\Gamma$ designa una función entera en el sentido de los enteros. Para dominios relativamente enteros enteros $\Gdom_{n\rho}$, con dominio de aplicación relativamente entero entero $\Gdom^*_{\rho\rho}$, el denominador $H(x)$ está dado por el coeficiente principal de la forma de involución correspondiente.}

\section*{\S\ 15. Dominios relativamente enteros enteros de primera especie y sistemas regulares.}

Para sistemas enteros, la base de integridad entera ocupa el lugar de la base de integridad; es decir, un número finito de polinomios enteros del sistema tal que todo polinomio entero del sistema puede representarse como combinación entera, en el sentido de los enteros, de ese número finito.

Los teoremas sobre la base de integridad entera corren esencialmente paralelos a los anteriores, y sólo las demostraciones requieren ligeras modificaciones.

Primero, la Definición V (\S\ 9) se transfiere directamente.

\emph{Definición V$'$: Una cantidad $\xi$ se llama algebraicamente entera y entera con respecto al dominio relativamente entero entero $\Gdom_{n\rho}$ si satisface alguna ecuación cuyo coeficiente principal es una unidad de $[\Omega]$ y cuyos restantes polinomios pertenecen a $\Gdom_{n\rho}$.}

La extensión del teorema de Gauss a polinomios con coeficientes algebraicamente enteros muestra, como en \S\ 9, que de esto puede obtenerse la definición por la ecuación irreducible.

De la Definición V$'$ se obtiene además la transferencia de la Definición VII.

\emph{Definición VII$'$: Un dominio relativamente entero entero $\Gdom_{n\rho}$ se llama de primera especie si, como dominio de aplicación, posee un dominio relativamente entero $\Gdom^*_{\rho\rho}$ tal que todo polinomio entero arbitrario en $x_1,\ldots,x_\rho$ es algebraicamente entero y entero sobre $\Gdom^*_{\rho\rho}$.}

Para mostrar la finitud de estos dominios de primera especie, observamos primero que, como en \S\ 10, se sigue de la definición que el coeficiente principal de la forma de involución es una unidad de $[\Omega]$.

Por el Teorema VI$'$ (\S\ 14), todo polinomio $F(x)$ de $\Gdom_{n\rho}$ admite por tanto una representación
\[
  F(x)=\frac{\Gamma(H_1(x),\ldots,H_\sigma(x))}{i}.
\]
Aplicando aquí al sistema formado por los polinomios enteros $\Gamma(H_1,\ldots,H_\sigma)$ el Teorema II de Hilbert sobre la base de módulo entera (\emph{Math. Ann.} 36), enunciado originalmente para coeficientes numéricos racionales enteros, en su extensión a polinomios algebraicamente enteros,\footnote{Si $w_1,\ldots,w_k$ designa un sistema fundamental de los enteros algebraicos de $\Omega$, y se pone
\[
  \Gamma(H_1,
\ldots,H_\sigma)=\Gamma_1(H)w_1+\cdots+\Gamma_k(H)w_k,
\]
entonces $\Gamma_1(H),\ldots,\Gamma_k(H)$ tienen coeficientes numéricos racionales enteros. La aplicación del Teorema II al sistema formado por las formas $v_1\Gamma_1(H)+\cdots+v_k\Gamma_k(H)$ muestra directamente la validez de la extensión.} se sigue la existencia de la base de integridad entera; es decir:

\emph{Teorema VII$'$: Todo dominio relativamente entero entero de primera especie posee una base de integridad entera.}

Transferimos además la definición de sistemas regulares:

\emph{Definición VIII$'$: Un sistema entero $\Ssys_{n\rho}$ se llama entero-regular si en un dominio de aplicación $\Jdom^*_{\rho\rho}$ del menor dominio de integridad entero que lo contiene existen $\rho$ polinomios tales que la resultante de los términos de dimensión más alta es una unidad de $[\Omega]$ ---$R_0=\eps$.}

Para transferir la demostración de existencia de la base de integridad, observamos que el lema auxiliar de \S\ 11 admite la siguiente formulación más precisa, que se sigue de las ecuaciones (5) y (6) de \S\ 11 y del hecho de que $A_1(\lambda),\ldots,A_\tau(\lambda)$ son funciones enteras, en el sentido de los enteros, de $\lambda_1,\ldots,\lambda_\rho$:

\emph{Una condición suficiente para que todo polinomio entero arbitrario en $x_1,\ldots,x_\rho$ dependa algebraicamente entera e integralmente sobre los enteros de $\rho$ polinomios enteros es que la resultante de los términos de dimensión más alta de estos polinomios sea una unidad de $[\Omega]$ ---$R_0=\eps$.}

De este lema auxiliar, exactamente como en \S\ 12, se sigue para todo polinomio $F(x)$ de un sistema entero-regular la representación (3) de \S\ 12, donde ahora $\Gamma_1(G_i),\ldots,\Gamma_k(G_i)$ son polinomios enteros en los $G_i$. Además, como en \S\ 12, la aplicación del Teorema II de Hilbert (\emph{Math. Ann.} 36) en su extensión a polinomios algebraicamente enteros, junto con la definición del menor dominio de integridad entero que lo contiene, da la existencia de la base de integridad entera; es decir:

\emph{Teorema IX$'$: Todo sistema entero-regular posee una base de integridad entera.}

Como ejemplo de dominio relativamente entero entero de primera especie, en analogía con el ejemplo 1 de \S\ 13, nos remitimos a la totalidad $\Gdom_{nn}$ de los polinomios enteros de un dominio de género de Lagrange. Que $\Gdom_{nn}$ es entero de primera especie se sigue de la identidad
\[
  x_k^n-\sigma_1(x)x_k^{n-1}+\cdots\pm\sigma_n(x)=0
  \qquad(k=1,2,
\ldots,n);
\]
puesto que la resultante de las funciones simétricas elementales tiene el valor $\pm1$, $\Gdom_{nn}$ es también un sistema entero regular.

Un ejemplo adicional está dado, por el lema auxiliar de \S\ 14, por la totalidad de las funciones simétricas enteras, en el sentido de los enteros, de $n$ filas de cantidades.

Erlangen, mayo de 1914.
\end{document}

